\documentclass[11pt,a4paper,oneside]{book}
\usepackage[T1]{fontenc}
\usepackage[utf8]{inputenc}
\usepackage[english]{babel}
\usepackage{lmodern,microtype}
\usepackage[margin=25mm,headheight=15pt]{geometry}
\usepackage{amsmath,amssymb,amsthm,mathtools,bm}
\usepackage{booktabs,longtable,tabularx,array}
\usepackage{graphicx,xcolor,enumitem,fancyhdr,titlesec,tocloft}
\setlength{\cftchapnumwidth}{4.5em}
\usepackage[unicode,colorlinks=true,linkcolor=blue!45!black,urlcolor=blue!55!black,citecolor=blue!45!black]{hyperref}
\hypersetup{pdftitle={META--CAUSAL: The Geometry of Predictive Memory and the Limits of Reconstructibility},pdfauthor={Andrei Preece and Boris Bateni}}
\definecolor{ink}{HTML}{173A5E}
\pagestyle{fancy}\fancyhf{}
\fancyhead[L]{\small\sffamily META--CAUSAL}
\fancyhead[R]{\small\sffamily Invariant Reconstruction}
\fancyfoot[C]{\thepage}
\renewcommand{\headrulewidth}{0.3pt}
\titleformat{\chapter}[hang]{\Large\bfseries\color{ink}}{\thechapter}{1em}{}
\titlespacing*{\chapter}{0pt}{10pt}{18pt}
\renewcommand{\thechapter}{\Roman{chapter}}
\renewcommand{\theequation}{\arabic{equation}}
\counterwithout{equation}{chapter}
\setlength{\parindent}{1.2em}\setlength{\parskip}{3pt}
\setlist{itemsep=3pt,topsep=5pt}
\setcounter{tocdepth}{0}
\emergencystretch=2em
\newcommand{\R}{\mathbb R}\newcommand{\C}{\mathbb C}\newcommand{\Z}{\mathbb Z}\newcommand{\N}{\mathbb N}\newcommand{\Q}{\mathbb Q}
\newcommand{\one}{\mathbf 1}\newcommand{\ii}{\mathrm i}\newcommand{\dd}{\mathrm d}
\DeclareMathOperator{\rank}{rank}\DeclareMathOperator{\Tr}{Tr}
\DeclareMathOperator{\spec}{spec}\DeclareMathOperator{\diag}{diag}
\DeclareMathOperator{\lcm}{lcm}\DeclareMathOperator{\spanR}{span_{\mathbb R}}
\newcommand{\statement}[3]{\par\medskip\noindent\textbf{#1 #2 (#3).}\ }
\newcommand{\note}[1]{\par\medskip\noindent\textbf{#1}\ }
\newcommand{\sourcepage}[1]{\par\smallskip\noindent{\scriptsize\color{gray}Source page #1}\par\smallskip}
\begin{document}
\begin{titlepage}
\centering
{\Large\sffamily\color{ink} META--CAUSAL\par}
\vspace{14mm}
{\Huge\bfseries The Geometry of Predictive Memory and the Limits of Reconstructibility\par}
\vspace{8mm}
{\Large Classical and Quantum Realizations, Spectral Invariants, and Divisor Hypergraphs\par}
\vspace{10mm}
{\large Research preprint\par}
\vspace{4mm}
{\large Andrei Preece and Boris Bateni\par}
{Independent Researchers\par Berlin, 2026\par}
\vfill
\begin{minipage}{0.94\textwidth}
\begin{center}\textbf{Abstract}\end{center}
This work studies the operational limits of predictive memory: linear, nonnegative, and normalized stochastic realizations; spectral and finite-time witnesses; exact uniform approximation of periodic signals; and information reconstruction under restricted observation. For a pure periodic spike, the classical minimax error reduces to a fractional problem on a divisor hypergraph. Later chapters establish exact staircases for three and four distinct prime divisors, a universal rational polytope, and the three upper structural levels for an arbitrary number of distinct prime divisors. Independent sections study finite-sample topology in open quantum systems, Berger spectral geometry, exact Hopf phase reduction and an integral of homogeneous Ricci flow, restricted Einstein--QFI--Hopf models, and specific identifiability questions. Rigorous results are distinguished from computational certificates, conditional physical models, and open problems. The compositional nature of several results and the absence of independent verification of worldwide priority are stated in the corresponding theorems.

\medskip\noindent\textbf{Scientific status and reading convention.} The complete chain of proofs is retained: statements, assumptions, counterexamples, and computational certificates form parts of a single investigation. Later corrections take precedence over historical conjectures that are explicitly refuted; numerical checks do not replace proofs. Geometric and gravitational bridges in which an action or coupling is postulated are not universal consequences of information theory. The manuscript has not undergone independent peer review or complete formal verification in Lean.
\end{minipage}
\vfill
\end{titlepage}
\tableofcontents
% 00_editorial.tex
\chapter*{Research Synthesis in English: Finite-State Memory Thresholds}
\addcontentsline{toc}{chapter}{Research synthesis: finite-state memory thresholds}
\begin{center}
{\small\sffamily EDITORIAL CORE · ENGLISH · SELF-CONTAINED RESEARCH SYNTHESIS · v37}\par\medskip
{\Large\bfseries Finite-State Memory Thresholds for Periodic Observations}\par\smallskip
Exact minimax Markov approximation, prime-power clocks and operational identification\par\smallskip
\textit{Integrated research edition, 13 October 2026. Not peer reviewed; priority not independently established.}
\end{center}
\note{Scope.} We study what finite homogeneous classical Markov systems can reproduce or approximate about prescribed periodic binary-output probability signals. The central result identifies the exact state budget at which the best infinite-horizon uniform approximation first improves beyond a constant, and gives the exact improvement at that budget. The complete technical record, translated from Russian, and historical corrections remain integrated below. Quantum examples, matrix factorizations and LLM probes are different comparison classes and are not consequences of the central approximation theorem.
\section*{1. The central scientific question}
Suppose we observe a calibrated binary experiment at equally spaced times. Its ideal probability is $f(k)\in[0,1]$. A classical surrogate with $m$ internal states has a single time-homogeneous transition matrix $M$, preparation $\mu$ and effect $e$:
\begin{equation}g(k)=e^{\mathsf T}M^k\mu,\qquad M\ge0,\quad\one^{\mathsf T}M=\one^{\mathsf T},\quad\mu\in\Delta_{m-1},\quad e\in[0,1]^m.\tag{1}\end{equation}
Our loss is the worst-case discrepancy over all nonnegative integer times, not a sample-average error:
\begin{equation}E_m(f)=\inf_{M,\mu,e}\sup_{k\ge0}|f(k)-g(k)|.\tag{2}\end{equation}
The challenge is not merely whether $E_m=0$, but exactly when an additional internal state starts to buy a strict improvement. This matters in automata theory, positive dynamical realizations, quantum--classical simulation and long-horizon system identification. It is not a general claim about arbitrary physical observers or Transformer compression.
\note{Operational boundaries.} The source experiment is prescribed, the hidden state count is finite, $M$ is time independent, and $e$ is a single bounded binary effect. We do not permit arbitrary time-dependent clock inputs, observation-conditioned interventions, reset operations, or external memory. Changing any of these assumptions changes the minimax problem.
\section*{2. Main theorem: the prime-power clock law}
Every integer $n\ge2$ has a unique factorization into pairwise coprime maximal prime powers
\begin{equation}n=\prod_{i=1}^r q_i,\qquad q_i=p_i^{\alpha_i},\quad p_i\ne p_j\ \text{for }i\ne j;\qquad K(n)=\sum_{i=1}^r q_i.\tag{3}\end{equation}
Take an $n$-periodic two-level spike signal
\begin{equation}f_{n,a,b}(k)=a+b\one_{\{n\mid k\}},\qquad a,b\ge0,\quad a+b\le1.\tag{4}\end{equation}
The integer $K(n)$ may be far smaller than $n$: for $n=210=2\cdot3\cdot5\cdot7$, $K(210)=17$.
\statement{Theorem}{1}{Exact first improvement threshold and its value} For (1)--(4) and every $n\ge2$, write $r=\omega(n)$ for the number of distinct prime divisors and $K=K(n)$. Then
\begin{equation}E_m(f_{n,a,b})=\frac b2\quad(1\le m<K),\qquad E_K(f_{n,a,b})=\max\left\{\frac{(r-1)b}{2r},\frac{(r-1)b-a}{2r-1}\right\}.\tag{5}\end{equation}
The infima are attained by finite homogeneous stochastic models. If $b>0$, the second expression is strictly smaller than $b/2$; thus $K(n)$ is the exact first budget allowing improvement. This extends the two-prime result in the technical section entitled ``Revision XXIX''.
\statement{Lemma}{2}{Cost of a recurrent clock} If the asymptotic period $L$ of a Markov response with at most $m$ internal states satisfies $n\mid L$, then $m\ge K(n)$. If $m=K(n)$, every recurrent class essential to $n\mid L$ is a deterministic cycle of length $q_i$ for one of the maximal prime powers, and there are no other states.
\begin{proof}Let $d_1,\ldots,d_s$ be the periods of the closed irreducible classes. Each class has at least $d_j$ states, so $m\ge\sum_jd_j$, and the response is asymptotically periodic with period dividing $L=\lcm(d_1,\ldots,d_s)$. If $n\mid L$, for each maximal prime power $q_i$ some class period $d_j$ is divisible by $q_i$. Assign each $q_i$ to one such class. For any assigned subset $B_j$, coprimality gives $d_j\ge\prod_{i\in B_j}q_i\ge\sum_{i\in B_j}q_i$. The last inequality follows from $uv\ge u+v$ for coprime integers $u,v\ge2$ (indeed $uv>u+v$ here) and induction. Summing over $j$ yields $m\ge\sum_iq_i=K(n)$. Equality forces one assigned maximal prime power per recurrent class, period exactly $q_i$, exactly $q_i$ states, no transients, and no additional classes. A period-$q_i$ irreducible chain on precisely $q_i$ states is a deterministic cycle.\end{proof}
\begin{proof}[Proof of Theorem 1]Let $g$ be any admissible response. Standard finite Markov-chain theory gives an asymptotically periodic $h$ with period $L=\lcm(d_1,\ldots,d_s)$ and $g(k)-h(k)\to0$. If $m<K$, Lemma 2 implies $n\nmid L$. Along arbitrarily large multiples of $L$ the target $f_{n,a,b}(k)$ takes both values $a$ and $a+b$: indeed $k=\ell L$ is divisible by $n$ for infinitely many $\ell$, but not for infinitely many others. Since $g(k)\to h(0)$ there, the uniform approximation error is at least $b/2$. The one-state constant response $g(k)=a+b/2$ attains this bound.

Now put $m=K$. If $n\nmid L$, the same $b/2$ bound applies and is no smaller than the claimed expression. Otherwise Lemma 2 shows the chain has exactly $r$ independent deterministic cycles, of respective lengths $q_i$. Consequently its response is an additive function of separate clock residues,
\[h(k)=\sum_{i=1}^r h_i(k\bmod q_i),\qquad h(k)\ge0.\]
Choose a nonzero residue $s_i\in\{1,\ldots,q_i-1\}$ for each $i$. Chinese remaindering realizes all $2^r$ combinations of residue $0$ or $s_i$. Denote the value with all zero residues by $H_*$, the value with all nonzero residues by $H_0$, and the value with all but the $i$th residue zero by $H_{-i}$. Additivity gives the exact identity
\begin{equation}\sum_{i=1}^rH_{-i}=(r-1)H_*+H_0.\tag{6}\end{equation}
If the approximation error is at most $\varepsilon$, then $H_*\ge a+b-\varepsilon$, $H_{-i}\le a+\varepsilon$ and $a-\varepsilon\le H_0$; also $H_0\ge0$. Applying these inequalities to (6) gives the two simultaneous lower bounds
\begin{equation}(r-1)b\le2r\varepsilon,\qquad (r-1)b-a\le(2r-1)\varepsilon.\tag{7}\end{equation}
They imply the right-hand side of (5).

To attain the bound, construct $r$ disjoint cycles of lengths $q_i$. Give each cycle initial probability weight $1/r$ concentrated at residue zero. For constants $t,u\ge0$, set its binary effect to $t+ru$ at the zero residue and to $t$ everywhere else. The resulting normalized classical Markov response is
\begin{equation}g(k)=t+u\sum_{i=1}^r\one_{\{q_i\mid k\}},\qquad0\le t\le t+ru\le1.\tag{8}\end{equation}
Write $c=(r-1)b/(2r)$. If $a\ge c$, choose $t=a-c$ and $u=b/r$, giving maximum error $c$. If $a<c$, choose $t=0$ and $u=(2a+b)/(2r-1)$, giving maximum error $((r-1)b-a)/(2r-1)$. All effects lie in $[0,1]$: in the first case $t+ru=a+b-c\le1$; in the second it equals $a+b-((r-1)b-a)/(2r-1)\le1$. At the target spike and at all nonspecial clock configurations these errors attain the stated values, completing the proof. For $r=1$ the second formula yields zero, as expected for a single exact $n$-cycle.\end{proof}
\statement{Corollary}{3}{Prior results recovered without separate cases} For prime $n=p$, $K=p$ and one $p$-cycle realizes the signal exactly; every $m<p$ has error $b/2$. For $n=pq$ with distinct primes, $K=p+q$ and (5) becomes
\[E_{p+q}=\max\{b/4,(b-a)/3\},\]
recovering the exact two-prime theorem in the technical section ``Revision XXIX''. For arbitrary numbers of distinct prime powers, the first improvement starts at their sum, not their product.
\section*{3. Exact cases and what the numbers mean}
\begin{center}\begin{tabular}{rrrrrrr}\toprule
$n$ & prime powers & $K$ & $r$ & $a$ & $b$ & $E_K$\\\midrule
7&7&7&1&0&1&0\\
21&3, 7&10&2&$4/25$&$21/25$&$17/75$\\
30&2, 3, 5&10&3&$4/25$&$21/25$&$38/125$\\
60&4, 3, 5&12&3&$4/25$&$21/25$&$38/125$\\
210&2, 3, 5, 7&17&4&0&1&$3/7$\\\bottomrule
\end{tabular}\end{center}
Before the first improvement $m<K$, all these cases have $E_m=b/2$.
\begin{center}\setlength{\unitlength}{1mm}
\begin{picture}(130,43)
\put(13,6){\vector(1,0){110}}\put(13,6){\vector(0,1){34}}
\put(0,34){\makebox(0,0)[l]{$0.42$}}\put(0,18){\makebox(0,0)[l]{$0.2267$}}
\linethickness{0.8pt}\put(20,34){\line(1,0){65}}\put(92,18){\circle*{1.6}}
\put(20,2){\makebox(0,0){1}}\put(49,2){\makebox(0,0){5}}\put(85,2){\makebox(0,0){9}}\put(114,2){\makebox(0,0){13}}
\put(92,10){\makebox(0,0){$K(21)=10$}}\put(123,4){\makebox(0,0)[l]{$m$}}
\put(10,40){\makebox(0,0){$E_m$}}\put(97,23){\makebox(0,0)[l]{\scriptsize strict gain at $K$}}
\put(24,38){\makebox(0,0)[l]{\scriptsize No improvement below the first compatible clock budget}}
\end{picture}\end{center}
The plotted step describes the proved part of the curve for $n=21$: the constant level through $m=9$ and the exact value at $m=10$. No assertion is made about the optimal curve for $11\le m<21$. This distinction matters when reporting resource tradeoffs to systems engineers.
\section*{4. Relation to spectral realization and quantum experiments}
The exact clock threshold is a result in stochastic dynamics, not in low-rank matrix approximation. For example, a qubit can produce
\[f_n(k)=\sin^2(\pi k/n)=\tfrac12-\tfrac14e^{2\pi\ii k/n}-\tfrac14e^{-2\pi\ii k/n}\]
with a three-dimensional linear predictive representation. A classical time-homogeneous realization preserving its primitive $n$th root-of-unity oscillation exactly requires at least $n$ recurrent states, because a finite stochastic matrix can carry that peripheral eigenvalue only through a recurrent class of period divisible by $n$. That is a statement about dynamical memory. In contrast, the nonnegative rank of a finite preparation--measurement table is a static factorization problem. These numbers must not be silently identified.

For the two-level Singer return signal
\[P_q(k)=\frac{q}{(q+1)^2}+\left(1-\frac{q}{(q+1)^2}\right)\one_{\{q^2+q+1\mid k\}},\]
Theorem 1 applies with $n=q^2+q+1$ regardless of whether that integer is prime. For the Singer example $q=4$, $n=21$, it gives an exact transition from error $21/50$ with fewer than ten classical states to $17/75$ with ten. The construction is a classical approximation to the signal, not a new physical model of geometric causation.
\begin{center}\small\begin{tabularx}{\textwidth}{@{}p{.22\textwidth}XX@{}}\toprule
Operational object&What it certifies&What it does not certify\\\midrule
Hankel rank&Linear predictive dimension for a consistent controlled process&Number of physical hidden states or neural parameters\\
Nonnegative rank&Size of a static nonnegative factorization&Normalized effects or time-homogeneous realization\\
$E_m$&Optimal uniform error in the fixed classical dynamical class&Accuracy at a finite observation horizon or under feedback\\
Projected log-probability SVD&Optimal low-rank output within a fixed feature space&Generalization of a retrained Transformer\\
Hopf holonomy / process witness&Specific phase or causal-order response under a declared apparatus&A universal theory of gravity\\\bottomrule
\end{tabularx}\end{center}
\section*{5. Finite experiments: where the infinite-time theorem stops}
A laboratory collects only finitely many samples at finitely many observation times. The theorem alone does not give a finite-data device-independent memory witness. It uses arbitrarily late residues and does not specify a mixing time for a competing Markov chain. A candidate finite-sample protocol must separately fix (i) accessible times $0\le k\le T$, (ii) number of Bernoulli repetitions per time, (iii) the competing model class including any mixing-time or spectral restrictions, and (iv) an explicit confidence level.

If an independent argument finds a time $k$ with a known probability gap $\Delta$ between two specified Bernoulli models, Hoeffding's inequality gives a sufficient number of samples on the order of $\Delta^{-2}\log(1/\alpha)$ for a fixed hypothesis test. This is neither an information-theoretic necessary sample bound nor a universal substitute for locating the relevant time. Therefore we distinguish \textbf{formal worst-case impossibility, constructive finite-state approximation, and statistical certification}.
\section*{8. Fibonacci transfer, delta-rule memories, SU(2) positional encodings and the Hopf sphere}
\addcontentsline{toc}{chapter}{New research module: Fibonacci, recurrent Transformer memory, qubits and $S^3$}
\note{Question and comparator.} Can a single exact algebraic object connect a gated recurrent attention operator, a qubit instrument and the Hopf geometry of normalized spinors, without confusing such a connection with an algorithmic or quantum advantage? We answer yes in a sharply limited two-dimensional model and prove the corresponding no-advantage control. RoPE [Su et al., 2024], delta-rule recurrences [Tumma et al., ICML 2026; Hatamizadeh et al., 2026] and Fibonacci-anyon fusion [Field--Simula, 2018] are antecedents, not results established here for the first time.
\statement{Theorem}{4}{Exact Fibonacci/delta/filter correspondence; finite-dimensional scope} Let $\phi=(1+\sqrt5)/2$, $F=\left(\begin{smallmatrix}1&1\\1&0\end{smallmatrix}\right)$, $K=F/\phi$ and $\lambda=-\phi^{-2}$. Let $v_-$ be a unit eigenvector of $K$ with eigenvalue $\lambda$. Then
\begin{equation}F^2=F+I,\qquad\spec(K)=\{1,\lambda\},\qquad K^2=I-(1-\phi^{-4})v_-v_-^{\mathsf T}.\tag{9}\end{equation}
Consequently the even-step linear recurrence $h_{t+1}=K^2h_t$ is exactly a rank-one delta-rule erasure with fixed key $v_-$ and strength $\beta=1-\phi^{-4}$, without a write input. In contrast $K$ itself is not a unitary qubit gate.
\begin{proof}The characteristic polynomial of $F$ is $x^2-x-1$ and $F$ is real symmetric; hence $K$ has an orthonormal eigenbasis with eigenvalues $1$ and $-\phi^{-2}$. Spectral projection onto the second eigenvector yields $K^2=v_+v_+^{\mathsf T}+\phi^{-4}v_-v_-^{\mathsf T}=I-(1-\phi^{-4})v_-v_-^{\mathsf T}$. As $\|K\|_2=1$ but one singular value is strictly less than one, $K$ is a contraction, not a unitary.\end{proof}
\statement{Proposition}{5}{Exact heralded qubit implementation} With $B=(I-K^2)^{1/2}$, the block operator
\begin{equation}U=\begin{pmatrix}K&B\\B&-K\end{pmatrix}\in SO(4)\subset SU(4)\tag{10}\end{equation}
is unitary. Preparing an ancilla in $|0\rangle$, applying $U$ and selecting ancilla $|0\rangle$ implements $\psi\mapsto K\psi$. The success probability is $\|K\psi\|^2\in[\phi^{-4},1]$. Conditional normalization maps the unit sphere $S^3\subset\C^2$ to itself and induces an invertible fractional linear (M\"obius) map on $\mathbb{CP}^1\simeq S^2$. Its topological degree is $+1$, independently of its contraction rate. It does not identify spacetime with $S^3$.
\begin{proof}The symmetric matrices $B$ and $K$ commute and satisfy $B^2+K^2=I$, giving $U^{\mathsf T}U=I$. For either eigenvector the restriction to its two-dimensional ancilla sector has determinant $-1$, so $\det U=1$. The postselection probability is the expectation of $K^2$ whose eigenvalues are $1$ and $\phi^{-4}$. Because $\det K\ne0$, its induced complex projective action is a bijective holomorphic M\"obius transformation of degree one.\end{proof}
\statement{Theorem}{6}{Predictive rank two and an exact classical countermodel} For the equal superposition of the two eigenvectors, the scalar unnormalized return amplitude is
\begin{equation}c_t=\tfrac12(1+\lambda^t),\qquad\det(c_{i+j})_{i,j=0}^1=\tfrac14(1-\lambda)^2>0,\qquad\rank(c_{i+j})_{i,j\ge0}=2.\tag{11}\end{equation}
Exactly the same sequence is the return probability of the two-state classical Markov chain
\begin{equation}M_\lambda=\tfrac12\begin{pmatrix}1+\lambda&1-\lambda\\1-\lambda&1+\lambda\end{pmatrix},\qquad c_t=(M_\lambda^t)_{11}.\tag{12}\end{equation}
Thus this Fibonacci--qubit--$S^3$ correspondence provides no quantum-memory advantage for the calibrated scalar passport in (11).
\begin{proof}The amplitude is a sum of two distinct exponential modes, hence its infinite Hankel rank is at most two; the displayed nonzero minor gives equality. The stochastic matrix $M_\lambda$ has eigenvalues $1$ and $\lambda\in(-1,0)$ and uniform stationary distribution, whence the indicated return probability.\end{proof}
\statement{Proposition}{7}{SU(2)-RoPE commutant and a two-pole phase kernel} For $SU(2)$ positional rotations $R_t=\diag(e^{-\ii t\theta/2},e^{\ii t\theta/2})$, one has $\langle R_sq,R_tk\rangle=\langle q,R_{t-s}k\rangle$. If $K$ and $R_t$ share an eigenbasis, $T=KR_1$ gives
\begin{equation}\langle\chi,T^t\chi\rangle=\tfrac12\bigl(e^{-\ii t\theta/2}+\lambda^te^{\ii t\theta/2}\bigr),\qquad\chi=(v_++v_-)/\sqrt2.\tag{13}\end{equation}
Its complex Hankel rank is exactly two. For a generic nontrivial real planar RoPE rotation, the real $2\times2$ commutant is $\{aI+bJ:a,b\in\R\}$, where $J=\left(\begin{smallmatrix}0&-1\\1&0\end{smallmatrix}\right)$. This explicitly limits allowed gauge changes of the positional basis, rather than asserting universal gauge symmetry.
\begin{proof}The relative-position identity is $R_s^*R_t=R_{t-s}$. Equation (13) follows from simultaneous diagonalization and has distinct poles because their absolute values are $1$ and $\phi^{-2}$. For the commutant solve $GJ=JG$, obtaining $G=aI+bJ$. The two-dimensional RoPE-gauge phenomenon is consistent with the richer Transformer gauge analysis of Wang--Wang (PMLR 282, 2026).\end{proof}
\note{A second, genuinely topological meaning of the same Fibonacci matrix.} The fusion-counting matrix of the Fibonacci anyon rule $\tau\otimes\tau=1\oplus\tau$ is $\left(\begin{smallmatrix}0&1\\1&1\end{smallmatrix}\right)$, which is permutation-conjugate to $F$. This counts fusion channels and has Perron--Frobenius eigenvalue $\phi$. It is not a physical braiding gate: the unitary braid representation requires separate fusion-associator ($F$-symbol) and $R$-symbol matrices. Numeric coincidence of $\phi$ is neither a model of quantum gravity nor a proof of Transformer universality.
\note{Catalog-based boundary audit.} OpenAI Math families 144 (simple Lebesgue Koopman spectrum), 146 (standard-map entropy), 148 (self-similar dimension), 151 (failure of a naive entropy inference), 229 (tree reconstruction) and 276 (qubit generalized amplitude damping) address adjacent but mathematically distinct scopes. None of them establishes an intrinsic Fibonacci gate in trained Transformers, a topological invariant of LLM memory or an $S^3$ cosmological origin. The available Lean scope documents must be checked separately; catalogue claims are not substitutes for proofs in the present operational passport.
\note{Falsifiable model proposal.} Compare a trainable delta gate $\beta\in[0,1]$ against $\beta=1-\phi^{-4}$ at matched state width, data, training and inference budget; hold the RoPE phase geometry fixed; report retrieval error, output KL, perplexity, parameter memory, latency and Hankel singular values on held-out prompts. No experiment on a publicly pretrained LLM is asserted here. The reproducible tests included in this edition verify exact identities, their qubit dilation, the relative-position gauge relation and deliberately unfavorable synthetic dynamics.
\section*{9. An operational finite-shot dimension witness for a damped qubit}
\addcontentsline{toc}{chapter}{Damped-qubit memory witnesses: four and five time points (v32)}
\note{Precise scope and previous literature.} The following algebraic identities are specializations of classical Hankel/linear-system realization criteria and are not asserted to be globally novel. The comparator is a single time-homogeneous two-state classical Markov kernel with a fixed preparation and a fixed binary effect. The experiment consists of distinct independently prepared binary trials at different evolution times; it does not involve consecutive projective measurements in one run, a changing classical external clock, or an unrestricted Transformer. Existing quantum temporal-memory dimension witnesses, including Vieira, Milz, Vitagliano and Budroni (\textit{Quantum} 8, 1224, 2024), are relevant precedents; their protocols are not assumed equivalent to ours.
\statement{Theorem}{8}{An exact four-time obstruction to two classical hidden states} Let $0<\rho\le1$ and $0<\theta<\pi$. A single-qubit phase-rotation/dephasing channel maps transverse Bloch coherence $w\mapsto\rho e^{\ii\theta}w$. Prepare $|+\rangle$ and measure $|-\rangle\langle-|$ after $k$ iterations, giving
\begin{equation}f_k=\frac{1-\rho^k\cos(k\theta)}2,\qquad k=0,1,2,\ldots.\tag{14}\end{equation}
Define the four-time witness
\begin{equation}J(f)=(f_2-f_1)^2-(f_1-f_0)(f_3-f_2).\tag{15}\end{equation}
Then
\begin{equation}J(f)=\frac{\rho^2\sin^2\theta}{4}(1+\rho^2-2\rho\cos\theta)>0.\tag{16}\end{equation}
By contrast, $J(g)=0$ for every classical response $g_k=e^{\mathsf T}M^k\mu$ with a fixed $2\times2$ stochastic matrix $M$, probability preparation $\mu$, and a fixed binary effect $0\le e\le1$.
\begin{proof}In the qubit channel the Bloch $x$ coordinate after $k$ steps is $\rho^k\cos(k\theta)$, so the $|-\rangle$ probability is (14). Set $c=\cos\theta$ and insert the Chebyshev identities $\cos(2\theta)=2c^2-1$, $\cos(3\theta)=4c^3-3c$ in (15); exact polynomial factorization yields $J=\rho^2(1-c^2)(1+\rho^2-2\rho c)/4$. Both nonnegative factors are strictly positive under the stated interior-angle assumptions. For the classical claim, the fixed stochastic matrix has an eigenvalue $1$ and its other eigenvalue $\lambda\in[-1,1]$; Cayley--Hamilton gives $g_{k+2}=(1+\lambda)g_{k+1}-\lambda g_k$. Hence the differences $\Delta_k=g_{k+1}-g_k$ satisfy $\Delta_{k+1}=\lambda\Delta_k$, whence $\Delta_1^2-\Delta_0\Delta_2=0$.\end{proof}
\statement{Theorem}{9}{Five-time Hankel certificate and exact factor relation} For the same response set $H_3(f)=[f_{i+j}]_{0\le i,j\le2}$ and $D(f)=-\det H_3(f)$. Then
\begin{equation}D(f)=\frac{\rho^2\sin^2\theta}{8}(1+\rho^2-2\rho\cos\theta)^2=\frac{1+\rho^2-2\rho\cos\theta}{2}J(f)>0.\tag{17}\end{equation}
The full Hankel rank of $f$ is exactly $3$; any response from a fixed two-state linear model has zero such determinant. The order-three witness requires only the probabilities at $k=0,1,2,3,4$.
\begin{proof}Expand $f_k=\frac12-\frac14(\rho e^{\ii\theta})^k-\frac14(\rho e^{-\ii\theta})^k$. The three roots $1,\rho e^{\ii\theta},\rho e^{-\ii\theta}$ are distinct. For $3\times3$ Hankel matrices the Vandermonde determinant identity yields $\det H_3=(\frac12)(-\frac14)^2\prod_{u<v}(z_u-z_v)^2=-\rho^2\sin^2\theta\,|1-\rho e^{\ii\theta}|^4/8$, which is (17). The rank of the infinite Hankel matrix is three by Vandermonde independence. For the classical case, Cayley--Hamilton forces linear Hankel rank at most two. The algebraic factor relation follows by dividing the explicit formulas for $D$ and $J$.\end{proof}
\statement{Corollary}{10}{Finite-window model-separation gaps} For any $[0,1]$-valued homogeneous two-state classical response $g$,
\begin{equation}\max_{0\le k\le3}|f_k-g_k|\ge\frac{J(f)}8,\qquad\max_{0\le k\le4}|f_k-g_k|\ge\frac{D(f)}9.\tag{18}\end{equation}
\begin{proof}For $x=(x_0,x_1,x_2,x_3)\in[0,1]^4$ the absolute partial derivatives of $J(x)$ are bounded, respectively, by $(1,3,3,1)$, so $|J(x)-J(y)|\le8\|x-y\|_\infty$. Every entry of the $3\times3$ Hankel block belongs to $[0,1]$ and each cofactor (a $2\times2$ determinant) has absolute value at most one. Along the straight segment between two such matrices, $|\det A-\det B|\le9\|A-B\|_{\max}$. The respective invariants vanish identically for $g$.\end{proof}
\statement{Corollary}{11}{Two independently testable level-$\alpha$ rules} For $0<\alpha<1$, estimate each of the required probabilities from $N$ independent Bernoulli trials, and write $\widehat J,\widehat D$ for the plug-in invariants. Under the two-state null, either threshold test has type-I error at most $\alpha$; under the specified qubit alternative its rejection probability is at least $1-\alpha$ if the corresponding sufficient per-time budget holds:
\begin{align}
\text{four points:}\quad&N\ge\frac{200}{J(f)^2}\log\frac8\alpha,\qquad\text{reject if }|\widehat J|>\frac{2J(f)}5;\tag{19}\\
\text{five points:}\quad&N\ge\frac{200}{D(f)^2}\log\frac{10}\alpha,\qquad\text{reject if }|\widehat D|>\frac{9D(f)}{20}.\tag{20}
\end{align}
The total trial budgets are $4N$ and $5N$, respectively. These are rigorous sufficient rather than optimized sample counts.
\begin{proof}Hoeffding and a union bound give $\Pr(\max_k|\widehat f_k-f_k|>\varepsilon)\le2qe^{-2N\varepsilon^2}$ when $q=4$ or $q=5$ probabilities are estimated. For four points choose $\varepsilon=J/20$ and use the Lipschitz constant $8$: on the concentration event the null yields $|\widehat J|\le8J/20=2J/5$, while the target gives $|\widehat J|\ge J-2J/5=3J/5$, strictly above threshold. The displayed budget ensures the required concentration. For five points choose $\varepsilon=D/20$ and repeat with constant $9$, producing threshold $9D/20$ and alternative lower value $11D/20$. No null-specific parameters are needed, but the alternative angle and damping must be calibrated in advance to choose these thresholds.\end{proof}
\statement{Proposition}{12}{Analytically optimal interrogation angles for the pure qubit} For $\rho=1$, the four-time contrast $J=(1-c)^2(1+c)/2$ is maximal at $c=-1/3$ with $J_{\max}=16/27$. The five-time contrast $D=(1-c)^3(1+c)/2$ is maximal at $c=-1/2$ with $D_{\max}=27/32$. The two protocols generally have different optimal phase angles.
\begin{proof}Direct differentiation on $c\in[-1,1]$; both contrasts vanish at both endpoints and the stated interior critical points are unique maxima.\end{proof}
\statement{Proposition}{13}{Fibonacci control and its limitation} The normalized Fibonacci stochastic cell from the previous section has the decaying nonstationary eigenvalue $-\phi^{-2}$. A physically allowed damped qubit can be calibrated to the same attenuation modulus $\rho=\phi^{-2}$, yet retains the strictly positive witnesses (16) and (17) at nontrivial phase. Equal modulus of the decay rate and equal nominal two-dimensionality therefore do not imply equal measured dynamical behavior. This is an operational comparison, not proof of Fibonacci topology inside a Transformer or the physical universe.
\note{Relation to current model architectures.} A Gated DeltaNet / Transformer layer uses input-dependent updates and is not ordinarily a fixed homogeneous two-state Markov process. The determinants certify only the specifically benchmarked autonomous surrogate class. The distinction between the algebraic witness, realistic noise calibration, and architectural generalization is part of the result.
\note{Evidence and provenance.} The self-contained \texttt{checks/verify\_v32\_finite\_witness.py} reproduces the symbolic identities and parameter sweeps; \texttt{results/V32\_CERTIFICATES.csv} provides numerical prescriptions. No Lean 4 formalization or kernel compilation is requested or claimed.
\note{New result in Revision XXXIII: a five-time witness against three hidden classical states.} For five successive response probabilities $x_0,\ldots,x_4$, set $\Delta_i=x_{i+1}-x_i$ ($0\le i\le3$) and
\[H=\Delta_0\Delta_2-\Delta_1^2,\quad A=\Delta_0\Delta_3-\Delta_1\Delta_2,\quad B=\Delta_1\Delta_3-\Delta_2^2.\]
The quartic polynomial
\begin{equation}T_3(x)=H^2-HA+A^2-3HB\ge0\tag{21}\end{equation}
is necessary for every response $x_k=e^{\mathsf T}M^k\mu$ of a fixed, homogeneous, three-state column-stochastic matrix, with arbitrary preparation and bounded binary effect. This follows exactly from $3\operatorname{tr}(M^2)\ge(\operatorname{tr}M)^2$ and Cayley--Hamilton. For the physical single-qubit response $x_k=[1-\rho^k\cos(k\theta)]/2$,
\[T_3=H^2\bigl[(1-\rho\cos\theta)^2-3\rho^2\sin^2\theta\bigr].\]
At $(\rho,\theta)=(1,\pi/2)$ it equals $-1/2$, disproving the entire three-state homogeneous classical comparison class on only five time points. A global Lipschitz bound yields the conservative robust gap $\max_{k\le4}|x_k-g_k|\ge1/384$ for any three-state comparison signal $g$. These facts do not follow from the rank-three Hankel determinant: the corresponding $3\times3$ difference Hankel determinant vanishes for both the qubit and every three-state model. The full derivation, finite-sample prescription, controls and source-provenance discussion are in Revision XXXIII below. The construction is not a device-independent quantum-certification claim, and its world-first priority has not been established.
\setcounter{equation}{21}


% 01_foundations.tex
\chapter{Axioms of Necessary Reconstruction}
Define an operational model as $\mathcal P=(S,A,E,K,\beta,C)$: a state space, permitted actions, a set of admissible effects, transition kernels, preparations, and admissible experimental contexts. No reconstruction result is meaningful without fixing these data. Realizations may admit coordinate and gauge transformations that preserve all observable probabilities.

For finite classical models, $M_a$ are column-stochastic $m\times m$ matrices, $\beta$ is a probability column, and $0\le\alpha\le\one^\top$. Words $w=a_1\cdots a_t$ are executed from left to right: $M_w=M_{a_t}\cdots M_{a_1}$. The outcome probability is $f(w)=\alpha M_w\beta$; the Hankel matrix is $H_{u,v}=f(uv)$, where $u$ is a prefix and $v$ a suffix.
\statement{Definition}{I.1}{Passport of invariant closure} For each assertion $T$, fix the passport
\[U(T)=(A,E,\sim,\mathcal C,\mathcal N,\mathcal V).\]
Here $\sim$ is the admissible gauge/operational equivalence, $\mathcal C$ the model class, $\mathcal N$ the error norm and horizon, and $\mathcal V$ a verifiable certificate. Calling a theorem universal without specifying the scope of its quantifiers is unjustified. This is a logical definition, not a physical law.
\statement{Theorem}{I.2}{Canonical predictive invariant; known result} For a finite linear realization, the minimal real predictive dimension equals $\rank H$. Coordinate changes $M_a\mapsto TM_aT^{-1}$, $\beta\mapsto T\beta$, $\alpha\mapsto\alpha T^{-1}$ preserve $H$. A probabilistic realization of dimension $m$ satisfies $\rank H\le m$; for a quantum model on $\C^d$ realized by linear instruments, $\rank H\le d^2$.
\begin{proof}The factorization $H_{u,v}=(\alpha M_v)(M_u\beta)$ passes through the state space. Minimizing through the reachable subspace and quotienting by the unobservable kernel proves the converse in the category of linear realizations. The quantum bound follows from the real dimension of the space of Hermitian operators.\end{proof}
\statement{Theorem}{I.3}{Exact Markov closure for a partition} Let $L$ be a deterministic coarse-graining of microstates into nonempty macroclasses: each column of $L$ contains exactly one unit entry. A column-stochastic $\overline M_a$ with $LM_a=\overline M_aL$ exists if and only if
\[Le_i=Le_j\quad\Longrightarrow\quad LM_ae_i=LM_ae_j\qquad\text{for all }i,j.\]
When this condition holds, $\overline M_a$ is uniquely determined by the columns $LM_ae_i$ for any representative of each macroclass.
\begin{proof}Necessity follows by multiplying the identity by $e_i,e_j$. For sufficiency, define the macrocolumn using a representative of its class. This is well-defined, and its components are nonnegative and sum to one.\end{proof}
\chapter{An Exact Distinction between Linear, Positive, and Quantum Memory (Revision 19)}
\note{Checking worldwide priority before stating a new theorem.} The family of single-qubit antidistinguishability problems with a probability matrix proportional to $\sin^2(\pi(i-j)/n)$ has already been published: T.~Heinosaari, O.~Kerppo, L.~Lepp\"aj\"arvi, M.~Pl\'avala, \textit{Physical Review A} 109, 032627 (2024), DOI \href{https://doi.org/10.1103/PhysRevA.109.032627}{10.1103/PhysRevA.109.032627}. This section makes no priority claim for the construction. Its purpose is to provide a self-contained, directly operational Hankel certificate, an explicit quantitative robustness bound in the maximum-error norm, and a reproducible finite-sample protocol. The historical novelty of these refinements has not been established. Ordinary $\rank H$ measures linear memory and therefore does not itself give the minimum number of stochastically positive states.
\statement{Definition}{II.1}{Passport of the preparation-and-measurement protocol} Fix an integer $n\ge3$. The permitted commands are preparations $R_i$ ($i\in\Z_n$), binary measurements $M_j$ with known outcome registration ($j\in\Z_n$), compositions of these commands, and the total number of registered outcomes. Quantum realizations consist of finite-dimensional completely positive instruments; classical realizations consist of time-homogeneous $m$-state stochastic instruments without an unlimited external shared-randomness resource. Conditional probabilities are compared for the same command set. The error norm for one series is $\|P-Q\|_{\max}=\max_{i,j}|P_{ij}-Q_{ij}|$.

Put $\theta_i=2\pi i/n$ and denote the standard Pauli matrices by $\sigma_x,\sigma_y$. Consider the qubit states and effects
\begin{equation}\rho_i=\frac{I+\cos\theta_i\,\sigma_x+\sin\theta_i\,\sigma_y}{2},\qquad E_j=I-\rho_j.\tag{22}\end{equation}
The exact probability of outcome 1 is
\begin{equation}W^{(n)}_{ij}=\Tr(\rho_iE_j)=\sin^2\frac{\pi(i-j)}n=\frac{1-\cos(\theta_i-\theta_j)}2.\tag{23}\end{equation}
\statement{Theorem}{II.2}{Constant full Hankel rank and unbounded classical memory} For every $n\ge3$, $W^{(n)}$ has ordinary rank 3 and nonzero spectrum
\begin{equation}\spec_{\ne0}W^{(n)}=\{n/2,-n/4,-n/4\}.\tag{24}\end{equation}
There is a physically admissible qubit instrument with commands $R_i,M_j$ and full operational Hankel rank 3 containing $W^{(n)}$ as a block. Every classical realization of all the same observables with $m$ hidden states must satisfy
\begin{equation}n\le\binom m{\lfloor m/2\rfloor}.\tag{25}\end{equation}
In particular, the minimal qubit Hilbert-space dimension is 2, whereas the necessary classical dimension grows without bound with $n$; at $n=6$, already $m=3$ states are insufficient despite Hankel rank 3.
\begin{proof}Let $\one=(1)_i$, $c=(\cos\theta_i)_i$, and $s=(\sin\theta_i)_i$. Then $W^{(n)}=\frac12\one\one^{\mathsf T}-\frac12(cc^{\mathsf T}+ss^{\mathsf T})$. For $n\ge3$, $\one,c,s$ are pairwise orthogonal and their squared norms are $n,n/2,n/2$. This yields the rank and spectrum (24).

Construct a full instrument acting on an arbitrary density matrix $X$:
\[R_i(X)=\Tr(X)\rho_i,\quad M_{j,1}(X)=\Tr(E_jX)\rho_0,\quad M_{j,0}(X)=\Tr((I-E_j)X)\rho_0.\]
These maps are completely positive; $R_i$ and $M_{j,0}+M_{j,1}$ are trace-preserving. All reachable unnormalized states lie in $V=\spanR\{I,\sigma_x,\sigma_y\}$, so the rank of the full operational Hankel matrix is at most 3. Conversely, prefixes $R_i$ and suffixes $M_{j,1}$ give the rank-3 block $W^{(n)}$, so the full rank is exactly 3.

For any classical model, after $R_i$ there is a probability vector $\mu_i\in\R_{\ge0}^m$, $\sum_s\mu_i(s)=1$, and the probability of outcome 1 under $M_j$ from state $s$ is $e_j(s)\in[0,1]$. Hence $W^{(n)}_{ij}=\sum_{s=1}^m\mu_i(s)e_j(s)$. Define support sets $S_i=\{s:\mu_i(s)>0\}$ and $T_i=\{s:e_i(s)>0\}$. The equality $W_{ii}=0$ gives $S_i\cap T_i=\varnothing$, while $W_{ij}>0$ for $i\ne j$ gives $S_i\cap T_j\ne\varnothing$. Therefore no $S_i$ is contained in any $S_j$ with $i\ne j$: otherwise $S_i\cap T_j=\varnothing$. The sets $S_i$ form an antichain in the Boolean lattice of subsets of $\{1,\ldots,m\}$, and the classical Sperner theorem gives (25). Since for $n\ge4$ the block rank is $3>1$, quantum dimension 1 is also impossible.\end{proof}
\statement{Theorem}{II.3}{Explicit robust bound against small classical memory} Let $n\ge3$, $m\ge1$, $n>\binom m{\lfloor m/2\rfloor}$, and $\Delta_n=\sin^2(\pi/n)$. Then every classical $m$-state model with probability table $Q$ satisfies
\begin{equation}\|Q-W^{(n)}\|_{\max}\ge\varepsilon_*(n,m):=\min\left\{\frac{\Delta_n}2,\frac{\Delta_n^2}{4m^2}\right\}.\tag{26}\end{equation}
This certificate requires no knowledge of the classical transition parameters and remains meaningful under arbitrary stochastic mixing within the $m$ hidden states.
\begin{proof}Suppose instead that $\varepsilon=\|Q-W^{(n)}\|_{\max}<\varepsilon_*$. Write $Q_{ij}=\sum_s\mu_i(s)e_j(s)$ and put $t=\Delta_n/(2m)$. Then $Q_{ii}<t^2$, so $S_i=\{s:\mu_i(s)>t\}$ and $T_i=\{s:e_i(s)>t\}$ are disjoint: a shared coordinate would make the diagonal sum strictly greater than $t^2$. For $i\ne j$, $Q_{ij}>\Delta_n/2=mt$, so some $s$ satisfies $\mu_i(s)e_j(s)>t$ (otherwise the sum would be at most $mt$). Since both factors are at most 1, both exceed $t$, giving $S_i\cap T_j\ne\varnothing$. Again the $S_i$ form an antichain of size $n$, contradicting Sperner's theorem.\end{proof}
\statement{Corollary}{II.4}{A rational counterexample: six settings} For $n=6$, $W^{(6)}$ is circulant with first row
\[\left(0,\tfrac14,\tfrac34,1,\tfrac34,\tfrac14\right).\]
Its ordinary rank is 3 and its singular values are $3,\frac32,\frac32$, but no three-state classical model realizes it exactly. Moreover,
\begin{equation}\inf_{Q\in\mathrm{Classical}(3)}\|Q-W^{(6)}\|_{\max}\ge\frac1{576}.\tag{27}\end{equation}
This is a certified lower bound, not the exact optimal distance to the three-state class.
\statement{Corollary}{II.5}{Finite-sample statistical protocol} Suppose the true model is $W^{(n)}$, every pair $(i,j)$ is measured independently $N$ times, and $\widehat W$ is the frequency matrix. For $0<\alpha<1$, put $\gamma=\varepsilon_*(n,m)/4$. It suffices to take
\begin{equation}N\ge\frac1{2\gamma^2}\log\frac{2n^2}\alpha=\frac8{\varepsilon_*^2}\log\frac{2n^2}\alpha\tag{28}\end{equation}
trials per pair to ensure $\|\widehat W-W^{(n)}\|_{\max}\le\gamma$ with probability at least $1-\alpha$. On this event no $m$-state classical model can satisfy $\|Q-\widehat W\|_{\max}\le\gamma$.
\begin{proof}Hoeffding's inequality and a union bound over the $n^2$ independent series give $\Pr(\|\widehat W-W\|_{\max}>\gamma)\le2n^2e^{-2N\gamma^2}$. If such a $Q$ existed, the triangle inequality would give $\|Q-W\|_{\max}\le2\gamma=\varepsilon_*/2$, contradicting (26).\end{proof}
\statement{Remark}{II.6}{Limits of physical inference and resources} Test (28) is conservative (for $n=6$, $m=3$, the gap is $\varepsilon_*=1/576$ and the number of measurements per pair is large); it does not report an experiment actually performed. It assumes honest random sampling, independent repetitions, no drift, and no unaccounted shared randomness between preparation and measurement. Unlimited jointly available randomness changes the comparison class, and the published work of Heinosaari--Kerppo--Lepp\"aj\"arvi--Pl\'avala specifically examines this qualification. Neither generalizing (26) to all quantum instruments nor deriving a physical law from it is admissible. OpenAI Math catalogue results on memory--sample tradeoffs (family 140), reconstruction on trees (229), and cones without semidefinite lifts (095) concern different problem classes; they cannot automatically replace the quantifiers proved here.
\note{Does a single rank answer the original question?} No. Already for $n=6$, the same operational protocol has full predictive rank 3, qubit physical memory of dimension 2, and admissible classical memory requiring at least 4 states. The separate formula (26) measures the inevitable cost of unjustifiably reducing classical memory. This three-way distinction is the verifiable part of the passport model, not a claim of worldwide first discovery.
\chapter{Rigorous Applicability to Artificial Intelligence: Predictive Memory, Errors, and Probing}
Let the alphabet of admissible commands and outcomes be finite, with a fixed reset regime, initial-condition distribution, and observation protocol. Write $f(w)$ for a probabilistic word $w$ and set $H_{u,v}=f(uv)$. An exact linear realization $f(w)=\alpha^{\mathsf T}T_w\beta$ of dimension $d$ satisfies $\rank H\le d$. When the rank $r$ is finite and word shifts are consistent, a minimal formal linear realization of dimension $r$ exists. A stochastic realization additionally requires nonnegativity, normalization, and closure under all admissible commands; these are substantially stronger conditions.
\statement{Proposition}{III.1}{Certificate of empirical predictive dimension} Let $H\in\R^{p\times q}$ be a selected subtable of the operational Hankel matrix, and let the frequency estimate $\widehat H$ satisfy $\|\widehat H-H\|_{\max}\le\varepsilon$. Then
\[\|\widehat H-H\|_2\le\sqrt{pq}\,\varepsilon,\qquad\sigma_r(H)\ge\sigma_r(\widehat H)-\sqrt{pq}\,\varepsilon.\]
Thus $\sigma_r(\widehat H)>\sqrt{pq}\,\varepsilon$ implies $\rank H\ge r$. Moreover every matrix $B$ of rank less than $r$ satisfies
\[\|H-B\|_{\max}\ge\frac{\sigma_r(\widehat H)-\sqrt{pq}\,\varepsilon}{\sqrt{pq}}.\]
The last bound is nontrivial only when the numerator is positive. It bounds reduction error on the selected tests; it does not guarantee the existence of physically admissible or neural-network compression of dimension $r$.
\begin{proof}The inequality $\|A\|_2\le\|A\|_F\le\sqrt{pq}\|A\|_{\max}$ and Weyl's inequality give the singular-value bound. By the Eckart--Young theorem, the spectral-norm distance to matrices of rank less than $r$ is at least $\sigma_r(H)$. Combining this with $\|H-B\|_2\le\sqrt{pq}\|H-B\|_{\max}$ gives the result.\end{proof}
\statement{Proposition}{III.2}{Safe error accumulation for a fixed command string} For every command $a$, let $P_a,Q_a$ be two probabilistic transitions on the same hidden state space, with $\sup_x\operatorname{TV}(P_a(x,\cdot),Q_a(x,\cdot))\le\delta_a$ and $\kappa_a:=\sup_{x,y}\operatorname{TV}(P_a(x,\cdot),P_a(y,\cdot))\le1$. For a command word $a_1\cdots a_T$ fixed in advance and a common initial law,
\[\operatorname{TV}(\mu P_{a_1}\cdots P_{a_T},\mu Q_{a_1}\cdots Q_{a_T})\le\sum_{t=1}^T\delta_{a_t}\prod_{j=t+1}^T\kappa_{a_j}.\]
If $\delta_a\le\delta$, $\kappa_a\le\kappa<1$, the final error is at most $\delta(1-\kappa^T)/(1-\kappa)$; for $\kappa=1$, the simple bound is $T\delta$. Without additional assumptions this conclusion does not extend to adaptively selected commands, incompatible hidden interfaces, changes in the context distribution, or open vocabularies.
\begin{proof}Writing $e_t=\operatorname{TV}(\mu P_{a_1}\cdots P_{a_t},\mu Q_{a_1}\cdots Q_{a_t})$, add and subtract $(\mu Q_{a_1}\cdots Q_{a_{t-1}})P_{a_t}$ to obtain $e_t\le\kappa_{a_t}e_{t-1}+\delta_{a_t}$. Iterate from $e_0=0$.\end{proof}
\note{Application scheme for an AI model.} Fix the query/response alphabet, condition distribution, and prefix and suffix sets; obtain independent estimates $\widehat H$; identify significant singular values using valid confidence bounds; separately attempt to construct a minimal admissible realization (stochastic, quantum, RNN, or state-space, according to the chosen passport); then test long-horizon error, calibration, and third-order correlations on held-out sequences. Not every Transformer, SSM, or diffusion model admits a finite exact realization of bounded rank, and data from a finite submatrix do not determine the global rank. The equation $\rank H=r$ does not prove that a network with $r$ neurons preserves quality.
\note{Essential boundary.} Propositions III.1 and III.2 apply known matrix and probability inequalities; no independent scientific priority is claimed. Section II demonstrates explicit physical inequivalence between linear and positive ranks and gives a quantitative lower bound that should supplement ordinary SVD heuristics when choosing a reduction class.
\chapter{An Autoregressive Model Passport: Finite Rank, Adaptive Risk, and the Limits of Extrapolation (Revision 20)}
\note{Purpose.} To specify precisely how rank and probability certificates apply to language models. Every assertion in this section concerns declared interfaces; none assumes universal Transformer compressibility or parameter reduction from rank alone.
\section*{Measurement passport and the correct Hankel object}
For a fixed finite alphabet $\Sigma$ (including EOS), decoding regime, initial prompt $c_0$, reset procedure, and horizon, define the consistent stochastic language
\[f(w):=\Pr\{X_1\cdots X_{|w|}=w\mid c_0\},\qquad w\in\Sigma^*.\]
It satisfies $f(\epsilon)=1$ and $\sum_{a\in\Sigma}f(wa)=f(w)$ if EOS is modeled as an absorbing termination. The full Hankel matrix $H_{u,v}=f(uv)$ consists of joint probabilities of a prefix and continuation. For a prefix with $f(u)>0$, the conditional continuation matrix $C_{u,v}=\Pr(v\mid u,c_0)$ satisfies
\begin{equation}H_{U,V}=\diag(f(u):u\in U)\,C_{U,V}.\tag{29}\end{equation}
Thus $\rank H_{U,V}=\rank C_{U,V}$ when all selected prefixes have positive probability, but singular values and errors change under scaling. If instead rows come from different externally supplied prompts $c_u$ that are not prefixes of one stochastic process, $C$ remains a measurement matrix but need not be a Hankel matrix of any consistent model. Its rank bounds the dimension of a possible linear factorization of these measurements, not unconditionally the memory of the original LLM.

Temperature 0 usually denotes deterministic token selection (greedy decoding), not the absence of logits: probabilities of the underlying softmax model must not be identified with the output distribution of a deterministic decoder. Positional indices, instruction roles, KV cache, context-truncation rules, stop tokens, random sources, and tools are part of the passport. For comparison models with a restricted window or open vocabulary, token mapping and the observation interface must be formalized separately.
\statement{Proposition}{IV.1}{Explicit finite-sample spectral certificate} Suppose $p\times q$ selected event probabilities $H_{uv}\in[0,1]$ are each estimated by the mean of $m$ independent Bernoulli trials after identical resets to the corresponding preparation. For $0<\gamma<1$, with probability at least $1-\gamma$, simultaneously
\begin{equation}\|\widehat H-H\|_{\max}\le\varepsilon_m:=\sqrt{\frac{\log(2pq/\gamma)}{2m}},\qquad\|\widehat H-H\|_2\le\sqrt{pq}\,\varepsilon_m.\tag{30}\end{equation}
In particular, $\sigma_r(\widehat H)>\sqrt{pq}\,\varepsilon_m$ implies $\rank H\ge r$ at the stated confidence level, and every $B$ with $\rank B<r$ satisfies
\begin{equation}\|H-B\|_{\max}\ge\max\left\{0,\frac{\sigma_r(\widehat H)}{\sqrt{pq}}-\varepsilon_m\right\}.\tag{31}\end{equation}
\begin{proof}For each event, Hoeffding gives $\Pr(|\widehat H_{uv}-H_{uv}|>\varepsilon)\le2e^{-2m\varepsilon^2}$. A union bound over $pq$ events proves the first inequality without requiring independence between different cells. Then $\|E\|_2\le\|E\|_F\le\sqrt{pq}\|E\|_{\max}$. The last conclusion follows from Weyl's inequality and the Eckart--Young theorem.\end{proof}
Limitation: selecting prompts after viewing results, correlated resets, and repeated selection of $r$ require confidence corrections. If event probabilities are computed exactly from available logprobs, Monte Carlo error can be replaced by justified numerical/model error; an arbitrary reported softmax probability cannot automatically be treated as a physical frequency.
\section*{Adaptive generation: two risk bounds}
Let $P,Q$ be two autoregressive laws on a common alphabet, with the same external causal control protocol (actions may depend on previously observed text). Let $p_t(\cdot\mid h),q_t(\cdot\mid h)$ be their conditional continuation kernels on the complete available context $h$.
\statement{Theorem}{IV.2}{TV certificate for adaptive continuation} If every reachable context $h$ of length $t-1$ satisfies $\operatorname{TV}(p_t(\cdot\mid h),q_t(\cdot\mid h))\le\delta_t<1$, then the joint laws of all generated tokens through horizon $T$ satisfy
\begin{equation}\operatorname{TV}(P_{1:T},Q_{1:T})\le1-\prod_{t=1}^T(1-\delta_t)\le\sum_{t=1}^T\delta_t.\tag{32}\end{equation}
Without a uniform bound, but for finite conditional laws, the weaker on-policy estimate is
\begin{equation}\operatorname{TV}(P_{1:T},Q_{1:T})\le\sum_{t=1}^T\mathbb E_{h\sim P_{<t}}\operatorname{TV}(p_t(\cdot\mid h),q_t(\cdot\mid h)).\tag{33}\end{equation}
If $P\ll Q$ on sequences and $D_t(h)=D_{\mathrm{KL}}(p_t(\cdot\mid h)\|q_t(\cdot\mid h))$, additionally
\begin{equation}\operatorname{TV}(P_{1:T},Q_{1:T})\le\sqrt{\frac12\sum_{t=1}^T\mathbb E_{h\sim P_{<t}}D_t(h)}.\tag{34}\end{equation}
\begin{proof}For (32), maximally couple the continuations while their histories agree. The conditional probability of maintaining agreement at step $t$ is at least $1-\delta_t$, yielding the product and the total-variation coupling bound. For (33), construct hybrid laws using $P$ for the first $t$ steps and $Q$ thereafter. When a single step is changed, its preceding prefix has law $P_{<t}$; applying the common future kernel cannot increase TV. A telescoping sum proves the result. Finally the relative-entropy chain rule and Pinsker's inequality give (34).\end{proof}
\statement{Remark}{IV.3}{Scope} The adaptive-generation estimate applies when both kernels are defined on the same space of complete histories and the same external controller's actions are explicitly included in the context. An error estimate only on predetermined teacher-forcing strings is insufficient for uniform control over all contexts. TV of the final hidden-state law must not be confused with TV of the entire text law.
\statement{Proposition}{IV.4}{Mixing of the final state does not protect the full text} Let $P$ generate independent Bernoulli$(1/2)$ variables and $Q$ independent Bernoulli$(1/2+\delta)$ variables, where $0<\delta<1/2$. Both transitions forget the entire history in one step and have Dobrushin coefficient 0. Nevertheless,
\begin{equation}\operatorname{TV}(P_{1:T},Q_{1:T})\ge1-2\exp(-T\delta^2/2).\tag{35}\end{equation}
Thus TV of the full generated text tends to 1 as $T\to\infty$, although TV of the final token equals $\delta$ independently of $T$.
\begin{proof}Take $A_T=\{T^{-1}\sum_{t=1}^TX_t>1/2+\delta/2\}$. Hoeffding gives $P(A_T)\le e^{-T\delta^2/2}$ and $Q(A_T)\ge1-e^{-T\delta^2/2}$; their difference bounds TV from below. The transitions do not depend on the previous state, so their Dobrushin coefficients are zero.\end{proof}
\section*{A global upper bound cannot be certified from finite tests}
\statement{Theorem}{IV.5}{A hidden delay destroys rank extrapolation} For every finite set of tested words $W\subset\{0,1\}^*$ and every integer $R$, there are two exactly finite-state consistent probabilistic processes with equal $f(w)$ for every $w\in W$, but global Hankel ranks 1 and at least $R$.
\begin{proof}Set $L=\max_{w\in W}|w|$ and choose $N>\max\{L,R\}$. The first process deterministically emits $000\ldots$ and the second $0^N111\ldots$. Their word probabilities agree for all words of length at most $L$. The first has Hankel rank 1. For the second, take $N+1$ prefixes $u_i=0^i$ and $N+1$ suffixes $v_j=0^j1$, $0\le i,j\le N$. Then
\[H_{u_i,v_j}=f(0^{i+j}1)=\one_{\{i+j=N\}},\]
so this submatrix is the reversed identity of order $N+1$, with rank $N+1>R$. The second process is realized by a finite automaton counting to $N$ and then emitting 1 forever.\end{proof}
\note{Implication for LLMs.} Even without noise, a finite table can certify a lower bound but not a general upper bound on predictive dimension without extra assumptions (such as bounded order, an identifiable family, or a proven finite realization).
\section*{Verifiable protocol v20}
\begin{enumerate}
\item Fix the observation object: joint word probabilities, conditional continuations, or arbitrary prompt--response pairs; do not automatically call the latter a Hankel matrix.
\item Before choosing data, specify $\Sigma$, the prompt/prefixes, suffixes, EOS, decoder, temperature, seed, context window, tools, and numerical format.
\item Estimate cells after independent resets and apply (30); when logprobs are accessible, specify their source and precision rather than assigning unnecessary Monte Carlo error.
\item Check nonnegative/stochastic realization separately from low linear rank. Test new prefixes and long contexts.
\item Compare distilled and original models on common contexts. Use contexts generated by the original model $P$ for on-policy estimates, and prove a uniform bound for a worst-case certificate.
\item For long generation, control the quality of the whole text and of the final token separately; one-step mixing does not guarantee stability of the text distribution.
\end{enumerate}
\note{Status.} The theorems in this section follow from classical tools of automata theory, maximal couplings, Hoeffding, and Pinsker. Their formulation as a unified passport for LLMs is practically useful, but worldwide priority, large-scale empirical advantages, and universal compressibility are not claimed. The archive contains no input logits from a real language model; subsequent tests are synthetic.
\chapter{An Exact Stochastic Family and Its Inverse Matrix}
Let $n\ge2$, $0\le a<1$, $q=1-a$, $\alpha=e_1^\top$, and $\beta=e_1$. Consider
\begin{align}Pe_1&=ae_1+qe_2,\quad Pe_j=e_{j+1}\ (2\le j<n),\quad Pe_n=e_2,\tag{36}\\
Qe_1&=Qe_2=e_1,\qquad Qe_j=e_j\ (j\ge3).\tag{37}\end{align}
Index the rows of $H_n(a)$ by suffixes $\varnothing,Q,PQ,\ldots,P^{n-2}Q$ and columns by prefixes $\varnothing,P,\ldots,P^{n-1}$. Then $r_{\{P\}}=r_{\{Q\}}=1$, $r_{\{P,Q\}}=n$, and $|\det H_n(a)|=q^{n-1}$. These assertions already appeared in earlier versions of META--CAUSAL.
\statement{Lemma}{V.1}{Explicit inverse matrix} For $n\ge3$, put $t=a^{n-2}$, $S=(1-t)/q$. The action of $C=H_n(a)^{-1}$ on $y=(y_0,\ldots,y_{n-1})^\top$ has coordinates
\begin{align}(Cy)_0&=(y_0-ay_1)/q,\quad(Cy)_1=(y_1-ay_{n-1}-(1-t)y_0)/q,\tag{38}\\
(Cy)_{n-j}&=(y_{j+1}-ay_j)/q\quad(2\le j\le n-2),\tag{39}\\
(Cy)_{n-1}&=(y_2-y_0)/q.\tag{40}\end{align}
For $n=2$, $C=q^{-1}\left(\begin{smallmatrix}1&-a\\-1&1\end{smallmatrix}\right)$.
\begin{proof}Direct multiplication by the probability matrix $H_n(a)$ verifies the result. For the recurrence, put $g_k=f(P^kQ)$. Then
\[g_k=a^k+q\sum_{\substack{r\ge0\\r(n-1)\le k-1}}a^{k-1-r(n-1)}.\]
The rows of $H_n$ are $(a^j)_{j=0}^{n-1}$ and shifts $(g_{j+k})_{j=0}^{n-1}$, $0\le k\le n-2$. Substitution into (38) eliminates adjacent shifts and gives $CH_n=I$. The case $n=2$ is computed directly.\end{proof}
\statement{Theorem}{V.2}{Exact invariant of entrywise robustness, new in series XI} For every $n\ge2$ and $0\le a<1$,
\begin{align}\|H_n(a)^{-1}\|_{\infty\to1}&=\frac{D_n(a)}{1-a},\tag{41}\\
D_n(a)&=(n-2)(1+a)+4-a^{n-2}+(-1)^{n-2}a,\tag{42}\\
\operatorname{dist}_{\max}(H_n(a),\{K:\rank K<n\})&=\rho_n(a)=\frac{1-a}{D_n(a)}.\tag{43}\end{align}
The distance is taken over all real matrices. For $n\ge3$, $\rho_n(0)=1/(n+2)$. At fixed $a<1$, $\rho_n(a)\sim(1-a)/[n(1+a)]$.
\begin{proof}For any invertible $A$, the general formula is $\operatorname{dist}_{\max}(A,\det=0)=1/\|A^{-1}\|_{\infty\to1}$. Indeed, if $(A+E)x=0$, then
\[\|Ax\|_\infty\le\|E\|_{\max}\|x\|_1,\qquad\|x\|_1\le\|A^{-1}\|_{\infty\to1}\|Ax\|_\infty.\]
Equality is attained by the rank-one matrix $E=-ys^\top/\|A^{-1}y\|_1$, where $y$ maximizes the norm and $s=\operatorname{sgn}(A^{-1}y)$.

For $C=H_n^{-1}$, the maximum occurs at a cube vertex $y_i\in\{\pm1\}$. By symmetry set $y_0=1$. In (38), with $y_{n-1}$ fixed, the sum of the absolute values of the first two coordinates is maximal at $y_1=-1$. For $n\ge3$ this gives
\[q\|Cy\|_1=3+a-t+ay_{n-1}+\sum_{j=2}^{n-2}(q+2a\one_{\{y_j\ne y_{j+1}\}})+2\one_{\{y_2=-1\}}.\]
Let $k$ count sign changes from $y_2$ to $y_{n-1}$. Then $y_{n-1}=y_2(-1)^k$, and $2ak+ay_2(-1)^k$ is nondecreasing in $k$, so one may choose $k=n-3$. Of the two choices for $y_2$, $y_2=-1$ is strictly better, with gain $2-2a(-1)^{n-3}>0$. Substitution gives $q\max\|Cy\|_1=D_n(a)$. For $n=2$, direct computation gives $(3+a)/q$, agreeing with (42).\end{proof}
\statement{Corollary}{V.3}{Dimension-uniform spectral stability} For $n\ge3$, $t=a^{n-2}$,
\[\|H_n^{-1}\|_1=\frac{3-t}q,\qquad\|H_n^{-1}\|_\infty=\frac{2+a-t}q,\]
and therefore
\[\sigma_{\min}(H_n(a))\ge\frac{1-a}{\sqrt{(3-t)(2+a-t)}}\ge\frac{1-a}{\sqrt{3(2+a)}}.\]
For fixed $a<1$, this is a positive lower bound independent of $n$.
\begin{proof}The absolute row and column sums follow from (38); then $\|C\|_2\le\sqrt{\|C\|_1\|C\|_\infty}$ and $\sigma_{\min}(H_n)=1/\|C\|_2$.\end{proof}
\chapter{The Law of Stable Causal Closure}
Consider a stochastic coarse-graining channel $L:\R^n\to\R^d$ and macrotransitions $\overline M_a$. Assume exact preservation of every used effect, $\alpha=\overline\alpha L$, $0\le\overline\alpha\le\one^\top$, and of the initial state, $\overline\beta=L\beta$. Define
\[\delta_a:=\tfrac12\|LM_a-\overline M_aL\|_{1\to1},\qquad\kappa_a:=\tfrac12\max_{i,j}\|\overline M_ae_i-\overline M_ae_j\|_1\le1.\]
\statement{Theorem}{VI.1}{Intervention-weighted telescoping} For every word $w=a_1\cdots a_t$, the error in every admissible measured probability is bounded by
\begin{equation}|f(w)-\overline f(w)|\le\Xi(w):=\sum_{j=1}^t\delta_{a_j}\prod_{k=j+1}^t\kappa_{a_k}.\tag{44}\end{equation}
If $\delta_a\le\delta$, $\kappa_a\le\kappa\le1$, then $|f(w)-\overline f(w)|\le\delta S_t(\kappa)$, where $S_t(\kappa):=\sum_{j=0}^{t-1}\kappa^j$. For $\kappa<1$, the error is uniformly bounded by $\delta/(1-\kappa)$ for arbitrarily long histories.
\begin{proof}Let $x_j=LM_{a_j}\cdots M_{a_1}\beta$ and $\overline x_j=\overline M_{a_j}\cdots\overline M_{a_1}L\beta$. With $e_j=\operatorname{TV}(x_j,\overline x_j)$, the triangle inequality and Dobrushin contraction give $e_j\le\delta_{a_j}+\kappa_{a_j}e_{j-1}$, $e_0=0$. Iteration gives $\Xi(w)$; a measurable effect cannot increase total variation.\end{proof}
\statement{Theorem}{VI.2}{Combined law of rank, mixing, and reduction cost} Let $H_n(a)$ be the table of Section II; any macrosystem has $d<n$ internal states and satisfies the assumptions above. For $T=2n-2$, $\delta:=\max(\delta_P,\delta_Q)$, and $\kappa:=\max(\kappa_P,\kappa_Q)$, necessarily
\begin{equation}\delta S_T(\kappa)\ge\rho_n(a)=\frac{1-a}{D_n(a)}.\tag{45}\end{equation}
In particular,
\[\kappa=1:\quad\delta\ge\frac{1-a}{(2n-2)D_n(a)},\qquad0\le\kappa<1:\quad\delta\ge\frac{(1-\kappa)(1-a)}{D_n(a)}.\]
\begin{proof}The macrotable $\overline H$ has rank at most $d<n$, hence $\|H_n(a)-\overline H\|_{\max}\ge\rho_n(a)$ by Theorem II.2. Every word used has length at most $T$, and (44) gives $\|H-\overline H\|_{\max}\le\delta S_T(\kappa)$.\end{proof}
\statement{Remark}{VI.3}{Scope} This is a necessary condition for this class of positive finite-dimensional reductions. It asserts neither sufficiency nor the existence of a reduction attaining the bound. It does not exclude nonlinear, nonstationary, non-Markovian, or quantum models of another dimension.
\chapter{A Noncommutative Order Invariant and a Dual Certificate}
In this family, the observable difference between two intervention orders is
\begin{equation}\Delta_{PQ}:=f(PQ)-f(QP)=1-a=q.\tag{46}\end{equation}
\statement{Theorem}{VII.1}{Cost of a commutative macrodescription} If $\overline P\overline Q=\overline Q\overline P$ while preserving the initial state and measurement, then
\[q\le(1+\kappa_Q)\delta_P+(1+\kappa_P)\delta_Q\le2(\delta_P+\delta_Q).\]
In particular $\delta_P+\delta_Q\ge(1-a)/2$, independently of the hidden state count $n$.
\begin{proof}The macropredictions for $PQ$ and $QP$ coincide. Their errors are at most $\kappa_Q\delta_P+\delta_Q$ and $\kappa_P\delta_Q+\delta_P$, respectively. Add and apply the triangle inequality.\end{proof}
Thus a single experimental model imposes two different independent requirements: high predictive dimension and preservation of intervention order. Low rank alone does not characterize the commutator, and the commutator does not determine full rank.
\chapter{Finitely Many Experiments: A Verifiable Certificate}
In each of the $n^2$ cells, measure $N$ independent Bernoulli outcomes. Let $\widehat H$ be the frequency table; fix $0<\gamma<1$ and
\[\eta=\sqrt{\frac{\log(2n^2/\gamma)}{2N}}.\]
Hoeffding's inequality gives $\Pr(\|\widehat H-H\|_{\max}>\eta)\le\gamma$.
\statement{Theorem}{VIII.1}{Robust verifiable rejection of all smaller models} Define $R(A)=\operatorname{dist}_{\max}(A,\{K:\rank K<n\})$; for invertible $A$, $R(A)=1/\|A^{-1}\|_{\infty\to1}$, and for singular $A$, $R(A)=0$. The rule
\[R(\widehat H)>\eta\quad\Longrightarrow\quad\text{reject all models of linear dimension below }n\]
has false-rejection probability at most $\gamma$ simultaneously for the entire composite class of smaller models. In the ideal family it rejects smaller models with probability at least $1-\gamma$ if $2\eta<\rho_n(a)$. This is a sufficient, not sample-optimal, bound.
\begin{proof}Distance to a fixed set is 1-Lipschitz in the original norm. Under any null hypothesis $R(H)=0$, so $R(\widehat H)\le\eta$ on the stated high-probability event. For the specified full family, $R(\widehat H)\ge\rho_n(a)-\eta$.\end{proof}
\chapter{The Algebraic Framework: Invariant Closure and a Singular Flag}
We first reintegrate the results of the earlier Revision IX into one principal manuscript. Let $W=\ker\one^\top$, $m=n-1$, $w_j=e_{j+1}-e_1$, and $u=w_1$. In the decomposition $\R e_1\oplus W$, $P$ and $Q$ have blocks
\[P=\begin{pmatrix}1&0\\qu&A\end{pmatrix},\qquad Q=\begin{pmatrix}1&0\\0&I-uu^\top\end{pmatrix}.\]
Here $u$ is the first basis vector in the internal coordinates of $W$; this $uu^\top$ is a matrix in the chosen coordinates, not an orthogonal projection in $\R^n$.
\statement{Theorem}{IX.1}{Two algebras with the same predictive rank} For $n\ge3$, the real unital algebra $\mathcal A_a=\operatorname{alg}_{\R}(I,P,Q)$ satisfies
\begin{align*}0<a<1:\quad&\mathcal A_a=\{T:\one^\top T=c_T\one^\top\},&&\dim\mathcal A_a=n^2-n+1;\\
a=0:\quad&\mathcal A_0=\{T:\one^\top T=c_T\one^\top,\ T\ell\subseteq\ell\},&&\dim\mathcal A_0=n^2-2n+3,
\end{align*}
where $\ell=\operatorname{span}(e_n-e_1)\subset W$. In both regimes $\rank H_n(a)=n$. In particular, the dimension of the operator algebra jumps by $n-2$.
\begin{proof}For $j=0,\ldots,m-1$, the vectors $A^ju$ form a basis of $W$; for $a>0$, the rows $u^\top A^j$ form a basis of $W^*$ (their matrix determinant has modulus $a^{m-1}$). At $a=0$, these rows vanish jointly on $\ell$ and span its annihilator. Thus the lower-right blocks of $P^i(I-Q)P^j$ span, respectively, $\operatorname{End}(W)$ or the ideal of operators annihilating $\ell$. At $a=0$, adjoining the identity generates all endomorphisms preserving $\ell$. Averaging powers of $P$ gives the projection $\Pi=\pi\one^\top$ onto the unique stationary distribution $\pi$ (uniform on the hidden cycle). The operators $(Q-I)\Pi$ and their left translates by powers of $P$ generate the lower-left block, while $\Pi$ separates the scalar block. This proves the algebra equalities; dimension counting completes the proof. To be precise, this is the real algebra of formal linear combinations, not physically realizable stochastic semigroup action.\end{proof}
\statement{Theorem}{IX.2}{Low-rank noncommutativity and full orbit controllability} For $n\ge3$ and $C:=QP-PQ$,
\[\rank C=2,\qquad C\beta=q(e_1-e_2),\qquad\left|\det[e_1,C\beta,PC\beta,\ldots,P^{n-2}C\beta]\right|=q^{n-1}.\]
Thus a single rank-two commutator cyclically generates all $n-1$ directions of the zero-total-mass space.
\begin{proof}$Q=I+(e_1-e_2)e_2^\top$ gives $C=(e_1-e_2)e_2^\top P-P(e_1-e_2)e_2^\top$, so $\rank C\le2$, with equality for $n\ge3$. Also $C\beta=-qw_1$, while the triangular Krylov matrix $[u,Au,\ldots,A^{m-1}u]$ has determinant of modulus 1. A factor $q$ in each of its $m$ columns gives the claimed determinant.\end{proof}
\chapter{One Probe: A Sparse Return Invariant and Period Identification}
For $n\ge3$, put $m=n-1$, $g_k=f(P^kQ)$, $g_0=1$. These quantities use one fixed terminal measurement probe $Q$ after $k$ repetitions of $P$.
\statement{Theorem}{X.1}{Sparse periodic defect: a new form of the canonical invariant} For every $k\ge0$ and $0\le a<1$,
\begin{equation}g_{k+1}-ag_k=q\one_{\{m\mid k\}},\qquad G(z):=\sum_{k\ge0}g_kz^k=\frac{1-z^m+qz}{(1-az)(1-z^m)}.\tag{47}\end{equation}
For $n\ge3$, $g_0=g_1=1$ and $g_2=a$; the least positive index $k$ for which $g_{k+1}-ag_k=q$ is $m$. Thus, within the prescribed family, $(g_0,\ldots,g_n)$ exactly recovers $a$ and $n$.
\begin{proof}Upon first leaving $e_1$, probability mass $q$ enters the hidden cycle $e_2\to\cdots\to e_n\to e_2$ of length $m$; the mass remaining at $e_1$ is multiplied by $a$ at every $P$ step. Measurement $Q$ returns to the registered state $e_1$ only mass at $e_2$ and mass already at $e_1$. Thus the stepwise increment $q$ occurs exactly after $k\equiv0\pmod m$. Multiplying the recurrence by $z^{k+1}$ and summing gives the rational formula. For $m\ge2$, the first two steps give $g_1=1$, $g_2=a$; the first subsequent positive increment has index $m$.\end{proof}
\statement{Theorem}{X.2}{Minimal order and full single-probe rank} For $n\ge3$, $G_n=(g_{i+j})_{0\le i,j<n}$ has
\begin{equation}\det G_n=(-1)^{m(m-1)/2}q^m(1-a^{m-1})\ne0,\tag{48}\end{equation}
and the minimal global linear recurrence of $g_k$ has order $n$ and characteristic polynomial $(x-a)(x^m-1)$. At $a=0$, the root $x=0$ represents a finite initial transient; it must not be discarded by considering only poles of $G(z)$.
\begin{proof}Factor $G_n=OV$ with $V=[P^j\beta]_{j<n}$ and $O=[\alpha QP^i]_{i<n}$. Reachability gives $\det V=q^m$. After a permutation of hidden columns, the first $m$ rows of $O$ form an identity block, and the last hidden row repeats the first; subtracting the first row from the last leaves coefficient $a^{m-1}-1$ in the first column. The permutation sign is $(-1)^{m(m-1)/2}$, giving (48). The recurrence from (47) has order at most $m+1$, and invertibility of $G_n$ excludes every smaller order.\end{proof}
\statement{Theorem}{X.3}{Stable period detection for a known model class} Suppose $\widehat g_k\in[0,1]$ and $|\widehat g_k-g_k|\le\varepsilon$ for $0\le k\le n$. Set $\widehat a=\widehat g_2$ and $\widehat r_k=\widehat g_{k+1}-\widehat a\widehat g_k$. If
\begin{equation}\varepsilon<\frac{1-a}{7},\tag{49}\end{equation}
classifying $k=1,\ldots,n-1$ by $\widehat r_k>(1-\widehat a)/2$ detects the first periodic increment at $k=n-1$ without error, recovering $n$ and $a$ with $|\widehat a-a|\le\varepsilon$. With independent Bernoulli repetitions, it suffices to take
\[N>\frac{49}{2(1-a)^2}\log\frac{2(n+1)}\gamma\]
repetitions of each of the $n+1$ words $P^kQ$, $0\le k\le n$, to guarantee this with probability at least $1-\gamma$.
\begin{proof}All estimated probabilities lie in $[0,1]$, so $|\widehat r_k-q\one_{\{m\mid k\}}|\le\varepsilon+a\varepsilon+\varepsilon\le3\varepsilon$. The threshold differs from $q/2$ by at most $\varepsilon/2$. If $7\varepsilon<q$, the shifted intervals cannot confuse 0 and $q$. Hoeffding and a union bound over the $n+1$ experimental words give the uniform error probability. This guarantee is conditional on membership in the $P,Q$ family; it does not itself reject an arbitrary smaller-dimensional alternative.\end{proof}
\chapter{An Isolated Gap for Deterministic Macrostates}
\statement{Theorem}{XI.1}{Exact minimax cost of discrete coarse-graining: a new structural result} Let $n\ge3$, $0\le a<1$, and let $L$ deterministically map $n$ microstates into $d<n$ nonempty macroclasses while preserving $e_1^\top=\overline\alpha L$. Macrooperations $\overline P,\overline Q$ may be arbitrary column-stochastic matrices. Then
\begin{equation}\min_{L,\overline P,\overline Q}\max_{A\in\{P,Q\}}\tfrac12\|LA-\overline AL\|_{1\to1}=\tfrac12,\tag{50}\end{equation}
where the minimum is over all such $L$ and $d<n$. Unlike general stochastic coarse-graining, this quantity does not decrease with $n$ and is independent of $a$.
\begin{proof}Exact preservation of $e_1^\top$ prevents merging $e_1$ with any hidden state: the effect must be constant within each macroclass. Since $d<n$, some distinct hidden states $e_i,e_j$, $i,j\ge2$, lie in the same class. If $LPe_i=LPe_j$ and $LQe_i=LQe_j$ held for every merged pair, there would be an exact simultaneous Markov coarse-graining for both actions, hence a macromodel of predictive rank below $n$, contradicting $\rank H_n=n$. Thus some merged pair and action $A$ satisfy $LAe_i\ne LAe_j$.

Because $Pe_i,Pe_j,Qe_i,Qe_j$ are deterministic for $i,j\ge2$, these unequal outputs are distinct basis vectors. But $\overline ALe_i=\overline ALe_j$ is one probability vector. By the triangle inequality, at least one TV distance from that vector to these two point masses is at least $1/2$, proving the lower bound.

For attainment at $n=3$, merge $e_2,e_3$: $P$ descends exactly, while the two incompatible $Q$ columns are replaced by their mean, giving maximum defect $1/2$. For $n\ge4$, merge $e_3,e_4$: $Q$ descends exactly, and for the sole ambiguous macrocolumn of $P$ use the mean of its two distinct target point masses. All other macrocolumns are exact. Again the defect is $1/2$.\end{proof}
\statement{Remark}{XI.2}{Scope} Formula (50) does not extend to arbitrary stochastic channels $L$, where separation into macroclasses may be blurred. Nor does it assert that every individual reduction attains $1/2$: this is the optimum over all admissible reductions.
\chapter{Quadratic Statistical Cost: A Matrix Improvement of the Single-Probe Certificate}
In the earlier Revision IX, entrywise Hoeffding and $\|E\|_2\le n\|E\|_{\max}$ gave a sufficient budget of order $n^3\log n$ binary trials. For one terminal probe this is not optimal: repeated antidiagonals require joint control.
\statement{Lemma}{XII.1}{Orthogonal geometry of antidiagonals} Let $J_k\in\R^{n\times n}$, $(J_k)_{ij}=\one_{\{i+j=k\}}$, $0\le k\le2n-2$, with $0\le i,j<n$. Then $J_k=J_k^\top$, $\|J_k\|_2=1$, and
\begin{equation}\sum_{k=0}^{2n-2}J_k^2=nI_n.\tag{51}\end{equation}
\begin{proof}Each $J_k$ is a partial permutation and $J_k^2$ is the diagonal projection onto rows containing a unit entry. For every row there are exactly $n$ pairs $(k,j)$ with $k=i+j$, $0\le j<n$. Hence the sum of squares equals $nI_n$.\end{proof}
\statement{Lemma}{XII.2}{Explicit uniform conditioning bound for the single-probe Hankel matrix} For $n=m+1\ge3$, write $G_n=OV$, where $V=[P^je_1]_{0\le j\le m}$ and $O=[e_1^\top QP^i]_{0\le i\le m}$. With $q=1-a$,
\[\|V^{-1}\|_2\le\frac{1+a}q,\qquad\|O^{-1}\|_2\le\frac\phi q\left(q+\sqrt{2+(1-a^2)^{-1}}\right),\qquad\phi=\frac{1+\sqrt5}2.\]
In particular, $\sigma_{\min}(G_n)\ge s_a$, where $s_a$, defined below, is independent of $n$.
\begin{proof}The reachability matrix factors as $V=D(I-aS)^{-1}$, where $D=\diag(1,q,\ldots,q)$ and $S$ is the upper one-step shift. Since $\|S\|_2\le1$, the first bound follows. Permuting the hidden columns of $O$ gives
\[\begin{pmatrix}I_m&v\\e_1^\top&a^{m-1}\end{pmatrix},\qquad v=(1,1,a,\ldots,a^{m-2})^\top\in\R^m.\]
Subtracting the first row from the last produces a triangular block whose final diagonal entry is $a^{m-1}-1$, of modulus at least $q$. The row-subtraction operator has norm $\phi$ (its nontrivial block is $\left(\begin{smallmatrix}1&0\\-1&1\end{smallmatrix}\right)$). Splitting the inverse triangular block into an identity block and its last column bounds its norm by $1+\sqrt{1+\|v\|_2^2}/q$. Finally $\|v\|_2^2\le1+(1-a^2)^{-1}$. Combine these inequalities with $G_n=OV$.\end{proof}
\statement{Theorem}{XII.3}{Improved universal test: $O(n^2\log n)$ trials} For each of the $2n-1$ experiments $P^kQ$, perform $N$ independent binary repetitions, independent across settings. Let $\widehat G_n=[\widehat g_{i+j}]$, $G_n=[g_{i+j}]$, and
\[L_\gamma=\log(2n/\gamma),\qquad t_{n,N,\gamma}=\sqrt{\frac{nL_\gamma}{2N}}+\frac{2L_\gamma}{3N}.\]
For any stationary binary process with independent repetitions,
\begin{equation}\Pr\{\|\widehat G_n-G_n\|_2>t_{n,N,\gamma}\}\le\gamma.\tag{52}\end{equation}
Thus $\sigma_{\min}(\widehat G_n)>t_{n,N,\gamma}$ falsely rejects any process of linear dimension below $n$ with probability at most $\gamma$. For this manuscript's family, put $\phi=(1+\sqrt5)/2$ and
\[s_a=\frac{q^2}{\phi(1+a)\bigl(q+\sqrt{2+(1-a^2)^{-1}}\bigr)}>0.\]
The preceding bound $\sigma_{\min}(G_n)\ge s_a$ implies power at least $1-\gamma$ if $2t_{n,N,\gamma}<s_a$. It suffices that
\begin{equation}N>\max\left\{\frac{8nL_\gamma}{s_a^2},\frac{8L_\gamma}{3s_a}\right\},\qquad(2n-1)N=O_a\bigl(n^2\log(n/\gamma)\bigr).\tag{53}\end{equation}
\begin{proof}Let $B_{kr}\in\{0,1\}$ be outcome $r$ at setting $k$, with $\mathbb EB_{kr}=g_k$. Then
\[\widehat G_n-G_n=\sum_{k=0}^{2n-2}\sum_{r=1}^N\frac{B_{kr}-g_k}{N}J_k.\]
The summands are independent, self-adjoint, centered, and of norm at most $1/N$. By (51), their matrix variance parameter satisfies $\|\sum_{k,r}\mathbb EX_{kr}^2\|\le(4N)^{-1}\|\sum_kJ_k^2\|=n/(4N)$. Tropp's matrix Bernstein inequality gives
\[\Pr(\|E\|_2>t)\le2n\exp\left[-\frac{t^2}{n/(2N)+2t/(3N)}\right].\]
The displayed $t_{n,N,\gamma}$ makes this at most $\gamma$. If $\rank G_n<n$, Weyl gives $\sigma_{\min}(\widehat G_n)\le\|E\|_2$. If $\sigma_{\min}(G_n)\ge s_a$, then $\sigma_{\min}(\widehat G_n)\ge s_a-\|E\|_2$. Finally (53) ensures $t<s_a/2$.\end{proof}
\statement{Remark}{XII.4}{Scope} This improves a stated sufficient budget; it is not a universal optimal information-theoretic lower bound. It uses the same repeatedly estimated probability at all positions on an antidiagonal. For nonstationary processes or temporally correlated repetitions, (52) requires a new probability estimate.
\setcounter{equation}{53}


% 02_spectral.tex
\chapter{A Spectral Bridge: Neutrino Flavors, Inertia, and Amplitudes}
This section specifies a separate conditional physical regime; it does not identify the finite stochastic family with cosmology. It uses standard vacuum-oscillation formulas (see the PDG review in the bibliography), developed here into an operational system of rank and amplitude checks.

Let $U\in U(3)$, $\alpha\ne\beta$, $c_i=U_{\beta i}\overline{U_{\alpha i}}$, $\lambda_i=m_i^2$, and $t=L/(2E)$ in units $\hbar=c=1$. For a fixed step $\tau>0$ and $k\ge0$, put
\begin{align}A_k&=\sum_{i=1}^3c_i e^{-\ii k\lambda_i\tau},\qquad p_k=|A_k|^2,\tag{54}\\
\omega_{ij}&=(\lambda_j-\lambda_i)\tau,\qquad b_{ij}=c_i\overline c_j\quad(i<j),\tag{55}\\
b_0&=\sum_i|c_i|^2,\qquad Z=\{1,e^{\pm\ii\omega_{12}},e^{\pm\ii\omega_{13}},e^{\pm\ii\omega_{23}}\}.\tag{56}\end{align}
Here $p_k$ is the ideal vacuum probability measured at different $L/E$, not a detector event count. The same $U,\lambda_i$ are required for all $k$, with no matter, coherence damping, or within-bin energy variation. Matter, an inhomogeneous potential, systematic errors, and energy averaging change the model class.
\statement{Theorem}{XIII.1}{Discriminant, inertia, and finite recurrence of a vacuum channel} Assume $c_1c_2c_3\ne0$ and that all seven points of $Z$ on the unit circle are pairwise distinct (no aliasing). For $H_r=(p_{i+j})_{0\le i,j<r}$,
\begin{align}\rank H_7&=7,\qquad\operatorname{Inertia}(H_7)=(4,3,0),\qquad\det H_8=0,\tag{57}\\
\det H_7&=b_0\prod_{i=1}^3|c_i|^4\prod_{1\le r<s\le7}(z_s-z_r)^2<0.\tag{58}\end{align}
Here $z_1,\ldots,z_7$ enumerate $Z$ arbitrarily; the squared Vandermonde product is independent of their ordering. The minimal shift polynomial of $p_k$ is
\begin{equation}F(X)=(X-1)\prod_{i<j}(X^2-2\cos\omega_{ij}X+1),\tag{59}\end{equation}
of degree seven. Fourteen consecutive exact samples $p_0,\ldots,p_{13}$ recover the coefficients of $F$ by solving a linear system with $H_7$, after which Prony's method recovers nodes and weights up to permutation and phase aliasing. For a transition $\alpha\ne\beta$, $p_0=0$ is known, so 13 nonzero baseline settings suffice.
\begin{proof}Expanding the square gives the real exponential sum
\[p_k=b_0+\sum_{i<j}\left[b_{ij}e^{\ii k\omega_{ij}}+\overline b_{ij}e^{-\ii k\omega_{ij}}\right].\]
Let $B_s$ be the coefficient corresponding to $z_s$, and $W_{rs}=z_s^r$, $0\le r\le6$. Then $H_7=W\diag(B_s)W^\top$. Distinct nodes make $W$ invertible. The Vandermonde formula gives (58), since $\prod_sB_s=b_0\prod_{i<j}|b_{ij}|^2=b_0\prod_i|c_i|^4$. To determine the sign, express each conjugate-frequency block in the real basis $(\cos k\omega,\sin k\omega)$. Its real symmetric form is
\[2\begin{pmatrix}\Re b&-\Im b\\-\Im b&-\Re b\end{pmatrix},\]
with eigenvalues $\pm2|b|$. The zero frequency contributes $b_0>0$, so Sylvester's law of inertia gives $(4,3,0)$ and a negative determinant. The exponential sum satisfies (59); invertibility of $H_7$ excludes a smaller recurrence. Coefficient recovery uses $\sum_{j=0}^7F_jp_{k+j}=0$, $0\le k\le6$, $F_7=1$. The same factorization with eight rows gives $\rank H_8\le7$.\end{proof}
\statement{Corollary}{XIII.2}{Counting coherent modes by signature} Suppose there are $M\ge2$ nondegenerate coherent eigenphases, all $M(M-1)/2$ positive pairwise differences are distinct up to sign modulo $2\pi$, and all $c_i$ are nonzero. For $p_k=|\sum_{i=1}^Mc_i e^{-\ii k\lambda_i\tau}|^2$, the minimal real linear dimension and inertia of its full Hankel matrix are
\[r=M(M-1)+1,\qquad\operatorname{Inertia}(H_r)=\left(1+\binom M2,\binom M2,0\right).\]
This conditionally counts phase modes; it does not identify an individual particle. Additional frequencies may have other physical or instrumental causes.
\begin{proof}Every pair $i<j$ contributes a conjugate pair of nodes with nonzero coefficients. The Vandermonde matrix is nonsingular by hypothesis, and every real pair block has one positive and one negative eigenvalue. The single constant term $\sum_i|c_i|^2$ is positive. Apply Sylvester's law.\end{proof}
\statement{Theorem}{XIII.3}{Exact boundary between finite quantum and finite classical memory} Suppose a real sequence $0\le p_k\le1$, $k\ge0$, has a finite purely nondissipative exponential representation
\[p_k=\sum_{z\in Z_*}b_z z^k,\qquad |z|=1,\quad b_z\ne0,\]
with distinct nodes and conjugate-paired weights making $p_k$ real. The following are equivalent:
\begin{enumerate}[label=(\alph*)]
\item A finite time-homogeneous classical model $p_k=\alpha M^k\beta$ exists, with column-stochastic $M$, probability vector $\beta$, and effect $0\le\alpha\le\one^\top$.
\item Every active node $z\in Z_*$ is a root of unity.
\item The sequence $p_k$ is strictly periodic from $k=0$.
\end{enumerate}
In particular, for the nondegenerate three-flavor vacuum channel of Theorem XIII.1, a nonzero coefficient at any frequency $\omega_{ij}/(2\pi)\notin\Q$ excludes every finite homogeneous classical Markov explanation of all $p_k$, although a quantum realization on $\C^3$ exists and the linear rank is seven. The quantifier ``all $k$'' is essential.
\begin{proof}Every unit-modulus eigenvalue of a finite nonnegative stochastic matrix is a root of unity: in Markov-chain normal form such eigenvalues arise only in closed irreducible classes with finite integer periods (the peripheral Perron--Frobenius theorem). Under (a), $\sum_{k\ge0}p_kt^k=\alpha(I-tM)^{-1}\beta$ has poles only among reciprocals of nonzero eigenvalues of $M$. The given exponential form has poles at $t=z^{-1}$ with nonzero residues for every distinct $z\in Z_*$. Hence all such $z$ are roots of unity, proving (b). Under (b), choose a common multiple $T$ of all node orders; then $p_{k+T}=p_k$, giving (c). Under (c), use $T$ classical states with $Me_j=e_{j+1\bmod T}$, initial $e_0$, and effect $\alpha e_j=p_j$. Since $0\le p_j\le1$, this is admissible and proves (a).\end{proof}
\statement{Corollary}{XIII.4}{Explicit rational three-level example with no finite classical realization} There exist an exact $3\times3$ unitary matrix $U$ with rational real and imaginary parts and phase nodes
\[z_{12}=\frac{3+4\ii}5,\qquad z_{23}=\frac{5+12\ii}{13},\qquad z_{13}=z_{12}z_{23}=\frac{-33+56\ii}{65}\]
such that the $e\to\mu$ channel has $\rank H_7=7$ and $\operatorname{Inertia}(H_7)=(4,3,0)$, but no finite time-homogeneous classical Markov model with fixed preparation and effect reproduces the whole sequence. The exact matrix certificate is in the accompanying \texttt{checks/verify\_neutrino.py}: $U=R_{23}R_{13}R_{12}$ with rational sine--cosine pairs $(4/5,3/5)$, $(12/13,5/13)$, $(8/17,15/17)$ and phase $\pi/2$ in $R_{13}$.
\begin{proof}All three nodes have unit modulus and, together with their inverses and 1, are distinct. Rational arithmetic verifies that all $c_j$ are nonzero in the stated example. But $z_{12}$ is not a root of unity: otherwise the algebraic integer $z_{12}+z_{12}^{-1}=6/5$ would be a rational integer, hence an ordinary integer, a contradiction. Apply Theorems XIII.1 and XIII.3.\end{proof}
\statement{Theorem}{XIII.5}{Triangular certificate for neutrino frequencies and complex amplitudes} Under Theorem XIII.1 and with physical difference labels $\omega_{ij}$, coefficients recovered from $p_k$ satisfy
\begin{align}\omega_{13}&\equiv\omega_{12}+\omega_{23}\pmod{2\pi},\tag{60}\\
R:=\frac{b_{12}b_{23}}{b_{13}}&=|c_2|^2\in\R_{>0},\tag{61}\\
b_0&=R+\frac{|b_{12}|^2+|b_{23}|^2}{R},\qquad b_0+2\Re(b_{12}+b_{13}+b_{23})=0.\tag{62}\end{align}
Conversely, for nonzero $b_{12},b_{13},b_{23}$ and $b_0>0$, (61)--(62) suffice for complex $c_1,c_2,c_3$ with $\sum c_i=0$, $b_{ij}=c_i\overline c_j$, and $b_0=\sum|c_i|^2$. This is scalar amplitude closure; without further unitary constraints it does not prove existence of a full physical mixing matrix $U$.
\begin{proof}The first equation expresses additivity of mass differences; the second is $c_1\overline c_2c_2\overline c_3/(c_1\overline c_3)=|c_2|^2$. Hence $|c_1|^2=|b_{12}|^2/R$ and $|c_3|^2=|b_{23}|^2/R$. Orthogonality of rows $U_\alpha,U_\beta$ gives $\sum_i c_i=0$, so $|\sum c_i|^2=b_0+2\Re\sum b_{ij}=0$. Conversely choose $c_2=\sqrt R$, $c_1=b_{12}/\sqrt R$, $c_3=\overline b_{23}/\sqrt R$. Then $c_1\overline c_3=b_{12}b_{23}/R=b_{13}$. The first part of (62) gives the norm, and the second gives $\sum c_i=0$.\end{proof}
\statement{Corollary}{XIII.6}{Even/odd components and standard CP asymmetry} For antineutrinos in vacuum, the $c_i$ are conjugated, so $\overline p_k=p_{-k}$. With $J_{\alpha\beta}:=\Im(c_1\overline c_2)$, the familiar three-flavor identity is
\begin{equation}p(t)-\overline p(t)=-16J_{\alpha\beta}\sin\frac{\Delta_{21}t}{2}\sin\frac{\Delta_{31}t}{2}\sin\frac{\Delta_{32}t}{2},\tag{63}\end{equation}
where $\Delta_{ij}=\lambda_i-\lambda_j$, $i>j$. The sign depends on the stated convention for $c_i$ and the channel direction. This identity is not new and does not extend unchanged to propagation through matter.
\begin{proof}For $i<j$, the neutrino--antineutrino difference of pair contributions is $-4\Im(b_{ij})\sin(\Delta_{ji}t)$. From $c_1+c_2+c_3=0$, $\Im b_{12}=\Im b_{23}=-\Im b_{13}=J_{\alpha\beta}$. Apply $\sin x-\sin(x+y)+\sin y=4\sin(x/2)\sin(y/2)\sin((x+y)/2)$.\end{proof}
\statement{Corollary}{XIII.7}{Parameter-free statistical rejection of three-flavor vacuum dynamics} Suppose $p_k$, $k=0,\ldots,14$, are measured at fixed $\tau$ using $N$ independent binary trials per setting. Put $\widehat H_8=(\widehat p_{i+j})_{0\le i,j<8}$ and
\[t_8=\sqrt{\frac{8\log(16/\gamma)}{2N}}+\frac{2\log(16/\gamma)}{3N}.\]
The rule $\sigma_{\min}(\widehat H_8)>t_8$ rejects the entire stationary coherent three-flavor vacuum model with false-rejection probability at most $\gamma$. It requires no estimation of $U$, masses, or CP phase. It also applies to degenerate frequencies because $\rank H_8\le7$ remains true. It does not ensure high power against arbitrary deviations and is not sufficient evidence for a new neutrino.
\begin{proof}The exponential decomposition gives $\rank H_8\le7$ even without nondegeneracy. Apply the antidiagonal matrix concentration theorem with $n=8$ and Weyl's smallest-singular-value bound.\end{proof}
\statement{Theorem}{XIII.8}{Operational boundary: reconstructing the antineutrino curve and nonidentifiability of the sign of $J$} In ideal vacuum, for any number of flavors and any pair $\alpha,\beta$, two logically distinct assertions hold.
\begin{enumerate}[label=(\roman*)]
\item $\overline p_{\alpha\beta}(t)=p_{\alpha\beta}(-t)$. Under the nondegenerate assumptions of XIII.1, fourteen exact neutrino samples $p_0,\ldots,p_{13}$ determine the full $\overline p_k$ sequence on the same discrete grid. In particular, if $F(X)=\sum_{j=0}^7F_jX^j$, $F_7=1$, then $F_0=-1$ and
\begin{equation}\overline p_1=p_{-1}=p_6+\sum_{j=1}^6F_jp_{j-1}.\tag{64}\end{equation}
\item For any real $C\ge\max_i\lambda_i$, the replacement
\begin{equation}U\longmapsto\overline U,\qquad\lambda_i\longmapsto C-\lambda_i\tag{65}\end{equation}
leaves all vacuum neutrino and antineutrino probabilities unchanged at every $t$, but sends $J_{\alpha\beta}=\Im(c_1\overline c_2)$ to $-J_{\alpha\beta}$. Thus without a spectral orientation and labeling convention, no vacuum probabilities identify the sign of $J$, while the physical CP asymmetry is preserved.
\end{enumerate}
\begin{proof}The first equality follows from $|\sum_i\overline c_i e^{-\ii\lambda_it}|^2=|\sum_i c_i e^{+\ii\lambda_it}|^2$. Since $F(0)=-1$, the recurrence extends uniquely from $k\ge0$ to $k<0$; at $k=-1$ it gives (64). For the second claim, each transition amplitude transforms as
\[A'_{\alpha\beta}(t)=\sum_i\overline c_i e^{-\ii(C-\lambda_i)t}=e^{-\ii Ct}\overline{A_{\alpha\beta}(t)}.\]
Taking the modulus preserves probability. The same calculation applies to antineutrinos. Separately $c_i\mapsto\overline c_i$, hence $J\mapsto-J$, while $\Delta_{ij}\mapsto-\Delta_{ij}$; the product on the right of (63) remains invariant.\end{proof}
\statement{Remark}{XIII.9}{Statistical meaning and interpretation limits} An arbitrary seven-phase signal can have the same rank pattern. Thus $\det H_8=0$ alone does not identify neutrinos; independent checks of (60)--(62) are substantially stronger. Under noise, zero and inertia tests are replaced by Weyl intervals for eigenvalues/singular values, with error probability controlled for the specific protocol. Including sterile states, matter interactions, cosmological curvature, or a dynamic environment requires a separate Hamiltonian derivation and does not follow from this theorem.
\chapter{An Arithmetic Invariant of Minimal Memory: An Exact Dimensional Separation}
The following assertions concern one scalar time sequence with fixed preparation and effect, not an arbitrary multistep process with interventions. ``Homogeneous'' and ``stationary'' are not interchangeable: the initial state need not be invariant under the transition.
\statement{Definition}{XIV.1}{Minimal classical memory} For $p=(p_k)_{k\ge0}\subset[0,1]$, define
\[\mu_{\mathrm{cl}}(p)=\inf\left\{n\in\N:\begin{array}{l}\exists M,\beta,\alpha:\ M\in\R_{\ge0}^{n\times n},\ \one^\top M=\one^\top,\\\beta\in\R_{\ge0}^n,\ \one^\top\beta=1,\ 0\le\alpha\le\one^\top,\\p_k=\alpha M^k\beta\quad\forall k\ge0\end{array}\right\},\]
with $\inf\varnothing=+\infty$. This is the size of a positive Markov realization and may strictly exceed real Hankel rank.
\statement{Theorem}{XIV.2}{Unbounded exact quantum--classical dimensional separation} For every integer $p\ge7$, put $\zeta_p=e^{2\pi\ii/p}$,
\[D_p=\diag(1,\zeta_p,\zeta_p^3),\qquad\psi=\tfrac1{\sqrt3}(1,1,1)^\top,\qquad s_p(k)=|\langle\psi,D_p^k\psi\rangle|^2.\]
Then
\begin{align}s_p(k)&=\tfrac13+\tfrac29\left(\cos\frac{2\pi k}p+\cos\frac{4\pi k}p+\cos\frac{6\pi k}p\right),\tag{66}\\
\rank H_7(s_p)&=7,\qquad\operatorname{Inertia}(H_7(s_p))=(4,3,0),\qquad\mu_{\mathrm{cl}}(s_p)=p.\tag{67}\end{align}
Thus minimal classical memory is unbounded on a fixed three-dimensional Hilbert space with fixed Hankel rank seven.
\begin{proof}The differences of exponents $0,1,3$ are $1,2,3$ up to sign. For $p\ge7$, the residues $0,\pm1,\pm2,\pm3$ are distinct. The exponential expansion therefore has seven distinct nodes $1,\zeta_p^{\pm1},\zeta_p^{\pm2},\zeta_p^{\pm3}$ with nonzero weights $1/3$ and $1/9$ at each nonzero frequency. Vandermonde factorization and Sylvester's law (XIII.1) give rank seven and inertia $(4,3,0)$. If an $n$-state classical realization exists, the nonzero coefficient at $\zeta_p$ gives a generating-function pole at $t=\zeta_p^{-1}$, hence $\zeta_p\in\spec M$. By peripheral Perron--Frobenius theory, it must arise from a closed irreducible Markov component of period $d$ with $\zeta_p^d=1$. Since $\zeta_p$ has exact order $p$, $p\mid d$. Such a component has at least $d$ states, so $n\ge d\ge p$. Conversely, the $p$-cycle $Me_j=e_{j+1\bmod p}$, $\beta=e_0$, $\alpha e_j=s_p(j)$ is admissible because $s_p\in[0,1]$ is $p$-periodic. The bounds coincide.\end{proof}
\statement{Corollary}{XIV.3}{An entirely rational result for $p=7$} For $p=7$, the six nonzero frequencies exhaust the nontrivial seventh roots of unity, so
\[s_7(k)=\begin{cases}1,&7\mid k,\\2/9,&7\nmid k.\end{cases}\]
The real rational matrix $H_7$ has
\begin{equation}\det H_7=-\frac{823543}{1594323},\qquad\spec(H_7)=\left\{\frac73,\left(\frac79\right)^{[3]},\left(-\frac79\right)^{[3]}\right\}.\tag{68}\end{equation}
In particular $\mu_{\mathrm{cl}}(s_7)=7$ without numerical approximation of roots of unity.
\begin{proof}$\sum_{j=0}^6\zeta_7^{jk}=7\one_{\{7\mid k\}}$ gives $s_7(k)=2/9+(7/9)\one_{\{7\mid k\}}$. The matrix $(\one_{\{7\mid i+j\}})_{0\le i,j\le6}$ is a permutation involution with eigenvalues $(+1)^{[4]},(-1)^{[3]}$. The all-ones matrix commutes with it and acts only on the direction $\one$, with eigenvalue seven. Thus one eigenvalue $7/9$ becomes $7/9+(2/9)7=7/3$, with the others unchanged. Their product gives (68).\end{proof}
\statement{Theorem}{XIV.4}{Barrier for finite continuous-time classical dynamics} Let $Q$ be a finite classical Markov generator: $Q_{ij}\ge0$ for $i\ne j$ and $\one^\top Q=0$. For any $\tau>0$, probability vector $\beta$, and effect $0\le\alpha\le\one^\top$, the sequence $q_k=\alpha e^{k\tau Q}\beta$ cannot agree for every $k\ge0$ with a nonconstant finite sum of pure oscillatory frequencies $s_k=\sum_{|z|=1}b_z z^k$ with distinct nodes and nonzero weights. In particular, even (66) admits no such finite continuous-time classical realization.
\begin{proof}$M=e^{\tau Q}$ is stochastic and every diagonal entry is strictly positive (there is positive probability of no jump in time $\tau$). Thus each closed irreducible component is aperiodic and its only possible peripheral eigenvalue is 1. A nonconstant finite oscillatory sum has an active node other than 1. Comparing generating-function poles as in XIII.3 gives a contradiction.\end{proof}
\statement{Theorem}{XIV.5}{No universal finite-sample rejection of unbounded classical memory} For every $N\ge0$ and numbers $p_0,\ldots,p_N\in[0,1]$, a finite homogeneous classical model with at most $N+1$ states reproduces all these numbers exactly. Consequently, in an experiment independently recording binary measurements of probability $p_k$ at each $0\le k\le N$, no test based only on these measurements uniformly separates a specified quantum curve from all finite-dimensional homogeneous classical models with unbounded state count.
\begin{proof}Use the cyclic shift $Me_j=e_{j+1\bmod(N+1)}$ on $N+1$ states, initial $e_0$, and effect $\alpha e_j=p_j$. Then $\alpha M^ke_0=p_k$ on the required horizon. Independent repetitions of the finite protocol have identical full joint observation laws, not merely identical expected probabilities. No test using this protocol can distinguish the models. Their continuations may differ beyond $N$.\end{proof}
\statement{Corollary}{XIV.6}{Three inequivalent complexity criteria} For every $p\ge7$, the same data satisfy
\[\dim\mathcal H=3,\qquad\rank H=7,\qquad\mu_{\mathrm{cl}}=p.\]
If an active eigenphase difference is irrational relative to $2\pi$ while the seven frequencies remain nondegenerate, the same pair $(3,7)$ occurs but $\mu_{\mathrm{cl}}=+\infty$ by XIII.3. Thus peripheral-spectrum arithmetic is a separate operational invariant not recoverable from rank or inertia alone.
\chapter{A Conditional Geometric Bridge to the Round $S^3$ and Oscillation Frequencies}
Let $S_R^3$ be the round sphere of radius $R$ and $-\Delta_{S_R^3}$ the positive scalar Laplace--Beltrami operator. Its standard spectrum is
\[\Lambda_\ell=\frac{\ell(\ell+2)}{R^2},\qquad\dim E_\ell=(\ell+1)^2,\qquad\ell=0,1,\ldots.\]
This is spectral geometry, not an independent evolution equation for our universe. Dynamics requires the additional explicit hypothesis of a one-particle Hamiltonian $H=\kappa(-\Delta_{S_R^3})$, $\kappa>0$, and independent justification of its physical applicability. Choose one normalized eigenvector in each of $E_0,E_1,E_2$ with equal initial amplitudes.
\statement{Theorem}{XV.1}{Arithmetic spectral barrier on the round $S^3$} Put $\theta=\kappa\tau/R^2$ and
\[g_\theta(k)=\tfrac19|1+e^{-3\ii\theta k}+e^{-8\ii\theta k}|^2=\tfrac13+\tfrac29\bigl(\cos(3\theta k)+\cos(5\theta k)+\cos(8\theta k)\bigr).\]
If the six nonzero nodes $e^{\pm3\ii\theta},e^{\pm5\ii\theta},e^{\pm8\ii\theta}$ are distinct and different from 1, then $\rank H_7(g_\theta)=7$ and its inertia is $(4,3,0)$. If $\theta/(2\pi)$ is irrational, $\mu_{\mathrm{cl}}(g_\theta)=\infty$. For prime $p\ge17$ and $\theta=2\pi/p$, all nodes are nondegenerate and $\mu_{\mathrm{cl}}(g_\theta)=p$.
\begin{proof}The first three eigenvalues are proportional to $0,3,8$, with differences $3,5,8$. The nondegenerate rank formula follows from Vandermonde factorization with weights $1/3$ and $1/9$. For irrational $\theta/(2\pi)$, nonzero integer multiples of $\theta$ do not generate roots of unity, so XIII.3 applies. For prime $p\ge17$, the residues $0,\pm3,\pm5,\pm8$ are distinct because all sums of two elements of $\{3,5,8\}$ are at most $16<p$; each nonzero frequency has exact order $p$. The Markov lower bound and cyclic upper construction of XIV.2 finish the proof.\end{proof}
\statement{Remark}{XV.2}{Physical boundary of $S^3$, neutrinos, and gravity} The established chain ``geometric spectrum $\to$ spectrum of a deliberately chosen Hamiltonian $\to$ classical memory'' is a rigorous conditional model. It does not imply that real neutrino mass splittings have ratio $3:5:8$, that $R$ determines the size of the universe, that a new CP phase exists, or that cosmological curvature is resolved. Physical identification requires a specific dynamical action, observables, metric, and data covariance model. OpenAI Math family 350 further illustrates why spectral-nodal properties of the round metric cannot be transferred without proof to all smooth perturbations of $S^3$.
\chapter{A Quadratic Memory Bound: Finite Geometry, Spectral Saturation, and Finite Certificates}
This section obtains an infinite family of exact dimension theorems new to this continuous manuscript. Its combinatorial starting point is the classical Singer construction; peripheral stochastic spectra and linear realization theory are also established results. The originality of the combination and worldwide priority have not been established. The results concern one experiment of the form ``preparation---$n$ repetitions of a fixed channel---binary effect'', not all multitime instrument correlations.
\statement{Lemma}{XVI.1}{Bound on the number of active frequencies of a unitary return} Let $U$ be unitary on $\C^d$, $\psi$ a unit vector, and $P(n)=|\langle\psi,U^n\psi\rangle|^2$. Then $\rank H(P)\le d(d-1)+1$. With pairwise distinct active eigenphases and nonzero coordinates of $\psi$, equality holds if and only if all ordered nonzero active-phase differences are distinct modulo $2\pi$. Here $H(P)=[P(i+j)]_{i,j\ge0}$ is the infinite linear Hankel matrix; for finite $d$, a sufficiently large finite submatrix determines its rank.
\begin{proof}In an eigenbasis of $U$, write $w_j=|\psi_j|^2>0$ on the $r\le d$ active phases $\theta_j$. Then $P(n)=\sum_{j,\ell=1}^r w_jw_\ell e^{\ii n(\theta_j-\theta_\ell)}$. Zero difference contributes one node and the others at most $r(r-1)$. All frequency coefficients are nonnegative and strictly positive for active pairs, so coincidences cannot cancel. Vandermonde factorization makes the Hankel rank of a finite exponential sum equal to its number of distinct active nodes. The maximum $d(d-1)+1$ is attained exactly under the stated conditions.\end{proof}
\statement{Theorem}{XVI.2}{Singer family: exact optimal quantum and classical memory} Let $q\ge2$ be a prime power, $d=q+1$, $p=q^2+q+1=d^2-d+1$. There is a set $D\subset\Z/p\Z$ of size $d$ such that every nonzero residue $r$ has exactly one ordered representation $r=a-b$, $a,b\in D$, $a\ne b$ (a cyclic planar Singer difference set). Put
\[U_D=\diag(\zeta_p^a:a\in D),\quad\psi_D=d^{-1/2}(1,\ldots,1)^\top,\quad P_q(n)=|\langle\psi_D,U_D^n\psi_D\rangle|^2,\quad\zeta_p=e^{2\pi\ii/p}.\]
For all $n\in\N_0$,
\begin{equation}P_q(n)=\frac{d-1}{d^2}+\frac p{d^2}\one_{\{p\mid n\}}.\tag{69}\end{equation}
Among all finite-dimensional quantum models $P(n)=\operatorname{tr}(E\Phi^n(\rho))$ with one fixed quantum channel, state, and effect $0\le E\le I$, and all finite homogeneous classical Markov realizations of XIV.1, the exact minimal dimensions are
\begin{equation}\mu_{\mathrm{qm}}(P_q)=d=q+1,\qquad\mu_{\mathrm{cl}}(P_q)=p=q^2+q+1=d^2-d+1.\tag{70}\end{equation}
The quantum minimum is not restricted to unitary channels.
\begin{proof}Existence of $D$ is Singer's classical result for the finite projective plane $PG(2,q)$, an explicit external premise rather than a new proof of finite geometry. The phase-difference formula gives
\[d^2P_q(n)=d+\sum_{r=1}^{p-1}\zeta_p^{rn}=\begin{cases}d+p-1=d^2,&p\mid n,\\d-1,&p\nmid n,\end{cases}\]
proving (69). In particular all $p$ frequencies are active because
\begin{equation}\widehat P_q(0)=\frac1d,\quad\widehat P_q(r)=\frac1{d^2}>0\ (1\le r<p),\quad\widehat P_q(r)=\frac1p\sum_{n=0}^{p-1}P_q(n)\zeta_p^{-rn}.\tag{71}\end{equation}
Thus $\rank H(P_q)=p$. In a quantum model of dimension $m$, the matrices $\Phi^n(\rho)$ belong to the real $m^2$-dimensional Hermitian space, so $\rank H\le m^2$. If $m\le d-1$, then $m^2\le(d-1)^2<d^2-d+1=p$ (here $d\ge3$). Hence every channel requires $m\ge d$, attained by the unitary construction.

For an $m$-state classical transition, Hankel rank is at most $m$, so $m\ge p$ without any peripheral-spectrum argument. The deterministic cycle $Me_j=e_{j+1\bmod p}$, $\beta=e_0$, $\alpha e_j=P_q(j)$ attains $m=p$. This proves both minima. This proof specifically uses full spectral saturation and does not apply to arbitrary $p$ in XIV.2, where rank stays fixed and the classical minimum requires a peripheral argument.\end{proof}
\statement{Corollary}{XVI.3}{Complete rational spectral signature} Let $H_p=[P_q(i+j)]_{0\le i,j<p}$ and $a=(d-1)/d^2$. Then $H_p=a\one\one^\top+(p/d^2)J_p$, where $(J_p)_{ij}=\one_{\{i+j\equiv0\pmod p\}}$. Since $p$ is odd, the exact spectral data are
\begin{align}\spec(H_p)&=\left\{\frac pd,\left(\frac p{d^2}\right)^{[(p-1)/2]},\left(-\frac p{d^2}\right)^{[(p-1)/2]}\right\},\tag{72}\\
\operatorname{Inertia}(H_p)&=\left(\frac{p+1}2,\frac{p-1}2,0\right),\quad\sigma_{\min}(H_p)=\frac p{d^2},\quad\kappa_2(H_p)=d,\tag{73}\\
\det H_p&=(-1)^{(p-1)/2}\frac{p^p}{d^{2p-1}}.\tag{74}\end{align}
\begin{proof}$J_p$ reflects indices $j\mapsto-j$; for odd $p$, one point is fixed and $(p-1)/2$ pairs are exchanged. Its eigenvalues are $(+1)^{[(p+1)/2]}$, $(-1)^{[(p-1)/2]}$. Along $\one$, the addition $a\one\one^\top$ changes $p/d^2$ to $p/d^2+ap=p/d$ and vanishes on its orthogonal complement. The formulas follow by multiplying these eigenvalues.\end{proof}
\statement{Theorem}{XVI.4}{The strict quadratic separation survives primitive mixing and depolarization} Under XVI.2, for $0<\eta\le1$ consider the completely positive trace-preserving channel
\[\Phi_\eta(\rho)=\eta U_D\rho U_D^\dagger+(1-\eta)\frac{I_d}d\operatorname{tr}\rho,\]
with $\rho_0=|\psi_D\rangle\langle\psi_D|$ and effect $E=\rho_0$. Its return probability is
\begin{equation}P_{q,\eta}(n)=\frac1d+\frac{\eta^n}{d^2}\sum_{r=1}^{p-1}\zeta_p^{rn}=\frac1d+\eta^n\left(P_q(n)-\frac1d\right).\tag{75}\end{equation}
For every strictly positive $\eta$,
\[\mu_{\mathrm{qm}}(P_{q,\eta})=d,\qquad\mu_{\mathrm{cl}}(P_{q,\eta})=p,\qquad\rank H(P_{q,\eta})=p.\]
For $0<\eta<1$, a strictly positive primitive aperiodic transition matrix with a unique stationary distribution and exponential mixing attains the classical minimum $p$.
\begin{proof}The depolarizing projection commutes with unitary conjugation, so $\Phi_\eta^n(\rho_0)=\eta^nU_D^n\rho_0(U_D^\dagger)^n+(1-\eta^n)I_d/d$, giving (75). Its $p$ distinct nodes $1,\eta\zeta_p,\ldots,\eta\zeta_p^{p-1}$ have nonzero weights $1/d,1/d^2,\ldots,1/d^2$, so the Hankel rank is $p$. The bounds $m^2\ge p$ and $m\ge p$ remain unchanged, and the given channel supplies quantum upper bound $d$.

For the classical upper bound, take $Ce_j=e_{j+1\bmod p}$, $\Pi=p^{-1}\one\one^\top$, $M_\eta=\eta C+(1-\eta)\Pi$, $\beta=e_0$, $\alpha e_j=P_q(j)$. Since $C\Pi=\Pi C=\Pi$, $M_\eta^n=\Pi+\eta^n(C^n-\Pi)$ and
\[\alpha M_\eta^n\beta=\frac1p\sum_{j=0}^{p-1}P_q(j)+\eta^n\left(P_q(n)-\frac1p\sum_{j=0}^{p-1}P_q(j)\right)=P_{q,\eta}(n),\]
because the period mean is $1/d$. For $\eta<1$, every entry of $M_\eta$ is at least $(1-\eta)/p>0$, making the chain primitive.\end{proof}
\statement{Corollary}{XVI.5}{Determinant and inertia with dissipation: an exact formula} For $0<\eta\le1$, define $H_p(\eta)=[P_{q,\eta}(i+j)]_{0\le i,j<p}$ and $S_p(\eta)=1+\eta+\cdots+\eta^{p-1}$. Then
\begin{equation}\det H_p(\eta)=(-1)^{(p-1)/2}\frac{p^{p-2}}{d^{2p-1}}\eta^{(p-1)(p-2)}S_p(\eta)^2.\tag{76}\end{equation}
Throughout $\eta\in(0,1]$, $\operatorname{Inertia}(H_p(\eta))=((p+1)/2,(p-1)/2,0)$. It is the exact nonzero determinant, not uniform numerical stability as $\eta\to0$, that preserves rank.
\begin{proof}Set $\lambda_0=1$, $\lambda_r=\eta\zeta_p^r$ for $1\le r<p$, $c_0=1/d$, $c_r=1/d^2$. With $V_{ir}=\lambda_r^i$, $H_p(\eta)=V\diag(c_r)V^\top$, hence
\[\det H_p(\eta)=\frac1{d^{2p-1}}\prod_{0\le r<s<p}(\lambda_s-\lambda_r)^2.\]
For the $p$th roots other than 1, the squared Vandermonde product of the monic polynomial $(x^p-1)/(x-1)$ is $(-1)^{(p-1)/2}p^{p-2}$. The factors involving 1 give $\prod_{r=1}^{p-1}(1-\eta\zeta_p^r)=S_p(\eta)$ because $p-1$ is even. The other differences contribute $\eta^{(p-1)(p-2)/2}$ before squaring, proving (76). Inertia at $\eta=1$ is given by XVI.3. Entries vary continuously and the determinant never vanishes on $(0,1]$, so no eigenvalue crosses zero.\end{proof}
\statement{Theorem}{XVI.6}{Two-state limit and failure of uniform testing at vanishing coherence} With $q,d,p$ fixed, as $\eta\downarrow0$ every exact $P_{q,\eta}(n)$ tends to
\[P_{q,0}(n)=\begin{cases}1,&n=0,\\1/d,&n\ge1.\end{cases}\]
This limit has a homogeneous classical two-state realization and Hankel rank two. Suppose independent binary measurements at times $k\ge0$ use finite trial counts $N_k$. Their joint laws $\mathbb P_\eta,\mathbb P_0$ satisfy
\begin{equation}\operatorname{TV}(\mathbb P_\eta,\mathbb P_0)\le\left(1-\frac1d\right)\sum_{k\ge1}N_k\eta^k\le\left(1-\frac1d\right)\eta\sum_{k\ge1}N_k.\tag{77}\end{equation}
Thus, at any prescribed finite trial budget, no test uniformly distinguishes exact classical memory $p$ for every $\eta>0$ from the two-state model with power bounded away from the trivial limit. The sum of its two errors is at least $1-\operatorname{TV}(\mathbb P_\eta,\mathbb P_0)$, tending to one. Exact dimension preservation in XVI.4 must not be confused with uniform experimental distinguishability.
\begin{proof}The two-state model $Me_0=e_1$, $Me_1=e_1$, $\beta=e_0$, $\alpha=(1,1/d)$ realizes $P_{q,0}$. From (75), for $k\ge1$,
\[|P_{q,\eta}(k)-P_{q,0}(k)|=\eta^k|P_q(k)-1/d|\le(1-1/d)\eta^k,\]
and the difference at zero vanishes. An optimal coupling of Bernoulli$(u)$ and Bernoulli$(v)$ has mismatch probability $|u-v|$. Coupling all independent repetitions and summing mismatch probabilities gives (77). The final claim is the standard two-point testing inequality.\end{proof}
\statement{Corollary}{XVI.7}{Rigidity of a flat spectrum and projective geometry} Let $p=d(d-1)+1$ and let $D\subset\Z/p\Z$ contain $d$ distinct exponents. For equal amplitudes $\psi_D$, the condition
\[\left|d^{-1}\sum_{a\in D}\zeta_p^{an}\right|^2=\frac{d-1}{d^2}\quad\text{for every }n=1,\ldots,p-1\]
is equivalent to every $r\ne0$ occurring exactly once among ordered differences $a-b$, $a\ne b$. In this cyclic equal-amplitude class, two-level spectral flatness is therefore equivalent to a perfect difference set.
\begin{proof}Let $N_r=|\{(a,b)\in D^2:a-b\equiv r\}|$. Then $d^2P(n)=\sum_{r=0}^{p-1}N_r\zeta_p^{rn}$, $N_0=d$. The stated values for $n\ne0$, together with $P(0)=1$, give $N_r=1$ for every $r\ne0$ by inverse discrete Fourier transform. The converse was proved in XVI.2.\end{proof}
\statement{Theorem}{XVI.8}{Finite-sample certificate for both minimal dimensions} Suppose $N$ independent binary trials with expectation $P(k)$ are available for each $k=0,1,\ldots,2p-2$; independence across different $k$ is unnecessary. Define $\widehat H_p=[\widehat P(i+j)]$ and reject rank below $p$ exactly when $\sigma_{\min}(\widehat H_p)>p/(2d^2)$. If
\begin{equation}N\ge8d^4\log\frac{2(2p-1)}\gamma,\qquad0<\gamma<1,\tag{78}\end{equation}
the false-rejection probability under any rank-below-$p$ null is at most $\gamma$, as is the false-acceptance probability for (69). In particular, this rejects all homogeneous classical models with $m<p$ and all quantum channels with $m<d$ levels at the stated confidence. It certifies a memory lower bound, not the absolute quantum nature of the process.
\begin{proof}For $\varepsilon=1/(4d^2)$, Hoeffding and a union bound give
\[\Pr\left\{\max_{0\le k\le2p-2}|\widehat P(k)-P(k)|>\varepsilon\right\}\le2(2p-1)e^{-2N\varepsilon^2}\le\gamma.\]
On its complement, $\|\widehat H_p-H_p\|_2\le p\varepsilon=p/(4d^2)$. Under the null, Weyl gives $\sigma_{\min}(\widehat H_p)\le p/(4d^2)$. Under (69), XVI.3 gives $\sigma_{\min}(\widehat H_p)\ge3p/(4d^2)$. The threshold strictly separates them. For $m\le d-1$, the quantum bound $\rank H\le m^2<p$ was proved above.\end{proof}
\statement{Proposition}{XVI.9}{A reversibility certificate for equilibrium classical autocorrelations} If $T$ is a finite stationary-chain stochastic matrix reversible under detailed balance with invariant distribution $\pi$, and $f$ is real on states, then $C(k)=\mathbb E_\pi[f(X_0)f(X_k)]$ satisfies
\begin{equation}C(k)\le C(2)\qquad\text{for every integer }k\ge2.\tag{79}\end{equation}
The Singer signal has $P_q(p)-P_q(2)=p/d^2>0$, so cannot be an autocorrelation of any such reversible equilibrium chain, regardless of dimension. If both estimates have absolute error at most $\delta<p/(2d^2)$, the difference stays positive. This obstruction concerns only that autocorrelation class, not all classical Markov effects $\alpha M^n\beta$.
\begin{proof}Detailed balance makes $T$ self-adjoint in $L^2(\pi)$ with real spectrum in $[-1,1]$. The spectral theorem gives $C(k)=\int_{-1}^1\lambda^k\,\dd\mu_f(\lambda)$ for a positive measure. Since $\lambda^k\le\lambda^2$ for $k\ge2$, integration proves (79). As $p\ge7$, (69) gives $P_q(2)=(d-1)/d^2$ and $P_q(p)=1$. Errors can reduce their difference by at most $2\delta$.\end{proof}
\statement{Remark}{XVI.10}{Interpretation, limits, and known antecedents} The Singer family is a proven mathematical construction, not a claim about real neutrino masses, new dimensions, fundamental constants, or cosmic topology. Its period $p$ is a model integer, not the size of the universe. Rank $p$ does not mean that every rank-$p$ system requires $p$ classical states; positive rank and minimal positive realization differ in general. With unbounded classical state count, finite-sample rejection is impossible by XIV.5; XVI.8 bounds only the explicitly restricted class $m<p$. Quantum theory permits a fixed channel of dimension $d$, but applying this to a physical system requires laboratory calibration of channel and measurement. Classical foundations include Singer (1938) and positive realization theory; direct studies of quantum--classical memory costs include Budroni and collaborators and Zonnios--Boyd--Binder (2024). Historical priority of the combined formulas has not been checked.
\begin{center}\begin{tabular}{rrrrl}\toprule
$q$&$d$&$p$&$P(n\not\equiv0\pmod p)$&$D$ (verified example)\\\midrule
2&3&7&$2/9$&$\{0,1,3\}$\\3&4&13&$3/16$&$\{0,1,3,9\}$\\4&5&21&$4/25$&$\{0,1,4,14,16\}$\\5&6&31&$5/36$&$\{0,1,3,8,12,18\}$\\\bottomrule
\end{tabular}\end{center}
The table was checked with exact arithmetic in \texttt{checks/verify\_singer\_memory.py}: all $p-1$ oriented differences, discrete frequencies, eigenvalues of $H_p^2$, the determinant sign, and the final rational case $p=7$ were verified. The general proof uses known Singer geometry and is not limited to these four computed examples.
\setcounter{equation}{79}


% 03_reconstruction.tex
\chapter{The Phase Obstruction to Inverse Reconstruction and Bispectral Recovery}
Single-probe memory theory now encounters a second, logically independent limit: invariance of observations at every time does not imply identifiability of internal spectral labels. The mathematical ingredients are cyclic Singer difference sets and the known completeness theorem for the bispectrum when the Fourier transform has no zeros (see Kakarala and earlier work on triple correlations). New to this manuscript is the exact combination of this boundary with the quadratic dimensional separation proved above, sparse one-row bispectral reconstruction, and an explicit stability certificate. Independent worldwide priority of this combination has not been established.

Let $G=\Z/p\Z$, $p=q^2+q+1$, let $D\subset G$ be a Singer set of size $d=q+1$, and $x_D(t)=\one_D(t)$. Define $F_D(k)=\sum_{t\in G}x_D(t)\zeta_p^{-kt}$,
\begin{align}C_D(r)&=\sum_{t\in G}x_D(t)x_D(t+r),\tag{80}\\
T_D(r,s)&=\sum_{t\in G}x_D(t)x_D(t+r)x_D(t+s),\tag{81}\end{align}
and $B_D(k,\ell)=F_D(k)F_D(\ell)\overline{F_D(k+\ell)}$.
\statement{Theorem}{XVII.1}{Strict ambiguity of the return signal and its resolution by a third-order correlator} Every $D$ from XVI.2 satisfies
\begin{align}C_D(0)&=d,\qquad C_D(r)=1\quad(r\ne0),\tag{82}\\
|F_D(0)|&=d,\qquad|F_D(k)|=\sqrt q\quad(k\ne0),\tag{83}\\
\sum_{r,s\in G}T_D(r,s)\zeta_p^{-kr-\ell s}&=B_D(k,\ell).\tag{84}\end{align}
Thus all difference sets with the same $(q,d,p)$ yield the same return signal $P_q(n)$ at every $n\in\Z$, regardless of the relative positions of spectral labels. Yet the complete $T_D$ determines $D$ uniquely up to cyclic translation. For $p=7$, $D=\{0,1,3\}$ and $D'=\{0,4,6\}=-D$ have equal returns for all $n$, but $D'$ is not a translate of $D$, and
\begin{equation}T_D(1,3)=1,\qquad T_{D'}(1,3)=0.\tag{85}\end{equation}
Repeating the scalar return measurement without changing the observables cannot detect this difference.
\begin{proof}There are $d$ ordered pairs with difference zero and one with every nonzero difference. This is $C_D(r)$, since a variable change merely reverses the index sign. Its discrete Fourier transform is $|F_D|^2$, proving (83). For indicators, $T_D(r,0)=C_D(r)$, so the full third-order table also determines the power spectrum. Substituting $u=t+r$, $v=t+s$ in the double sum proves (84).

For real arrays $x,y$ with nonzero Fourier coefficients, equal Fourier moduli, and equal bispectra, set $u_k=F_y(k)/F_x(k)\in\mathbb T$. Since $F_x(0)=F_y(0)=d>0$, $u_0=1$. Bispectral equality gives $u_{k+\ell}=u_ku_\ell$. Every homomorphism $\Z/p\Z\to\mathbb T$ is $u_k=\zeta_p^{tk}$ for some $t\in G$. Fourier inversion gives $y(s)=x(s+t)$, up to the translation-sign convention. Here both arrays are equal-size indicators, so $F(0)=d$.

For the stated pair, inspection of the seven translates shows that $D'$ is not a translate of $D$. Only $t=0$ gives $(t,t+1,t+3)\subset D$, and no such triple lies in $D'$. Equation (81) gives (85).\end{proof}
\statement{Theorem}{XVII.2}{A one-dimensional bispectral slice suffices; explicit stability} Let $x=x_D$ and let $y\in\R^p$ satisfy $F_y(0)=F_x(0)=d$ and $|F_y(k)|=|F_x(k)|$ for every $k$. The equalities
\begin{equation}B_y(k,1)=B_x(k,1)\qquad(1\le k\le p-2)\tag{86}\end{equation}
suffice for $y(s)=x(s+t)$ for some $t\in G$. Thus, given the two-point spectrum, at most $p-2$ complex bispectral values replace the full $p\times p$ third-order table.

More quantitatively, suppose $|B_y(k,1)-B_x(k,1)|\le\delta$ for $1\le k\le p-2$, $0\le\delta\le2q^{3/2}$. Put $\varepsilon=2\arcsin(\delta/(2q^{3/2}))$. Then some $t\in G$ satisfies
\begin{equation}\|y-x(\cdot+t)\|_{\ell^2(G)}\le2\sqrt d\,(p-2)\varepsilon.\tag{87}\end{equation}
If $y$ is also a size-$d$ indicator and the right-hand side is strictly below $\sqrt2$, it is exactly a cyclic translate of $x$.
\begin{proof}Set $u_k=F_y(k)/F_x(k)\in\mathbb T$. Equal moduli give $|B_y(k,1)|=|B_x(k,1)|=q^{3/2}$. Choose $e_k\in[-\varepsilon,\varepsilon]$ with $u_ku_1\overline u_{k+1}=e^{\ii e_k}$. Recursion yields
\[u_k=u_1^k\exp\left(-\ii\sum_{j=1}^{k-1}e_j\right),\qquad1\le k\le p-1.\]
Reality gives $u_{p-1}=\overline u_1$, hence $u_1^p=e^{\ii S}$, $S=\sum_{j=1}^{p-2}e_j$. Thus $u_1=\zeta_p^te^{\ii S/p}$ for some $t$, and
\[|u_k-\zeta_p^{tk}|\le\left|\frac{kS}p-\sum_{j=1}^{k-1}e_j\right|\le\left[\frac{k(p-2)}p+(k-1)\right]\varepsilon<2(p-2)\varepsilon.\]
Parseval and $\|x\|_2=\sqrt d$ give (87); translation orientation does not affect the minimum. For $\delta=0$ this is exact equivalence. Distinct equal-size indicators differ at at least two positions, hence their $\ell^2$ distance is at least $\sqrt2$, proving the final assertion.\end{proof}
\statement{Proposition}{XVII.3}{Exact sample cost of an additional third invariant in the test case} For $p=7$, $D=\{0,1,3\}$, $D'=-D$, assume an additional calibrated oracle independently samples uniform $t\in G$ and returns $Z=x_D(t)x_D(t+1)x_D(t+3)\in\{0,1\}$, or the corresponding result for $D'$. Under $D$, $Z\sim\operatorname{Bernoulli}(1/7)$; under $D'$ it is identically zero. Accepting $D$ if at least one of $N$ outcomes is 1, and $D'$ otherwise, has zero error under $D'$ and miss probability $(6/7)^N$ under $D$. For one-sided error at most $\gamma\in(0,1)$, $N\ge\lceil\log\gamma/\log(6/7)\rceil$ suffices.
\begin{proof}Exactly one $t$ yields a product of 1 for $D$, and none for $D'$. Independent repetitions give the all-zero probability.\end{proof}
\statement{Remark}{XVII.4}{Operational boundary and antecedents} This oracle does not follow automatically from the ordinary binary experiment $P_q(n)=|\langle\psi,U_D^n\psi\rangle|^2$. Implementing it physically requires additional controlled spectral operations or a separate three-point interferometric channel. Without that access, XVII.1 establishes mathematical distinguishability, not a laboratory demonstration. The pair $D,-D$ is distinct at fixed spectral-label orientation; if reflection $t\mapsto-t$ is an allowed gauge equivalence, an additional reference is needed to regard it as physically distinct. Bispectral completeness for nonvanishing Fourier transforms was known earlier, including Kakarala's triple-correlation work (1992, 2012) and astronomical-interferometry phase reconstruction. This manuscript supplies refined formulas for the special family and their conditional integration with Hankel memory; their global historical priority is not proved.
\begin{center}\begin{tabular}{rrrrr}\toprule
$p$&$d$&$|F(k\ne0)|$&$|B(k,1)|$&$\#\{(r,s):T_D(r,s)\ne T_{-D}(r,s)\}$\\\midrule
7&3&$\sqrt2$&$2\sqrt2$&12\\13&4&$\sqrt3$&$3\sqrt3$&48\\21&5&2&8&120\\31&6&$\sqrt5$&$5\sqrt5$&240\\\bottomrule
\end{tabular}\end{center}
For all rows, integer computations checked difference sets, the full cyclic correlation, the identity between the bispectrum and the double Fourier transform of the triple correlator, and the difference between $D$ and $-D$. For $p=7$, all 14 size-3 Singer sets (two translation orbits) were exhaustively checked; see \texttt{checks/verify\_bispectral\_identifiability.py}. These computations verify examples and do not replace the theorem for every prime power.
\chapter{An Interferometric Phase Witness: A Finite Protocol from Mathematical Ambiguity to a Measurable Effect}
The preceding chapter establishes nonidentifiability of spectral labels from one return-probability function. Here an additional experiment with a phase-reference channel distinguishes two specific quantum processes using two temporal depths while remaining invariant under the choice of energy zero. This is new within the present construction, not a new natural law or a report of a laboratory experiment. The interferometric Hadamard test and classical bispectral phase reconstruction are known; the purpose is their exact integration with the minimal Singer pair, confidence bounds, and a specified gauge group.
\statement{Definition}{XVIII.1}{Physically specified data and calibration} Let $\omega=e^{2\pi\ii/7}$, $D=\{0,1,3\}$, and $\mathcal H=\C^3$ with orthonormal basis $|0\rangle,|1\rangle,|3\rangle$. Define
\[|\psi\rangle=\frac{|0\rangle+|1\rangle+|3\rangle}{\sqrt3},\quad U_\pm=\sum_{a\in D}\omega^{\pm a}|a\rangle\langle a|,\]
\[A_\pm(n)=\langle\psi|U_\pm^n|\psi\rangle=\frac13\sum_{a\in D}\omega^{\pm an},\qquad R_\pm(n)=|A_\pm(n)|^2.\]
The available phase-reference experiment prepares $|\psi\rangle$, executes controlled-$U_\pm^n$ for $n=1,2$, and reads both ancilla-qubit quadratures. The common energy-zero shift $U\mapsto e^{\ii\chi}U$ is a gauge transformation. Its cancellation in a multistep protocol requires $U^2$ to mean two physically consistent evolutions, not an independently arbitrarily phased gate.
\statement{Theorem}{XVIII.2}{A two-time gauge-invariant orientation witness} All return probabilities of $U_+,U_-$ coincide exactly:
\begin{equation}R_+(n)=R_-(n)=\begin{cases}1,&7\mid n,\\2/9,&7\nmid n.\end{cases}\tag{88}\end{equation}
However, the complex phase closure
\begin{equation}B_\pm=A_\pm(1)^2\overline{A_\pm(2)},\qquad W_\pm=\Im B_\pm\tag{89}\end{equation}
has exact values
\begin{equation}B_\pm=\frac{-1\pm\ii\sqrt7}{27},\qquad W_\pm=\pm\frac{\sqrt7}{27}.\tag{90}\end{equation}
It is preserved by $U_\pm\mapsto e^{\ii\chi}U_\pm$, including cyclic shifts $D\mapsto D+t$. Under $U\mapsto U^\dagger$, $B\mapsto\overline B$. Distinguishability requires calibrated time orientation; if time reversal identifies systems, the sign of $W$ is not a physical invariant of the equivalence class.
\begin{proof}$U_-=U_+^\dagger$ implies $A_-(n)=\overline{A_+(n)}$, hence equal returns. The set $D$ is a cyclic $(7,3,1)$ difference set, so $|\sum_{a\in D}\omega^{an}|^2=2$ when $7\nmid n$, and 9 otherwise. Put $z=\omega$. Using $z^7=1$ and reducing modulo $1+z+\cdots+z^6=0$ gives
\[27B_+=(1+z+z^3)^2(1+z^{-2}+z^{-6})=2(z+z^2+z^4).\]
For $S=z+z^2+z^4$, $S+\overline S=-1$, $S\overline S=2$, and $\Im S>0$, hence $2S=-1+\ii\sqrt7$. The other sign is its conjugate. A common phase multiplies $A(n)$ by $e^{\ii n\chi}$ and contributes $e^{2\ii\chi}e^{-2\ii\chi}=1$ to $B$.\end{proof}
\statement{Proposition}{XVIII.3}{Two depths are necessary among scalar phase-reference probes} If only one complex return $A(n)$ at a fixed $n\ne0$ is accessible and $U\mapsto e^{\ii\chi}U$ is an unknown gauge, no gauge-invariant function of that response alone distinguishes the pair in XVIII.2. Depths $n=1,2$ suffice through (89).
\begin{proof}As $\chi$ varies, $A(n)\mapsto e^{\ii n\chi}A(n)$ traces its entire constant-modulus circle. Its orbit therefore depends only on $|A(n)|$, which is equal for the two systems. Sufficiency follows from XVIII.2. This minimality concerns only the stated scalar probes, not arbitrary additional observables, preparations, or multitime instruments.\end{proof}
\statement{Theorem}{XVIII.4}{Independent finite-sample interferometric certificate} For each $n\in\{1,2\}$, suppose two independent series of ancilla $\{\pm1\}$ measurements have expectations $\Re A(n)$ and $\Im A(n)$, with $N$ trials per series (total $4N$). Construct $\widehat A(n)$ from their means and project onto the unit complex disk if necessary. Set $\widehat W=\Im[\widehat A(1)^2\overline{\widehat A(2)}]$ and $s=\sqrt7/27$. For $\gamma\in(0,1)$,
\begin{equation}N\ge\left\lceil\frac{144}{s^2}\log\frac8\gamma\right\rceil=\left\lceil\frac{104976}{7}\log\frac8\gamma\right\rceil\tag{91}\end{equation}
guarantees $\Pr(|\widehat W-W_\pm|\le s/2)\ge1-\gamma$. Thus $\operatorname{sign}(\widehat W)$ correctly identifies orientation with this confidence under the stated assumptions. At $\gamma=0.05$, $N=76111$ per quadrature, or 304444 measurements total, suffices; this bound is far from optimal.
\begin{proof}Hoeffding bounds error exceeding $t$ in one $\pm1$ mean by $2e^{-Nt^2/2}$, hence in any of four means by $8e^{-Nt^2/2}$. On the complementary event, each complex error is at most $\sqrt2t$. Disk projection cannot increase distance to $A(n)$, since $|A(n)|\le1$. Telescoping the triple products with factors of modulus at most one gives
\[|\widehat W-W|\le|\widehat B-B|\le2|\widehat A(1)-A(1)|+|\widehat A(2)-A(2)|\le3\sqrt2t.\]
Set $t=s/(6\sqrt2)$ to obtain error at most $s/2$; the union bound gives (91), and $W_\pm=\pm s$ fixes the correct sign.\end{proof}
\statement{Corollary}{XVIII.5}{Certificate with finite systematic error} If quadrature expectation biases combine at each $n$ into $|\widetilde A(n)-A(n)|\le\kappa$, where $0\le\kappa<s/6$, the same $s/2$ certificate holds when
\[N\ge\left\lceil\frac{144}{(s-6\kappa)^2}\log\frac8\gamma\right\rceil.\]
Physical consistency of $U$ and $U^2$ is also required: unknown phase drift must be included in $\kappa$, not hidden in a gauge transformation.
\begin{proof}Each total complex error is at most $\kappa+\sqrt2t$. Choose $t=(s-6\kappa)/(6\sqrt2)>0$, so $3(\kappa+\sqrt2t)=s/2$, and apply the previous theorem.\end{proof}
\note{Reproducible algorithm.} Prepare $|\psi\rangle$. At each $n=1,2$, run separate Hadamard tests estimating the real and imaginary parts of $\langle\psi|U^n|\psi\rangle$ through ancilla $X$ and $Y$ measurements. Compute $\widehat B$ and the sign of $\widehat W$. Both devices must share a phase convention, controlled evolution powers, and unchanged time orientation. Symbolic checks use \texttt{verify\_phase\_closure.py}; synthetic Hadamard-distribution tests use \texttt{simulate\_phase\_protocol.py}. Synthetic samples do not replace hardware data.
\note{Separation from neutrino physics.} T2K/NOvA data measure transition probabilities and constrain PMNS parameters; they do not measure the closure $A(1)^2\overline{A(2)}$ of the artificial three-level $U_D$. A controlled three-level processor could confirm this quantum model, but that alone would not discover a particle, prove physical $S^3$ cosmology, or establish a new natural law.
\note{Historical priority.} Bispectral phase reconstruction predates this study, including triple-correlation work from 1984 and 1992 and Holambe--Ray--Basu (1996). Hadamard measurement of complex overlaps and gauge-invariant third-order phase products are also known. Thus (90) is an exact verifiable integration of known tools with a specific Singer pair, not a claimed world-first law. A global-priority claim would require broader bibliographic research, independent specialist assessments, and a publicly dated preprint.
\chapter{Conditional Phase Closure: Complete Reconstruction from One Row of Triple Correlation}
We strengthen the preceding bispectral result in a strictly limited geometric setting. For a perfect difference set, one row with a fixed nonzero offset suffices instead of the full two-dimensional third-order correlator. Fourth order also obeys an exact rank law extending from $\lambda=1$ to arbitrary cyclic $(v,k,\lambda)$ difference sets. These are mathematical proofs, not new experimental claims about neutrinos or cosmology.

Let $G=\Z/v\Z$, $D\subset G$, $|D|=k$, and $x(t)=\one_D(t)$. Suppose $D$ is a cyclic $(v,k,\lambda)$ difference set, $0<\lambda<k<v$, so $|D\cap(D-r)|=\lambda$ for every $r\ne0$. Fix such an $r$ and define
\begin{align}A_r&=\{t\in G:x(t)x(t+r)=1\},\qquad|A_r|=\lambda,\tag{92}\\
T_r(s)&=\sum_{t\in G}x(t)x(t+r)x(t+s),\tag{93}\\
K_r(s,u)&=\sum_{t\in G}x(t)x(t+r)x(t+s)x(t+u).\tag{94}\end{align}
\statement{Theorem}{XIX.1}{Exact spectrum of the anchored fourth-order correlator} For every fixed $r\ne0$, $K_r\in\R^{v\times v}$ is symmetric positive semidefinite, has rank $\lambda$, and nonzero spectrum
\begin{equation}\spec_+(K_r)=\{k+\lambda(\lambda-1)\text{ once},\quad k-\lambda\text{ with multiplicity }\lambda-1\}.\tag{95}\end{equation}
For $\lambda=1$, the second family is absent. The spectrum is independent of the nonzero offset and reveals $\lambda$ without using the coordinates of $D$.
\begin{proof}For each $t\in A_r$, take $v_t\in\R^v$ with $(v_t)_s=x(t+s)$, and assemble their $\lambda$ columns into $V_r$. By definition $K_r=V_rV_r^{\mathsf T}$, while $\langle v_t,v_{t'}\rangle$ equals $k$ if $t=t'$ and $\lambda$ otherwise. Hence $V_r^{\mathsf T}V_r=(k-\lambda)I_\lambda+\lambda\one_\lambda\one_\lambda^{\mathsf T}$, with the stated eigenvalues. They are positive because $k>\lambda$, so $V_r$ has rank $\lambda$; the nonzero spectra of $V_rV_r^{\mathsf T}$ and $V_r^{\mathsf T}V_r$ coincide.\end{proof}
\statement{Theorem}{XIX.2}{One-row reconstruction and exact factorization of all anchored moments} If $\lambda=1$, in particular for any Singer $(q^2+q+1,q+1,1)$ set, each $r\ne0$ has a unique $t_r$ with $t_r,t_r+r\in D$. Then
\begin{equation}T_r(s)=x(t_r+s)\qquad(s\in G).\tag{96}\end{equation}
Thus $\{s:T_r(s)=1\}=D-t_r$: one row recovers all of $D$ up to translation. For every integer $m\ge1$ and $s_1,\ldots,s_m\in G$,
\begin{equation}\sum_{t\in G}x(t)x(t+r)\prod_{j=1}^m x(t+s_j)=\prod_{j=1}^mT_r(s_j).\tag{97}\end{equation}
The fourth-order matrix is a scaled projector:
\begin{equation}K_r=T_rT_r^{\mathsf T},\qquad K_r^2=kK_r,\qquad\rank K_r=1.\tag{98}\end{equation}
\begin{proof}$A_r=\{t_r\}$, so each sum has exactly one nonzero contribution, the corresponding product of $x(t_r+s_j)$. This proves (96)--(97). The vector $T_r$ has exactly $k$ unit entries, so $\|T_r\|_2^2=k$, yielding (98).\end{proof}
\statement{Corollary}{XIX.3}{Complete recovery of phase closures with additional spectral access} For a Singer set, put $A_D(n)=k^{-1}\sum_{a\in D}e^{2\pi\ii na/v}$. From one row compute
\[\widetilde A_r(n)=\frac1k\sum_{s\in G}T_r(s)e^{2\pi\ii ns/v}=e^{-2\pi\ii nt_r/v}A_D(n).\]
Every integer phase closure is therefore determined without knowing $t_r$:
\begin{equation}A_D(n)A_D(m)\overline{A_D(n+m)}=\widetilde A_r(n)\widetilde A_r(m)\overline{\widetilde A_r(n+m)}.\tag{99}\end{equation}
For the quantum model XVI.2, simultaneously $\rank H_{\mathrm{temporal}}=v=q^2+q+1$ and $\rank K_r=1$ for every $r\ne0$. These ranks concern different operational tables; this is neither a contradiction nor a universal identity for arbitrary quantum processes.
\begin{proof}Insert (96), substitute $a=s+t_r$, and cancel the phase factors in the triple product. The last claim follows from the established temporal Hankel rank and (98).\end{proof}
\statement{Proposition}{XIX.4}{Threshold stability and one-sided sample reconstruction} For $\lambda=1$, if $\max_s|\widehat T_r(s)-T_r(s)|<1/2$, thresholding at $1/2$ exactly recovers $D-t_r$. Alternatively, suppose an extra ideal oracle, for specified $s$, independently samples uniform $t\in G$ and returns $Z_s=x(t)x(t+r)x(t+s)$. Then $\Pr(Z_s=1)=T_r(s)/v$, either 0 or $1/v$. Running
\begin{equation}N=\left\lceil v\log\frac k\gamma\right\rceil\tag{100}\end{equation}
trials per $s\in G$ and declaring $T_r(s)=1$ exactly when at least one 1 occurs recovers the full slice with probability at least $1-\gamma$, using $vN=O(v^2\log(k/\gamma))$ oracle calls.
\begin{proof}Each true coordinate is 0 or 1, so error below $1/2$ preserves its threshold label. Zero coordinates cannot be false positives. At each of the $k$ positive coordinates, the miss probability is $(1-1/v)^N\le e^{-N/v}$; a union bound gives $ke^{-N/v}\le\gamma$.\end{proof}
\statement{Remark}{XIX.5}{Applicability and experimental obstruction} The oracle accesses spectral-label occupancy and correlated shifts $t+r,t+s$. Scalar return probabilities $|A_D(n)|^2$ do not provide this, and standard neutrino detectors do not assume it. Complete reconstruction concerns an expanded observable set. With possible false positives, (100) requires a new error model. At $\lambda>1$, (95) replaces rank-one closure; a sum of $\lambda$ columns cannot be replaced by one. Difference-set theory and bispectral inversion predate this work (Singer 1938; Bartelt et al., 1984; Bendory et al., 2018). The contribution within this manuscript is the explicit joint description of anchored fourth-order rank, one-row recovery, all higher conditional moments, and quantum--classical rank separation. Independent global novelty of the combination is unproved.
\note{Exact test examples.} For $(v,k,\lambda)=(7,3,1)$, $D=\{0,1,3\}$, and $r=1$, $T_r=\one_{\{0,1,3\}}$ and $\spec_+(K_r)=\{3\}$. For $D'=-D=\{0,4,6\}$, the same row has ones at $\{0,1,5\}$, distinguishing the return-indistinguishable processes. For the classical $(11,5,2)$ set $D=\{1,3,4,5,9\}$, every $r\ne0$ gives $\spec_+(K_r)=\{7,3\}$: $\lambda=2$ appears as a second nonzero eigenvalue. The code \texttt{checks/verify\_anchored\_closure.py} also checks $(13,4,1)$, $(21,5,1)$, $(31,6,1)$, and $(7,4,2)$ over the integers; error-free output only verifies these examples.
\chapter{The Critical Tree Boundary: The Same Hankel Rank, Opposite Reconstructibility}
OpenAI Math family 229 studies source reconstruction from deep generations of a symmetric three-state tree process. Its selected Lean formalization confirms the supercritical side, not every negative-side claim in the full preprint. The self-contained counterexample below disproves inference of tree reconstructibility from the rank of a one-point temporal correlation alone; it needs no unverified threshold-theorem cases.
\statement{Theorem}{XX.1}{Hankel rank does not determine reconstruction under branching} Let $b\ge2$ be an integer. Each edge of the full $b$-ary tree carries the same symmetric three-state transition
\[M_\lambda=\lambda I_3+(1-\lambda)\frac{\one\one^\top}{3},\qquad0\le\lambda<1,\]
and the root is uniform on $\{1,2,3\}$. Let $f$ be nonzero and real with $\mathbb E_\pi f=0$, $\mathbb E_\pi f^2=1$. The single-line autocorrelation is $C_\lambda(k)=\lambda^k$ for $k\ge0$, with $0^0=1$, so
\[\rank[C_\lambda(i+j)]_{i,j\ge0}=1\qquad(0\le\lambda<1).\]
At $\lambda=0$, every generation $n\ge1$ is independent of the root, making reconstruction impossible. Conversely, if $\theta=b\lambda^2>1$ and $\mathcal F_n$ contains the entire $n$th generation, then for every $n\ge1$,
\begin{equation}\mathbb E\left[\mathbb E(f(X_\varnothing)\mid\mathcal F_n)^2\right]\ge\frac1{1+\dfrac{b-1}{b(1-\theta^{-1})}}>0.\tag{101}\end{equation}
For $b=2$, $\lambda=3/4$, the bound is $2/11$; for $\lambda=0$, the left side is identically zero for all $n\ge1$, although both Hankel ranks are one.
\begin{proof}$M_\lambda f=\lambda f$ and stationarity give $C_\lambda(k)=\mathbb E_\pi[f(X_0)f(X_k)]=\lambda^k$. At $\lambda=0$, only $C_0(0)$ is nonzero; at $\lambda>0$, the Hankel matrix is the outer product of $(\lambda^i)_{i\ge0}$ with itself. Both ranks are one. At $\lambda=0$, each child is uniform independently of its parent, hence the whole generation is independent of the root.

For $\lambda>0$, put $S_n=\sum_{|v|=n}f(X_v)$, $Y_n=S_n/(b\lambda)^n$, $R=f(X_\varnothing)$. Then $\mathbb E(RY_n)=1$. Two distinct depth-$n$ leaves whose most recent common ancestor has depth $h$ have correlation $\lambda^{2(n-h)}$ by conditional subtree independence. There are $(b-1)b^{2n-h-1}$ ordered pairs of this type and $b^n$ diagonal pairs. Therefore
\[\mathbb E(Y_n^2)=\theta^{-n}+\frac{b-1}b\sum_{h=0}^{n-1}\theta^{-h}\le1+\frac{b-1}{b(1-\theta^{-1})}=:K_b(\theta).\]
As $Y_n$ is $\mathcal F_n$-measurable, Cauchy--Schwarz gives $1=\mathbb E[\mathbb E(R\mid\mathcal F_n)Y_n]\le\sqrt{\mathbb E[\mathbb E(R\mid\mathcal F_n)^2]\mathbb E(Y_n^2)}$, proving (101). For $b=2$, $\lambda=3/4$, $\theta=9/8$ and $K_b=1+(1/2)/(1-8/9)=11/2$. This example does not use family 229's unformalized threshold side.\end{proof}
\statement{Remark}{XX.2}{A necessary new coordinate in the causal-reconstruction passport} For branching processes, the straight-line sequence $C(k)$ and its rank are insufficient. One must at least specify propagation geometry, here $b$, and transfer coefficient $\lambda$. The invariant $b\lambda^2$ is relevant to this symmetric family, not a universal criterion for asymmetric channels or quantum instruments. Equation (101) asserts positive asymptotic $L^2$ information about the root, not successful operation of a specific physical detector.
\chapter{A Conditional Computational Boundary of Causal Diagnostics: Forced Flows}
OpenAI Math family 376 describes an effective construction of smooth forced incompressible Navier--Stokes flows: a pair (finite machine, input) determines a smooth force and a distinguished material trajectory entering a fixed open region exactly when the machine halts. Its Lean description includes selected classes of these encodings. Unlike the unforced global Clay problem, the external force is tailored to the machine; these classes must not be substituted for one another.
\statement{Proposition}{XXI.1}{Conditional undecidability of an exact reachability certificate} Suppose an effective map $\mathcal C$ from pairs ``finite machine $M$, input $w$'' to smooth data for a forced three-dimensional Navier--Stokes system has the equivalence proved in family 376:
\[M\text{ halts on }w\quad\Longleftrightarrow\quad\text{the distinguished trajectory enters the detector}.\]
No algorithm terminating on every description in the image of $\mathcal C$ can always correctly decide whether the distinguished trajectory ever reaches the detector.
\begin{proof}Composing such an algorithm with the computable compiler $\mathcal C$ would give an everywhere-terminating correct halting algorithm for arbitrary finite machines and inputs, contradicting undecidability of the halting problem.\end{proof}
\statement{Remark}{XXI.2}{Reachability undecidability is not loss of regularity} This is a logical consequence of an explicitly accepted catalogue premise, not independent verification of the complete hydrodynamic construction. It concerns future particle entry in a forced flow. It proves no finite-time singularity, irregularity, undecidability of every numerical approximation, or negative answer to the Navier--Stokes Millennium Problem. It separates existence of a smooth solution from universal algorithmic decidability of its unbounded-time trajectory.
\chapter{A Unified Mathematical Boundary of Transferability: OpenAI Math Sources}
Together the theorems define a conditional closure passport $\mathcal C=(A,E,R,L,N,P)$, where $A$ is the algebra of permitted operations and its physically admissible semigroup; $E$ the effects; $R$ the reachable and observable predictive space; $L$ allowed factors/gauges; $N$ the perturbation structure; and $P$ certificates and a verification protocol. This is an architectural definition, not an experimental natural law.
\statement{Theorem}{XXII.1}{Boundary of one-sided correlation reconstruction} Suppose a unitary $U$ on an infinite-dimensional Hilbert space has a vector $\phi$ such that $\{U^k\phi:k\in\Z\}$ is an orthonormal basis (a smooth dynamical construction of this kind is claimed in OpenAI Math family 144). Then $c_k=\langle\phi,U^k\phi\rangle=\one_{\{k=0\}}$, and
\[\rank[c_{i+j}]_{i,j\ge0}=1,\qquad\rank[c_{i+j}]_{i,j\in\Z}=\infty.\]
\begin{proof}The first table has only its $(0,0)$ entry nonzero. The second acts on $\ell^2(\Z)$ as the permutation $x_j\mapsto x_{-j}$ and has infinite rank. These are correlation tables, not the full operational Hankel matrix for every intervention.\end{proof}
Three claim classes must remain distinct: (i) finite-process rank lower bounds; (ii) invariant spectral gaps in Hamiltonians; and (iii) global regularity of nonlinear PDEs. OpenAI Math provides substantive results in each, but transfer between them needs additional maps and estimates. For example, 269 (the Laughlin gap problem) does not establish a mass gap for four-dimensional Yang--Mills theory; 376 concerns forced Navier--Stokes flows and does not prove smoothness for the unforced Clay problem; 350 limits transfer of nodal bounds between smooth metrics on $S^3$; and 365 separates smooth inverse reconstruction from a counterexample for bounded measurable conductivities.
\note{Source-check status.} The README, CONTENTS index, and thematic Lean descriptions in \url{https://github.com/openai/math} were examined for families 003, 095, 144, 229, 269, 287, 350, 362, 365, 374, and 376. Family 229's Lean description covers supercritical reconstruction, not the whole claimed negative threshold result. For 269, a partial gap proof must not be presented as verification of all stability under external potentials; for 365, the described formal scope concerns conductivity nonuniqueness. The original proofs of all these families were not fully reproduced. A register of exact URIs and applicability limits is included in the accompanying archive.
\note{Methodological map of specific transfers (Revision 18).} Three operationally different boundaries cannot be removed by a single spectral norm:
\[\rank H\ \not\Rightarrow\ \text{identifiability of spectral labels},\qquad\rank H\ \not\Rightarrow\ \text{reconstruction under branching}.\]
The first nonimplication is certified by XVII.1 in the explicit Singer family; XX.1 shows that rank alone does not determine reconstructibility. This does not assert independence from every conceivable complete invariant set. Computational reachability in forced PDEs is a further separate level (XXI.1).
\note{Open problems.} General optimal bounds for random stochastic coarse-graining, lower sample complexity of single-probe tests, and full structural classification of hypergraph optima for arbitrary $r$ remain undetermined. Physical transfers require separate checks of experimental covariances, interactions, and stability. No particular spectral or computational example automatically solves a Clay problem.
\chapter{Interdisciplinary Logical Bridges}
\begin{enumerate}
\item \textbf{Statistical and gauge indistinguishability.} If an involution $g$ preserves the distribution of all accessible data and reverses a quantity $J$, its sign is unidentifiable from those data. If $\operatorname{TV}(P_x,P_{gx})\le\xi$, optimal correct-sign probability under equal priors is at most $(1+\xi)/2$. This applies to conditional CP-even channels but establishes no new neutrino physics.
\item \textbf{Hydrodynamics and finite sensors.} For any finite family of linear sensors on an infinite-dimensional space of smooth solenoidal modes, nonzero fields invisible to all sensors exist. This does not automatically extend to infinite-dimensional boundary-measurement operators.
\item \textbf{Quantum spectrum.} A positive gap $\Delta$ alone does not ensure stability under the sum of $N$ small local perturbations: their total operator norm may grow as $N\varepsilon$. Stability requires further restrictions on the perturbation class.
\item \textbf{Geometry and inverse problems.} Coefficient regularity and gauge equivalence must be specified explicitly; agreement on a finite measurement series does not establish uniqueness of an infinite-dimensional geometry.
\item \textbf{Computability.} Agreement of $A$-actions with operational facts in an effective model does not guarantee an algorithm for finite exact closure in all cases. The catalogue's forced Navier--Stokes flows encoding halting limit such universal compilers but do not solve the unforced Clay problem.
\end{enumerate}
\chapter{A Universal System: Necessary, Sufficient, and Unverified Levels}
\begin{center}\small\begin{tabularx}{\textwidth}{@{}p{.18\textwidth}XX@{}}\toprule
Level&Verifiable invariant&Permissible conclusion\\\midrule
Logical&Intervention/effect class, equivalence&What counts as one process\\
Algebraic&$\rank H$, operator algebra, commutator&Minimum linear memory and prohibited simplifications\\
Geometric&Coarse-graining kernel, commuting diagram&Existence of an exact intervention factor\\
Dynamic&$\delta_a,\kappa_a,\Xi(w)$&Cost of approximate closure across scales\\
Statistical&$R(\widehat H),\eta,\gamma$&Frequency-based test of the claimed dimension\\
Physical&Admissible apparatus, locality, positivity, causality&Whether mathematics transfers to experiment\\\bottomrule
\end{tabularx}\end{center}
The combined criterion for this stochastic family is
\begin{equation}\begin{cases}
\rank H_n(a)=n,\\
\rho_n(a)=(1-a)/D_n(a),\\
\delta S_{2n-2}(\kappa)\ge\rho_n(a)&(d<n),\\
(1+\kappa_Q)\delta_P+(1+\kappa_P)\delta_Q\ge1-a&(\overline P\overline Q=\overline Q\overline P),\\
R(\widehat H)>\eta\ \Longrightarrow\ \rank H\ge n\text{ with confidence }1-\gamma.
\end{cases}\tag{102}\end{equation}
This combines rank, dynamic, commutator, and statistical invariants. No row alone derives a fundamental universal law: each has an explicit domain of validity.
\chapter{Measurable Language-Model Compression: A Closed Protocol and Results (Revision 21)}
\note{Experimental status.} For the first time in this manuscript, this section reports an executed experiment on a trained causal Transformer with an unchanged base model and several compressed output layers. It is not a test of an external large-LLM checkpoint: the computing environment could not download \texttt{roneneldan/TinyStories-1M} weights. Instead, a local character decoder was created and the actually trained weights from three runs were saved. Conclusions are limited to a 105232-parameter Transformer, this manuscript's corpus, and the output layer. The original experiments are reproducible from the archive: \texttt{experiments/benchmark.py}, \texttt{experiments/results\_*.json}, \texttt{experiments/model\_*.pt}, logs, and the passport.
\section*{XXV.1 A rank certificate that actually applies to soft outputs}
Fix finite contexts $u_1,\ldots,u_N$ and a $V$-token alphabet. Let $p_{ia}=p(a\mid u_i)>0$ be exact conditional probabilities. Set $P_N=I_N-\one\one^\top/N$, $P_V=I_V-\one\one^\top/V$, and
\begin{equation}C(p)=P_N[\log p_{ia}]_{i,a}P_V.\tag{103}\end{equation}
This is a doubly centered matrix of log conditional probabilities, not the joint Hankel matrix $\Pr(uv)$; the distinction matters for LLMs.
\statement{Theorem}{XXV.1}{Lower bound on unavoidable rank error in softmax outputs} Suppose any candidate on the same contexts outputs
\[q(a\mid u_i)=\operatorname{softmax}_a(z_iB^\top+b),\qquad B\in\R^{V\times d},\quad\rank B\le r,\]
with arbitrary $z_i$ and arbitrary retuning of preceding layers. Then $\rank C(q)\le r$ and
\begin{equation}\|C(p)-C(q)\|_{\max}\ge\frac{\sigma_{r+1}(C(p))}{\sqrt{NV}}.\tag{104}\end{equation}
If an estimated $\widehat C$ satisfies $\|\widehat C-C(p)\|_2\le\eta$, replace the numerator validly by $(\sigma_{r+1}(\widehat C)-\eta)_+$. Without such assumptions, noise certification is impossible.
\begin{proof}The identity $\log q_{ia}=z_iB_a^\top+b_a-\log\sum_j\exp(z_iB_j^\top+b_j)$ contains a row-constant last term removed by right multiplication by $P_V$, while the $i$-constant bias is removed by left multiplication by $P_N$. Thus $C(q)=P_NZB^\top P_V$ has rank at most $r$. Eckart--Young gives $\|C(p)-C(q)\|_F\ge\sigma_{r+1}(C(p))$; use $\|A\|_F\le\sqrt{NV}\|A\|_{\max}$. The noisy case follows from Weyl.\end{proof}
\statement{Remark}{XXV.2}{Certificate limitations} Equation (104) concerns centered log-probability error, not directly accuracy, perplexity, or TV. If $p_{ia},q_{ia}\ge m>0$ and rowwise $\operatorname{TV}(p_i,q_i)\le\delta$, additionally $\|C(p)-C(q)\|_{\max}\le4\delta/m$, but $m$ may be extremely small for a large vocabulary. The certificate neither asserts whole-network compressibility nor upper-bounds global Hankel dimension; nonlinear output decoders require a new argument.
\section*{XXV.2 Executed protocol and verifiable data}
The source was \texttt{META\_CAUSAL\_MASTER\_v20.tex}, deterministically stripped of markup and mathematical environments: 102933 characters, alphabet $V=144$. The autoregressive decoder-only Transformer had 2 causal attention blocks, 4 heads, width 64, context 48, and 105232 parameters. The first 78\% of the corpus was used for training and the remaining 22\% for position-independent testing, with a 200-character gap. AdamW performed 240 training steps; fixed pseudorandom generators selected 512 calibration and 512 test contexts. Three independent runs used seeds 20261012, 20261013, and 20261014. This is a small-model local demonstration, not an experiment on modern multibillion-parameter LLM properties.

Only the final matrix $W\in\R^{144\times64}$ was compressed. Comparisons were (A) weight truncated SVD; (B) activation-covariance-whitened SVD; (C) approximate diagonal output weighting by mean probabilities (Fisher-inspired, not the exact Fisher matrix); and (D) 90 equal low-rank KL-distillation steps from different initializations. All allowed the same free-bias correction. Direct SVD of the conditional-probability table was also tested; it was substantially worse and is not presented as an improvement.

Each rank-$r$ factored layer stores $r(64+144)+144$ real numbers and requires $r(64+144)$ multiplications per token, excluding bias and the rest of the network. The original layer stores $64\cdot144+144=9360$ numbers and requires 9216 multiplications. At $r=16$, layer memory falls by 62.9\%, but whole-model memory falls by only about 5.6\% with other parameters unchanged.
\note{Result.} Probability weighting outperformed raw weight SVD at all five tested ranks and all three seeds, but not always the established activation-SVD baseline. At $r=16$, KL-distilled ordinary SVD was best in all three reported runs. Thus no advantage over a strong existing method with retraining was proved, and the result cannot be interpreted as a new state-of-the-art LLM compression method.
\begin{table}[htbp]\centering\small
\caption{Test $D_{\mathrm{KL}}(p_{\mathrm{teacher}}\|p_{\mathrm{student}})$ at $r=16$; lower is better. Separate trainings, equal operation counts in the compressed final layer.}
\begin{tabularx}{\textwidth}{r*{4}{>{\raggedleft\arraybackslash}X}}\toprule
Seed&Weight SVD&Activation SVD&Probability-weighted SVD&Weight SVD + 90 KL steps\\\midrule
20261012&0.07194&0.03056&0.02577&0.01486\\
20261013&0.08332&0.04238&0.04464&0.02263\\
20261014&0.07592&0.04376&0.04082&0.02109\\\bottomrule
\end{tabularx}\end{table}
For the first run at $r=16$, $\sigma_{17}(C(p))\approx18.0135$. Equation (104) certifies unavoidable centered log-probability maximum error of at least 0.06634 for every rank-16 linear softmax output on these $512\times144$ tests. This lower bound does not promise that any particular SVD attains it.
\note{Computation time.} On a single-thread CPU with batch 2048, measured median time for the final layer alone was 437.45 versus 216.85 microseconds for rank-16 factorization. At batch 16 it was 7.82 versus 8.11 microseconds, a slowdown from overhead. With four threads and batch 2048 it was 210.43 versus 155.96 microseconds. These are not end-to-end model speedups; fitting and calibration also have nonzero cost.
\section*{XXV.3 Antecedents and transfer limits}
Activation-calibrated and probability-weighted algorithms are known: Fisher-Weighted SVD (Hsu et al., 2022, \url{https://arxiv.org/abs/2207.00112}), ASVD (Yuan et al., 2023, \url{https://arxiv.org/abs/2312.05821}), SVD-LLM (Wang et al., ICLR 2025, \url{https://arxiv.org/abs/2403.07378}), and LEMS/KFAC-SVD (ICML 2026, \url{https://github.com/lems-svd/lems}). The implementations tested here are not equivalent to their full methods, so this comparison does not prove superiority over them.

Families 140 (memory--sample tradeoff), 229 (causal reconstruction limits), 095 (positive realization limits), and 376 (algorithmic prediction limits) were selected from \url{https://github.com/openai/math}. They motivate certificates mathematically rather than supply ready-made LLM compression procedures. Family 140 uses another statistical setting; 229 formalizes only a selected direction in Lean; 376 concerns specially forced Navier--Stokes dynamics and does not solve the general Clay problem. See documents 140, 229, and 376 for scope.
\note{Direct execution on a public LLM.} The archive includes the standalone \texttt{experiments/run\_pretrained\_llm.py} for a real independently supplied checkpoint and corpus, with \url{https://huggingface.co/roneneldan/TinyStories-1M} specified as a candidate. It was not run in this environment because weights and the \texttt{transformers} library were unavailable. A publication-level comparison needs at least an independent large checkpoint, multiple corpora, matched end-to-end latency, ASVD/SVD-LLM/LEMS baselines, bootstrap intervals, and independent replication. Until then, global superiority and a link to Millennium Problem solutions are unestablished.
\note{Reproducibility of earlier numerical checks.} \texttt{checks/verify\_master.py} tests $H_n(a)$, explicit inverses, determinants, commutators, and finite stochastic coarse-grainings. For $2\le n\le9$ and four $a$ values, 32 rational checks and eight control reductions were performed; these do not prove the general case. Worldwide priority of the combination is unestablished, and physical interpretation requires independent measurement models.
\setcounter{equation}{104}


% 04_certificates.tex
\chapter{Forty Verifiable Consequences for Predictive Compression and Invariant Reconstruction (Revision 22)}
\note{Scientific qualification.} These forty exact statements, consequences, and computable certificates are assembled here for the first time within this manuscript. They are not forty established worldwide discoveries or forty solved open problems. Most elementary tools---SVD, Frobenius optimality, Weyl, Hoeffding, Pinsker, automaton factorization, and information-channel theory---are classical. New to this program is their composition with precisely specified passports, particularly the two-sided softmax rank certificate. Independent historical assessment of each consequence is incomplete. Catalogue assertions are deliberately not declared proved when only fragments of their Lean code have been checked.
\section*{XXVI.1 A unified passport and the central two-sided theorem}
Let $N,V\ge2$, $p_{ia}>0$, $\sum_a p_{ia}=1$, $L_{ia}=\log p_{ia}$, $P_N=I_N-\one_N\one_N^\top/N$, $P_V=I_V-\one_V\one_V^\top/V$, and $C=P_NLP_V$. Competitors $\mathcal Q_r$ are all panel models $q_{ia}=\operatorname{softmax}_a(z_iB_a^\top+b_a)$ with $\rank B\le r$ and arbitrary free $z_i,b_a$, sharing the finite vocabulary and the same $N$ contexts. This permits independent features at the test contexts; it is not automatically implementable by one parametrically compact Transformer on new contexts.

For $r<\min(N-1,V-1)$, put $s_r=(\sum_{j>r}\sigma_j(C)^2)^{1/2}$, $\tau_r=\sigma_{r+1}(C)$, and $\varepsilon_r=s_r$.
\statement{Theorem}{XXVI.1}{Exact oracle interval for softmax compression on a finite panel} Under this passport,
\begin{equation}\inf_{q\in\mathcal Q_r}\|C(p)-C(q)\|_F=s_r,\qquad\frac{s_r}{\sqrt{NV}}\le\inf_{q\in\mathcal Q_r}\|C(p)-C(q)\|_{\max}\le s_r.\tag{105}\end{equation}
An explicit $q_r\in\mathcal Q_r$ also satisfies
\begin{equation}\max_i\operatorname{TV}(p_i,q_{r,i})\le\tanh(s_r),\qquad\max_iD_{\mathrm{KL}}(p_i\|q_{r,i})\le2s_r.\tag{106}\end{equation}
The max norm means largest absolute entry, not an operator norm; all right-hand sides vanish for $r\ge\rank C$.
\begin{proof}Set $c_a=N^{-1}\sum_iL_{ia}$. Double centering gives $\operatorname{softmax}(C_{i\cdot}+c)=p_{i\cdot}$. Let $C_r$ be the truncated SVD of $C$. It is doubly centered because its nonzero singular vectors are orthogonal to the all-ones directions. Set $q_{r,i}=\operatorname{softmax}(C_{r,i\cdot}+c)$. Then $C(q_r)=C_r$, realizable by a rank-at-most-$r$ linear output with free contextual features. Eckart--Young--Mirsky gives the exact Frobenius minimum because every $C(q)$ has rank at most $r$. The inequalities $\|A\|_F\le\sqrt{NV}\|A\|_{\max}\le\sqrt{NV}\|A\|_F$ give the interval. Since $\|C-C_r\|_{\max}\le s_r$, row logits differ by at most $s_r$. After softmax normalization, log probability ratios lie in $[-2s_r,2s_r]$. The bounded-density-ratio inequality gives TV at most $\tanh s_r$, and $D_{\mathrm{KL}}(p\|q)=\mathbb E_p\log(p/q)$ gives $2s_r$.\end{proof}
\statement{Remark}{XXVI.2}{Scope} This does not provide cheap realizable computation of $z_i$ for arbitrary new prompts. The quantity $s_r$ is an exact oracle panel limit in a chosen norm, not the limit of architecture-constrained end-to-end compression or a promise of lower perplexity.
\section*{XXVI.2 Register of 40 mathematical certificates}
\subsection*{I. Log-probability invariants (01--08)}
\begin{enumerate}
\item \textbf{Algebraic invariance under double centering.} For any $a\in\R^N$, $b\in\R^V$, $P_N(L+a\one_V^\top+\one_Nb^\top)P_V=C$. Each added term is annihilated by a projector. A column shift changes $p$ and is not a physical softmax gauge; it is discarded only because a separate vocabulary bias is allowed.
\item \textbf{Complete gauge separation.} $L=C+u\one_V^\top+\one_Nc^\top$, with $c_a=N^{-1}\sum_iL_{ia}$ and suitable $u_i$. Decompose with respect to $P_N,P_V$, absorbing the overall constant into $u$.
\item \textbf{Exact minimal panel output dimension.} With free context features and vocabulary bias, exact reproduction of $p$ requires minimal rank $\rank C$. Necessity follows from centered softmax; sufficiency from $C=ZB^\top$ by SVD and addition of $c$.
\item \textbf{Exact Frobenius compression cost.} $\inf_{q\in\mathcal Q_r}\|C(p)-C(q)\|_F=s_r$, by the $q_r$ construction and Eckart--Young--Mirsky.
\item \textbf{Unavoidable maximum error.} Every $q\in\mathcal Q_r$ has $\|C(p)-C(q)\|_{\max}\ge s_r/\sqrt{NV}$, because its Frobenius error is at least $s_r$.
\item \textbf{Attainable maximum error.} The stated truncated-SVD realization has $\|C(p)-C(q_r)\|_{\max}\le s_r$, since an entry's modulus is bounded by the Frobenius norm.
\item \textbf{Upper TV certificate.} For the same $q_r$, $\max_i\operatorname{TV}(p_i,q_{r,i})\le\tanh s_r$, since normalized probability ratios lie between $e^{-2s_r}$ and $e^{2s_r}$.
\item \textbf{Upper KL certificate.} $\max_iD_{\mathrm{KL}}(p_i\|q_{r,i})\le2s_r$, by the log-ratio bound after softmax normalization.
\end{enumerate}
\subsection*{II. Noise, observability, and confidence (09--16)}
\begin{enumerate}[start=9]
\item \textbf{Entrywise to spectral noise.} If $\|\widehat C-C\|_{\max}\le\epsilon$, then $\|\widehat C-C\|_2\le\sqrt{NV}\epsilon$, using $\|A\|_2\le\|A\|_F\le\sqrt{NV}\|A\|_{\max}$.
\item \textbf{Certified rank lower bound.} If $\|\widehat C-C\|_2\le\eta$ and $\sigma_j(\widehat C)>\eta$, then $\rank C\ge j$ by Weyl.
\item \textbf{Frobenius-cost lower bound from a noisy table.} $\operatorname{dist}_F(C,\operatorname{Rank}\le r)\ge(s_r(\widehat C)-\|\widehat C-C\|_F)_+$, by the triangle inequality and infimum over low-rank matrices.
\item \textbf{TV error to logarithms.} If $p_{ia},q_{ia}\ge m>0$ and $\delta=\max_i\operatorname{TV}(p_i,q_i)$, then $\|C(p)-C(q)\|_{\max}\le4\delta/m$. Use $|p_{ia}-q_{ia}|\le\operatorname{TV}(p_i,q_i)$, $|\log p-\log q|\le|p-q|/m$, and factor at most four from double centering.
\item \textbf{TV lower bound for every rank-$r$ competitor.} Under item 12 and $q\in\mathcal Q_r$, $\delta\ge ms_r/(4\sqrt{NV})$, combining items 5 and 12.
\item \textbf{Finite-sample confidence window.} With $K$ independent token draws per context, $\max_{ia}|\widehat p_{ia}-p_{ia}|\le t$ with probability at least $1-2NV e^{-2Kt^2}$. If $p_{ia}\ge m>t$, then $\|C(\widehat p)-C(p)\|_2\le4\sqrt{NV}t/(m-t)$. Apply Hoeffding, a cellwise union bound, the mean-value theorem for logarithms, and items 9 and 12 with entrywise error.
\item \textbf{Valid low-rank exclusion test.} On item 14's event, $\sigma_{r+1}(C(\widehat p))>4\sqrt{NV}t/(m-t)$ excludes exact reproduction by every $q\in\mathcal Q_r$, by item 10.
\item \textbf{Simultaneous protection of many ranks.} All item-15 tests for $r=0,\ldots,\min(N,V)-1$ are valid on the same event without an extra rank-count correction: one spectral event controls all singular values through Weyl.
\end{enumerate}
\subsection*{III. Weight matrices, activations, and budgets (17--24)}
\begin{enumerate}[start=17]
\item \textbf{Exact activation-SVD optimum.} For $W\in\R^{V\times d}$, $X\in\R^{d\times K}$, the minimum of $\|(W-\widehat W)X\|_F$ over $\rank\widehat W\le r$ is $s_r(WX)$. The choice $\widehat W=(WX)_rX^\dagger$ has rank at most $r$ and $\widehat WX=(WX)_r$, because $\operatorname{row}(WX)\subseteq\operatorname{row}(X)$. Eckart--Young supplies the converse.
\item \textbf{Covariance optimum.} For $G=XX^\top/K$, $\widehat W_r=(WG^{1/2})_rG^{\dagger/2}$ attains $\min_{\rank\widehat W\le r}\operatorname{tr}((W-\widehat W)G(W-\widehat W)^\top)=s_r(WG^{1/2})^2$, by item 17 and spectral factorization of $G$.
\item \textbf{Transfer under bounded covariance shift.} If $G_{\mathrm{test}}\preceq\kappa G_{\mathrm{cal}}$, then $\operatorname{tr}(EG_{\mathrm{test}}E^\top)\le\kappa\operatorname{tr}(EG_{\mathrm{cal}}E^\top)$ for any $E$, by positive semidefiniteness of $\kappa G_{\mathrm{cal}}-G_{\mathrm{test}}$.
\item \textbf{Failure without shift control.} If $v\in\ker G_{\mathrm{cal}}$ but $v^\top G_{\mathrm{test}}v>0$, then $E=uv^\top$ has zero calibration error and positive test error for every $u\ne0$, by direct quadratic-form evaluation.
\item \textbf{Condition for parameter savings.} Replacing a $V\times d$ layer by $V\times r$ and $r\times d$ factors stores fewer weights exactly when $r(d+V)<dV$. This follows by counting; bias $V$ is the same on both sides.
\item \textbf{Whole-model savings.} If total original parameters are $P$ and all other parameters stay fixed, the saving fraction is $[dV-r(d+V)]/P$ by subtraction and division by $P$.
\item \textbf{Globally optimal rank allocation with equal costs in an additive quadratic surrogate.} For layer losses $E_\ell(r)=\sum_{j>r}\sigma_{\ell j}^2$, budget $\sum_\ell r_\ell=R$, and unit cost per rank, choose the $R$ largest $\sigma_{\ell j}^2$ while retaining an index prefix in each layer. Gains are nonincreasing within a layer, so an exchange argument proves optimality. This does not guarantee global network NLL.
\item \textbf{Exact allocation with unequal integer costs.} For allowed rank costs $c_\ell(r)\in\N$, dynamic programming minimizes the additive surrogate:
\[D_\ell(B)=\min_{r:c_\ell(r)\le B}\{D_{\ell-1}(B-c_\ell(r))+E_\ell(r)\}.\]
This follows from optimal substructure at the last layer. For an exact budget, $D_0(0)=0$, $D_0(B>0)=+\infty$; for a budget upper bound, $D_0(B\ge0)=0$.
\end{enumerate}
\subsection*{IV. Long and adaptive sequences (25--32)}
\begin{enumerate}[start=25]
\item \textbf{Adaptive whole-text bound.} If $\sup_{h:|h|=t-1}\operatorname{TV}(P_t(\cdot\mid h),Q_t(\cdot\mid h))\le\delta_t$, then $\operatorname{TV}(P_{1:T},Q_{1:T})\le1-\prod_{t=1}^T(1-\delta_t)$, by successive maximal coupling while histories agree.
\item \textbf{Exact KL chain rule.} Under conditional absolute continuity,
\[D_{\mathrm{KL}}(P_{1:T}\|Q_{1:T})=\sum_{t=1}^T\mathbb E_{h\sim P_{<t}}D_{\mathrm{KL}}(P_t(\cdot\mid h)\|Q_t(\cdot\mid h)),\]
by taking the logarithm of the product of conditional densities.
\item \textbf{TV guarantee from teacher KL measurements.} $\operatorname{TV}(P_{1:T},Q_{1:T})\le\sqrt{\frac12\sum_t\mathbb E_{P_{<t}}D_{\mathrm{KL}}(P_t\|Q_t)}$, by item 26 and Pinsker.
\item \textbf{Guarantee for every tested event.} For any text property $A$, $|P_{1:T}(A)-Q_{1:T}(A)|\le\operatorname{TV}(P_{1:T},Q_{1:T})$, by the variational definition.
\item \textbf{Current-state contraction.} If $\sup_x\operatorname{TV}(M_t(x,\cdot),\widetilde M_t(x,\cdot))\le\delta_t$ and $M_t$ has Dobrushin coefficient $\kappa_t$, state-law distance obeys $e_t\le\kappa_te_{t-1}+\delta_t$. Add and subtract the action of $M_t$, then use TV contraction.
\item \textbf{Geometric reduction cost.} If $e_0=0$, $\kappa_t\le\kappa<1$, $\delta_t\le\delta$, then $e_T\le\delta(1-\kappa^T)/(1-\kappa)$ by iterating item 29.
\item \textbf{Counterexample to replacing whole-text events by final-token events.} Let $P$ always emit 0 and $Q$ independently emit 1 with probability $\delta\in(0,1)$. Both Dobrushin coefficients vanish and final-token TV is $\delta$, but whole-text TV is $1-(1-\delta)^T\to1$, since agreement means no 1 occurred.
\item \textbf{Stopping and EOS.} If stopping uses absorbing EOS and length is at most $T$, any stopped-result test obeys item 25: mapping a full trajectory to the stopped one cannot increase TV.
\end{enumerate}
\subsection*{V. Hankel matrices, realizability, and phase information (33--40)}
\begin{enumerate}[start=33]
\item \textbf{Every observed panel lower-bounds full rank.} A finite $B=(f(u_iv_j))_{ij}$ satisfies $\rank B\le\rank H$ for $H_{u,v}=f(uv)$, including infinite rank: restricting rows and columns cannot increase rank.
\item \textbf{Verifiable shift consistency.} For every word $a$, $H_{ua,v}=H_{u,av}$, since both equal $f(uav)$. A data violation rejects that joint-function passport.
\item \textbf{Conditional and joint ranks for genuine prefixes.} If $f(u)>0$ and $f(uv)=f(u)g(v\mid u)$, fixed joint and conditional panels have equal rank: the former is the latter multiplied by invertible $\diag(f(u_i))$. Arbitrary LLM resets do not thereby become one process.
\item \textbf{Linear realization rank bound.} If $f(w)=\alpha M_w\beta$, $M_w\in\R^{d\times d}$, then $\rank H\le d$, by $H_{u,v}=(\alpha M_v)(M_u\beta)$ with consistent action order.
\item \textbf{Positive realization cost.} For a nonnegative finite submatrix $B$ of an $m$-state stochastic model, $\rank_+(B)\le m$, hence $m\ge\rank_+(B)\ge\rank B$. Factor through nonnegative reachable states and effects.
\item \textbf{Quantum linear bound.} A complete instrument system on a $d$-dimensional Hilbert space has $\rank H\le d^2$, since each sequence acts linearly on the real $d^2$-dimensional Hermitian space. This inequality alone does not determine positive realization cost.
\item \textbf{Phase nonidentifiability from autocorrelation.} Singer sets with the same $(v,k,1)$ have identical $R_D(r)=\sum_t x_D(t)x_D(t+r)$: $R_D(0)=k$, $R_D(r\ne0)=1$. Second order therefore cannot distinguish translation-inequivalent sets. This follows from the perfect-difference-set definition and restates a boundary from Revisions 17--18, not a new universal law.
\item \textbf{Single-slice phase reconstruction.} For the same $D$ and $r\ne0$, $T_r(s)=\sum_t x_D(t)x_D(t+r)x_D(t+s)=x_D(t_r+s)$ at the unique $t_r$ with $t_r,t_r+r\in D$. If $\sup_s|\widehat T_r(s)-T_r(s)|<1/2$, threshold $1/2$ recovers $D$ up to translation. Uniqueness of the pair, binary occupancy, and threshold stability prove this. The identity belongs to the earlier proved program; the explicit deterministic error tolerance is the refinement here.
\end{enumerate}
\section*{XXVI.3 Numerical experiment on saved local neural networks}
The three Revision-21 checkpoints were used without retraining on a new panel of 192 held-out contexts and 144 output characters. For each checkpoint and five $r$ values, $C_r$ and oracle $q_r$ were constructed directly, verifying $C(q_r)=C_r$ and the exact Frobenius cost. The following three-run averages use a different panel from v21 and must not be compared as head-to-head method training.
\begin{center}\begin{tabular}{rrrr}\toprule
Rank $r$&Mean KL&Mean TV&Mean $s_r$\\\midrule
8&0.14026&0.18320&63.858\\16&0.02684&0.07798&29.324\\24&0.00507&0.03456&15.396\\32&0.00115&0.01666&8.915\\48&0.00010&0.00517&2.906\\\bottomrule
\end{tabular}\end{center}
This verifies the theorem and the oracle approximation cost on known contexts; it is not a working compact language model. At these $s_r$, the theoretical bounds $\tanh s_r$ are almost one and too weak to guarantee useful rank-16 quality. The directly measured KL/TV and exact centered-logarithm equality are informative. Source code, the deterministic panel, 15 result rows, and saved checkpoints are included in the ZIP.
\section*{XXVI.4 Logical map: what forty means}
Items 01--08 establish exact oracle softmax-rank geometry; 09--16 finite-sample testing; 17--24 realizable activation approximation and surrogate optimization; 25--32 valid and invalid transfer to long generation; and 33--40 compatibility with the initial causal-Hankel and phase problems. The main benefit is one consistent system of gauges, norms, classes, and verifiable restrictions, not the count of named discoveries.
\section*{XXVI.5 OpenAI Math: 18 target families and transfer limits}
As of 12 October 2026, \url{https://github.com/openai/math} was reviewed through \texttt{overview.tex}, \texttt{CONTENTS.md}, \texttt{README.md}, Lean-scope documents, and correction history. The following connections are research directions, not borrowed proofs of our assertions. The catalogue publicly acknowledges differing levels of formal verification and records withdrawal of three manuscripts for a sign error on 7 October 2026.
\small
\begin{longtable}{@{}p{.07\textwidth}p{.32\textwidth}p{.55\textwidth}@{}}
\toprule ID&Catalogue topic&Permissible use and limitation\\\midrule\endhead
095&Nonspectrahedral hyperbolicity cones&Not every realizability constraint admits a simple SDP; not a theorem about arbitrary LLMs.\\
104&Stochastic and parity games&Algorithmic comparison classes for control inputs; a different complexity category.\\
116&Identity testing for noncommutative formulas&Motivates symbolic operator-composition certificates, not verification of trained weights.\\
119&Information from noisy bits&Motivates channel information loss, without unconditional transfer.\\
122&Reconstruction from deletion traces&Examples of finite-observation lower bounds, not a universal tokenization law.\\
130&Accelerated discrete Fourier transform&Algorithmic motivation for spectral analysis, not a Transformer speedup claim.\\
140&Memory--sample bounds for Gaussian regression&Model memory must be distinguished from learning/sample complexity; a different problem.\\
144&Simple Lebesgue spectrum in smooth dynamics&Separates spectral multiplicity and observability for a fixed operator.\\
145&Multiple mixing&Warns against replacing higher correlations by pairwise ones without assumptions.\\
178&Ramanujan graphs&Graph spectral gaps and contractions, only for specially constructed chains.\\
221&M\'ezard--Parisi formula&Variational-energy limits; no direct LLM compression follows.\\
229&Three-/four-state channel reconstruction&Branching causal-reconstruction threshold; Lean documentation limits formal coverage to the supercritical regime.\\
265&Area law for two-dimensional gapped systems&Spectrum and tensor-state approximation cost under a gap; no guarantee of a cheap LLM.\\
269&Laughlin gap and stability&Distinguish proved from conjectured stable spectral properties.\\
284&Quantum--classical query separation&Realization classes can have different costs; not equivalent to a quantum LLM.\\
350&Nodal sets and Yau bounds&Nodal bounds on $S^3$ do not transfer unconditionally between geometries.\\
365&Inverse boundary problems for metrics/connections&Gauge ambiguity and identifiability limits; the settings differ.\\
376&Computation by forced Navier--Stokes flows&Algorithmic prediction limits for constructed forced systems, not a Clay solution.\\\bottomrule
\end{longtable}\normalsize
\note{Passport and audit sources.}
\url{https://github.com/openai/math/blob/main/README.md};
\url{https://github.com/openai/math/blob/main/overview.tex};
\url{https://github.com/openai/math/blob/main/CONTENTS.md};
\url{https://github.com/openai/math/blob/main/history.md};
\url{https://github.com/openai/math/blob/main/lean/docs/140.md};
\url{https://github.com/openai/math/blob/main/lean/docs/229.md};
\url{https://github.com/openai/math/blob/main/lean/docs/376.md}.
Outside the catalogue, ASVD, SVD-LLM, and LEMS (ICML 2026) are relevant compression techniques; superiority over them has not been established.
\note{Proof boundary of Revision 22.} No external public LLM checkpoint was tested; the three earlier small trained networks were retained unchanged. Independent peer review and full Lean formalization of these forty consequences were not performed; some statements reorganize known results. No Millennium Problem or new natural law is proved. The next objectively falsifiable test is an end-to-end comparison against SVD-LLM/LEMS on an open checkpoint with matched FLOPs, perplexity, and training cost.
\setcounter{equation}{106}
\chapter{Four Structural Barriers and a Corrected End-to-End Certificate}
This revision critically examines four proposed claims about stochastic coarse-graining, nonnegative factorization, GKLS spectra, and output-layer compression. Each receives a fully verifiable form or a counterexample to its original proof. Earlier sections remain; this revision does not establish historical scientific priority.
\section*{XXVII.1 Cyclic-barrier passport and a rigorous finite-size-dependent bound}
Let $e_1,\ldots,e_n$ be standard probability columns, $n\ge3$. The action $P$ fixes $e_1$ and cycles $e_2,\ldots,e_n$; $Q$ fixes all states except $e_2$, with $Qe_2=e_1$. The effect is $\alpha=e_1^{\mathsf T}$. Let $L\in\R_{\ge0}^{d\times n}$ be column-stochastic, $\overline\alpha\in[0,1]^d$, $\overline\alpha L=e_1^{\mathsf T}$, and let $\overline P,\overline Q$ be column-stochastic. Define $\varepsilon=\max_{A\in\{P,Q\}}\frac12\|LA-\overline AL\|_{1\to1}$.
\statement{Theorem}{XXVII.1}{Rigorously proved finite-size-dependent barrier} If $2\le d<n$, every admissible blurred channel and macrotransition satisfy
\begin{equation}\varepsilon\ge\frac{\sqrt{n-d}}{2(n-1)^2}>0.\tag{107}\end{equation}
This gives positive error for each fixed $n$, but does not prove the proposed stronger constant $1/[2(d-1)]$.
\begin{proof}The equalities $\overline\alpha Le_1=1$, $\overline\alpha Le_j=0$ for $j\ge2$ put the support of $Le_1$ on coordinates with $\overline\alpha_k=1$ and all hidden supports on those with $\overline\alpha_k=0$, of which there are at most $d-1$. Set $m=n-1$. For each $i\ge2$ choose $k_i\in\{0,\ldots,m-1\}$ with $P^{k_i}e_i=e_2$. Then $\alpha QP^{k_i}e_j=\delta_{ij}$ for hidden $i,j$. The matrix $G_{ij}=\overline\alpha\overline Q\overline P^{k_i}Le_j$ factors through a coordinate face of dimension at most $d-1$, so $\rank G\le d-1$. Telescoping norm-one operators gives $\|LQP^{k_i}-\overline Q\overline P^{k_i}L\|_{1\to1}\le2(k_i+1)\varepsilon\le2m\varepsilon$, hence $|G_{ij}-\delta_{ij}|\le2m\varepsilon$. Eckart--Young--Mirsky gives $\sqrt{n-d}\le\|I_m-G\|_F\le2m^2\varepsilon$.\end{proof}
\statement{Remark}{XXVII.2}{Status of the strengthening} The input-text claim $\varepsilon\ge1/[2(d-1)]$ is not established by this reasoning: support separation and one distance of 2 after $Q$ do not imply it for all $n$ without another lemma. A million random trials do not prove a global infimum. The parameter $a$ in that formulation did not enter $P,Q$, so had no verifiable effect.
\section*{XXVII.2 Sperner's lower bound: necessity is not exact attainment}
For $W^{(n)}_{ij}=\sin^2(\pi(i-j)/n)$, let $m_S(n)=\min\{m:\binom m{\lfloor m/2\rfloor}\ge n\}$, and let $\mu_{\mathrm{cl}}(W)$ be the minimal full preparation--measurement instrument dimension under Section II's passport.
\statement{Proposition}{XXVII.3}{Static and dynamic memory: a verified chain} For $n\ge3$,
\begin{equation}\rank W^{(n)}=3,\qquad m_S(n)\le\rank_+(W^{(n)})\le\mu_{\mathrm{cl}}(W^{(n)})\le n.\tag{108}\end{equation}
\begin{proof}Ordinary rank was proved in (24). Every nonnegative factorization $W=UV$ with $m$ internal components defines supports $S_i=\{k:U_{ik}>0\}$ and $T_j=\{k:V_{kj}>0\}$. Zero diagonal and positive off-diagonal entries give $S_i\cap T_i=\varnothing$, $S_i\cap T_j\ne\varnothing$ for $i\ne j$, so the $S_i$ form an antichain and $m\ge m_S(n)$. Every classical instrument supplies such a factorization, giving $\rank_+W\le\mu_{\mathrm{cl}}$. The upper bound is attained by storing label $i$ under $R_i$, outputting 1 with probability $W_{ij}$ under $M_j$, and resetting to a fixed label. These commands also realize every later sequence of the reset instrument.\end{proof}
\statement{Remark}{XXVII.4}{Unproved equality} The antichain condition controls only zeros, not the positive $\sin^2$ entries, and does not imply exact factorization through $m$ nonnegative rays. Equality of the three proposed minima requires a constructive factorization with normalized probability effects and compatible transitions. Even the numerical assertion $\rank_+=4$ at $n=6$ is not certified here.
\section*{XXVII.3 GKLS: only observable peripheral modes}
Let $\mathcal L$ be a finite-dimensional GKLS generator, $\rho_0$ a density matrix, $0\le E\le I$, and $\Pi_\lambda$ its spectral projection. Boundedness of the CPTP semigroup $e^{t\mathcal L}$ makes its purely imaginary eigenvalues semisimple.
\statement{Theorem}{XXVII.5}{Condition for a visible undamped mode} For $\omega>0$, $p(t)=\Tr(Ee^{t\mathcal L}(\rho_0))$ has a nonzero undamped coefficient at frequency $\omega$ exactly when
\begin{equation}\ii\omega\in\spec\mathcal L\quad\text{and}\quad\Tr(E\Pi_{\ii\omega}\rho_0)\ne0.\tag{109}\end{equation}
If $M$ distinct strictly positive undamped frequencies occur, then $d^2\ge2M+1$ and $d\ge\lceil\sqrt{2M+1}\rceil$.
\begin{proof}Finite-dimensional spectral decomposition gives decaying modes and terms $e^{\ii\omega t}\Tr(E\Pi_{\ii\omega}\rho_0)$. Boundedness makes the imaginary modes semisimple and their coefficients are unique. Hermiticity preservation and real $p(t)$ pair $\pm\ii\omega$. A CPTP semigroup also has a stationary state, hence eigenvalue zero. Thus at least $2M+1$ distinct complex eigen-directions occur in the $d^2$-dimensional operator space.\end{proof}
\statement{Remark}{XXVII.6}{Counterexample to the original sufficiency claim} For $H=(\omega/2)\sigma_z$, $\rho_0=E=|0\rangle\langle0|$, and no jump operators, $\mathcal L=-\ii[H,\cdot]$ has eigenvalues $\pm\ii\omega$ but $p(t)=1$. Peripheral eigenvalues alone do not suffice. The proposed proof that ``$H=0$ excludes oscillations'' incorrectly assumed $\langle X,\mathcal LX\rangle_{\mathrm{HS}}\le0$ without extra conditions: for jump $L=|0\rangle\langle1|$, $X=\diag(2,1)$, the form equals $+1$. The universal $H=0$ conclusion is therefore excluded from this proof; stronger spectral or detailed-balance assumptions suffice.
\section*{XXVII.4 An exact end-to-end certificate for a fixed encoder}
On $N$ fixed contexts, let $p>0$, $L_{ia}=\log p_{ia}$, and $C=P_NLP_V$. Let $Z\in\R^{N\times d}$ be fixed encoder features, $S=\operatorname{Im}(P_NZ)$, and $\Pi_S$ the orthogonal projector. A candidate is $q_{ia}=\operatorname{softmax}_a(z_iAB^{\mathsf T}+b)$, with $A\in\R^{d\times r}$, $B\in\R^{V\times r}$ and free $b\in\R^V$.
\statement{Theorem}{XXVII.7}{Exact architecture-conditional error replacing a false decomposition} For fixed $Z$, unrestricted bias, and $r\ge0$,
\begin{equation}\inf_{A,B,b}\|C(p)-C(q)\|_F^2=\|(I-\Pi_S)C\|_F^2+\sum_{j>r}\sigma_j(\Pi_SC)^2.\tag{110}\end{equation}
A rank-at-most-$r$ factorization attains the optimum. For trainable $Z_\theta$, take the infimum of the right side over admissible $\theta$; it cannot be declared free independently of architecture.
\begin{proof}After centering removes bias, the attainable matrices are precisely $Y$ with $\rank Y\le r$, $\operatorname{col}Y\subseteq S$, and $Y\one_V=0$. Orthogonal projection onto fixed $S$ gives $\|C-Y\|_F^2=\|(I-\Pi_S)C\|_F^2+\|\Pi_SC-Y\|_F^2$. Truncated SVD of $\Pi_SC$ minimizes the second term. Its nonzero right singular vectors are orthogonal to $\one_V$, and left vectors lie in $S$, so it factors through $P_NZAB^{\mathsf T}P_V$. A suitable $b$ and softmax realize the corresponding centered logits, proving attainment.\end{proof}
\statement{Corollary}{XXVII.8}{Stability under measurement error} Let $e_r(C;S)$ be the square root of (110). For fixed $S$, $|e_r(C;S)-e_r(\widehat C;S)|\le\|C-\widehat C\|_F$.
\begin{proof}Distance to the same admissible set is 1-Lipschitz: apply the triangle inequality twice and take infima.\end{proof}
\statement{Remark}{XXVII.9}{Explicit counterexample to proposed Theorem 4} Let $u_1=(1,-1,0)^{\mathsf T}/\sqrt2$, $u_2=(1,1,-2)^{\mathsf T}/\sqrt6$, and $C=2u_1u_1^{\mathsf T}+u_2u_2^{\mathsf T}$. This is doubly centered. At $r=1$, $S=\operatorname{span}(u_2)$, the optimum $Y=u_2u_2^{\mathsf T}$ gives squared error 4. The proposed lower bound using $C_r=2u_1u_1^{\mathsf T}$ demanded $1+4=5$ and is false. It incorrectly applied Pythagoras to $C_r$ and arbitrary low-rank $Y$ without orthogonality. Equation (110) repairs it.
\note{Methodological result.} The four ideas yield two exact constructive theorems (finite barrier and architecture-conditional optimum), an observability-qualified spectral theorem, and a rigorous chain of necessary nonnegative-memory bounds. The stronger universal coarse-graining constant and exact nonnegative-rank claim remain unproved; an explicit counterexample rejects the false end-to-end formula. Code checks reproduce the valid formula on several random matrices and both counterexamples. This is mathematical revision, not experimental physical verification.
\note{External sources and priority.} OpenAI Math families 095, 140, 229 and formalization documents, \url{https://github.com/openai/math}; Baumgartner--Narnhofer, \textit{Analysis of quantum semigroups with GKS--Lindblad generators II}, \url{https://arxiv.org/abs/0806.3164}; Bu\v ca et al., \textit{Non-stationary coherent quantum many-body dynamics through dissipation}, Nature Communications (2019), \url{https://www.nature.com/articles/s41467-019-09757-y}; Hrube\v s, \textit{On the nonnegative rank of distance matrices} (2012), \url{https://doi.org/10.1016/j.ipl.2012.02.009}. The latter also illustrates why geometrically plausible nonnegative-rank equalities need separate proof.
\chapter{Relational Causality, Hopf Holonomy, and the Limits of Observable Memory}
\section*{XXVIII.1 Subject and status of the initial claims}
Moving beyond an external observer and linear time separates into independent questions: local-instrument correlations without a single causal order; holonomy of a nontrivial $U(1)$ bundle over $S^2$; exact stochastic-channel reduction; and conditional dynamics of an internal quantum observer. Each has a concrete operational model; their combination is not an established theory of quantum gravity. The results below do not identify Singer projective geometry with the Hopf fibration.
\statement{Definition}{XXVIII.1}{Relational operational passport} Define $\mathcal P=(A,B,I,G,H,N)$: allowed local instrument families $A,B$; accessible joint statistics $I$; operationally indistinguishable gauges $G$; admissible time slices and conditional clocks $H$; and an experimental error norm/criterion $N$. It may include an internal observer and requires no imagined device outside the entire closed system.
\section*{XXVIII.2 Indefinite order: a correlation matrix, not an automatic Hankel matrix}
For finite-dimensional local instruments, process-matrix theory writes $p(i,j)=\Tr[W(M_i^A\otimes M_j^B)]$. Besides $W\succeq0$, linear normalization constraints for admissible local CPTP instruments are required; positivity alone is insufficient. The instrument sets must be fixed in the passport. Here $K=(p(i,j))$ is a response matrix for instrument pairs, with no automatically defined prefixes and suffixes.
\statement{Proposition}{XXVIII.2}{Rank invariance of a local response matrix} For a bilinear form $B$ on finite-dimensional instrument spaces and finite bases $a_i,b_j$, let $K_{ij}=B(a_i,b_j)$. Invertible coordinate changes give $K'=XKY^{\mathsf T}$, hence $\rank K'=\rank K$. Restricting accessible instruments to a proper subspace may reduce rank.
\begin{proof}Invertible left/right multiplication preserves image dimension. Restriction selects a submatrix or composes with a noninjective linear map, which need not preserve rank.\end{proof}
\statement{Remark}{XXVIII.3}{Causal classification is not rank classification} Causal nonseparability asks whether $W$ belongs to the convex set of causally separable processes and can be tested by a causal witness. A quantum switch is causally nonseparable in its appropriate setting but does not by itself violate causal inequalities without extra restrictions. Coherent order, causal nonseparability, and device-independent causal-inequality violation are distinct claims.
\section*{XXVIII.3 One measurable witness of order coherence}
For unitary $A,B$ and normalized $|\psi\rangle$, the coherently controlled order outputs
\[|\Psi_{\mathrm{sw}}\rangle=\frac{|0\rangle\otimes BA|\psi\rangle+|1\rangle\otimes AB|\psi\rangle}{\sqrt2}.\]
Write $|u\rangle=BA|\psi\rangle$, $|v\rangle=AB|\psi\rangle$, $z=\langle u|v\rangle$.
\statement{Theorem}{XXVIII.4}{Operational phase-order witness} Control-qubit measurements in Pauli $X,Y$ bases give $p_{X+}=(1+\Re z)/2$, $p_{Y+}=(1+\Im z)/2$. If control is discarded, every target-only measurement has the same probabilities for the coherent switch and the equal classical mixture of orders. For $A=\sigma_x$, $B=\sigma_z$, $z=-1$ for every $|\psi\rangle$, so $p_{X+}=0$, whereas an equal incoherent mixture has $p_{X+}=1/2$.
\begin{proof}Tracing out the target gives $\rho_c=\frac12\left(\begin{smallmatrix}1&\overline z\\z&1\end{smallmatrix}\right)$; tracing against $(I+X)/2$ and $(I+Y)/2$ gives the probabilities. Tracing out control removes cross terms and leaves $\frac12(|u\rangle\langle u|+|v\rangle\langle v|)$. Pauli anticommutation gives $u=-v$.\end{proof}
\statement{Remark}{XXVIII.5}{Scope} This witnesses order interference with a trusted switch architecture. A single $p_{X+}$ without extra assumptions is not a universal device-independent certificate of causal nonseparability.
\section*{XXVIII.4 The Hopf bundle: a topological quantity and an observable signal}
On $S^3=\{z\in\C^2:\|z\|=1\}$, the $U(1)$ fibration is $S^3\to\mathbb{CP}^1\simeq S^2$. Choose the north-chart section
\[z(\theta,\varphi)=\begin{pmatrix}\cos(\theta/2)\\e^{\ii\varphi}\sin(\theta/2)\end{pmatrix}.\]
With fixed connection-sign convention,
\[A=-\ii z^\dagger\dd z=\frac{1-\cos\theta}{2}\dd\varphi,\qquad F=\dd A=\tfrac12\sin\theta\,\dd\theta\wedge\dd\varphi.\]
\statement{Theorem}{XXVIII.6}{Topological number and full return-signal rank are independent} For the standard Hopf bundle, $c_1=(2\pi)^{-1}\int_{S^2}F=1$. For a positively oriented latitude $C_\theta$, $\gamma=\oint_{C_\theta}A=\pi(1-\cos\theta)$. A two-path interferometer whose one path acquires phase $\gamma$ has after $n$ loops
\[p_n=\left|\frac{1+e^{\ii n\gamma}}2\right|^2=\frac{1+\cos(n\gamma)}2,\qquad
\rank[p_{i+j}]_{i,j\ge0}=\begin{cases}1,&\gamma\in2\pi\Z,\\2,&\gamma\in\pi+2\pi\Z,\\3,&\gamma\notin\pi\Z.\end{cases}\]
Thus fixed $c_1=1$ coexists with ranks 1, 2, and 3.
\begin{proof}The curvature integral is $\int_0^{2\pi}\int_0^\pi\frac12\sin\theta\,\dd\theta\dd\varphi=2\pi$, and latitude phase is $2\pi\sin^2(\theta/2)$. The expansion $p_n=\frac12+\frac14e^{\ii n\gamma}+\frac14e^{-\ii n\gamma}$ has respectively one, two, or three distinct nonzero exponential components. Vandermonde nonsingularity makes their count the Hankel rank.\end{proof}
\statement{Corollary}{XXVIII.7}{Return probabilities do not determine orientation} $p_n(\gamma)=p_n(-\gamma)$ for all $n$. A calibrated coherent ancilla can instead measure $\sin\gamma$ and distinguish signs when it is nonzero.
\begin{proof}Cosine is even. For $(|0\rangle+e^{\ii\gamma}|1\rangle)/\sqrt2$, measuring $Y$ gives $(1+\sin\gamma)/2$, changed by orientation reversal.\end{proof}
\statement{Proposition}{XXVIII.8}{Local data do not determine the global Chern class} Let $U\subset S^2$ be contractible, with an open disk in its complement. Consistent holonomy and curvature measurements confined to $U$ do not determine $c_1$ without assumptions outside $U$: for every $k\in\Z$, a $U(1)$ bundle and connection of $c_1=k$ can match those local data.
\begin{proof}Every $U(1)$ bundle is trivial over contractible $U$. Extend the specified local connection over a finite open cover and choose a degree-$k$ clutching function on a disk in the complement. Standard complex-line-bundle gluing gives $c_1=k$ without changing the connection on $U$.\end{proof}
\statement{Remark}{XXVIII.9}{What remains unproved} Singer projective sets $PG(2,q)$ have parameters $(q^2+q+1,q+1,1)$, while the Hopf bundle has an integer Chern class. Matching dimensions or the shared word ``projective'' does not create a canonical geometric link. Neither classical-memory separation as a consequence of $SU(2)$ noncommutativity nor reflection $D\mapsto-D$ as measurement of a Chern--Simons invariant has been proved.
\section*{XXVIII.5 A constructive correspondence: Singer phases as Hopf holonomy}
\statement{Theorem}{XXVIII.10}{Compiling finite cyclic phases into geometric holonomy} Let $p\ge3$, $\varnothing\ne D\subset\Z_p$, $d=|D|$, and $\theta_p=\arccos(1-2/p)$. The standard Hopf connection has latitude holonomy $h_p=e^{2\pi\ii/p}$. On $\C^d$, define
\[|s_D\rangle=d^{-1/2}\sum_{j\in D}|j\rangle,\qquad U_D=\sum_{j\in D}h_p^j|j\rangle\langle j|.\]
For every integer $n\ge0$,
\[|\langle s_D|U_D^n|s_D\rangle|^2=\frac1{d^2}\left|\sum_{j\in D}e^{2\pi\ii jn/p}\right|^2.\]
For a Singer $(p,d,1)$ set with $p=q^2+q+1$, $d=q+1$, this is 1 when $p\mid n$ and $q/(q+1)^2$ otherwise.
\begin{proof}$\gamma=\pi(1-\cos\theta_p)=2\pi/p$, so $h_p=e^{\ii\gamma}$. The diagonal operator gives the amplitude. For $p\nmid n$,
\[\left|\sum_{j\in D}h_p^{jn}\right|^2=\sum_{j,k\in D}h_p^{n(j-k)}=d+\sum_{x\in\Z_p\setminus\{0\}}h_p^{nx}=d-1=q,\]
using the $\lambda=1$ difference condition and root-of-unity summation. For $p\mid n$, every phase is 1 and the squared unnormalized amplitude is $d^2$.\end{proof}
\statement{Corollary}{XXVIII.11}{Geometric encoding does not determine global topology} Every phase is a power of one loop's holonomy, connecting the earlier Singer family to an explicit two-level geometric-phase module. But the $d$-path interferometer need not physically be one spin or fibration: phase powers are independently controlled operations. Local holonomy does not determine global $c_1$. This is a rigorous realization bridge, not a derivation of $SU(2)$ noncommutativity, quantum gravity, or cosmological parameters from difference sets.
\statement{Proposition}{XXVIII.12}{Finite-sample loop-orientation test} For calibrated orientations $\gamma,-\gamma$ with $\sin\gamma\ne0$, independent ancilla-$Y$ trials have probabilities $(1\pm\sin\gamma)/2$. A sample-frequency threshold test distinguishes the sign with error at most $\alpha$ when $N\ge2\log(1/\alpha)/\sin^2\gamma$. This conservative sufficient bound assumes ideal independent trials and known bases; it omits drift and correlated noise.
\begin{proof}Hoeffding bounds crossing of $1/2$ from true mean $(1+|\sin\gamma|)/2$ by $\exp[-2N(|\sin\gamma|/2)^2]$. The opposite sign is symmetric.\end{proof}
\section*{XXVIII.6 Zero-defect coarse-graining exists without a fractal}
\statement{Proposition}{XXVIII.13}{Exact reduction through a stochastic right-inverse lift} Let column-stochastic $L\in\R^{d\times n}$ and $R\in\R^{n\times d}$ satisfy $LR=I_d$. For any macrotransitions $\overline P_a$, set $P_a=R\overline P_aL$. Then these are stochastic and
\[LP_a=\overline P_aL,\qquad LP_{a_T}\cdots P_{a_1}=\overline P_{a_T}\cdots\overline P_{a_1}L.\]
This is exact zero-defect reduction for compatible transitions, not removal of the obstruction for all microdynamics.
\begin{proof}$LP_a=LR\overline P_aL=\overline P_aL$; iterate along the word. Compositions of stochastic maps are stochastic.\end{proof}
\statement{Proposition}{XXVIII.14}{Normalization excludes nontrivial state scaling} If $\Lambda$ is trace-preserving, $\rho$ a normalized density matrix, and $\Lambda(\rho)=\lambda\rho$, then $\lambda=1$. Thus $LM_aL^\dagger=\lambda M_a$ with $\lambda\ne1$ cannot by itself mean preservation of normalized physical states or all probabilities. Unnormalized eigenoperators may scale, but require a separate definition.
\begin{proof}$1=\Tr\Lambda(\rho)=\lambda\Tr\rho=\lambda$.\end{proof}
\statement{Remark}{XXVIII.15}{Dimension and rank} The self-similar dimension $\log N/\log(1/r)$ concerns spatial similarity scaling under standard separation assumptions, not observability-operator rank. Finite-dimensional linear rank is integer-valued; ``fractional Hankel rank'' needs a different explicitly defined regularization or spectral dimension.
\section*{XXVIII.7 An internal observer and finite conditional time}
\statement{Proposition}{XXVIII.16}{Finite relational time without an external dynamical parameter} Let $U$ be unitary on $\mathcal H_S$, $U^N=I$, and $T_C|t\rangle=|t+1\bmod N\rangle$ on $\mathcal H_C=\C^N$. The history state
\[|\Omega\rangle=N^{-1/2}\sum_{t=0}^{N-1}|t\rangle_C\otimes U^t|\psi_0\rangle_S\]
is normalized, invariant under $V=T_C\otimes U$, and has conditional system state $U^t|\psi_0\rangle$ at clock outcome $t$. This discrete relational-time model does not prove that all physical time or causality arises solely from entanglement.
\begin{proof}Clock orthogonality gives normalization. $V$ shifts the $t$ term to $|t+1\rangle U^{t+1}|\psi_0\rangle$, and $U^N=I$ closes the sum. Projection onto $\langle t|$ gives the conditional normalized vector.\end{proof}
\statement{Theorem}{XXVIII.17}{Global-state indistinguishability under local access} Bell states $|\Phi^\pm\rangle=(|00\rangle\pm|11\rangle)/\sqrt2$ give equal probabilities for all measurements confined to the first subsystem, although they are orthogonal and distinguishable jointly.
\begin{proof}$\Tr_B|\Phi^\pm\rangle\langle\Phi^\pm|=I_2/2$, so every local POVM agrees, while $\langle\Phi^+|\Phi^-\rangle=0$.\end{proof}
\statement{Remark}{XXVIII.18}{Access, not a dimension slogan} This indistinguishability comes from the restricted observable algebra, not a purported general formula $\dim\mathcal H_{\mathrm{apparatus}}<\dim\mathcal H_{\mathrm{universe}}$. No-cloning of arbitrary unknown states follows from linearity/unitarity even with a large ancillary space. An internal observer's required predictive coordinates still depend on the operational passport.
\section*{XXVIII.8 A joint invariant: a defined alternative to ``topological rank''}
For specified geometric and operational structures, introduce $\mathcal J_{\mathcal P}=(c_1(L),\rank H_{\mathcal P},\Im z_{\mathrm{ctrl}},A,B)$. Its entries concern the global bundle, the passport's action/measurement order, and calibrated interferometry. It is a tuple of distinct invariants, not one new scalar law. They are compatible without determining each other: fixed $c_1$ admits three ranks, and the local no-go result explains why global coverage is needed to measure $c_1$.
\begin{center}\small\begin{tabularx}{\textwidth}{@{}p{.22\textwidth}XX@{}}\toprule
Object&Established&Does not follow without more assumptions\\\midrule
Causal switch&Order interference with trusted instruments&Causal-inequality violation\\
Hopf holonomy&$c_1=1$, phase response, ranks 1/2/3&Global topology from one local loop\\
Stochastic reduction&Zero defect under exact lumpability&Removal of the obstruction for all transitions\\
Internal clock&Conditional discrete dynamics&Universal quantum-gravitational causality\\
Local rank&Instrument-basis invariant&Topological ``subject--object'' obstruction\\\bottomrule
\end{tabularx}\end{center}
\section*{XXVIII.9 Reproducibility, independent priority, and empirical limits}
The revision's tests check exact Hankel strata at $\gamma=0,\pi,\pi/2,\pi/3$, Pauli controlled-order interference, zero-error stochastic reduction, Bell-state local indistinguishability, and finite-history-state invariance. Code is in the archive's \texttt{checks/} directory. These are finite-model computations, not laboratory holonomy measurements or cosmological data.

Sources and priority: Oreshkov--Costa--Brukner (2012), \url{https://www.nature.com/articles/ncomms2076}; Ara\'ujo et al., \textit{Witnessing causal nonseparability} (2015), \url{https://arxiv.org/abs/1506.03776}; Wootters--Zurek (1982), \url{https://doi.org/10.1038/299802a0}; discussion of Page--Wootters, \url{https://ora.ox.ac.uk/objects/uuid:e1a7308e-9830-4924-9ab7-9df0dc4abcd2}; Singer's construction, \url{https://personal.ntu.edu.sg/bernhard/Publications/pub/Singer.pdf}. A general fundamental law would require operational independence of passport conventions, a nontrivial new theorem linking these structures, falsifiable predictions, independent experiments, and worldwide novelty assessment. These steps are not claimed complete.


% 05_positive_memory.tex
\chapter{Conic Impossibility, Closed Clocks, and Exact Feature Geometry}
\section*{XXIX.1 Verification passport and verdict map}
This revision examines four proposals following Sections XXIII--XXIV: a six-label trigonometric nonnegative factorization; a causal phase transition attributed to superposed internal clocks; a spectral identity between an irreducible $SU(2)$ spin orbit and a Singer difference set; and extension of softmax compression certificates to deep-network feature-space families. Operator definitions, dimensions, matrix products, and infimum domains must be checked before accepting a claim. Nonnegative rank is treated by exact algebra and geometric stability by norm inequalities. No worldwide novelty is claimed.
\begin{center}\small\begin{tabularx}{\textwidth}{@{}p{.3\textwidth}X@{}}\toprule
Proposed claim&Rigorous revision\\\midrule
$\rank_+(W^{(6)})=4$&The proposed factorization is false; $\rank_+(W^{(6)})\ge5$ is proved.\\
Phase threshold $\theta_c$&It does not follow from the process definition. An invertible allowed local control-phase rotation preserves the causal class.\\
Irreducible spin reproduces Singer&False for standard $J_z$: Hankel rank is $2d-1$, not $d^2-d+1$. A general spectral gate $f(J_z)$ can realize Singer.\\
Geometric layer error&The fixed-encoder exact optimum of XXIII remains; a valid Lipschitz comparison under subspace drift is obtained.\\\bottomrule
\end{tabularx}\end{center}
\section*{XXIX.2 Six-label antidistinguishability needs more than four positive rays}
Let $W^{(6)}_{ij}=\sin^2[\pi(i-j)/6]$, $0\le i,j<6$. The proposed nonnegative $U\in\R^{6\times4}$, $V\in\R^{4\times6}$ do not satisfy $UV=W^{(6)}$: exact multiplication gives $(UV)_{02}=1/2\ne3/4=W^{(6)}_{02}$ and $(UV)_{03}=1/4\ne1=W^{(6)}_{03}$. Moreover $\rank(UV)=4$, whereas $\rank W^{(6)}=3$. Relabeling supports does not repair this.
\statement{Theorem}{XXIX.1}{Six-label barrier for matrices of linear rank three} If $A\in\R_{\ge0}^{6\times6}$ has $A_{ii}=0$, $A_{ij}>0$ for $i\ne j$, and $\rank A\le3$, then $\rank_+(A)\ge5$. In particular,
\[\rank W^{(6)}=3,\qquad5\le\rank_+(W^{(6)})\le\mu_{\mathrm{cl}}(W^{(6)})\le6.\]
\begin{proof}Suppose $A=UV$ through four nonnegative components, padding smaller factorizations by zeros. Let $S_i=\{a:U_{ia}>0\}$, $T_i=\{a:V_{ai}>0\}$. Zero diagonal gives $S_i\cap T_i=\varnothing$; positive off-diagonal gives $S_i\cap T_j\ne\varnothing$ for $i\ne j$. Containment $S_i\subseteq S_j$ would contradict the latter, so the six $S_i$ form an antichain on four elements. The Lubell--Yamamoto--Meshalkin inequality $\sum_i\binom4{|S_i|}^{-1}\le1$, together with $\binom4k\le6$, forces every $|S_i|=2$. Thus the supports are the six edges of $K_4$. For fixed $j$, if some $a\in S_j^c$ were absent from $T_j\subseteq S_j^c$, the edge $\{a,b\}$ with $b\in S_j$ would miss $T_j$, contradicting positivity. Hence $T_j=S_j^c$ and the columns also have all six two-element supports.

If $Ux=0$, each edge $\{a,b\}$ gives $u_ax_a+u_bx_b=0$ with positive coefficients. A zero coordinate propagates through connected $K_4$ to $x=0$. If none is zero, signs must be opposite across every edge, impossible on a triangle. Thus $\rank U=4$. The same argument for column supports gives $\rank V=4$. Since $U$ is injective and $V$ surjective onto $\R^4$, $\rank(UV)=4$, contradicting $\rank A\le3$. For $W^{(6)}$, the independent modes $1,\cos(\pi i/3),\sin(\pi i/3)$ give rank three; $W^{(6)}=W^{(6)}I_6$ gives the upper bound six.\end{proof}
\statement{Remark}{XXIX.2}{Boundary of the result} The theorem excludes rank four; it does not attain rank five. Static nonnegative rank is not automatically the memory of a controlled instrument preserving every intervention. Relevant antecedents include Lin--Chu and Hrube\v s on Euclidean distance matrices; independent historical novelty assessment was not performed.
\section*{XXIX.3 Cyclic quantum clocks: the correct sector and no universal causal threshold}
Let $\mathcal H_C=\C^N$, let $\mathcal H_S$ be finite-dimensional, and let $U_t$ be unitary. Define
\[S=\sum_{t=0}^{N-1}|t+1\bmod N\rangle\langle t|\otimes U_t,\qquad J=I-S.\]
Although $S$ is unitary, $J$ is generally not Hermitian. Hence $e^{\ii\alpha J}$ is not generally a unitary gauge group, and $\int_0^{2\pi}e^{\ii\alpha J}\rho e^{-\ii\alpha J}\,\dd\alpha$ is not a physical averaging operation without further conditions.
\statement{Proposition}{XXIX.3}{Existence condition for a cyclic history} Put $M=U_{N-1}\cdots U_1U_0$. Then $\ker(I-S)\simeq\ker(I-M)$. For $\psi_0\in\ker(I-M)$, a normalized history is
\[|\Omega\rangle=N^{-1/2}\sum_{t=0}^{N-1}|t\rangle\otimes U_{t-1}\cdots U_0|\psi_0\rangle,\qquad S|\Omega\rangle=|\Omega\rangle,\]
with empty product the identity. If $\ker(I-M)=\{0\}$, no nonzero stationary history of this form exists.
\begin{proof}Coordinatewise, $S\Psi=\Psi$ says $\psi_{t+1}=U_t\psi_t$. A full circuit requires $M\psi_0=\psi_0$; conversely any such vector gives a consistent cycle. Unitarity supplies normalization.\end{proof}
\statement{Proposition}{XXIX.4}{Control phase alone does not define a causal-order critical point} For normalized $|u\rangle,|v\rangle$, let $|\Psi_\theta\rangle=(|0\rangle|u\rangle+e^{\ii\theta}|1\rangle|v\rangle)/\sqrt2$. Then
\[\Tr_C|\Psi_\theta\rangle\langle\Psi_\theta|=\tfrac12(|u\rangle\langle u|+|v\rangle\langle v|),\]
independent of $\theta$. Retaining control, its phase shift can be undone locally, so phase alone does not change a causally separable class into a nonseparable one. A fixed $X$ detector has $p_{X+}(\theta)=\frac12[1+\Re(e^{\ii\theta}\langle u|v\rangle)]$.
\begin{proof}Partial trace removes $|0\rangle\langle1|$ cross terms; the control amplitude in the $|\pm\rangle$ basis gives the final formula. The phase is local unitary conjugation of the joint process in its control sector, preserving the physical class under an allowed inverse. This presumes an initially correctly normalized process operator and an admissible control transformation.\end{proof}
Consequently $\operatorname{dist}_{\mathrm{sep}}W(\theta)=\max\{0,(|\cos\theta|-\kappa_{\mathrm{crit}})/2\}$ does not follow from process-matrix axioms. It requires a specific admissible family, fixed norm, and solved convex optimization. Cross terms contain both sine and cosine; replacing coherence by $\cos\theta$ alone is not a general identity.
\section*{XXIX.4 Spectral rigidity of an irreducible $SU(2)$ orbit and a valid Singer gate}
\statement{Theorem}{XXIX.5}{Rank of a uniform spin orbit} Let $d=q+1\ge3$, $p=q^2+q+1=d^2-d+1$, and let $J_z$ be the irreducible spin operator with eigenvalues $m=-j,-j+1,\ldots,j$, $j=(d-1)/2$. For equal weights $|c_m|^2=1/d$,
\[P_{\mathrm{spin}}(n)=\left|\sum_m\frac1d e^{2\pi\ii mn/p}\right|^2=\frac1{d^2}\left[\frac{\sin(d\pi n/p)}{\sin(\pi n/p)}\right]^2\quad(p\nmid n),\]
with value 1 for $p\mid n$, and $\rank H(P_{\mathrm{spin}})=2d-1<p$. Thus ordinary $e^{2\pi\ii J_z/p}$ does not reproduce the $p$-mode Singer return sequence.
\begin{proof}The geometric series yields the sine ratio. Its squared-modulus Fourier coefficients are $(d-|k|)/d^2>0$ for $-(d-1)\le k\le d-1$. Since $p>2(d-1)$, these $2d-1$ frequencies are distinct modulo $p$; Vandermonde factorization gives the rank. A Singer set has each nonzero ordered difference once, whereas consecutive spin labels have difference 1 with multiplicity $d-1\ge2$.\end{proof}
\statement{Proposition}{XXIX.6}{Valid spectral compilation of Singer without topological identification} For any Singer $D=\{a_1,\ldots,a_d\}\subset\Z_p$, choose a real interpolation polynomial $f$ of degree at most $d-1$ with $f(m_i)=a_i$. The gate $U=\exp(2\pi\ii f(J_z)/p)$ on the equal-weight spin state gives
\[|\langle\psi|U^n|\psi\rangle|^2=d^{-2}\left|\sum_{a\in D}e^{2\pi\ii an/p}\right|^2=P_D(n),\qquad\rank H(P_D)=p.\]
For $q=2$, $d=3$, $D=\{0,1,3\}$, use $f(x)=(x^2+3x+2)/2$ at $m=-1,0,1$.
\begin{proof}Interpolation assigns the desired eigenvalues without mixing eigenvectors. The return amplitude is the normalized phase sum, and perfect differences give positive weights at all $p$ frequencies. Nonlinear $f(J_z)$ is not a standard spin-rotation generator, so this does not identify $S^3$ topology with Singer combinatorics.\end{proof}
\section*{XXIX.5 Grassmannian subspaces: a valid comparison instead of a false loss sum}
For doubly centered $C\in\R^{N\times V}$, $S\subseteq\one_N^\perp$, projector $\Pi_S$, and output rank $r$, define
\[e_r(C;S)=\min\{\|C-Y\|_F:\operatorname{col}Y\subseteq S,\ \rank Y\le r,\ Y\one_V=0\}.\]
Theorem XXIII.4 is equivalent to
\begin{equation}e_r(C;S)^2=\|C\|_F^2-\sum_{j=1}^r\sigma_j(\Pi_SC)^2=\|(I-\Pi_S)C\|_F^2+\sum_{j>r}\sigma_j(\Pi_SC)^2.\tag{111}\end{equation}
When all network weights are optimized, $S$ changes with the encoder. An arbitrary $S_L$ cannot appear on the right of an inequality with $\inf_{W_{1:L}}$ on the left unless the attainable set of $S_L$ is first fixed.
\statement{Theorem}{XXIX.7}{Lipschitz stability in feature space} For $S,T\subseteq\one_N^\perp$,
\[|e_r(C;S)-e_r(C;T)|\le\|\Pi_S-\Pi_T\|_2\|C\|_F.\]
For measured layer spaces $S_1,\ldots,S_L$,
\[e_r(C;S_L)\ge e_r(C;S_1)-\|C\|_F\sum_{\ell=1}^{L-1}\|\Pi_{S_{\ell+1}}-\Pi_{S_\ell}\|_2.\]
The analogous telescoping bound holds for projection errors $\|(I-\Pi_{S_\ell})C\|_F$. Neither forces a positive sum of layer errors: a next layer may increase, preserve, or decrease error.
\begin{proof}Take the truncated-SVD minimizer $Y_S$ and set $Y_T=\Pi_TY_S$. It has rank at most $r$, columns in $T$, and zero row sums. Thus
\[e_r(C;T)\le\|C-\Pi_TY_S\|_F\le e_r(C;S)+\|(\Pi_S-\Pi_T)Y_S\|_F\le e_r(C;S)+\|\Pi_S-\Pi_T\|_2\|C\|_F,\]
since $\|Y_S\|_F\le\|C\|_F$. Exchange $S,T$ and telescope.\end{proof}
\statement{Remark}{XXIX.8}{Errors in the proposed Grassmannian version} (1) For equal-dimensional spaces, $\|\Pi_S-\Pi_T\|_F^2=2\sum_i\sin^2\theta_i$ is squared chordal distance, not squared standard Riemannian geodesic length $\sum_i\theta_i^2$. (2) The proposed projection-error bound incorrectly used $\|C\|_2$ rather than $\|C\|_F$. Take orthonormal $u_1,\ldots,u_4\in\one_N^\perp$, $S_1=\operatorname{span}(u_3,u_4)$, $S_L=\operatorname{span}(u_1,u_2)$, $C=[u_1\ u_2]$, adding columns and an orthogonal embedding for right centering if needed. The left side is zero, but the proposed right side is $\sqrt2-1>0$. Use $\|C\|_F$ or the sharper $\|(\Pi_S-\Pi_T)C\|_F$. (3) The proposed lower bound using the singular-value tail of $\Pi_{S_*}C$ cannot improve the general oracle residual: $\sigma_j(\Pi_{S_*}C)\le\sigma_j(C)$, so its tail is bounded above by the already universal rank-$r$ lower bound $\sum_{j>r}\sigma_j(C)^2$.
\section*{XXIX.6 Exact register of checks and open problems}
\texttt{checks/verify\_structural\_v25.py} independently reproduces rational factorization failures, matrix ranks, spin-frequency examples, clock nonperiodicity, phase erasure, the exact fixed-encoder optimum, and projector perturbations. The lower bound five is analytic; random checks support but do not replace it. The full protocol is in \texttt{checks/RESULTS\_v25.json}.

The exact $\rank_+(W^{(6)})$, unsettled in Revision XXV, is proved in Section XXVI. Open questions remain physical implementation of the nonlinear Singer gate, a causally nonseparable family with a specified metric, trainable-encoder optimization for $e_r(C;S)$, and independent priority review.

Antecedents: Singer difference sets and finite projective planes, \url{https://personal.ntu.edu.sg/bernhard/Publications/pub/Singer.pdf}; Lin--Chu, \textit{On the Nonnegative Rank of Euclidean Distance Matrices}, Linear Algebra Appl. 433 (2010), DOI 10.1016/j.laa.2010.03.038; Hrube\v s, \textit{On the nonnegative rank of distance matrices}, Information Processing Letters 112 (2012), DOI 10.1016/j.ipl.2012.02.009. Indefinite causal order and the quantum switch are discussed in XXIV; this revision claims no new physical result.
\chapter{Exact Conic Complexity and the Separation of Static and Dynamic Memory}
\section*{XXX.1 Resolving the question of five positive rays}
Revision XXV established $\rank_+(W^{(6)})\ge5$ through support structure and a four-component rank contradiction. Here a constructive upper bound closes the question. For $0\le i,j\le5$,
\[W^{(6)}_{ij}=\sin^2\frac{\pi(i-j)}6,\qquad W^{(6)}=\frac14\begin{pmatrix}0&1&3&4&3&1\\1&0&1&3&4&3\\3&1&0&1&3&4\\4&3&1&0&1&3\\3&4&3&1&0&1\\1&3&4&3&1&0\end{pmatrix}.\]
\statement{Theorem}{XXX.1}{Exact nonnegative factorization dimension of six labels} $\rank W^{(6)}=3$ and $\rank_+(W^{(6)})=5$. An explicit upper certificate is $W^{(6)}=AD/4$, where
\[A=\begin{pmatrix}1&0&1&0&0\\1&0&0&0&1\\0&0&0&1&2\\0&1&0&0&1\\0&1&1&0&0\\0&0&2&1&0\end{pmatrix},\qquad D=\begin{pmatrix}0&0&1&3&3&1\\3&3&1&0&0&1\\0&1&2&1&0&0\\1&1&0&1&1&0\\1&0&0&0&1&2\end{pmatrix}.\]
\begin{proof}All entries are nonnegative integers, and multiplication over $\Q$ gives $W^{(6)}$, so nonnegative rank is at most five. The opposite bound is the earlier theorem: a six-element antichain in a four-component factorization forces both factors and their product to have rank four, contradicting ordinary rank three. The latter also follows from $\sin^2t=(1-\cos2t)/2$ and independence of $1,\cos(2\pi j/6),\sin(2\pi j/6)$.\end{proof}
\statement{Proposition}{XXX.2}{Exact origin of the five-factor construction} Put
\[B=\begin{pmatrix}0&0&0&1&2&1\\1&2&1&0&0&0\\0&0&1&1&0&0\\0&1&0&0&1&0\\1&0&0&0&0&1\end{pmatrix},\qquad S=AB.\]
Up to the chosen scale and indexing, $S$ is the regular hexagon's standard slack matrix: $S_{ij}=0$ when $j=i$ or $i+1\pmod6$, and 1 or 2 otherwise. With $G_{kj}=\frac14(\delta_{kj}+\delta_{k,j+1\bmod6})$,
\[W^{(6)}=SG=A(BG)=\tfrac14AD.\]
\begin{proof}Rational multiplication verifies $D=4BG$ and $ABG=W^{(6)}$. Geometrically, the tangent constraint at a vertex is a positive combination of its two adjacent facet slacks. The five-factor construction for $S$ goes back to Gouveia--Parrilo--Thomas, Example 3.4; no global priority is claimed for that lift.\end{proof}
\section*{XXX.2 Exact and asymptotic separation of three complexities}
\statement{Theorem}{XXX.3}{Static factorization versus dynamic stochastic memory} For every integer $n\ge3$, let $p_n(k)=\sin^2(\pi k/n)$ and $W^{(n)}_{ij}=p_n(i-j)$, $i,j\in\Z_n$. Under passport $\mathcal P_n$---exact probabilities, every $k\in\N_0$, one homogeneous classical transition and one fixed finite quantum unitary dynamics---
\[\rank H(p_n)=\rank W^{(n)}=3,\qquad m_S(n)\le\rank_+(W^{(n)})\le2\lceil\log_2n\rceil,\]
\[\mu_{\mathrm{qm}}^{\mathrm{unitary}}(p_n)=2,\qquad\mu_{\mathrm{cl}}^{\mathrm{homogeneous}}(p_n)=n,\]
where $m_S(n)=\min\{m:\binom m{\lfloor m/2\rfloor}\ge n\}$. Hence static nonnegative rank is $\Theta(\log n)$ whereas exact dynamic classical memory is $n$.
\begin{proof}\textit{Linear ranks.} Expansion gives three distinct active harmonics $0,\pm2\pi/n$, so both finite ordinary rank and infinite Hankel rank are 3. \textit{Static lower bound.} Zero diagonal and positive off-diagonal entries imply the support-antichain bound $m_S(n)=\Omega(\log n)$. \textit{Static upper bound.} Define
\[v_i=(\cos(2\pi i/n),\sin(2\pi i/n)),\quad w_{j+1/2}=(\cos(2\pi(j+1/2)/n),\sin(2\pi(j+1/2)/n)).\]
The regular-polygon slack $S^{(n)}_{i,j}=\cos(\pi/n)-v_i\cdot w_{j+1/2}\ge0$ satisfies
\[W^{(n)}_{ij}=\frac{S^{(n)}_{i,j-1}+S^{(n)}_{i,j}}{4\cos(\pi/n)}.\]
Thus $W^{(n)}=S^{(n)}G^{(n)}$ with $G^{(n)}\ge0$, so $\rank_+W^{(n)}\le\rank_+S^{(n)}\le2\lceil\log_2n\rceil$ by known regular-polygon extended formulations (Fiorini--Rothvoss--Tiwary; refinement by Vandaele--Gillis--Glineur).

\textit{Quantum dynamics.} Take $|\pm\rangle=(|0\rangle\pm|1\rangle)/\sqrt2$, $U=\diag(1,e^{2\pi\ii/n})$, $E=|-\rangle\langle-|$. Then $\langle+|U^{-k}EU^k|+\rangle=p_n(k)$. One level cannot produce a nonconstant signal. \textit{Classical dynamics.} The expansion $p_n(k)=\frac12-\frac14(e^{2\pi\ii k/n}+e^{-2\pi\ii k/n})$ has a primitive order-$n$ root. Any finite homogeneous realization must have it in its peripheral spectrum; some closed irreducible class then has period divisible by $n$ and at least $n$ states. An $n$-cycle with effect $\alpha_j=p_n(j)$ attains the bound.\end{proof}
\section*{XXX.3 Probability normalization does not follow from conic factorization}
\statement{Proposition}{XXX.4}{The five factors cannot be diagonally normalized} No positive diagonal scaling $A\mapsto A\diag(s)$, $(D/4)\mapsto\diag(s)^{-1}(D/4)$ simultaneously makes the first factor's rows probability distributions and the second factor's columns effects in $[0,1]$.
\begin{proof}$As=\one$ forces $s=(1-t,1-t,t,1-2t,t)^\top$. Admissible effects require $s_k\ge\max_j(D/4)_{kj}$. The first row demands $1-t\ge3/4$, or $t\le1/4$, while the third demands $t\ge1/2$, a contradiction.\end{proof}
This does not exclude another normalized five-state factorization. The valid bound remains $5=\rank_+(W^{(6)})\le\mu_{\mathrm{cl}}^{\mathrm{static}}(W^{(6)})\le6$. Exact normalized static cost and independent sequential-intervention requirements remain separate questions.
\section*{XXX.4 Historical limits and numerical reproducibility}
The contribution is an exact integration of known polygonal lifts with distinctions between dynamic and predictive memory. Five-factor regular-hexagon lifts were known at least from Gouveia--Parrilo--Thomas (2011/2013), and logarithmic upper bounds from regular-polygon studies. Historical priority of this particular fact is not established. The integrated three-problem passport is new within the manuscript; its historical originality needs separate study.

Rational products, factor ranks, $SG=W$, and diagonal-normalization impossibility were checked symbolically; the executable test is in \texttt{checks/}, with results in \texttt{RESULTS\_v26.json}. Sources: Gouveia, Parrilo, Thomas, ``Lifts of Convex Sets and Cone Factorizations'', \textit{Mathematics of Operations Research} 38 (2013), 248--264, DOI 10.1287/moor.1120.0575, Example 3.4; Vandaele, Gillis, Glineur, ``On the linear extension complexity of regular $n$-gons'', \textit{Linear Algebra and its Applications} 521 (2017), 217--239, DOI 10.1016/j.laa.2016.12.023. The value five is established here by exact recomputation plus Section XXV's lower bound; global novelty is not claimed.
\chapter{Joint Predictive Memory, the Exact Cost of a Shared Latent Space, and OpenAI Math Boundaries}
\section*{XXXI.1 Passport and the multi-interface problem}
Let $N$ experimental contexts serve $K$ independent conditional output interfaces with alphabet sizes $V_a$. On this fixed panel, strictly positive probabilities $p^{(a)}_{iv}$ are known. Define
\[C_a=P_N[\log p^{(a)}_{iv}]P_{V_a},\quad P_m=I_m-\one_m\one_m^\top/m,\quad C=[\sqrt{w_1}C_1\mid\cdots\mid\sqrt{w_K}C_K],\quad w_a>0.\]
This $C$ is the observable object; double centering removes arbitrary row-logit shifts and output biases. Let $Z\in\R^{N\times d}$ be frozen encoder activations, $S=\operatorname{Im}(P_NZ)\subset\one_N^\perp$, and $\Pi_S$ its orthogonal projector. One shared adapter $A\in\R^{d\times r}$ serves all interfaces, each with its own $B_a\in\R^{V_a\times r}$ and bias $b_a$. The predictions $q_i^{(a)}=\operatorname{softmax}(z_iAB_a^\top+b_a)$ are valid conditional distributions. The loss is the weighted sum of $\|C_a(p)-C_a(q)\|_F^2$, not direct KL or inference time.
\statement{Theorem}{XXXI.1}{Exact joint spectral certificate} For $0\le r\le\min(d,\sum_aV_a)$,
\begin{equation}\inf_{A,(B_a,b_a)}\sum_{a=1}^K w_a\|C_a(p)-C_a(q)\|_F^2=\|(I-\Pi_S)C\|_F^2+\sum_{j>r}\sigma_j(\Pi_SC)^2.\tag{112}\end{equation}
A finite-dimensional linear shared adapter attains the infimum. This concerns only frozen $Z$ and these interfaces, not all-layer compression or superiority in KL, latency, or parameters over modern algorithms.
\begin{proof}The centered joint prediction is
\[Y=P_NZA[\sqrt{w_1}B_1^\top P_{V_1}\mid\cdots\mid\sqrt{w_K}B_K^\top P_{V_K}],\]
so $\operatorname{col}Y\subset S$ and $\rank Y\le r$. Conversely any such $Y$ with separately centered output blocks admits this factorization: choose its common left space within $S$, express a basis through $P_NZ$, and split the right factor into blocks. Truncated SVD of $\Pi_SC$ has precisely these properties because each block is right-centered. The orthogonal identity $\|C-Y\|_F^2=\|(I-\Pi_S)C\|_F^2+\|\Pi_SC-Y\|_F^2$ and Eckart--Young--Mirsky prove the claim.\end{proof}
\statement{Corollary}{XXXI.2}{Cost of sharing a representation} Let $E_{\mathrm{shared}}$ be (112), and let
\[E_{\mathrm{separate}}=\sum_a w_a\left[\|(I-\Pi_S)C_a\|_F^2+\sum_{j>r}\sigma_j(\Pi_SC_a)^2\right]\]
be the sum of minima with independent rank-$r$ adapters. The exact sharing-cost invariant $\Delta_r=E_{\mathrm{shared}}-E_{\mathrm{separate}}\ge0$ vanishes exactly when one subspace of dimension at most $r$ in $S$ simultaneously realizes a Frobenius-optimal rank approximation to every $\Pi_SC_a$. This is not a parameter-matched comparison: the independent construction permits $K$ adapters.
\begin{proof}Shared factorizations form a subclass of independent ones. Equality requires one left space to attain every nonnegative individual gap simultaneously.\end{proof}
\statement{Proposition}{XXXI.3}{Strictly positive sharing cost} Two rank-one interfaces and a two-dimensional observed encoder can have $E_{\mathrm{separate}}=0$, $E_{\mathrm{shared}}=1$ for shared rank $r=1$ and equal weights.
\begin{proof}Choose orthonormal $u,v\in\one_N^\perp$ and unit centered output vectors $a,b$. Set $C_1=ua^\top$, $C_2=vb^\top$, $S=\operatorname{span}\{u,v\}$. Each interface is separately exact at rank one; the concatenation has two singular values equal to one. Apply (112).\end{proof}
\statement{Theorem}{XXXI.4}{Stability under measurement and dataset noise} Let $\widehat C=C+E$, $\|E\|_F\le\eta$, with fixed $S$. For $e_r=\sqrt{E_{\mathrm{shared}}}$,
\[|e_r(\widehat C;S)-e_r(C;S)|\le\eta,\qquad|\widehat\Delta_r-\Delta_r|\le4\eta\|C\|_F+2\eta^2.\]
\begin{proof}Distance to a fixed nonempty set is 1-Lipschitz. The shared feasible set contains zero, so $e_r(C)\le\|C\|_F$; the squared costs differ by at most $\eta(2\|C\|_F+\eta)$. For separate tasks, blockwise application and Cauchy--Schwarz give the same total bound. Add the two estimates.\end{proof}
\section*{XXXI.2 Comparison using actual frozen representations}
Fifteen checks used the three previously trained small Transformers, not an external large LLM: three saved independent-initialization checkpoints, each at ranks $4,8,12,16,24$. The passport used the same 320 calibration contexts and separate 320 held-out contexts. Three conditional output interfaces partitioned the character vocabulary by index modulo three; frozen encoders supplied all features. The Frobenius optimum was attained exactly on the training panel; held-out performance was mean teacher/student KL over the three conditional vocabularies.
\begin{center}\begin{tabular}{rrrr}\toprule
$r$&Mean $\Delta_r$ (Frobenius$^2$)&KL, shared adapter&KL, three adapters\\\midrule
4&1833.8814&0.28546&0.25429\\8&1596.0038&0.13649&0.09691\\12&998.2802&0.06446&0.03457\\16&557.5252&0.02805&0.01265\\24&203.9798&0.00695&0.00212\\\bottomrule
\end{tabular}\end{center}
All 15 calibration identities, the orthogonal decomposition, and nonnegative sharing cost were checked in \texttt{verify\_multi\_passport\_v27.py}; results were saved as CSV and JSON. These measure the quality tradeoff of a narrow shared adapter, but separate adapters have different parameter/operation costs. The numbers do not prove superiority of a new deployable compression algorithm.
\section*{XXXI.3 Systematic selection from the full OpenAI Math catalogue}
Descriptions of all 372 repository families were considered. Detailed scope and logical-barrier mappings for 42 priorities are in \texttt{CATALOG\_RELEVANCE\_v27.csv}, spanning inverse identification, conic complexity, spectral dynamics, quantum information, statistical memory, PDE computability, and geometry. This is not independent proof of 719 manuscripts.

Specific limits: 095 and 126 motivate studying admissible positive-lift classes; 140 lower-bounds sample cost for a precisely specified streaming regression problem; 229 concerns regular or Poisson-tree reconstruction, with selected Lean coverage narrower than all claimed directions. Families 350/362/365/376 have explicit geometric, physical, and computational assumptions. Linking them to a particular Millennium Problem needs an additional bridge proof absent here.

Key verifiable documents and scopes:
\begin{itemize}
\item \url{https://github.com/openai/math/blob/main/README.md}: catalogue and general incomplete-formalization qualification.
\item \url{https://github.com/openai/math/blob/main/history.md}: three withdrawals and fourteen corrections dated 7 October 2026.
\item \url{https://github.com/openai/math/blob/main/lean/docs/095.md}: positive cones and PSD-lift restrictions.
\item \url{https://github.com/openai/math/blob/main/lean/docs/126.md}: selected formalization does not cover the full claimed exponential lower bound.
\item \url{https://github.com/openai/math/blob/main/lean/docs/140.md}: narrowly specified streaming regression.
\item \url{https://github.com/openai/math/blob/main/lean/docs/229.md}: supercritical reconstruction.
\item \url{https://github.com/openai/math/blob/main/lean/docs/269.md}: claimed stability differs from the selected proof fragment.
\item \url{https://github.com/openai/math/blob/main/lean/docs/350.md}: nodal restrictions and metric dependence.
\item \url{https://github.com/openai/math/blob/main/lean/docs/362.md}: a single-species kinetic model.
\item \url{https://github.com/openai/math/blob/main/lean/docs/365.md}: the formal result is not the full smooth-metric reconstruction claim.
\item \url{https://github.com/openai/math/blob/main/lean/docs/376.md}: forced Navier--Stokes flows, not the unforced Clay problem.
\end{itemize}
\note{Worldwide novelty status.} Matrix analysis yields the exact spectral expression, not an independent natural law. The joint problem, integration into operational passports, and reproducible multi-interface comparison are new within this manuscript. Historical priority for joint SVD/multitarget factorization theorems is not asserted without outside bibliographic expertise. Experimental superiority over modern LLM compressors, complete quantum gravity, and Millennium Problem solutions have not been achieved.
\chapter{Exact Robustness of Spectral Memory against Classical Approximation}
\section*{XXXII.1 Operational passport and central problem}
Sections XI and XXVI separated exact reproduction of periodic quantum responses from positive realization. But absence of an exact $m$-state model alone does not lower-bound approximation error: horizon, norm, and admissible transitions must be fixed. This section closes that gap for a class of periodic signals.

Fix $\mathcal P_{n,m}^{\infty}$: discrete $k\in\N_0$, one fixed binary effect, preparation $\mu\in\Delta_{m-1}$, any homogeneous column-stochastic $M\in\R_{\ge0}^{m\times m}$, effect $e\in[0,1]^m$, and $g(k)=e^\top M^k\mu$. Error is $\|f-g\|_\infty=\sup_{k\ge0}|f(k)-g(k)|$. Unlimited external memory, $k$-dependent transitions, and autonomous clocks with at least $n$ states are excluded. Preparation and effect may be reselected for each model. We prove optimal approximate as well as exact bounds in this class.
\statement{Lemma}{XXXII.1}{Asymptotic periodicity of finite Markov realizations} Every admissible response has an integer $L\ge1$ and an $L$-periodic $h(k)$ with $g(k)-h(k)\to0$. One may choose $L$ as the least common multiple of periods of closed irreducible classes. Each period is at most its class's state count and hence at most $m$.
\begin{proof}Canonical Markov form separates transient and closed irreducible classes. Transients decay. A closed period-$d$ class has only $d$th roots of unity on the unit circle, with semisimple peripheral spectrum. All other modes have modulus below one and decay, including polynomial factors from Jordan blocks. The remaining sum has period dividing $L$. This is standard Perron--Frobenius theory applied only to finite stochastic matrices.\end{proof}
\statement{Theorem}{XXXII.2}{Exact uniform cost of insufficient memory at prime period} Let prime $p\ge3$ and any $p$-periodic $f:\Z\to[0,1]$ have range $R_f=\max_{0\le j<p}f(j)-\min_{0\le j<p}f(j)$. Then over finite homogeneous classical realizations with $1\le m<p$,
\begin{equation}\inf_{\substack{1\le m<p,\ M,\mu,e\\M\text{ stochastic}}}\sup_{k\ge0}|f(k)-e^\top M^k\mu|=\frac{R_f}2.\tag{113}\end{equation}
The lower bound includes irreducible, periodic, or nonergodic competitors; the one-state midpoint response attains it.
\begin{proof}Choose $L$ from XXXII.1. No class period can be divisible by $p$, since it is at most $m<p$, so $\gcd(L,p)=1$. For each $a\bmod L$ and $j\bmod p$, the Chinese remainder theorem gives arbitrarily large $k$ with those residues. Along them $g(k)\to h(a)$ and $f(k)=f(j)$. If uniform error is $\varepsilon$, then $|f(j)-h(a)|\le\varepsilon$ for every $j$. Any real value is at distance at least $R_f/2$ from one extreme. The one-state model $M=(1)$, $\mu=(1)$, $e=(\max f+\min f)/2$ attains equality.\end{proof}
\statement{Corollary}{XXXII.3}{Qubit signal: limiting error tends to one half} For odd prime $p$, $f_p(k)=\sin^2(\pi k/p)$ has a qubit realization, infinite Hankel matrix of rank 3, and exact homogeneous classical memory $p$. Smaller classical models satisfy
\begin{equation}\inf_{m<p}\|f_p-g\|_\infty=\tfrac12\cos^2\frac\pi{2p}\longrightarrow\tfrac12\quad\text{as primes }p\to\infty.\tag{114}\end{equation}
For binary measurements, TV is numerically the same error: $\operatorname{TV}(\operatorname{Ber}(f),\operatorname{Ber}(g))=|f-g|$.
\begin{proof}The minimum is zero and the maximum $\cos^2(\pi/(2p))$. Apply XXXII.2; the exact quantum and classical realizations were proved in XXVI.\end{proof}
\statement{Corollary}{XXXII.4}{Prime Singer periods: a robust quadratic separation} If $q$ is a prime power, $p=q^2+q+1$ is prime, $d=q+1$, and $P_q(k)$ equals 1 at multiples of $p$ and $q/(q+1)^2$ elsewhere, then its dimension-$d$ quantum unitary realization coexists with
\begin{equation}\inf_{m<p}\|P_q-g\|_\infty=\frac12\left(1-\frac q{(q+1)^2}\right).\tag{115}\end{equation}
For $(q,p)=(2,7),(3,13),(5,31)$, the values are $7/18,13/32,31/72$.
\begin{proof}Apply XXXII.2 to the two-level range. The rational values follow directly. Arbitrary composite periods require extra hypotheses.\end{proof}
\section*{XXXII.2 Composite periods: a universal primitive-harmonic barrier}
\statement{Theorem}{XXXII.5}{Fourier certificate of uniform error} Let bounded $n$-periodic $f:\Z\to\R$, $n\ge3$, have $\gcd(a,n)=1$, $\zeta=e^{2\pi\ii a/n}$, and $\widehat f(a)=n^{-1}\sum_{j=0}^{n-1}f(j)\zeta^{-j}$. Every finite homogeneous classical realization with $m<n$ obeys
\begin{equation}\|f-g\|_\infty\ge|\widehat f(a)|.\tag{116}\end{equation}
In particular, $\inf_{m<n}\|\sin^2(\pi k/n)-g\|_\infty\ge1/4$ for every $n\ge3$.
\begin{proof}Consider $T^{-1}\sum_{k=0}^{T-1}g(k)\zeta^{-k}$. Peripheral eigenvalue orders are bounded by their closed class sizes and hence by $m<n$, so none equals the primitive order-$n$ root $\zeta$. Their modulated averages tend to zero, as does the decaying part. For periodic $f$, the limit is $\widehat f(a)$. The triangle inequality gives (116). For sine squared, the three distinct frequencies have coefficients $1/2,-1/4,-1/4$, hence $\widehat f_n(1)=-1/4$.\end{proof}
\section*{XXXII.3 Continuous time, finite horizons, and verifiable data}
\statement{Theorem}{XXXII.6}{Continuous-time classical barrier} For a continuous nonconstant periodic $f:[0,\infty)\to[0,1]$ of range $R_f$, all finite-dimensional homogeneous continuous-time classical generators and admissible $e,\mu$ satisfy
\begin{equation}\inf_{m,Q,\mu,e}\sup_{t\ge0}|f(t)-e^\top e^{tQ}\mu|=R_f/2.\tag{117}\end{equation}
For $f(t)=\sin^2(\omega t)$, $\omega\ne0$, the optimal error is $1/2$, with no bound on the finite classical state count.
\begin{proof}Finite continuous-time closed classes have no nonzero imaginary generator eigenvalues, since $e^{tQ}$ is aperiodic for $t>0$. Thus the response tends to a constant. A continuous periodic signal reaches both extrema infinitely often, forcing error at least half its range. The one-state midpoint model attains it.\end{proof}
\statement{Proposition}{XXXII.7}{Exact one-state approximation cost on a finite horizon} For $f_n(k)=\sin^2(\pi k/n)$ and $T\in\N_0$,
\begin{equation}\inf_{c\in[0,1]}\max_{0\le k\le T}|f_n(k)-c|=\tfrac12\sin^2\left(\frac{\pi\min(T,\lfloor n/2\rfloor)}n\right).\tag{118}\end{equation}
For $T\le n/2$, this is at most $\pi^2T^2/(2n^2)$.
\begin{proof}The prefix minimum is $f_n(0)=0$ and its maximum occurs at $\min(T,\lfloor n/2\rfloor)$. The best constant is the range midpoint, with half-range error. Use $\sin x\le x$ for the last bound.\end{proof}
\statement{Remark}{XXXII.8}{Infinite certificates are not finite experiments} XXXII.2 and XXXII.6 control all future times. Their gaps do not imply the same power for a universal finite experiment without competitor mixing-time information. For $n\gg T$, one state approximates the first $T$ values with error of order $(T/n)^2$. If two known Bernoulli models have gap $\Delta>0$ at a specified time, $N\ge2\Delta^{-2}\log(4/\alpha)$ independent trials per model suffice by Hoeffding for both estimates to be within $\Delta/2$ with probability at least $1-\alpha$. Existence of a suitable unknown time does not supply an algorithm to find it.
\section*{XXXII.4 The remaining conic question and an exact geometric passport}
Revision XXVI proved $\rank_+W^{(6)}=5$, but normalized probability memory for six preparations and six binary effects may differ. The following geometric reduction avoids identifying it heuristically with nonnegative rank.
\statement{Proposition}{XXXII.9}{Exact cube embedding of six probability columns} Let $w_j=(W^{(6)}_{0j},W^{(6)}_{1j},W^{(6)}_{2j})^\top$, $j\in\Z_6$. A normalized $m$-state classical model reproduces all six binary measurements exactly if and only if some $x_1,\ldots,x_m\in[0,1]^3$ satisfy
\begin{equation}\{w_0,\ldots,w_5\}\subseteq\operatorname{conv}\{x_1,\ldots,x_m\}.\tag{119}\end{equation}
Explicitly,
\begin{align*}w_0&=(0,\tfrac14,\tfrac34),&w_1&=(\tfrac14,0,\tfrac14),&w_2&=(\tfrac34,\tfrac14,0),\\w_3&=(1,\tfrac34,\tfrac14),&w_4&=(\tfrac34,1,\tfrac34),&w_5&=(\tfrac14,\tfrac34,1).
\end{align*}
They lie in the cube plane $x-y+z=1/2$.
\begin{proof}Given preparations $\mu_j$ and first-three effect rows $a_i\in[0,1]^m$, set $x_k=(a_{0k},a_{1k},a_{2k})$. Then $w_j=\sum_k(\mu_j)_kx_k$. Conversely, convex weights $t_{kj}$ define preparations, and coordinates of $x_k$ define the first effects. Set $a_{i+3}=1-a_i$ for $i=0,1,2$; these reproduce all remaining rows because $W^{(6)}_{i+3,j}=1-W^{(6)}_{ij}$. Coordinates and the plane equation follow by substitution.\end{proof}
\statement{Remark}{XXXII.10}{Proved and unproved} The established chain $5\le\rank_+W^{(6)}\le\mu_{\mathrm{normalized}}\le6$ reduces the question to whether five cube points can contain these six rational points in their convex hull. Nonexistence would prove normalized cost six; an explicit rational embedding would prove five. Uncertified nonlinear optimization proves neither.
\section*{XXXII.5 Sources and historical novelty}
The formulas use classical finite-chain periodicity and Fourier analysis; independent global priority of the particular minimax claims is unestablished. Spectral inputs may be checked in \textit{A Tutorial on the Spectral Theory of Markov Chains}, \url{https://direct.mit.edu/neco/article/35/11/1713/117578/A-Tutorial-on-the-Spectral-Theory-of-Markov-Chains}; \textit{Minimal positive realizations: A survey}, \url{https://www.sciencedirect.com/science/article/pii/S0005109822002758}; Chistikov, Kiefer, Maru\v si\'c, Shirmohammadi, Worrell, \textit{On Restricted Nonnegative Matrix Factorization}, \url{https://arxiv.org/abs/1605.07061}; and OpenAI Math families 140, 229, 376, \url{https://github.com/openai/math}, whose hypotheses have not been transferred here without proof.

New within this monograph are the exact uniform small-classical-memory approximation bound, the specific cube formulation of normalized factorization, and physically interpretable limits of finite measurements. None alone proves a new natural law, quantum gravity, or a Millennium Problem.
\chapter{Exact Memory Thresholds for Composite Periods and an Audit of Mathematical Transfers}
\section*{XXXIII.1 Passport and problem}
For distinct primes $p,q$, put $n=pq$ and $f(k)=a+b\one_{\{n\mid k\}}$, $a,b\ge0$, $a+b\le1$, $k\in\N_0$. Compare $g(k)=e^\top M^k\mu$ with a column-stochastic matrix of order at most $m$, probability column $\mu$, and $e\in[0,1]^m$, in uniform error over all $k\ge0$. Finite samples do not imply these infinite-horizon bounds.
\statement{Theorem}{XXXIII.1}{Exact threshold for two coprime clocks} For $n=pq\ge6$,
\begin{equation}\varepsilon_m(f)=\begin{cases}b/2,&1\le m<p+q,\\\max\{b/4,(b-a)/3\},&m=p+q.\end{cases}\tag{120}\end{equation}
Finite homogeneous Markov models attain both values; first improvement over the best constant requires exactly $p+q$ states.
\begin{proof}Every finite-chain response tends to a periodic $h$ of period $L$, the lcm of closed-class periods. If $n\nmid L$, along multiples of $L$ the residues modulo $n$ include zero and nonzero values infinitely often, so $f$ takes both $a$ and $a+b$ while $g\to h(0)$. Thus error is at least $b/2$, attained by $a+b/2$.

If $n\mid L$, class periods must include a multiple of $p$ and of $q$. If in one class, at least $pq$ states are needed; if in two, at least $p+q$. Since $pq>p+q$ for distinct primes, $m<p+q$ is impossible. At $m=p+q$, the model must consist of deterministic $p$- and $q$-cycles with no transients. Thus $h(k)=u(k\bmod p)+v(k\bmod q)$, $u,v\ge0$.

By Chinese remaindering, choose values $h_{00},h_{01},h_{10},h_{11}$ at $(0,0),(0,j),(i,0),(i,j)$ with $i,j\ne0$. Additivity gives $h_{00}+h_{11}=h_{01}+h_{10}$. Error at most $\varepsilon$ implies $h_{00}\ge a+b-\varepsilon$, $h_{11}\ge a-\varepsilon$, $h_{01},h_{10}\le a+\varepsilon$, and $h_{11}\ge0$. Hence $b\le4\varepsilon$ and $b-a\le3\varepsilon$. The earlier $b/2$ bound when $n\nmid L$ is no weaker.

For attainment, give each independent cycle half the initial mass at residue zero. Set ordinary effects to $t$ and zero-residue effects to $t+2u$, giving
\begin{equation}g(k)=t+u(\one_{\{p\mid k\}}+\one_{\{q\mid k\}}).\tag{121}\end{equation}
For $b\le4a$, take $t=a-b/4$, $u=b/2$, yielding error $b/4$. For $b\ge4a$, take $t=0$, $u=(2a+b)/3$, yielding $(b-a)/3$. The assumptions ensure $0\le t\le t+2u\le1$, so the effects and transitions are physically normalized.\end{proof}
\statement{Corollary}{XXXIII.2}{Composite Singer period 21} For $f(k)=4/25+(21/25)\one_{\{21\mid k\}}$, $\varepsilon_9=21/50=0.42$ and $\varepsilon_{10}=17/75\approx0.226667$. A ten-state realization is $g(k)=(29/75)(\one_{\{3\mid k\}}+\one_{\{7\mid k\}})$, from cycles of lengths 3 and 7, initial masses $1/2$, special effects $58/75$, and all others zero. Its uniform error is $17/75$; periodicity reduces verification to 21 consecutive times.
\statement{Remark}{XXXIII.3}{A discrete threshold, not a linear memory budget} Moving from 9 to 10 states lowers error from 0.42 to 0.2267 without changing the original quantum signal. This concerns one two-level periodic probability over an infinite horizon, not universal LLM compression, an experimentally discovered physical law, or a Millennium Problem.
\section*{XXXIII.2 Correction to the cube passport of XXVIII}
In XXXII.9, the candidate five-state extension's vertices belong to the entire cube $[0,1]^3$. The six target points lying in $x-y+z=1/2$ does not force enclosing-polytope vertices into that plane. A proof based only on pentagons in that plane would be incomplete. The bound $5\le\mu_{\mathrm{normalized}}\le6$ persists; numerical optimization without a global certificate does not settle it.
\section*{XXXIII.3 Transferable theorems versus inadmissible physical inferences}
\begin{itemize}
\item \textbf{095.} A particular hyperbolicity cone lacking a finite LMI representation does not imply that finite-dimensional density-matrix sets lack semidefinite representations; they have one directly.
\item \textbf{140.} Memory-limited noiseless Gaussian-regression sample bounds are not universal LLM-compression limits. A fabricated general ``Shannon--Witsenhausen threshold'' cannot replace their hypotheses.
\item \textbf{351.} The catalogue concerns extension of Ricci flow on closed four-manifolds under bounded scalar curvature, not a proof of three-dimensional surgery theory.
\item \textbf{365.} Joint recovery of a smooth metric and connection from boundary data is distinct from nonidentifiability of bounded measurable scalar conductivities. Neither infers Planck-scale metric defects.
\item \textbf{265, 269.} Lattice/spherical spectral gaps do not automatically establish a four-dimensional continuum Yang--Mills mass gap.
\item \textbf{350.} Smooth-surface nodal-measure upper bounds do not extend unconditionally to every dimension; the catalogue separately claims counterexamples in dimensions three and four.
\item \textbf{Statistical passport.} A sample count obtained directly from Hoeffding is sufficient for a specified test, not a universal necessary lower bound.
\item \textbf{376.} Computation by specially forced flow does not solve unforced Clay regularity.
\end{itemize}
Here catalogue status follows family descriptions and stated formalization scope; inclusion of a preprint is not independent scientific review. See \url{https://github.com/openai/math} and \url{https://github.com/openai/math/blob/main/history.md}.
\section*{XXXIII.4 Status and reproducibility}
Theorem XXXIII.1 does not rely on formal catalogue claims. The archive includes exact rational checks for periods 6, 15, 21, 35 and a linear-programming check of the four-point certificate. The latter is numerical corroboration, not a substitute for the analytic proof. Historical novelty has not been independently established against the full literature.
\chapter{A Unified Memory Threshold for Arbitrary Composite Periods}
\note{Principal result of this revision.} Sections XXVIII--XXIX established prime and two-prime bounds. The following treats every $n\ge2$, including products of three or more prime powers. Its complete self-contained English proof appears at the beginning of the manuscript; the translated version of the Russian restatement is retained here for mathematical continuity.
\statement{Theorem}{XXXIV.1}{First exact approximation-memory threshold} Let $n=\prod_{i=1}^r q_i$, where $q_i=p_i^{\alpha_i}$ are maximal powers of distinct primes, and $K(n)=\sum_iq_i$. For $a,b\ge0$, $a+b\le1$, $f(k)=a+b\one_{\{n\mid k\}}$, and normalized homogeneous $m$-state competitors,
\[\inf_g\sup_{k\ge0}|f(k)-g(k)|=\begin{cases}b/2,&1\le m<K(n),\\\displaystyle\max\left\{\frac{(r-1)b}{2r},\frac{(r-1)b-a}{2r-1}\right\},&m=K(n).\end{cases}\]
Both infima are attained. For $b>0$, $m=K(n)$ is exactly the first budget permitting error below $b/2$.
\begin{proof}A finite-chain response is asymptotically periodic with period dividing $L=\lcm(d_1,\ldots,d_s)$, where $d_j$ are closed-class periods and $m\ge\sum_jd_j$. If $n\mid L$, assign each maximal prime power $q_i$ to a class whose period it divides. Grouping assigned coprime powers gives $\sum_jd_j\ge\sum_iq_i=K(n)$, since the product of assigned coprime integers at least two is at least their sum. Thus for $m<K(n)$, $n\nmid L$. Along multiples of $L$, the target takes both levels infinitely often while the competitor tends to one value, forcing half-range error. The midpoint constant attains it.

At $m=K(n)$ and $n\mid L$, equality in the state bound forces exactly $r$ disjoint deterministic cycles of lengths $q_i$. Their response is additive in independent residues. Restrict each residue to 0 or 1 and denote all-zero value by $H_*$, all-one by $H_0$, and all-zero-except-$i$ by $H_{-i}$. Chinese remaindering realizes these values and additivity yields $\sum_iH_{-i}=(r-1)H_*+H_0$. Error $\varepsilon$ gives $H_*\ge a+b-\varepsilon$, $H_{-i}\le a+\varepsilon$, $H_0\ge\max\{0,a-\varepsilon\}$, hence the two asserted lower bounds. When $n\nmid L$, the stronger half-range bound applies.

For attainment, give each cycle mass $1/r$, ordinary effect $t$, and zero-vertex effect $t+ru$, yielding $g(k)=t+u\sum_i\one_{\{q_i\mid k\}}$. If $a\ge(r-1)b/(2r)$, choose $t=a-(r-1)b/(2r)$, $u=b/r$; otherwise choose $t=0$, $u=(2a+b)/(2r-1)$. All effects lie in $[0,1]$ and substitution gives the claimed maximum.\end{proof}
\statement{Remark}{XXXIV.2}{Strictly limited meaning} The theorem concerns one fixed binary observation, homogeneous classical Markov dynamics, and the entire infinite horizon. It does not imply a finite-time experimental advantage, universal LLM compression, or a new physical equation. Periodic methods have antecedents in unary automata, including Bianchi--Mereghetti--Palano--Pighizzini (2011) and Bianchi--Pighizzini (2012). Independent priority of the exact minimax formula is unestablished.
\note{Examples.} At $a=4/25$, $b=21/25$, $n=21=3\cdot7$ gives $K=10$, $\varepsilon_{10}=17/75$; $n=30=2\cdot3\cdot5$ gives $K=10$, $\varepsilon_{10}=38/125$. For $n=210=2\cdot3\cdot5\cdot7$, $a=0$, $b=1$, $K=17$ and $\varepsilon_{17}=3/7$.
\note{Reproducibility.} \texttt{checks/verify\_v30\_general\_crt.py} performs 64 independent linear optimizations and exact rational checks for integer periods, saved in \texttt{RESULTS\_GENERAL\_CRT\_v30.csv}. These check special cases, not the general analytic proof.
\setcounter{equation}{121}


% 06_fibonacci_witnesses.tex
\chapter{The Fibonacci Matrix, Recurrent Transformers, the Qubit, and the Geometry of $S^3$}
\note{Scientific setting.} Modern Gated DeltaNet, Kimi Delta Attention, and preconditioned recurrent mechanisms treat the memory operator as a controlled sequence of low-rank corrections. Rotary Position Embedding (RoPE) encodes relative positional phases by rotations. We construct one exact two-dimensional example connecting the Fibonacci matrix to a delta rule, an admissible qubit filter, and the Hopf fibration. This is neither an identity for an arbitrary Transformer nor an empirical result on a pretrained LLM.
\statement{Definition}{XXXV.1}{Specification of the Fibonacci recurrent cell}
Fix $F=\left(\begin{smallmatrix}1&1\\1&0\end{smallmatrix}\right)$, $\varphi=(1+\sqrt5)/2$, $K=F/\varphi$, discrete time $t\in\N_0$, and the state space $\C^2$. For the real recurrent cell, the operator $G=K^2$ is admissible; in the quantum branch, $K$ is applied conditionally, with the successful filter outcome explicitly recorded. The permitted scalar observations and norms are specified in each theorem.
\statement{Theorem}{XXXV.2}{Two Fibonacci steps are an exact delta rule}
Let $v_-$ be a unit eigenvector of $K$ with eigenvalue $-\varphi^{-2}$. Then
\begin{equation}F^2=F+I,\qquad K^2=I-\beta v_-v_-^\top,\qquad \beta=1-\varphi^{-4}\simeq0.8541019663.\tag{122}\end{equation}
In particular, the update $h_{t+1}=K^2h_t+\beta v_-x_t$ is an exact two-dimensional realization of a static delta correction with key $v_-$ and scalar write input $x_t$.
\begin{proof}
The characteristic polynomial $\det(xI-F)=x^2-x-1$ gives eigenvalues $\varphi$ and $-\varphi^{-1}$. The normalized symmetric operator $K$ has spectrum $\{1,-\varphi^{-2}\}$, and its square has spectrum $\{1,\varphi^{-4}\}$. The spectral theorem gives the asserted identity. The expression with an external input is the definition of a delta-rule step with a fixed key; it does not establish equivalence of entire Gated DeltaNet or Transformer models.
\end{proof}
\statement{Proposition}{XXXV.3}{Exact qubit embedding: two instrument branches}
The normalized operator $K$ is nonunitary, but $K^\dagger K\preceq I$. Set $B=(I-K^2)^{1/2}$ and
\[U=\begin{pmatrix}K&B\\B&-K\end{pmatrix}.\]
Then $U\in SO(4)\subset SU(4)$. Postselecting the ancillary qubit on $|0\rangle$ implements $K$ with success probability
\[\varphi^{-4}\le p_{\rm succ}(\psi)=\langle\psi,K^2\psi\rangle\le1.\]
The complete measurement must also include the unsuccessful branch with operator $B$; without it, the operation is not trace preserving. $K$ must not be presented as an ordinary Bloch-sphere rotation.
\begin{proof}
The symmetric matrices $B,K$ commute and satisfy $B^2+K^2=I$, so $U^2=I$. On each of the two spectral subspaces, the corresponding two-dimensional ancilla block has determinant $-1$; the full determinant is $+1$. The probability bounds are the extreme eigenvalues of $K^2$.
\end{proof}
\statement{Proposition}{XXXV.4}{The spinor sphere and the absence of a cosmological deduction}
Normalized pure one-qubit state vectors form $S^3=\{(z_1,z_2)\in\C^2:|z_1|^2+|z_2|^2=1\}$. Identifying global phases gives $\mathbb{CP}^1\cong S^2$; the standard projection $S^3\to S^2$ is the Hopf map. The postselected dynamics $\psi\mapsto K\psi/\|K\psi\|$ induces a homeomorphism, indeed a holomorphic automorphism of $\mathbb{CP}^1$, of topological degree $+1$, since $K$ is invertible. Its attracting spectral mode is nevertheless not a unitary rotation. The Chern class of the Hopf bundle does not determine the decay rate $\varphi^{-2}$.
\statement{Theorem}{XXXV.5}{Hankel rank and a strict classical negative control}
Write $\lambda=-\varphi^{-2}$ and $\chi=(v_++v_-)/\sqrt2$. The scalar amplitude $c_t=\langle\chi,K^t\chi\rangle$ satisfies
\begin{equation}c_t=\tfrac12(1+\lambda^t),\quad \rank(c_{i+j})_{i,j\ge0}=2,\quad
\det\begin{pmatrix}c_0&c_1\\c_1&c_2\end{pmatrix}=\frac{(1-\lambda)^2}{4}>0.\tag{123}\end{equation}
There is a homogeneous two-state classical stochastic matrix
\begin{equation}M=\frac12\begin{pmatrix}1+\lambda&1-\lambda\\1-\lambda&1+\lambda\end{pmatrix}\tag{124}\end{equation}
with $c_t=(M^t)_{11}$. Thus, under this scalar specification, the quantum filter has no memory advantage over two classical states.
\begin{proof}
The spectral theorem gives a sum of two exponentials and hence Hankel rank at most two; the nonzero minor proves the lower bound. $M$ is nonnegative and stochastic with eigenvalues $1,\lambda$. Decomposing $e_1$ into symmetric and antisymmetric modes yields $e_1^\top M^te_1=(1+\lambda^t)/2$.
\end{proof}
\statement{Proposition}{XXXV.6}{Joint RoPE--$SU(2)$ phase spectrum}
Let $R_t=\diag(e^{-it\theta/2},e^{it\theta/2})\in SU(2)$ act in the eigenbasis of $K$. Then $\langle R_sq,R_tk\rangle=\langle q,R_{t-s}k\rangle$, and the two-mode phase signal
\begin{equation}C_t=\langle\chi,(KR_1)^t\chi\rangle=\frac{e^{-it\theta/2}+\lambda^te^{it\theta/2}}2\tag{125}\end{equation}
has complex Hankel rank exactly two. This is a controlled extension of RoPE's relative-phase principle, not a claim that all RoPE Transformers are qubits.
\begin{proof}
The group identity $R_s^\dagger R_t=R_{t-s}$ gives the first assertion. Simultaneous diagonalization gives two spectral exponentials with distinct pole moduli $1$ and $\varphi^{-2}$, so Vandermonde nondegeneracy gives rank two.
\end{proof}
\statement{Proposition}{XXXV.7}{Gauge restriction imposed by positional rotation}
For a nontrivial real two-dimensional rotation $R_\theta$, $\sin\theta\ne0$, every real matrix commuting with it has the form $G=aI+bJ$, where $J=\left(\begin{smallmatrix}0&-1\\1&0\end{smallmatrix}\right)$. With invertibility imposed, this is the group $\C^\times$. Positional encoding therefore restricts the admissible basis gauge; corresponding group restrictions have been studied for large Transformers (Wang--Wang, PMLR, 2026).
\begin{proof}
Since $R_\theta=\cos\theta I+\sin\theta J$, its commutant equals that of $J$. For $G=\left(\begin{smallmatrix}a&b\\c&d\end{smallmatrix}\right)$, $GJ=JG$ immediately gives $a=d$, $b=-c$.
\end{proof}
\statement{Proposition}{XXXV.8}{Fibonacci anyons: a common counting operator, different physical actions}
The fusion rule $\tau\otimes\tau=1\oplus\tau$ gives the channel-counting matrix $N_\tau=\left(\begin{smallmatrix}0&1\\1&1\end{smallmatrix}\right)=PFP$, with $P=\left(\begin{smallmatrix}0&1\\1&0\end{smallmatrix}\right)$. Its Perron--Frobenius radius is $\varphi$. However, $N_\tau$ is not a unitary anyon-braiding operator: the latter requires separate unitary categorical $F$- and $R$-matrices. Equality of counting matrices does not imply that a language model realizes anyonic topology.
\section*{OpenAI Math catalog directions and their verifiable status}
\begin{longtable}{p{.10\textwidth}p{.42\textwidth}p{.39\textwidth}}
\toprule Family&Substantive connection&Inference not justified\\\midrule
144&Koopman spectrum and dynamical components invisible to scalar observation&A simple spectrum does not imply finite Transformer memory.\\
146, 151&Dynamical entropy and counterexamples to transferring spectral information to entropy&The golden ratio alone does not determine model entropy.\\
148&Dimension of self-similar measures under specific hypotheses&There is no automatic fractional ``Hankel rank.''\\
229&Reconstruction on branching trees&Rank along one ray does not determine collective reconstructibility.\\
276&Classical capacity of qubit amplitude-damping channels&Physical channel capacity is not software-layer accuracy.\\
095, 116&Conic representations and formal verification of noncommutative identities&No universal compression theorem for all Transformers follows.\\\bottomrule
\end{longtable}
\section*{Checks and an engineering experiment}
The executable \path{checks/verify_fibonacci_transformer_bridge_v31.py} checks symbolic identities, 20 zero determinants of $3\times3$ Hankel matrices, two-state classical simulation over 65 time steps, 300 arbitrary pure qubit states, the $4\times4$ unitary embedding, $SU(2)$ relative-phase relations, and negative controls. These are mathematical and synthetic tests, not an experiment on an external pretrained LLM.

A simple counterexample directly refutes universality of a golden forgetting rate: if the true decay coefficient $\rho\ne\varphi^{-4}$, a learned coefficient $\rho$ achieves zero error on a linear test sequence, whereas the fixed golden coefficient generally does not. A real experimental architecture should therefore learn or calibrate memory times and be compared with Mamba, Gated DeltaNet-2, standard RoPE, and ordinary low-rank SSMs under equal computational budgets.

\note{Module references and priority.} Su et al., \emph{RoFormer}, Neurocomputing 568 (2024), \url{https://arxiv.org/abs/2104.09864}; Tumma--Loo--Rus, \emph{Preconditioned DeltaNet}, ICML 2026, \url{https://proceedings.mlr.press/v306/tumma26a.html}; Hatamizadeh--Choi--Kautz, \emph{Gated DeltaNet-2}, 2026, \url{https://arxiv.org/abs/2605.22791}; Wang--Wang, \emph{Gauge Symmetries in Transformer Architectures}, PMLR 282 (2026), \url{https://proceedings.mlr.press/v282/wang26a.html}; Field--Simula, \emph{Introduction to topological quantum computation}, \url{https://arxiv.org/abs/1802.06176}; Malhotra et al., \emph{Tracking Representation Dynamics in LLMs with Persistent Homology}, PMLR 334 (2026), \url{https://proceedings.mlr.press/v334/malhotra26a.html}; OpenAI Math, \url{https://github.com/openai/math}. All exact identities in this section are derived within the chosen specification; independent priority has not been established.

\chapter{Four- and Five-Sample Memory Witnesses without Lean}
\note{Scientific specification.} One qubit with repeated channel $w\mapsto\rho e^{i\theta}w$, independent preparation of $|+\rangle$ before each run, binary effect $|-\rangle\langle-|$, and a homogeneous classical two-hidden-state comparator with a fixed effect. No external clock, changing conditions between steps, or adaptive interventions are allowed. This text does not prove an advantage for a commercial LLM.
\statement{Theorem}{XXXVI.1}{Four times suffice for a dimension test}
Let $0<\rho\le1$, $0<\theta<\pi$, and $f_k=[1-\rho^k\cos(k\theta)]/2$. The quantity
\[J=(f_2-f_1)^2-(f_1-f_0)(f_3-f_2)=\frac{\rho^2\sin^2\theta}{4}(1+\rho^2-2\rho\cos\theta)>0\]
is identically zero for every homogeneous classical two-state response $g_k=e^TM^k\mu$. This follows from the Cayley--Hamilton recurrence for a $2\times2$ stochastic matrix and direct factorization of $J(f)$.
\statement{Corollary}{XXXVI.2}{Finite-horizon simulation cost}
Every admissible classical comparator has at least one $k\in\{0,1,2,3\}$ with $|f_k-g_k|\ge J/8$. With five measured times, define $D=-\det[f_{i+j}]_{0\le i,j\le2}$. Then
\[D=\frac{\rho^2\sin^2\theta}{8}(1+\rho^2-2\rho\cos\theta)^2=\frac{1+\rho^2-2\rho\cos\theta}{2}J>0.\]
Consequently $\max_{0\le k\le4}|f_k-g_k|\ge D/9$. The proofs use Lipschitz bounds 8 and 9 for the polynomial witnesses on the probability cube.
\statement{Corollary}{XXXVI.3}{A sufficient statistical sample}
At confidence $1-\alpha$, it suffices to perform $N\ge200J^{-2}\log(8/\alpha)$ independent binary trials at each of four times, or $N\ge200D^{-2}\log(10/\alpha)$ at each of five times. Reject the two-state hypothesis when $|\widehat J|>2J/5$ or, respectively, $|\widehat D|>9D/20$. These conditions allow arbitrary probabilities for the alternative Markov model but require pre-calibrated $\rho,\theta$ for the target qubit. The budgets are conservative, not optimal.
\statement{Proposition}{XXXVI.4}{Noiseless phase-experiment optimization}
For $\rho=1$, $J$ is maximal at $\cos\theta=-1/3$, with $J_{\max}=16/27$; $D$ is maximal at $\cos\theta=-1/2$, with $D_{\max}=27/32$. Thus the optimal phase depends on the number of measured times and the statistic chosen.
\note{The golden ratio, geometry, and AI.} With $\rho=\varphi^{-2}$, the damped qubit and the Fibonacci Markov matrix have the same modulus of the decaying eigenvalue but different $J$. Equal decay rates therefore do not imply equal operational memory. The $SU(2)$ connection describes the geometry of qubit phase rotation; it proves neither new laws of the Universe nor unconditional compressibility of modern Transformers.
\note{Independent verification.} The calculations are checked by \path{checks/verify_v32_finite_witness.py}. This revision uses neither Lean 4, Lake, nor formal compilation certificates: its mathematical status rests on the stated proofs and exact symbolic arithmetic. The protocol, trial counts, and restrictions are available in the appendix CSV.

\chapter{Five-Sample Constraints on Positive Memory and the Spectral Geometry of Three States}
\note{Specification.} The same source is reset before every run and measured at times $k=0,\ldots,4$. The null class consists of all $g_k=e^TM^k\mu$ with $M\in\R_{\ge0}^{3\times3}$, $\one^TM=\one^T$, $\mu\in\Delta_2$, and $e\in[0,1]^3$. The transition is homogeneous; arbitrary random transitions, degenerate spectra, and initial distributions are allowed. The quantum protocol prepares a fresh qubit, applies the same unitary--dephasing transformation $k$ times, then measures one binary effect. This is not a test of arbitrary nonstationary or controlled classical machines, nor a universal LLM certificate.
\section{A decisive null control for the proposed six-point test}
For $0<\rho\le1$, $0<\theta<\pi$, write $z=\rho e^{i\theta}$ and
\begin{equation}f_k=\frac{1-\rho^k\cos(k\theta)}2=\frac12-\frac14z^k-\frac14\bar z^{\,k}.\tag{126}\end{equation}
The differences $d_k=f_{k+1}-f_k$ contain exactly the exponentials $z^k$ and $\bar z^k$. Therefore
\begin{equation}\det[d_{i+j}]_{0\le i,j\le2}=0.\tag{127}\end{equation}
Differences of every homogeneous classical three-state response also have Hankel rank at most two, since a stochastic matrix has eigenvalue 1. Thus the six-point $3\times3$ difference determinant does not distinguish these classes; its proposed positive signal is refuted by (127). Transition-matrix positivity, rather than linear rank, is needed.
\section{Three universal polynomial constraints from five times}
For arbitrary $x_0,\ldots,x_4\in[0,1]$, set $d_j=x_{j+1}-x_j$ and
\begin{equation}H=d_0d_2-d_1^2,\qquad A=d_0d_3-d_1d_2,\qquad B=d_1d_3-d_2^2.\tag{128}\end{equation}
All three expressions use five observed probabilities, without reconstructing the hidden matrix.
\statement{Theorem}{XXXVII.1}{Finite-sample certificates for three-state positive realizations}
Every homogeneous stochastic response $x_k=e^TM^k\mu$ with three hidden states satisfies
\begin{align}T_1&=H(H+A)\ge0,\tag{129}\\T_2&=H^2+A^2-2HB\ge0,\tag{130}\\T_3&=H^2-HA+A^2-3HB\ge0.\tag{131}\end{align}
These constraints also hold at $H=0$, without diagonalizability or invertibility assumptions on $M$.
\begin{proof}
Factor $\det(tI-M)=(t-1)(t^2-\tau t+\delta)$, where $\tau=\Tr M-1$. For $\nu=(M-I)\mu$, Cayley--Hamilton gives $(M^2-\tau M+\delta I)\nu=0$, hence $d_{k+2}=\tau d_{k+1}-\delta d_k$. Direct elimination yields
\begin{equation}A=\tau H,\qquad B=\delta H.\tag{132}\end{equation}
Entrywise nonnegativity implies $\Tr M=1+\tau\ge0$ and $\Tr(M^2)=1+\tau^2-2\delta\ge0$. Furthermore,
\[\Tr(M^2)=\sum_iM_{ii}^2+\sum_{i\ne j}M_{ij}M_{ji}\ge\sum_iM_{ii}^2\ge(\Tr M)^2/3,\]
so $1-\tau+\tau^2-3\delta\ge0$. Multiplying these three inequalities by $H^2$ and using (132) gives the certificates (129)--(131). At $H=0$, the identities force $A=B=0$, so all certificates vanish. The proof uses neither numerical optimization nor Lean.
\end{proof}
\statement{Theorem}{XXXVII.2}{A strict five-point quantum--classical gap}
For (126), with $c=\cos\theta$,
\begin{align}H&=-\frac{\rho^2\sin^2\theta}{4}(1+\rho^2-2\rho c),\tag{133}\\A&=2\rho cH,\qquad B=\rho^2H,\tag{134}\\T_3(f)&=H^2[(1-\rho c)^2-3\rho^2\sin^2\theta].\tag{135}\end{align}
If $(1-\rho\cos\theta)^2<3\rho^2\sin^2\theta$, no homogeneous classical three-state model reproduces even the first five probabilities. At $(\rho,\theta)=(1,\pi/2)$,
\begin{equation}(f_0,f_1,f_2,f_3,f_4)=(0,\tfrac12,1,\tfrac12,0),\quad(H,A,B)=(-\tfrac12,0,-\tfrac12),\quad T_3=-\tfrac12.\tag{136}\end{equation}
Every classical three-state model $g$ then obeys the robust bound
\begin{equation}\max_{0\le k\le4}|f_k-g_k|\ge\frac1{384}.\tag{137}\end{equation}
\begin{proof}
Substitution of $\cos(k\theta)=T_k(c)$ and polynomial factorization gives (133)--(135). For (136), use $\cos(k\pi/2)\in\{1,0,-1\}$. The negative value contradicts Theorem XXXVII.1. For robustness, put $\eta=\max_k|f_k-g_k|$. Then $|d_j(f)-d_j(g)|\le2\eta$ and $d_j\in[-1,1]$. Each of $H,A,B$ changes by at most $8\eta$ and has modulus at most 2. Comparing terms gives
\begin{equation}|T_3(f)-T_3(g)|\le(32+32+32+96)\eta=192\eta.\tag{138}\end{equation}
Since $T_3(g)\ge0$ and $T_3(f)=-1/2$, (137) follows. The estimate is conservative and is not claimed to be the exact minimax value.
\end{proof}
\section{Reconstructing eigenfrequencies and the trace hierarchy}
When $H\ne0$, five observations determine the sum and product of the comparator's two nonconstant spectral modes:
\begin{equation}\tau=A/H,\qquad\delta=B/H.\tag{139}\end{equation}
Set $s_0=2$, $s_1=\tau$, and $s_{k+2}=\tau s_{k+1}-\delta s_k$. Every three-state stochastic realization must satisfy infinitely many conditions reconstructible from finite-dimensional data:
\begin{equation}1+s_k=\Tr(M^k)\ge0\qquad(k=1,2,\ldots).\tag{140}\end{equation}
These necessary tests all depend on the same five values. Algebraic calculation of $s_k$ from five probabilities must not be confused with stable experimental determination at a distant time: coefficient errors may be amplified as $H\to0$ or $k$ grows.
\statement{Corollary}{XXXVII.3}{Pure rotations not realizable with three classical states}
For $\rho=1$, $0<\theta<\pi$, and $\theta\ne2\pi/3$, no single homogeneous three-state chain exactly reproduces $f_0,\ldots,f_4$. At $\theta=2\pi/3$, a realization exists: a three-cycle initialized at state 0 with effect $(0,3/4,3/4)$.
\begin{proof}
Here $H<0$, so the five observations uniquely determine $\tau=2\cos\theta$, $\delta=1$. A stochastic realization would have spectrum $\{1,e^{i\theta},e^{-i\theta}\}$, requiring $1+2\cos(k\theta)\ge0$ for every $k$. For irrational $\theta/(2\pi)$, density of the rotation orbit gives a $k$ with cosine below $-1/2$. For rational angles the root of unity has order $q\ge3$. If $q\ge4$, one of the $q$ equally spaced phases lies within $\pi/q$ of $\pi$ and has cosine below $-1/2$. Order 3 gives only $\theta=2\pi/3$ in the specified interval and the stated cycle. Orders 1 and 2 are excluded by the angular range.
\end{proof}
\section{Statistical protocol and negative controls}
\statement{Proposition}{XXXVII.4}{A nonoptimal but valid five-point budget}
For the prespecified alternative $(\rho,\theta)=(1,\pi/2)$, measure $f_k$ by independent Bernoulli series of $N$ trials at each of five times and calculate $\widehat T_3$. Reject the three-state hypothesis if $\widehat T_3<-1/4$. Both error probabilities are at most $\alpha$ provided
\begin{equation}N\ge\tfrac12(1536)^2\log\frac{10}{\alpha}\qquad\text{at each time}.\tag{141}\end{equation}
For $\alpha=0.05$, $N=6\,250\,150$ per point suffices, or $5N$ independent measurements in total.
\begin{proof}
Set $\eta=1/1536$. Hoeffding's inequality and a union bound give $\Pr(\max_k|\widehat f_k-f_k|>\eta)\le10e^{-2N\eta^2}\le\alpha$. On the concentration event, (138) bounds the witness change by $192\eta=1/8$. Under the null, its true value is nonnegative and the estimate is at least $-1/8$. Under the alternative, the true value is $-1/2$ and the estimate is at most $-3/8$. The threshold $-1/4$ separates them. The sample size is deliberately extremely conservative; simulations with smaller $N$ do not replace the guarantee.
\end{proof}
\note{Exact controls.} At qubit phase $2\pi/3$, a three-state cycle reproduces the signal: an arbitrary quantum model need not violate the certificate. At $\pi/2$, the first trace $1+\tau=1$ reveals no contradiction, but the second trace $1+\tau^2-2\delta=-1$ and the strengthened JLL inequality detect it using five points. Any particular nonstationary or input-controlled comparator needs a new specification; the theorem does not transfer automatically.
\begin{figure}[ht]
\centering
\setlength{\unitlength}{.8mm}
\begin{picture}(140,85)
\put(50,40){\circle{60}}\put(10,40){\vector(1,0){85}}\put(50,3){\vector(0,1){78}}
\put(35,14){\line(0,1){52}}\put(35,66){\line(5,-3){45}}\put(35,14){\line(5,3){45}}
\put(50,70){\circle*{2}}\put(35,66){\circle*{2}}
\put(96,39){$\operatorname{Re}\lambda$}\put(52,79){$\operatorname{Im}\lambda$}
\put(53,70){$e^{i\pi/2}$: violates $T_3\ge0$}\put(10,67){$e^{2\pi i/3}$}
\put(42,36){\small admissible}\put(43,31){\small $3\times3$ spectra}\put(73,10){$\rho=1$}
\end{picture}
\caption{Necessary spectral region for the complex conjugate pair of a three-state stochastic matrix.}
\end{figure}
\note{Replication, historical status, and AI.} The script \path{checks/verify_v33_three_state.py} verifies the symbolic identities in SymPy, 64 exact rational and 1200 random stochastic null controls, 250 quantum-protocol simulations, and 250 classical controls with 1200 tosses at each time; results are stored in CSV and JSON. A numerical run is not a physical experiment. The method belongs to classical nonnegative-matrix theory, Cayley--Hamilton theory, and trace inequalities; compare work on Karpelevich eigenvalue regions, Johnson--Loewy--London inequalities, and temporal dimension witnesses (Vieira--Milz--Vitagliano--Budroni, \emph{Quantum} 8, 1224 (2024); the 2026 review of stochastic-matrix spectra). Priority for this particular polynomial has not been established. Application to recurrent networks or Transformers requires an explicit reduction to a single homogeneous finite stochastic model; input-dependent attention and delta rules generally do not meet this assumption. Lean is not used.
\setcounter{equation}{141}


% 07_reversibility_trees.tex
\chapter{A Critical Register of 42 Claims and a Spectral Specification of Reversibility}
\note{Purpose.} The register supplied by the user contains exactly 42 items. ``Law,'' ``invariant,'' and ``new fundamental discovery'' are not interchangeable: a mathematical identity may have long been known, and a proof within a specified model class does not establish physical universality. The register below maps admissible conclusions; it does not assert 42 independent discoveries of worldwide novelty. Some formulas from the original list were absent from the supplied text. We do not invent them, but compare their meaning with the full manuscript preceding this revision.
\section{Classification of the forty-two specifications}
Notation: A denotes a proved mathematical core within an explicitly restricted model class; C a conditional or special assertion whose full specification must be recovered before application; R a formulation in the supplied list requiring mathematical correction. Even A does not imply independent priority verification. For historically known facts (for example, Sch\"utzenberger--Fliess, Perron--Frobenius, and the spectral theorem), the number of new discoveries is not the number of rows.
\begingroup\small
\begin{longtable}{p{.04\textwidth}p{.29\textwidth}p{.04\textwidth}p{.53\textwidth}}
\toprule No.&Invariant / claim&Status&Specification or necessary restriction\\\midrule\endhead
1&Linear Hankel realization rank&A&Finite-rank rational series; complete word Hankel; linear rather than positive memory.\\
2&Static nonnegative factorization&A&Probability table only; stochastic preparation/effect normalization is a separate condition.\\
3&Sperner bound for sine-squared matrices&A&Exact diagonal zeros, positive off-diagonal entries; finite $n$.\\
4&Robust orthant gap&C&Requires specified $n,m,\delta$ and maximum-entry error; not a universal constant.\\
5&Six-point nonnegative rank five&A&Static rank only; no five-state normalized model established.\\
6&Asymptotic linear/positive/dynamic split&C&Separate operational comparisons and the $n\to\infty$ regime.\\
7&Observable-algebra dimension jump&C&Specific canonical $P,Q$ and branching parameter only.\\
8&Cyclic commutator rank&C&Canonical transition pair; noncommutation need not imply an observable effect.\\
9&Exact stochastic lumpability&A&Deterministic quotient and action-wise intertwining; not for arbitrary $L$.\\
10&Distance to rank-deficient matrices&C&Fixed finite matrix, explicit entrywise norm, and algebraic assumptions.\\
11&Uniform spectral stability&C&Only the specified canonical Hankel family; not arbitrary systems.\\
12&Dobrushin telescoping&C&Shared hidden-state interface, specified kernels, and output horizon.\\
13&Rank/mixing joint constraint&C&Requires complete finite tests and robust-rank separation.\\
14&Order-sensitive commutator witness&C&A specified preparation/effect pair and calibrated intervention order.\\
15&Deterministic coarse-graining barrier&C&Fixed microscopic processes and preserved effect; no global claim.\\
16&Stochastic coarse-graining barrier&C&Dimensional $n,d$ bound in the canonical family; no uniform gap.\\
17&Stochastic lift with zero defect&A&Constructed compatible microscopic $P=R\bar P L$; not fixed $P$.\\
18&Trace-preserving scaling no-go&A&Normalized quantum/classical states and a linear trace-preserving map.\\
19&Neutrino Hankel inertia&C&Nondegenerate frequencies, nonzero mixing amplitudes, and specified sampling.\\
20&Peripheral spectrum of classical chains&C&For purely periodic finite-response signals; the root-of-unity condition is necessary, and period cycles suffice for valid periodic probabilities.\\
21&Neutrino coefficient triangle&C&Rank-one amplitudes, known labels, and phase convention.\\
22&CP antiunitary ambiguity&C&Vacuum model and simultaneous orientation transformation; not generic matter propagation.\\
23&Singer quadratic memory scale&C&Specified unitary return with a Singer difference set; not a universal quantum advantage.\\
24&Depolarizing primitive memory&C&Homogeneous depolarizing channel with $\eta>0$ and $p$ visible distinct modes.\\
25&Vanishing-coherence nonuniformity&C&Specified finite horizon and sample resources; distinguish from exact rank.\\
26&Reversible-chain correlation monotonicity&R&False at every lag: a two-state periodic chain has $C(k)=(-1)^k$; valid for even lags and alternating finite differences.\\
27&Bispectrum resolves reflection&C&Requires a nonzero orientation-sensitive imaginary component and fixed chirality gauge.\\
28&Two-time interferometric closure&C&Controlled complex amplitudes and a common energy-zero phase; not raw return probabilities.\\
29&Anchored fourth-order rank&A&Difference set $(v,k,\lambda)$ and anchored nonzero shift; rank $\lambda$, equal to 1 only if $\lambda=1$.\\
30&Hopf Chern charge versus Hankel rank&A&Fixed Hopf bundle, variable chosen trajectory and observable.\\
31&Quantum-switch phase witness&C&Trusted control and instruments; no device-independent causal violation implied.\\
32&Closed-clock monodromy criterion&A&Unitary cycle and eigenvalue 1 of the full monodromy.\\
33&Relational conditional time&C&Compatible cyclic history and clock conditionalization; not proof of cosmological Wheeler--DeWitt dynamics.\\
34&Prime-period minimax limit&A&Homogeneous classical chain, $m<p$, and uniform infinite time.\\
35&Continuous-time oscillation bound&A&Finite homogeneous CTMC, uniform infinite time, and periodic target.\\
36&Primitive Fourier obstruction&A&Frequency coprime to $n$ and $m<n$; explicit Fourier convention.\\
37&Prime-power first useful budget&C&Specific spike signal $a+b\one_{n\mid k}$ and infinite uniform loss; not an arbitrary periodic signal.\\
38&Fibonacci delta-rule identity&A&Static $2\times2$ memory slice; neither a learned Transformer nor physical necessity.\\
39&Two-state qubit short-window witnesses&C&Known damping/phase specification and a fixed homogeneous two-state classical comparator.\\
40&Oracle low-rank softmax error&A&Finite prompt panel; freely chosen latent features, not a deployable network.\\
41&Fixed-encoder optimum&A&Frozen context feature matrix and centered log-probabilities.\\
42&Shared multi-interface rank penalty&C&Specified common rank budget and task weights; parameter/FLOP parity is not automatic.\\\bottomrule
\end{longtable}\endgroup
\section{Correction: reversibility does not imply correlation monotonicity at every step}
\statement{Proposition}{XXXVIII.1}{An exact counterexample to lag monotonicity}
Let $P=\left(\begin{smallmatrix}0&1\\1&0\end{smallmatrix}\right)$, $\pi=(1/2,1/2)$, and $h=(1,-1)^T$. The chain satisfies detailed balance with respect to $\pi$, but its stationary autocorrelation is $C(k)=\mathbb E_\pi[h(X_0)h(X_k)]=(-1)^k$, which is not monotone in $k$. Item 26 of the register, without restrictions on the class and time parity, is false.
\begin{proof}
$P=P^T$ and $\pi_0P_{01}=\pi_1P_{10}$. Since $Ph=-h$, $P^kh=(-1)^kh$, and $C(k)=\langle h,P^kh\rangle_\pi=(-1)^k\langle h,h\rangle_\pi=(-1)^k$.
\end{proof}
\statement{Theorem}{XXXVIII.2}{Complete discrete monotonicity at even lags}
Let $P$ be the transition of a finite stationary reversible chain with invariant probability $\pi$ positive on its state space, and let $h$ be any real observable. For $a_k=C(2k)=\langle h,P^{2k}h\rangle_\pi$, all $j,k\ge0$ satisfy
\begin{equation}(-1)^j\Delta^ja_k\ge0,\qquad\Delta a_k=a_{k+1}-a_k.\tag{142}\end{equation}
More precisely, there are finitely many $t_i\in[0,1]$ and weights $w_i\ge0$ such that
\begin{equation}a_k=\sum_iw_it_i^k,\qquad(-1)^j\Delta^ja_k=\sum_iw_it_i^k(1-t_i)^j.\tag{143}\end{equation}
Conversely, every finite representation of this form is the even-lag sequence of a finite stationary reversible chain with a suitable real observable.
\begin{proof}
Detailed balance makes $P$ self-adjoint on $L^2(\pi)$. As a Markov operator it is a contraction, with real eigenvalues $\lambda_i\in[-1,1]$. Expanding $h=\sum_ic_iu_i$ in an orthonormal eigenbasis gives $C(k)=\sum_i|c_i|^2\lambda_i^k$. Set $w_i=|c_i|^2$, $t_i=\lambda_i^2$; the finite-difference formula follows from the binomial identity.

Conversely, for $W=\sum_iw_i>0$, take a disjoint union of reversible two-state blocks
\[P_i=\frac12\begin{pmatrix}1+\sqrt{t_i}&1-\sqrt{t_i}\\1-\sqrt{t_i}&1+\sqrt{t_i}\end{pmatrix},\]
with block mass $w_i/W$, equally distributed over its two vertices. On block $i$, let $h=\sqrt W(1,-1)^T$. Its stationary even-lag correlation is $\sum_iw_it_i^k$. For $W=0$, use the zero observable. Irreducibility of the whole chain is unnecessary.
\end{proof}
\statement{Corollary}{XXXVIII.3}{Universal cost of simulating a return spike by a stationary reversible chain}
Let $n\ge3$, $f(k)=a+b\one_{\{n\mid k\}}$, $a,b\ge0$, and $a+b\le1$. Compare with all finite stationary reversible Markov autocorrelations $g(k)=\mathbb E_\pi[h(X_0)h(X_k)]$, for real $h$ and $g(k)\in[0,1]$. Independently of the number of hidden states,
\begin{equation}\inf_{P,\pi,h}\sup_{k\ge0}|f(k)-g(k)|=\frac b2.\tag{144}\end{equation}
For the Singer return signal $a=q/(q+1)^2$, $b=1-a$, $n=q^2+q+1$ (in particular, $n\ge7$), the lower bound is $\tfrac12[1-q/(q+1)^2]$. Unlike for unrestricted Markov chains, this obstruction persists as the state count grows, as long as reversibility and stationary autocorrelation are retained.
\begin{proof}
Since $n\nmid2n-2$, $f(2n-2)=a$, $f(2n)=a+b$, while (142) gives $g(2n)\le g(2n-2)$. If the uniform error is $\varepsilon$, then $b\le2\varepsilon+g(2n)-g(2n-2)\le2\varepsilon$. A one-state reversible chain with constant observable $h=\sqrt{a+b/2}$ has $g(k)=a+b/2$ and attains error $b/2$.
\end{proof}
\statement{Proposition}{XXXVIII.4}{A finite-sample reversibility witness}
For a prespecified $n\ge3$, conduct $N$ independent Bernoulli measurements at each of $t_-=2n-2$, $t_+=2n$. Let $b>0$ be the known alternative gap and $\widehat\delta=\widehat f(t_+)-\widehat f(t_-)$. Reject the stationary reversible autocorrelation model if $\widehat\delta>b/2$. Both error probabilities are at most $\alpha$ if
\begin{equation}N\ge\left\lceil\frac8{b^2}\log\frac4\alpha\right\rceil\qquad\text{at each time}.\tag{145}\end{equation}
This is sufficient, not optimal. All hypotheses must concern one experimentally realizable binary output admitting the specified autocorrelation representation under the null. The result does not apply without an additional condition to arbitrary responses $e^TP^k\mu$ with independent $e,\mu$.
\begin{proof}
Each empirical frequency differs from its true value by at least $b/4$ with probability at most $2\exp[-2N(b/4)^2]$. The union probability is at most $4e^{-Nb^2/8}\le\alpha$. Outside this event, the difference of the two errors has modulus strictly below $b/2$. The true difference is nonpositive under the null and exactly $b$ under the alternative. The strict threshold separates them.
\end{proof}
\note{The concrete $(7,3,1)$ model.} For $n=7$, $a=2/9$, $b=7/9$, $f(12)=2/9$ and $f(14)=1$. Every stationary reversible classical autocorrelator has uniform error at least $7/18$, regardless of its state count. At $\alpha=0.05$, a sufficient measurement plan is $N\ge\lceil(648/49)\log80\rceil=58$ samples per time, 116 in total, if the autocorrelation specification is experimentally accessible. This example neither detects quantumness device independently nor excludes arbitrary irreversible classical processes: a seven-state cycle reproduces the signal exactly.
\section{The substantive boundary for AI specialists}
Analyses of recurrent networks and Transformers must distinguish (i) rank of a temporal function, (ii) realization by a general homogeneous stochastic chain, (iii) a stationary reversible autocorrelation representation, and (iv) an input-dependent neural network, which need satisfy none of these specifications. Theorem XXXVIII.2 supplies necessary statistical tests for a stationary reversible proxy, not a universal certificate of language-model incompressibility. Separate finite-panel theorems XXIII and XXVII apply to SVD and multi-interface compression.
\note{Scientific priority.} Discrete spectral methods and complete monotonicity of moment sequences are classical mathematics: reversible Markov chains, the spectral theorem, and the Hausdorff moment problem. This chapter's value is to separate a correct condition from the false claim of lag monotonicity, give an exact applied minimum for periodic return, and specify a falsifiable statistical protocol. No separate proof of worldwide priority is provided. Compare Aldous--Fill, \emph{Reversible Markov Chains and Random Walks on Graphs}, Chapter 3, and Berstel--Reutenauer, \emph{Rational Series and Their Languages}, for the classical equality between Hankel rank and minimal linear dimension.

\chapter{Complementary Symmetrization, Memory Geometry, and Branching Reconstruction}
\note{Scientific specification.} This chapter examines two transferable mathematical mechanisms: normalized classical realizations of pairwise complementary binary effects, and reconstruction of hidden information in branching stochastic channels despite low Hankel rank of a one-dimensional observation. Results apply only to explicitly fixed model classes; they do not solve universal Transformer compression, quantum gravity, or Millennium Problems. Historical priority of the formulas has not been established.
\section{Exact symmetrization of complementary observables}
\statement{Theorem}{XXXIX.1}{Complementarization without additional hidden states}
Suppose $B\in[0,1]^{N\times2q}$ has a normalized classical realization $B_{ij}=\sum_{s=1}^m\mu_i(s)e_j(s)$ with $\mu_i\in\Delta_{m-1}$ and $e_j(s)\in[0,1]$. For $0\le j<q$, define
\[e'_j(s)=\frac{1+e_j(s)-e_{j+q}(s)}2,\qquad e'_{j+q}(s)=1-e'_j(s).\]
Then $e'_j(s)\in[0,1]$, the state count does not increase, paired effects sum to 1 componentwise, and
\begin{equation}B'_{ij}=\tfrac12(B_{ij}+1-B_{i,j+q}),\qquad B'_{i,j+q}=1-B'_{ij}.\tag{146}\end{equation}
If $B_{ij}+B_{i,j+q}=1$ for all $i,j$, then $B'=B$. If $\max_{i,j<q}|B_{ij}+B_{i,j+q}-1|\le\eta$, then $\|B'-B\|_{\max}\le\eta/2$. This is attained already by one numerical pair and is the smallest $\ell^\infty$ correction making a pair complementary.
\begin{proof}
The effect bounds imply $(1+e_j-e_{j+q})/2\in[0,1]$. New effect pairs sum to one at every hidden state. Normalization $\sum_s\mu_i(s)=1$ gives (146). Exact complementarity implies equality of the tables. For any corrected pair $x'+y'=1$, $|x'-x|+|y'-y|\ge|1-x-y|$, hence the largest correction is at least $|1-x-y|/2$. The symmetric correction attains this bound and stays within $[0,1]$.
\end{proof}
\statement{Corollary}{XXXIX.2}{Normalized rank equals a cube-covering number}
For $B=[C\mid1-C]$, $C\in[0,1]^{N\times q}$, the minimum number of normalized hidden states equals the least $m$ for which points $z_1,\ldots,z_m\in[0,1]^q$ satisfy $C_{i,:}\in\operatorname{conv}\{z_1,\ldots,z_m\}$ for every $i$. For $W^{(6)}_{ij}=\sin^2[\pi(i-j)/6]$, the first three effects determine the remaining three. Thus
\[\mu_{\rm normalized}(W^{(6)})=\min\{m:\{w_i\}_{i=0}^5\subseteq\operatorname{conv}(z_1,\ldots,z_m),\ z_s\in[0,1]^3\},\]
where $(w_i)_j=W^{(6)}_{ij}$ for $j=0,1,2$. These six points lie in the plane $x-y+z=1/2$, but the unknown $z_s$ need not lie there. Testing only planar pentagons therefore cannot establish impossibility of a normalized five-state model.
\begin{proof}
Theorem XXXIX.1 allows effects to be made complementary without loss. Their first $q$ values form $z_s$, and $C_{i,:}=\sum_s\mu_i(s)z_s$ is exactly the required convex representation. Conversely, convex coefficients give probabilities $\mu_i$, while $z_s$ and $1-z_s$ give both halves of the effect table. For $W^{(6)}$, direct calculation gives $W_{i,j+3}=1-W_{ij}$ and $x-y+z=1/2$.
\end{proof}
\statement{Lemma}{XXXIX.3}{The open cube interior is unused}
Each of the six points $w_i$ lies in the relative interior of a different face of $[0,1]^3$. In an exact realization, every $z_s\in(0,1)^3$ has zero weight in all six convex combinations. It therefore suffices to search for $z_s$ on the cube boundary. Restriction to the cube's eight vertices is not justified.
\begin{proof}
Each $w_i$ has exactly one coordinate equal to 0 or 1; its other coordinates are strictly between them. An average of nonnegative numbers is zero only if all positively weighted terms are zero, and an average of numbers at most 1 is one only if all positively weighted terms are one. An interior point cannot participate in any of the six representations and can be removed.
\end{proof}
\note{Current open problem.} The result $\rank_+(W^{(6)})=5$ concerns unnormalized nonnegative factors. Corollary XXXIX.2 currently gives only $5\le\mu_{\rm normalized}(W^{(6)})\le6$. The multistart nonlinear optimization in the checking code does not certify global impossibility with five effects.
\section{Branching reconstructs information at Hankel rank one}
\statement{Theorem}{XXXIX.4}{Exact variance and information certificates for a binary-spin tree}
Let a regular $b$-ary tree, $b\ge2$, have a uniform root spin $X_\varnothing\in\{-1,+1\}$. Each edge carries an independent binary symmetric channel $X_v=X_u\xi_{uv}$ with $\mathbb E\xi_{uv}=\lambda\in(0,1]$. Set $S_t=\sum_{|v|=t}X_v$ and $Z_t=S_t/(b\lambda)^t$. Then
\begin{align}\mathbb E[Z_t\mid X_\varnothing]&=X_\varnothing,\tag{147}\\V_t:=\operatorname{Var}(Z_t\mid X_\varnothing)&=(1-\lambda^2)\sum_{j=1}^t(b\lambda^2)^{-j}.\tag{148}\end{align}
For the conditional distributions $P_t^+,P_t^-$ of all leaves under the two root values,
\begin{equation}\operatorname{TV}(P_t^+,P_t^-)\ge\frac1{1+V_t}.\tag{149}\end{equation}
If $b\lambda^2>1$, at every depth
\begin{equation}\operatorname{TV}(P_t^+,P_t^-)\ge\frac{b\lambda^2-1}{(b-1)\lambda^2}>0.\tag{150}\end{equation}
The optimal root-guessing success probability with equal priors is at least $1/2+1/[2(1+V_t)]$. Yet the correlation along any single branch is $c_k=\lambda^k$, of infinite Hankel rank one.
\begin{proof}
Conditional on depth-$t$ spins, $S_{t+1}$ has mean $b\lambda S_t$ and variance $b^{t+1}(1-\lambda^2)$. Thus $Z_t$ is a martingale conditional on the root, with $V_{t+1}-V_t=(1-\lambda^2)(b\lambda^2)^{-(t+1)}$ and $V_0=0$.

Let $Q=(P_t^++P_t^-)/2$ and $h=(dP_t^+-dP_t^-)/(dP_t^++dP_t^-)$ on the support of the sum. Then $|h|\le1$, $\mathbb E_Q[Z_th]=1$, $\mathbb E_QZ_t^2=1+V_t$, and $\mathbb E_Q|h|=\operatorname{TV}(P_t^+,P_t^-)$. Cauchy--Schwarz gives
\[1\le\sqrt{(1+V_t)\mathbb E_Qh^2}\le\sqrt{(1+V_t)\operatorname{TV}(P_t^+,P_t^-)}.\]
This proves (149). For $b\lambda^2>1$, $V_t\le(1-\lambda^2)/(b\lambda^2-1)$, giving (150). Along one branch, independence gives mean $\lambda^k$, and $[c_{i+j}]$ is the outer product of the two nonzero vectors $(\lambda^i)$ and $(\lambda^j)$.
\end{proof}
\statement{Corollary}{XXXIX.5}{Robustness to independent terminal detector noise}
Suppose the observed leaves are $Y_v=X_v\zeta_v$, with independent signs $\zeta_v$, independent also of the tree, and $\mathbb E\zeta_v=\nu\in(0,1]$. The estimator $\widetilde Z_t=(b\lambda)^{-t}\nu^{-1}\sum_{|v|=t}Y_v$ remains conditionally unbiased, with
\begin{equation}\operatorname{Var}(\widetilde Z_t\mid X_\varnothing)=V_t+\frac{1-\nu^2}{\nu^2(b\lambda^2)^t}.\tag{151}\end{equation}
The TV distinguishability from noisy leaves is at least the reciprocal of one plus the right-hand side. For $b\lambda^2>1$ and fixed $\nu>0$, the additional variance vanishes exponentially. The conclusion does not hold at $\nu=0$.
\begin{proof}
Given the true leaves, $\mathbb E[Y_v\mid X_v]=\nu X_v$ and $\operatorname{Var}(Y_v\mid X_v)=1-\nu^2$. Independence gives sum variance $b^t(1-\nu^2)$. Total variance and Theorem XXXIX.4 prove the formula; the same Cauchy--Schwarz argument applied to $\widetilde Z_t$ gives the TV bound.
\end{proof}
\statement{Remark}{XXXIX.6}{Relation to OpenAI Math family 229}
Family 229 in \path{openai/math} studies reconstruction thresholds for three- and four-state broadcasting on trees. Its companion \path{lean/docs/229.md} says the formalization contains the supercritical direction for three states. This does not make our binary bound a new worldwide theorem: the binary threshold and second-moment method are classical; see Y. Peres and S. Roch, \emph{Reconstruction on Trees: Exponential Moment Bounds for Linear Estimators} (2009), arXiv:0908.2056. Here we give a self-contained elementary formula linking the tree to the limitation of a one-dimensional Hankel witness.
\begin{figure}[ht]\centering
\fbox{\parbox{.8\textwidth}{Original figure available in the complete source archive: \path{results/V35_TREE_BOUNDS.png}.}}
\caption{Exact analytical lower bounds on TV distinguishability of a root spin from leaves of a regular tree at $\lambda=3/4$, with independent readout visibility $\nu$. Four specifications are shown: $b\in\{2,3\}$, $\nu\in\{1,0.8\}$. These curves are proved bounds, not laboratory measurements or Transformer quality scores.}
\end{figure}
\section{A map of the OpenAI Math catalog's uses and limits of transfer}
The \path{openai/math} repository contains 372 research families. The table groups them by operational connections rather than advertised novelty. The number of accompanying manuscripts and formalizations can change across revisions; specific families and domains of applicability are listed for reproducibility.
\begin{longtable}{p{.13\textwidth}p{.32\textwidth}p{.46\textwidth}}
\toprule Families&Actual direction&Direct admissible transfer to META--CAUSAL\\\midrule
095, 126&Restrictions on conic lifts&Motivation for a precisely specified family of positive factors; no exclusion of all quantum models.\\
116&Noncommutative identity verification&Symbolic checking of operator polynomials within degree and field restrictions.\\
119, 140&Noisy information; memory bits in Gaussian regression&Information limits in the corresponding statistical problems, not a universal LLM compression bound.\\
144, 148&Koopman spectrum; entropy and dimension of self-similar measures&Counterexamples to simplistic interpretations of rank and fractal dimension.\\
229&Reconstruction in tree channels&The branching threshold $b\lambda^2>1$; distinguish spatial observation from one trajectory.\\
265, 269, 276&Gap, quantum correlations, qubit noise&Transfer only when Hamiltonian or channel classes match.\\
327, 365&Markov metric types; inverse conductivity&Geometric stability and nonidentifiability conditions, not a cosmological theory.\\
376&Universal computation in forced flows&Predictability limits of specified forced PDEs, not a solution of a Clay problem.\\\bottomrule
\end{longtable}
\note{Reading formalizations precisely.} The documents \path{lean/docs/229.md} and \path{lean/docs/365.md} restrict their Lean assertions to supercritical reconstruction for a particular three-state channel and nonuniqueness for measurable scalar conductivity, respectively. They do not validate every claim suggested by the families' overview titles. Lean was not used in this work.
\section{Verification protocol}
The script \path{checks/verify_v35_cube_tree.py} reproduces pairwise complementarization, hexagon geometric invariants, conditional variances, and the statistical TV bound. Rigorous checks use rational arithmetic and exact enumeration of small trees. Diagnostics for the open five-state problem use only heuristic optimization, explicitly labeled non-probative. No numerical result is claimed as proof of global optimality for five hidden states.
\setcounter{equation}{151}


% 08_information.tex
\chapter{An Operational Reconstructibility Invariant and the Exact Cost of Compression}
\note{Context and precise formulation.} Previous chapters used three fundamentally different quantities: Hankel rank (a linear predictive space), nonnegative rank (an orthant representation), and normalized memory (probabilistic states with admissible effects). None alone determines reconstructibility of a hidden bit from several noisy sensors. We introduce a fourth operational invariant that also gives the exact cost of linear feature compression. The construction uses known Mahalanobis quadratic geometry and the Ky Fan principle; it is not a claim of new universal physics.
\statement{Definition}{XL.1}{Binary-observation specification}
Let $X$ be uniform on $\{-1,+1\}$ and let the observed real vector $Y\in\R^d$ satisfy
\begin{equation}\mathbb E[Y\mid X=x]=xm,\qquad\operatorname{Cov}(Y\mid X=x)=\Sigma\succ0\quad(x=\pm1).\tag{152}\end{equation}
Write the conditional measures as $P_+,P_-$. Fix the accessible observation algebra, dimension $d$, conditional covariance, and measurement norm. Define
\begin{equation}\mathcal I(m,\Sigma)=m^\top\Sigma^{-1}m\ge0.\tag{153}\end{equation}
This measures linearly accessible statistical reconstructibility, not the full mutual information $I(X;Y)$ or neural-network dimension. For singular $\Sigma\succeq0$, a component of $m$ in $\ker\Sigma$ gives an error-free linear discriminator; otherwise use the Moore--Penrose inverse on the support of $\Sigma$.
\statement{Theorem}{XL.2}{Gauge invariance and a universal TV certificate}
Under Definition XL.1,
\begin{align}\mathcal I(Gm,G\Sigma G^\top)&=\mathcal I(m,\Sigma)\quad\text{for all }G\in GL_d(\R),\tag{154}\\\operatorname{TV}(P_+,P_-)&\ge\frac{\mathcal I}{1+\mathcal I}.\tag{155}\end{align}
The bound requires neither normality, independent coordinates, nor a special tree structure.
\begin{proof}
Invariance follows from $(G\Sigma G^\top)^{-1}=G^{-\top}\Sigma^{-1}G^{-1}$. Put $Q=(P_++P_-)/2$ and $h=(dP_+-dP_-)/(dP_++dP_-)$ on its support. Then $|h|\le1$, $\mathbb E_Q|h|=\operatorname{TV}(P_+,P_-)$, $\mathbb E_Q[Yh]=m$, and $\mathbb E_Q[YY^\top]=\Sigma+mm^\top$. For any $w$, Cauchy--Schwarz and $h^2\le|h|$ give
\[(w^\top m)^2\le\bigl(w^\top\Sigma w+(w^\top m)^2\bigr)\operatorname{TV}(P_+,P_-).\]
The maximum of $(w^\top m)^2/(w^\top\Sigma w)$ is $m^\top\Sigma^{-1}m$, attained at $w=\Sigma^{-1}m$.
\end{proof}
\statement{Proposition}{XL.3}{Negative control: incompleteness of the quadratic invariant}
$\mathcal I$ need not determine full distinguishability. For one observation $Y\in[-1,1]$, take
\[P_+=\tfrac12\delta_{-1/2}+\tfrac12\delta_{1/2},\qquad P_-=\tfrac18\delta_{-1}+\tfrac34\delta_0+\tfrac18\delta_1.\]
Both distributions have mean 0 and variance $1/4$, so $m=0$ and $\mathcal I=0$, but their supports are disjoint and TV equals 1. The invariant is a guaranteed resource of linear observation, not a complete information characterization or a quantumness criterion.
\statement{Theorem}{XL.4}{Exact irreversible information reduction under compression}
Let $A\in\R^{r\times d}$ have full row rank, and let $Z=AY$ be the compressed instrument's sole observation. Write
\[\mathcal I_A=(Am)^\top(A\Sigma A^\top)^{-1}Am,\qquad\Pi_A=\operatorname{proj}_{\operatorname{row}(A\Sigma^{1/2})}.\]
Then
\begin{equation}\mathcal I-\mathcal I_A=\|(I-\Pi_A)\Sigma^{-1/2}m\|_2^2\ge0.\tag{156}\end{equation}
Equality holds exactly when $\Sigma^{-1}m\in\operatorname{Im}(A^\top)$. In particular, the single tailored statistic $z=(\Sigma^{-1}m)^\top Y$ preserves all of $\mathcal I$. It need not preserve the full distribution of $Y$ or work for an arbitrary new task.
\begin{proof}
Put $u=\Sigma^{-1/2}m$, $T=A\Sigma^{1/2}$. Then $\mathcal I=\|u\|^2$ and $\mathcal I_A=u^\top T^\top(TT^\top)^{-1}Tu=\|\Pi_Au\|^2$. Pythagoras proves (156). Equality means $u\in\operatorname{row}(T)$, equivalent to the stated condition.
\end{proof}
\statement{Theorem}{XL.5}{The exact joint spectrum of multitask memory}
Tasks $a=1,\ldots,K$ have means $\pm m_a$, common conditional covariance $\Sigma\succ0$, and prescribed weights $\omega_a\ge0$. All use one rank-$r$ linear instrument $A$. Define
\begin{equation}G=\Sigma^{-1/2}\left(\sum_{a=1}^K\omega_am_am_a^\top\right)\Sigma^{-1/2}\succeq0.\tag{157}\end{equation}
For eigenvalues $\gamma_1\ge\cdots\ge\gamma_d\ge0$,
\begin{equation}\max_{\rank A=r}\sum_{a=1}^K\omega_a\mathcal I_A(m_a)=\sum_{j=1}^r\gamma_j,\qquad L_r=\sum_{j>r}\gamma_j.\tag{158}\end{equation}
The loss is invariant under any common invertible feature-coordinate change. This optimum concerns linear observability with fixed common covariance, not full Transformer reconfiguration or inference cost.
\begin{proof}
The preceding transformation gives $\sum_a\omega_a\mathcal I_A(m_a)=\Tr(\Pi_AG)$. As full-row-rank $A$ varies, $\Pi_A$ runs through all rank-$r$ orthogonal projectors: for a projector with orthonormal image basis $U$, choose $A=U^\top\Sigma^{-1/2}$. Ky Fan's variational principle gives the sum of the largest $r$ eigenvalues. Invariance follows because the operational objective is unchanged under $m_a\mapsto Tm_a$, $\Sigma\mapsto T\Sigma T^\top$, $A\mapsto AT^{-1}$.
\end{proof}
\begin{figure}[ht]\centering\includegraphics[width=.85\textwidth]{assets/page_116_1.png}
\caption{Exact decrease of optimal loss $L_r$ with the rank of a shared linear adapter on three synthetic tasks. This is not a measurement of LLM quality.}\end{figure}
\statement{Theorem}{XL.6}{A robust certificate with inexact measurements}
Suppose $\|m-\widehat m\|_2\le\eta$ and $\|\Sigma-\widehat\Sigma\|_2\le\delta$. For any $w$ with positive denominator,
\begin{equation}L(w):=\frac{(|w^\top\widehat m|-\eta\|w\|_2)_+^2}{w^\top\widehat\Sigma w+\delta\|w\|_2^2}\le\mathcal I,\qquad\operatorname{TV}(P_+,P_-)\ge\frac{L(w)}{1+L(w)}.\tag{159}\end{equation}
This remains valid for a data-dependent choice of $w$, provided the norm bounds hold simultaneously.
\begin{proof}
The triangle inequality gives $|w^\top m|\ge(|w^\top\widehat m|-\eta\|w\|)_+$, and the operator-norm bound gives $w^\top\Sigma w\le w^\top\widehat\Sigma w+\delta\|w\|^2$. Maximize the Rayleigh quotient defining $\mathcal I$ and apply Theorem XL.2.
\end{proof}
\statement{Corollary}{XL.7}{Finite sampling without hidden feature selection}
If $Y_i\in[-1,1]$ and there are $N$ independent preparations with $X=+1$, all $d$ means and $d^2$ raw second moments have absolute error at most $\epsilon$ with probability at least $1-\alpha$ when
\begin{equation}N\ge\frac2{\epsilon^2}\log\frac{2(d+d^2)}\alpha.\tag{160}\end{equation}
On this event, use $\eta=\sqrt d\,\epsilon$ and $\delta=d(3\epsilon+\epsilon^2)$ in (159), with $\widehat\Sigma=\widehat K-\widehat m\widehat m^\top$. Independence concerns repeated preparations, not coordinates within one vector $Y$.
\begin{proof}
Hoeffding gives $\Pr(|\widehat\mu-\mu|>\epsilon)\le2e^{-N\epsilon^2/2}$ for each bounded variable $Y_i$ or $Y_iY_j$. A union bound proves (160). Mean-vector error is at most $\sqrt d\epsilon$. Entrywise covariance error is at most $\epsilon+2\epsilon+\epsilon^2$, so its operator norm is at most $d(3\epsilon+\epsilon^2)$. The simultaneous bounds permit adaptive choice of a linear sensor on the same data, not an arbitrary new nonlinear model.
\end{proof}
\section{A tree with heterogeneous channels: exact moments and an example}
\statement{Proposition}{XL.8}{A binary-spin tree with independent noisy detectors}
Let a binary root $X=\pm1$ propagate on a finite tree via independent edge multipliers $\xi_e=\pm1$, $\mathbb E\xi_e=\lambda_e$. An independent detector at each observed leaf $v$ multiplies its spin by $\zeta_v=\pm1$, $\mathbb E\zeta_v=\nu_v$. Then
\begin{align}m_v&=\nu_v\prod_{e\in\operatorname{path}(\varnothing,v)}\lambda_e,\tag{161}\\
K_{uv}&=\begin{cases}1,&u=v,\\\nu_u\nu_v\displaystyle\prod_{e\in\operatorname{path}(u,v)}\lambda_e,&u\ne v,\end{cases}\tag{162}\\
\Sigma&=K-mm^\top.\tag{163}\end{align}
When $\Sigma\succ0$, Theorem XL.2 bounds root distinguishability and gives optimal leaf weights $w\propto\Sigma^{-1}m$. On a homogeneous regular tree, permutation symmetry makes $w$ uniform, so the leaf-sum theorem XXXV is a special case.
\begin{proof}
Conditional means are products of independent multiplier means. In $Y_uY_v$, shared root-path edges cancel, leaving the path $u\leftrightarrow v$; diagonally $Y_v^2=1$. Weight optimality follows from the Rayleigh quotient in XL.2. On a homogeneous tree, automorphism-induced leaf permutations preserve $m,\Sigma$, forcing the unique optimal weights to be constant.
\end{proof}
\note{A fully rational checking example.} Consider a six-edge tree with $\lambda=(3/5,4/5,2/3,3/4,4/5,5/6)$ and four independent detectors $\nu=(5/6,3/4,2/3,7/8)$. Enumeration of all $2^{10}=1024$ edge/detector sign configurations gives
\begin{align}\mathcal I&=\frac{2174372499}{2591339288}\approx0.8391,\tag{164}\\
\operatorname{TV}(P_+,P_-)&\ge\frac{2174372499}{4765711787}\approx0.4563,\tag{165}\\
\operatorname{TV}(P_+,P_-)&=\frac{109627}{172800}\approx0.6344\quad\text{(exact full enumeration).}\tag{166}\end{align}
The optimal conditional variance of an unbiased linear estimate is $\mathcal I^{-1}=2591339288/2174372499$, while uniform leaf weighting gives the larger variance $5663031/4068289$. Exact $m,\Sigma$, the script, and checking protocol are stored in \path{checks/RESULTS_v36.json}.
\section{Invariant interpretation and relation to research}
The chapter provides a verifiable hierarchy:
\[\text{observation specification}\longrightarrow(m,\Sigma)\longrightarrow\mathcal I\longrightarrow\mathcal I_A\longrightarrow L_r\longrightarrow\text{finite-sample lower bound}.\]
$\mathcal I$ is neither Hankel rank, a topological Chern class, nor a new natural charge. It is an operational quadratic invariant for specified conditional distributions. Its direct AI application is quantitative choice of shared linear adapters under known or bounded covariance. Separate tests of token-prediction quality, actual latency, energy, and generalization to new contexts are also needed.
\note{Priority and methodological sources.} The proofs use standard Mahalanobis geometry, linear discriminant analysis, Rayleigh quotients, Ky Fan's variational principle, and classical broadcasting reconstruction on trees. See Peres--Roch (2009), Hopkins--Li (COLT 2024), and \path{openai/math:229}. Catalog works 119, 140, 144, 327, and 365 have different mathematical specifications and do not assert universal Transformer compression. Unification of these certificates is part of this research program; worldwide priority and the status of a new fundamental law of nature are not established.
\note{Verification.} \path{checks/verify_v36_information_geometry.py} performs independent exact rational enumeration and checks gauge invariance, irreversible compression loss, optimal multitask projections, the robust bound, and Hoeffding's inequality using SymPy, NumPy, and SciPy. Lean is not used. Numerical tests do not replace the proofs above.

\chapter{Distinguishability Invariants and the Exact Cost of Stochastic Observation (Revision v37)}
\note{Scientific status and starting resources.} We extend the preceding Mahalanobis-invariant results and connect them to sequential reconstruction, feature compression, and binary information channels. Relevant OpenAI Math directions are 119 (binary-channel information compression), 229 (supercritical tree reconstruction), 140 (memory and sample size), 144 (the Koopman spectral operator), and 365 (inverse-problem nonuniqueness). None proves the following theorems for us. The immediate tools are classical comparison of statistical experiments (Blackwell--Le Cam), conditional expectation, triangular discrimination, and the Ky Fan principle. We give their joint operationally testable consequence; no new worldwide priority is claimed.
\section{Specification and the full binary-distinguishability invariant}
On a finite alphabet $\mathcal Y$, let $p=P(Y=\cdot\mid B=+1)$ and $q=P(Y=\cdot\mid B=-1)$ be arbitrary distributions, with equal priors for $B=\pm1$. Define the mixture and posterior contrast
\begin{equation}r_y=\frac{p_y+q_y}2,\qquad h_y=\begin{cases}(p_y-q_y)/(p_y+q_y),&p_y+q_y>0,\\0,&p_y+q_y=0.\end{cases}\tag{167}\end{equation}
Then $h_y=\mathbb E[B\mid Y=y]$, $|h_y|\le1$, and $\mathbb E_rh=0$. Define
\begin{equation}\mathcal J(Y)=\mathbb E_rh^2=\frac12\sum_{y\in\mathcal Y}\frac{(p_y-q_y)^2}{p_y+q_y},\tag{168}\end{equation}
with $0/0=0$. This is half the standard triangular discrimination. It is invariant under relabeling observations, not a topological charge, a new entropy, or a standalone universal physical law.
\statement{Theorem}{XLI.1}{Exact statistical distinguishability bounds}
For every admissible $p,q$,
\[\mathcal J(Y)\le\operatorname{TV}(p,q)\le\sqrt{\mathcal J(Y)},\qquad R^*_{\rm Bayes}=\frac{1-\operatorname{TV}(p,q)}2.\]
Both bounds between $\mathcal J$ and TV are attainable.
\begin{proof}
$\operatorname{TV}(p,q)=\tfrac12\sum_y|p_y-q_y|=\mathbb E_r|h|$. Apply $h^2\le|h|$ and Cauchy--Schwarz. The lower equality is attained by $p=(1-t,t,0)$, $q=(1-t,0,t)$, $0\le t\le1$, where $\mathcal J=\operatorname{TV}=t$. For the upper equality, take $r=(1/2,1/2)$, $h=(c,-c)$, $0\le c\le1$; then TV is $c$ and $\mathcal J=c^2$. The Bayes-risk formula is the equal-prior binary-testing identity.
\end{proof}
\section{Exact loss under an arbitrary stochastic channel}
The observer receives $Z$ from $Y$ through any column-stochastic channel $T_{z\mid y}$ independent of $B$. This includes a quantized label, random aggregation, discretization, a neural encoder as a deterministic input function, and measurement noise. Quantum channels are not assumed.
\statement{Theorem}{XLI.2}{Exact dissipation of posterior contrast}
For $h_Z(z)=\mathbb E_r[h_Y(Y)\mid Z=z]$,
\begin{equation}\mathcal J(Y)-\mathcal J(Z)=\mathbb E_r[(h_Y-h_Z)^2]=\mathbb E_r[\operatorname{Var}(h_Y\mid Z)]\ge0.\tag{169}\end{equation}
The loss vanishes exactly when $h_Y$ is almost surely a function of $Z$. For sequential channels $Y\to Z\to U$, losses add exactly:
\[\mathcal J(Y)-\mathcal J(U)=[\mathcal J(Y)-\mathcal J(Z)]+[\mathcal J(Z)-\mathcal J(U)].\]
\begin{proof}
Bayes' rule gives $\mathbb E[B\mid Z]=\mathbb E[\mathbb E[B\mid Y]\mid Z]$. Conditional expectation is an orthogonal projection in $L^2(r)$. Pythagoras and total variance give the squared-norm decomposition. Equality holds when the residual square vanishes. Sequential channels give a telescoping difference.
\end{proof}
\statement{Corollary}{XLI.3}{Limits on claims about minimal AI dimension}
For one fixed binary task, one real feature $h(Y)$ preserves all statistical distinguishability $\mathcal J$ and TV. Thus no universal lower bound on nonlinear representation dimension can follow from the number of original observations without restrictions on arithmetic precision, task family, memory horizon, and architecture class.
\begin{proof}
The likelihood ratio $p_y/q_y$, including boundary cases, is uniquely determined by $h_y\in[-1,1]$. This is a sufficient statistic for the binary experiment, preserving the conditional label probability and optimal test.
\end{proof}
\section{Linear features, Mahalanobis geometry, and a shared adapter}
Let $\phi(Y)\in\R^d$ be arbitrary features with finite second moments, $G=\operatorname{Cov}_r(\phi)$, and $v=\mathbb E_r[h(\phi-\mathbb E_r\phi)]$.
\statement{Theorem}{XLI.4}{A bridge between full and linear distinguishability}
With the pseudoinverse $G^+$,
\begin{equation}\mathcal J_\phi=v^\top G^+v=\|\Pi_{\operatorname{span}\{\phi_i-\mathbb E_r\phi_i\}}h\|_{L^2(r)}^2\le\mathcal J(Y).\tag{170}\end{equation}
If conditional feature means are $\pm m$ and their common conditional covariance is $\Sigma\succ0$, then
\begin{equation}\mathcal J_\phi=\frac{\mathcal I}{1+\mathcal I},\qquad\mathcal I=m^\top\Sigma^{-1}m.\tag{171}\end{equation}
The earlier TV bound is therefore a special case of Theorem XLI.1.
\begin{proof}
Projection onto centered linear feature combinations has coefficients satisfying $Ga=v$. Since $v$ is orthogonal to $\ker G$, its squared norm is $v^\top G^+v$. Total covariance gives $G=\Sigma+mm^\top$ and $v=m$. Sherman--Morrison gives $m^\top(\Sigma+mm^\top)^{-1}m=\mathcal I/(1+\mathcal I)$.
\end{proof}
\statement{Theorem}{XLI.5}{Exact spectral cost of a shared multitask feature}
Several binary experiments have common mixture measure $r$, posterior contrasts $h_1,\ldots,h_K\in L^2_0(r)$, and weights $w_a\ge0$. For a shared observable subspace $S\subseteq L^2_0(r)$ of dimension at most $d$,
\begin{equation}\inf_{\dim S\le d}\sum_{a=1}^Kw_a\|h_a-\Pi_Sh_a\|_2^2=\sum_{j>d}\lambda_j(\mathsf K),\qquad\mathsf K=\sum_{a=1}^Kw_ah_a\otimes h_a.\tag{172}\end{equation}
This concerns a freely chosen feature space under a fixed specification, not an automatically deployable low-rank Transformer architecture.
\begin{proof}
Total energy is $\Tr\mathsf K$; energy retained by $S$ is $\Tr(\Pi_S\mathsf K)$. Ky Fan's principle maximizes this by the sum of the $d$ largest eigenvalues of the positive finite-rank operator. Subtraction gives the spectral tail.
\end{proof}
\section{Finite samples, robustness, and repeated observations}
\statement{Theorem}{XLI.6}{A continuous finite-sample certificate}
Let $s=|\mathcal Y|$ and estimate $p,q$ independently from $N$ identically distributed measurements each. With empirical frequencies $\widehat p,\widehat q$, $0<\alpha<1$, set $\epsilon_N=\sqrt{\log(4s/\alpha)/(2N)}$, $\widehat{\mathcal J}=\mathcal J(\widehat p,\widehat q)$. With probability at least $1-\alpha$,
\begin{equation}\operatorname{TV}(p,q)\ge\mathcal J(p,q)\ge\max\{0,\widehat{\mathcal J}-3s\epsilon_N\}.\tag{173}\end{equation}
No positive lower bound on cell masses is required.
\begin{proof}
For arbitrary distributions,
\[\mathcal J(p,q)=\sup_{|g_y|\le1}\sum_y\left[(p_y-q_y)g_y-\frac{p_y+q_y}{2}g_y^2\right],\]
with maximizer $g_y=(p_y-q_y)/(p_y+q_y)\in[-1,1]$. Thus
\[|\mathcal J(p,q)-\mathcal J(p',q')|\le\tfrac32(\|p-p'\|_1+\|q-q'\|_1),\]
since each perturbation coefficient has modulus at most $3/2$. Hoeffding and a union bound over $2s$ frequencies give coordinate errors at most $\epsilon_N$ with probability at least $1-\alpha$. Their total $L^1$ error is at most $2s\epsilon_N$; apply Theorem XLI.1.
\end{proof}
\statement{Corollary}{XLI.7}{Independent repetitions amplify distinguishability}
For $N$ conditionally independent repetitions with the same conditional means $\pm m$ and covariance $\Sigma\succ0$,
\[\operatorname{TV}(P_+^{\otimes N},P_-^{\otimes N})\ge\frac{N\mathcal I}{1+N\mathcal I},\qquad R^*_{{\rm Bayes},N}\le\frac1{2(1+N\mathcal I)}.\]
This does not automatically extend to adaptive or correlated probes.
\begin{proof}
The concatenated vector has block covariance $\diag(\Sigma,\ldots,\Sigma)$ and mean vector $(m,\ldots,m)$, giving $\mathcal I_N=N\mathcal I$. Apply Theorem XLI.4 and the TV lower bound.
\end{proof}
\section{Exact checks and negative controls}
For uniform $r$ on six elements and contrasts
\[h_1=\tfrac14(3,-3,2,-2,1,-1),\quad h_2=\tfrac16(2,2,-3,-3,1,1),\quad h_3=\tfrac14(2,-2,-1,1,0,0),\]
define $p_a=r(1+h_a)$, $q_a=r(1-h_a)$. A three-sensor channel merging pairs $(0,1)$, $(2,3)$, $(4,5)$ completely destroys $\mathcal J(h_1)$ and $\mathcal J(h_3)$ but preserves $\mathcal J(h_2)$ exactly. These rational examples demonstrate why a joint multitask specification is needed; they are not a standalone LLM experiment.
\begin{figure}[ht]\centering\includegraphics[width=.95\textwidth]{assets/page_121_1.png}
\caption{Exact stochastic-observation cost in finite binary experiments. The identity is proved analytically; the plot shows computed rational examples, not an LLM experiment.}\end{figure}
The program \path{checks/verify_v37_posterior_geometry.py} checks exact rational equalities, 3500 random stochastic trials, and independent numerical tests of multitask spectral optimality, robustness, and the empirical bound at $N=10^5$ measurements per hypothesis. Numerical tests do not replace analytical proofs; Lean was not used. Programs and exact results are included in the complete archive.
\setcounter{equation}{173}
\note{Scientific limitations and existing results.}
\begin{enumerate}
\item $\mathcal J$ is a known symmetric $f$-divergence; the loss formula is total variance for posterior contrast. It does not establish a previously unknown fundamental law of nature.
\item OpenAI Math family 119 studies particular information inequalities for binary symmetric noise; we do not attribute the general formula (169) to it.
\item Family 229 establishes root reconstructibility in a special tree-broadcasting model, not compressibility of Transformer architectures. Results from 140 cannot be transferred to arbitrary networks without checking hypotheses.
\item Families 144 and 365 concern separate spectral and inverse problems; they do not assert quantum gravity, universal holography, or a new physical law.
\item For one binary task, a sufficient nonlinear statistic is scalar. Conclusions about required hidden-network dimension need encoder-class, parameter, and multitask restrictions.
\end{enumerate}


% 09_cube_horizons.tex
\chapter{Cubic Rigidity and the Exact Cost of Physical Normalization}
\note{The six-label control set.} To make the geometric proof reproducible, the six observed points in the central cube section $x-y+z=1/2$ are specified exactly:
\[(0,1/4,3/4),\ (1/4,0,1/4),\ (3/4,1/4,0),\ (1,3/4,1/4),\ (3/4,1,3/4),\ (1/4,3/4,1).\]
These are the original six-label witness data, not additional assumptions about physical measurements.

\note{New central assertion.} Revisions XXVI--XXXVII proved $\rank_+W^{(6)}=5$, but bounded normalized classical hidden-state dimension only between 5 and 6. This chapter resolves the question: $\mu_{\rm norm}(W^{(6)})=6$. The proof uses the geometry of the entire three-dimensional effect cube, not an unjustified requirement that hidden states lie in the plane of observed preparations.
\section{Physical specification and the exact distinction among three ranks}
For $n=6$, put
\begin{equation}W_{ji}=\sin^2\frac{\pi(j-i)}6,\qquad i,j\in\{0,\ldots,5\}.\tag{174}\end{equation}
A normalized classical static realization of size $m$ is a factorization
\begin{equation}W_{ji}=\sum_{s=1}^m\mu_j(s)e_i(s),\qquad\mu_j\in\Delta_{m-1},\quad0\le e_i(s)\le1.\tag{175}\end{equation}
All preparations share one hidden-state space, and every effect is physically bounded between zero and one. This is stricter than an orthant factorization $W=UV$, $U,V\ge0$, with no probabilistic factor constraints.
\statement{Lemma}{XLII.1}{Complementary symmetrization without additional memory}
If $W_{j,i+3}=1-W_{ji}$, every normalized realization (175) can be replaced by an equivalent realization of the same size with $e_{i+3}(s)=1-e_i(s)$ for $i=0,1,2$.
\begin{proof}
Set $e'_i(s)=[1+e_i(s)-e_{i+3}(s)]/2$ and $e'_{i+3}=1-e'_i$. Both remain in $[0,1]$. Averaging under any $\mu_j$ gives $[1+W_{ji}-W_{j,i+3}]/2=W_{ji}$ and $W_{j,i+3}$. No hidden index is added.
\end{proof}
Thus only the first three effects are needed. Set $v_s=(e_0(s),e_1(s),e_2(s))\in C=[0,1]^3$. The six preparations give
\begin{equation}\begin{aligned}
p_0&=(0,\tfrac14,\tfrac34),&p_1&=(\tfrac14,0,\tfrac14),&p_2&=(\tfrac34,\tfrac14,0),\\
p_3&=(1,\tfrac34,\tfrac14),&p_4&=(\tfrac34,1,\tfrac34),&p_5&=(\tfrac14,\tfrac34,1).
\end{aligned}\tag{176}\end{equation}
Let $\Pi=\{(x,y,z):x-y+z=1/2\}$ and $H=C\cap\Pi$. This is a hexagonal cube section, and each $p_j$ lies strictly inside a different side. Since $p_j=\sum_s\mu_j(s)v_s$, an $m$-state realization is equivalent to
\begin{equation}\{p_0,\ldots,p_5\}\subseteq\operatorname{conv}\{v_1,\ldots,v_m\},\qquad v_s\in C.\tag{177}\end{equation}
The $v_s$ are arbitrary points of the three-dimensional cube.
\begin{figure}[ht]\centering\includegraphics[width=.65\textwidth]{assets/figure_123.pdf}
\caption{The planar section $H$ of the three-dimensional cube and six points in the interiors of its sides. Arbitrary hidden vertices $v_s\in C$ cannot be replaced by points $v_s\in\Pi$ without a separate proof.}\end{figure}
\section{Two independent geometric lemmas}
\statement{Lemma}{XLII.2}{Six contacts require six section vertices}
Let $H$ be a nondegenerate hexagon, with $p_j$, $0\le j<6$, in the relative interior of its respective side. If a convex polygon $Q$ satisfies $\{p_0,\ldots,p_5\}\subseteq Q\subseteq H$, it has at least six vertices. If it has exactly six, then either $Q=\operatorname{conv}\{p_j\}$ or $Q=H$.
\begin{proof}
Let $S$ be the vertices of $H$ that are also vertices of $Q$, and $c=|S|$. Let $e(S)$ count sides of $H$ with both endpoints in $S$. On each of the other $6-e(S)$ sides, $Q$ must have a relative-interior vertex: $p_j$ lies on a supporting line and must be a convex combination of vertices of $Q$ on that line. If neither endpoint is present, an interior vertex is necessary; if only one is present, it is still necessary. Hence $\#\operatorname{vert}(Q)\ge c+6-e(S)$. The side graph is $C_6$, so every proper nonempty $S$ satisfies $e(S)\le c-1$, forcing at least seven vertices. If $S=\varnothing$, equality at six means one vertex on each side, necessarily $p_j$. If $S$ includes all six outer vertices, $Q=H$. These exhaust equality cases.
\end{proof}
\statement{Lemma}{XLII.3}{Sections of a convex hull of five points}
Let $P=\operatorname{conv}\{v_1,\ldots,v_m\}\subset\R^3$, $m\le5$, and suppose a plane $\Pi$ cuts it in a two-dimensional polygon $Q$. If $Q$ has at least six vertices, then $m=5$, no $v_s$ lies on $\Pi$, and the plane divides them into groups of sizes 2 and 3. All six segments joining these groups meet $\Pi$ in distinct vertices of $Q$.
\begin{proof}
Let $k_0$ vertices lie on the plane, with $a,b$ on its two sides; $a+b+k_0=m\le5$. Each section vertex is either one of the $k_0$ coplanar points or lies on a segment between opposite sides. Thus $\#\operatorname{vert}(Q)\le k_0+ab$. For $m\le5$, this reaches six only at $(k_0,a,b)=(0,2,3)$ or $(0,3,2)$. All six segments must then produce distinct section vertices. No prescribed configuration of the five original points is assumed.
\end{proof}
\section{Main theorem: $\mu_{\rm norm}=6$}
\statement{Theorem}{XLII.4}{Exact normalized cost of the hexagonal probability table}
For (174),
\begin{equation}\rank W^{(6)}=3,\qquad\rank_+W^{(6)}=5,\qquad\mu_{\rm norm}(W^{(6)})=6.\tag{178}\end{equation}
The table has a one-qubit realization. Physical normalization therefore incurs a strict additional cost beyond nonnegative rank.
\begin{proof}
Six deterministic preparations $\mu_j(s)=\delta_{js}$ with effects $e_i(s)=W_{si}\in[0,1]$ give the upper bound. If $m\le5$ were possible, (177) and Lemma XLII.1 would give a polytope $P\subseteq C$ whose section $Q=P\cap\Pi$ contains all $p_j$. The two lemmas force $m=5$, a $2+3$ split across $\Pi$, and either $Q=\operatorname{conv}\{p_j\}$ or $Q=H$. Denote the two vertices on one side by $u_1,u_2$ and the three on the other by $z_1,z_2,z_3$. Their six joining segments meet $\Pi$ in the six section vertices.

Case 1: $Q=\operatorname{conv}\{p_j\}$. Each $p_j$ lies in the relative interior of a different cube face. If an interior point of a segment in the cube lies on a supporting face, both endpoints lie on that face, by the coordinate equality 0 or 1. Thus each $u_a$ lies on three different cube faces and is a cube vertex. The two face triples are disjoint, so $u_1,u_2$ are opposite cube vertices. But $\Pi$ passes through the center and no cube vertex lies on it, so opposite vertices are on opposite sides---a contradiction.

Case 2: $Q=H$. The six section vertices lie in the relative interiors of six cube edges crossed by $\Pi$. A segment with an interior point on a cube edge must have both endpoints on that edge. Thus each $u_a$ would be incident to three distinct crossed edges. Direct enumeration of $(x,y,z)\in\{0,1\}^3$ shows that crossed edges join levels 0 and 1 of $x-y+z$, and no cube vertex is incident to three of them; the maximum incidence is two. This is another contradiction.

Hence $m\le5$ is impossible and six is attainable. The decomposition $W_{ij}=1/2-\tfrac12\cos[\pi(i-j)/3]$ has three independent real modes, giving rank three. Chapters XXV--XXVI proved nonnegative rank five by an exact rational five-factor decomposition and a combinatorial lower bound. Equatorial qubit-phase preparation followed by an orthogonal-projector measurement gives (174).
\end{proof}
\section{Extension to a family of hexagonal sections}
\statement{Theorem}{XLII.5}{The six-face cube barrier}
Let $\Pi$ pass through the center of $C=[0,1]^3$ but no cube vertex. Suppose $H=C\cap\Pi$ is a hexagon whose sides lie in six distinct cube faces. For arbitrary points $p_j$ in the relative interiors of these sides, the minimum $m$ such that $\{p_j\}\subseteq\operatorname{conv}\{v_s\}$ with $v_s\in C$ is exactly six.
\begin{proof}
The preceding proof uses only the six different boundary faces and the cube-face location of each side. Its first case forces opposite cube vertices onto the same side of a central plane. In the second case, each section vertex lies on a crossed cube edge. No cube vertex is incident to more than two crossed edges: three crossings from a vertex would put all three neighbors on the other side, which is impossible for a central hexagonal section, as can be checked using the eight sign vectors or central-section symmetry. Thus five points are again impossible; $v_s=p_s$ gives six.
\end{proof}
\statement{Corollary}{XLII.6}{The zero-noise threshold and computable robustness}
For every particular rational table $(p_0,\ldots,p_5)$ in Theorem XLII.5, there exists $\varepsilon_*>0$ such that every normalized five-state model has maximum-coordinate error at least $\varepsilon_*$. For rational points this threshold is a positive real algebraic number, computable in principle by quantifier elimination over real closed fields. No explicit numerical threshold is claimed here.
\begin{proof}
Five hidden points and six stochastic mixtures form a compact product of cubes and simplices. Their continuous image in the space of six observed triples is compact and excludes the target by Theorem XLII.5, so the distance is positive. Polynomial minimization with rational coefficients over a compact rational semialgebraic set has an algebraic minimum and is decidable by real algebraic geometry.
\end{proof}
\section{Verification, sources, and limits of novelty}
\begin{itemize}
\item Exact certificate: \path{checks/verify_v38_cube_section.py} checks rational coordinates, section equality, all six crossed edges, incidence degrees, 64 outer-hexagon vertex subsets, 21 distributions of five latent points across the plane, and the exact five-factor decomposition.
\item Critical negative control: replacing three-dimensional hidden effects by points of the original plane is invalid. The proof first allows any $v_s\in[0,1]^3$; the plane enters only when intersecting the polytope with the observed-preparation region.
\item Predecessors: A. Padrol and J. Pfeifle, \emph{Polygons as Sections of Higher-Dimensional Polytopes}, Electron. J. Combin. 22 (2015), \#P1.24, \url{https://doi.org/10.37236/4315}; D. Chistikov et al., \emph{On Restricted Nonnegative Matrix Factorization}, ICALP 2016, \url{https://arxiv.org/abs/1605.07061}. A hexagon by itself can be a section of a three-dimensional five-vertex polytope; the decisive restriction is that vertices belong to the cube.
\item Status: a self-contained deduction for this operational problem and its geometric family. Worldwide novelty has not been exhaustively checked; this is no claim to solve a Millennium Problem, establish a new gravitational theory, or outperform an LLM.
\end{itemize}

\chapter{Finite-Horizon Memory, Fibonacci Approximation, and Exact Specifications of Open Problems}
\note{Scientific setting.} This revision uses four independently maintained resources: \url{https://www.proofatlas.ai/open-problems/}, \url{https://github.com/jgmcm/problems/tree/main/catalogue/problems}, \url{https://jgmcm.github.io/problems/}, and \url{https://github.com/openai/math}. ProofAtlas (7 October 2026 revision) contains 500 ranked questions; the ranking is an assessment by language models, not expert consensus. J. McM's catalog contains 84 physics problems with July 2026 statuses. OpenAI Math lists 372 result families in 719 manuscripts (6 October 2026 release), with heterogeneous verification and some works lacking independent audit. Catalogs provide precise formulations and possible methods, not certificates of every included proof.
\section{Main theorem: the same quantum phase gives two answers depending on the horizon}
Let $\varphi=(1+\sqrt5)/2$, $\alpha=\varphi^{-2}=(3-\sqrt5)/2$, and
\begin{equation}f_\alpha(k)=\sin^2(\pi k\alpha),\qquad k=0,1,\ldots.\tag{179}\end{equation}
One qubit realizes it: prepare $|+\rangle$, apply $U=\diag(1,e^{2\pi i\alpha})$ $k$ times, then measure $|-\rangle\langle-|$.
\statement{Theorem}{XLIII.1}{Fibonacci cost of finite-horizon classical simulation}
For every integer $T\ge1$ and $0<\varepsilon\le1/2$, a deterministic homogeneous classical Markov chain with $m$ states and effect $e\in[0,1]^m$ approximates (179) on $\{0,\ldots,T\}$ with
\begin{equation}\max_{0\le k\le T}|f_\alpha(k)-g(k)|\le\varepsilon,\qquad m<2\sqrt{\frac{\pi T}{\varepsilon}}.\tag{180}\end{equation}
Yet every fixed finite homogeneous stochastic model, with any state count, satisfies
\begin{equation}\sup_{k\ge0}|f_\alpha(k)-g(k)|\ge\frac12.\tag{181}\end{equation}
This is sharp, attained by the one-state model $g\equiv1/2$.
\begin{proof}
Let $F_0=0,F_1=1$ be the Fibonacci numbers. The fractions $p_s/q_s=F_{s-2}/F_s$, $s\ge4$, are convergents of $\alpha=[0;2,1,1,\ldots]$, and the standard convergent identity gives
\begin{equation}|\alpha-p_s/q_s|<q_s^{-2}.\tag{182}\end{equation}
Set $R=\sqrt{\pi T/\varepsilon}>2$ and choose the first Fibonacci denominator $q_s\ge R$; then $R\le q_s<2R$. Use a deterministic $q_s$-cycle, $M|j\rangle=|j+1\pmod{q_s}\rangle$, initial state $|0\rangle$, and effect $e(j)=\sin^2(\pi p_sj/q_s)$. Its output is $g(k)=\sin^2(\pi kp_s/q_s)$ for all $k$. Since $\sin^2x$ is 1-Lipschitz,
\[|f_\alpha(k)-g(k)|\le\pi k|\alpha-p_s/q_s|<\pi T/q_s^2\le\varepsilon.\]
For the second assertion, Perron--Frobenius theory gives an integer-periodic sequence $h(k)$ of period $L$ with $g(k)-h(k)\to0$. On every residue subsequence $k=a+L\ell$, irrationality of $L\alpha$ makes the phases dense on the circle. Thus $f_\alpha(k)$ approaches both zero and one while $g(k)\to h(a)$, forcing limiting maximum error at least $\max\{h(a),1-h(a)\}\ge1/2$.
\end{proof}
\statement{Remark}{XLIII.2}{Physical and computational limits}
No topological quantum computer is required; $S^3$ here is only the normalized qubit-amplitude sphere. The Fibonacci matrix from Chapter XXXI is not used as a physical unitary. Fibonacci numbers instead serve as denominators of rational approximations to an irrational angle. The $O(\sqrt{T/\varepsilon})$ estimate is a constructive upper bound, not a proved lower or optimal bound. Real trained Transformers and DeltaNet were not tested by this result.
\begin{center}\begin{tabular}{rrrr}\toprule $T$&$\varepsilon$&Cycle size $m$&Verified bound $\pi T/m^2$\\\midrule10&0.05&34&$\approx0.02717$\\100&0.05&89&$\approx0.03966$\\1000&0.05&377&$\approx0.02210$\\\bottomrule\end{tabular}\end{center}
\begin{figure}[ht]\centering\includegraphics[width=.9\textwidth]{assets/figure_127.pdf}
\caption{At horizon $T=100$, maximum error of the Fibonacci classical cycle and the rigorous analytical bound from the convergent denominator. More states reduce fixed-window error, while the infinite-horizon gap remains $1/2$.}\end{figure}
The numerical values are illustrations checked in \path{checks/verify_v39_fibonacci_and_witness.py}; the proof does not rely on a computational experiment.
\section{An explicit strengthening of the three-state finite-sample witness}
Chapter XXXIII proved a polynomial necessary condition for every autonomous three-state stochastic model. For $x=(x_0,\ldots,x_4)$, define
\[d_j=x_{j+1}-x_j,\quad H=d_0d_2-d_1^2,\quad A=d_0d_3-d_1d_2,\quad B=d_1d_3-d_2^2,\quad T_3=H^2-HA+A^2-3HB.\]
Every admissible homogeneous three-state model has $T_3\ge0$. The qubit at phase $\pi/2$ gives $f_*=(0,1/2,1,1/2,0)$, $(H,A,B)=(-1/2,0,-1/2)$, and $T_3(f_*)=-1/2$.
\statement{Theorem}{XLIII.3}{A stronger analytical distance from all three-state chains}
Every homogeneous classical three-state Markov chain, with arbitrary initial distribution and fixed binary effect, producing $g=(g_0,\ldots,g_4)$ satisfies
\begin{equation}\max_{0\le k\le4}|f_{*,k}-g_k|>\frac1{29}.\tag{183}\end{equation}
This strictly improves the previous guarantee $1/384$. The exact optimal gap is not calculated here.
\begin{proof}
If $\|x-f_*\|_\infty\le\varepsilon$, each difference changes by at most $2\varepsilon$. Since the original differences have modulus $1/2$, product perturbation bounds give $|\delta H|,|\delta A|,|\delta B|\le b:=4\varepsilon+8\varepsilon^2$. Expanding at $(-1/2,0,-1/2)$,
\[T_3(x)=-\tfrac12+\tfrac12\delta H+\tfrac12\delta A+\tfrac32\delta B+(\delta H)^2-\delta H\delta A+(\delta A)^2-3\delta H\delta B,\]
so
\begin{equation}T_3(x)\le-\tfrac12+\tfrac52b+6b^2.\tag{184}\end{equation}
At $\varepsilon=1/29$, $b=124/841$ and $1/2-(5b/2+6b^2)=1349/1414562>0$. Hence every point of the closed radius-$1/29$ cube has $T_3<0$, excluding all three-state models.
\end{proof}
\statement{Corollary}{XLIII.4}{A finite-sample test controlling both errors}
Take independent repeated binary measurements at each of five times, with frequencies $\widehat f_k$. Test the simple alternative $f=f_*$ against the composite null of all homogeneous three-state classical chains. Reject the null when $\|\widehat f-f_*\|_\infty\le1/58$. If each time has
\[N\ge\left\lceil1682\log\frac{10}{\alpha}\right\rceil\]
independent repetitions, the worst-case type-I error and the type-II error at $f_*$ are at most $\alpha$. At $\alpha=0.05$, $N=8912$ per time, or 44,560 binary outcomes, suffices.
\begin{proof}
Under the alternative, failure to accept occurs only when the estimation error exceeds $1/58$. Under a null model $g$, Theorem XLIII.3 gives $\|g-f_*\|_\infty>1/29$, so acceptance of the alternative implies $\|\widehat f-g\|_\infty>1/58$. Hoeffding and a union bound give $\Pr(\|\widehat f-f\|_\infty>u)\le10e^{-2Nu^2}$. Use $u=1/58$. The protocol does not cover arbitrary time-inhomogeneous hidden models or devices retaining memory between independent trials.
\end{proof}
\section{Checking exact target matches across four catalogs}
The comparison is by mathematical object, not similarity of titles. The machine-readable register is \path{sources/FOUR_CATALOGUES_CROSSWALK_v39.csv}.
\begin{longtable}{p{.22\textwidth}p{.16\textwidth}p{.53\textwidth}}
\toprule JGMCM&OpenAI Math&Verifiable relation between formulations\\\midrule
001 Haldane&268&Claimed result concerns the same periodic pure spin-1 chain.\\
004 Area Law&265&A general 2D area law is claimed; ProofAtlas marks it ``source-reported proved.''\\
021 Bose gas&267&A separate zero-temperature manuscript result is listed; do not confuse it with positive temperature.\\
029 Kohn--Sham&278&A counterexample is claimed in the same local-potential class.\\
044 $O(3)$&215&Exponential correlation decay is claimed for every finite $\beta$.\\
003 Laughlin&269&No match: bosonic $V_0$ at $\nu=1/2$ versus fermionic $V_1$ at $\nu=1/3$.\\
012 Anderson&261&No match: Gaussian versus uniform potential.\\
017 lattice&090&Minimality and uniqueness are different claims.\\
058 LTI--SDP&---&No exact match to a hierarchy-completeness theorem found.\\
066 HK inverse&365&Calder\'on's inverse boundary problem is not the inverse-density problem.\\\bottomrule
\end{longtable}
\note{Priority qualification.} A catalog issued before October 2026 should not automatically retain an ``open'' label after new manuscripts appear. However, an author's solution claim is not independent expert checking, peer review, or community acceptance. This is only a map of relevant assertions; no prior problem status is independently changed here.
\section{Exact result list and limits of extrapolation}
\begin{enumerate}
\item Proved here: Theorem XLIII.1 gives finite-horizon simulation of an irrational qubit with budget $O(\sqrt{T/\varepsilon})$ and infinite-horizon barrier $1/2$. The former specializes familiar convergent facts; worldwide originality is not established.
\item Proved here: Theorem XLIII.3 gives the stronger explicit distance $1/29$ for a five-sample test against three classical states, with the valid statistical guarantee XLIII.4. The constant is not globally optimal.
\item Not claimed: solutions of Yang--Mills, unforced Navier--Stokes, or P versus NP; general quantum Transformer advantage; or verified worldwide priority.
\item Applied meaning: real recurrent architectures should be compared by admissible state count, horizon $T$, error $\varepsilon$, per-step cost, external-noise robustness, and statistical confidence. Finite-window classical qubit simulation can be unexpectedly cheap; an advantage certificate without a horizon is inapplicable to an experiment.
\end{enumerate}
\note{Sources.} ProofAtlas, \url{https://www.proofatlas.ai/open-problems/} (7 October 2026); J. McM Physics Open Problems, \url{https://github.com/jgmcm/problems/tree/main/catalogue/problems} (20--21 July 2026 statuses); OpenAI Math, \url{https://github.com/openai/math/blob/main/overview.tex} (6 October 2026), families 090, 215, 261, 265, 267--269, 278. Standard methods include continued fractions, Kronecker's density theorem for irrational rotations, Perron--Frobenius theory, and Hoeffding's inequality. These predecessors are not renamed ``newly discovered laws.''
\setcounter{equation}{184}


% 10_stress_topology.tex
\chapter{Invariant Theorems after a Rigorous Stress Test of Four New Hypotheses}
\note{Purpose and status.} Three proposals continuing revision XXXIX (incoming files \path{USER_PROPOSAL_v40_A/B/C.txt}) contain sound research directions but also false universal formulas. This revision preserves those sources as a history of the questions and supplies explicit mathematical counterexamples. No false result is subsequently treated as established.
\section{The exact combinatorial vertex count of a hypercube section}
\statement{Theorem}{XLIV.1}{Generalized central section and exact vertex count}
For odd $n=2r+1\ge3$, let
\[\Pi_n=\left\{x\in\R^n:\sum_{j=0}^{n-1}(-1)^jx_j=\tfrac12\right\},\qquad H_n=[0,1]^n\cap\Pi_n.\]
Then $\dim H_n=n-1$, and
\begin{equation}f_0(H_n)=n\binom{n-1}{(n-1)/2}.\tag{185}\end{equation}
Thus $f_0(H_3)=6$, $f_0(H_5)=30$, and $f_0(H_7)=140$. The claim that $H_n$ is always a $2n$-gon is already false at $n=5$.
\begin{proof}
No cube vertex lies on $\Pi_n$, because the linear form is integer valued there. Section vertices therefore arise on cube edges joining form values 0 and 1. Fix an edge direction. For coefficient $+1$, the remaining $2r$ coefficients comprise $r$ positives and $r$ negatives. The initial sum must be zero, so the number of such edges is $\sum_{a=0}^r\binom ra^2=\binom{2r}r$. For coefficient $-1$, the remaining signs comprise $r+1$ positives and $r-1$ negatives; their fixed sum must be 1. Vandermonde's identity gives $\sum_{a=0}^{r-1}\binom{r+1}{a+1}\binom{r-1}a=\binom{2r}r$. There are $n$ directions. Distinct edges have distinct interior intersections, accounting for all vertices. The cube center belongs to $\Pi_n$, giving full affine dimension $n-1$.
\end{proof}
\begin{figure}[ht]\centering\includegraphics[width=.85\textwidth]{assets/figure_130.pdf}
\caption{Exact vertex count of a hypercube section: already at $n=5$ the count is 30, not 10. Comparison on a logarithmic scale.}\end{figure}
\statement{Remark}{XLIV.2}{Limits of the hexagon theorem XXXVIII}
At $n=3$, the section is two-dimensional and the six observations lie in relative interiors of its sides. At $n=5$, it is four-dimensional with 30 vertices; observations from $W^{(2n)}$ lie on codimension-one faces, not on one-dimensional sides of a $2n$-gon. The six-contact argument of Chapter XXXVIII cannot be transferred by simply replacing 3 with $n$. Neither $\mu_{\rm norm}(W^{(2n)})=2n$ for $n\ge5$ nor a constant normalization premium follows.
\section{Exact minimax barrier for a multidimensional irrational phase}
\statement{Theorem}{XLIV.3}{A quasiperiodic law of finite stochastic memory}
Suppose $1,\alpha_1,\ldots,\alpha_d$ are linearly independent over $\Q$, $b_j>0$, $a\ge0$, and $a+\sum_jb_j\le1$. Set
\begin{equation}f(k)=a+\sum_{j=1}^db_j\sin^2(\pi k\alpha_j),\qquad k\in\N_0.\tag{186}\end{equation}
For all finite homogeneous classical Markov systems $g(k)=e^\top M^k\mu$, with stochastic $M$, $\mu\in\Delta_{m-1}$, $e\in[0,1]^m$,
\begin{equation}\inf_{m<\infty,M,\mu,e}\sup_{k\ge0}|f(k)-g(k)|=\frac12\sum_{j=1}^db_j.\tag{187}\end{equation}
Meanwhile, $f$ has an exact $d$-qubit realization with one admissible binary effect, and ordinary Hankel rank $1+2d$.
\begin{proof}
Peripheral eigenvalues of a finite stochastic matrix are roots of unity and are semisimple because $M^k$ is bounded. Hence $g$ approaches an $L$-periodic response $h$. Kronecker--Weyl implies that within every residue class $k=s+L\ell$, the phase vector is dense in $\mathbb T^d$. Thus $f$ approaches both $a$ and $a+B$, $B=\sum_jb_j$, while $g\to h(s)$, forcing error at least $B/2$. The one-state effect $a+B/2$ attains it.

For the quantum realization, prepare $|+\rangle^{\otimes d}$, apply $U_j=\diag(1,e^{2\pi i\alpha_j})$ on each qubit at each step, and measure $E=aI+\sum_jb_jP_-^{(j)}$. Since $0\le E\le(a+B)I\le I$, this is admissible and gives $f$. Rational independence makes its $1+2d$ nonzero exponential modes distinct, giving the Hankel rank.
\end{proof}
\statement{Remark}{XLIV.4}{Normalization is mandatory}
The proposed unnormalized sum $\sin^2(\pi k\alpha)+\sin^2(\pi k\beta)$ is dense in $[0,2]$ and is not a binary-effect probability. The physical two-qubit probability is half this sum; its exact minimax gap is $1/2$, not 1. The unnormalized sum can be an expectation of an observable in $[0,2]$, which is a different specification.
\section{Exact operator cost of quantum measurement: correcting normalization}
\statement{Proposition}{XLIV.5}{Failure of the proposed inner product}
The proposed identity involving $J_{\rm qm}=\Tr(\Delta\sigma^{-1}\Delta)$ and a squared Frobenius norm already fails for $\rho_+=|0\rangle\langle0|$, $\rho_-=|1\rangle\langle1|$, with the projective measurement in this basis.
\begin{proof}
Here $\sigma=I/2$, $\Delta=\diag(1,-1)$. The proposed formulas give $J_{\rm qm}=4$, $J_{\rm cl}=1$, but the operator inside the right-hand Frobenius norm is identically zero. The equality would require $3=0$.
\end{proof}
\statement{Theorem}{XLIV.6}{An operator-projection identity with Naimark dilation}
Let $\rho_\pm$ be states, $\sigma=(\rho_++\rho_-)/2\succ0$, and $A=(\rho_+-\rho_-)/2$. There is a unique Hermitian $L$ solving
\begin{equation}(\sigma L+L\sigma)/2=A.\tag{188}\end{equation}
For any POVM $\{E_y\}$, put $p_y=\Tr(\sigma E_y)$, $a_y=\Tr(AE_y)$, $J_Q=\Tr(\sigma L^2)$, and $J_C=\sum_{y:p_y>0}a_y^2/p_y$. Let $V:\mathcal H\to\mathcal K$ be a Naimark isometry and $P_y$ orthogonal projectors with $V^\dagger P_yV=E_y$. Write $\widetilde\sigma=V\sigma V^\dagger$, $\widetilde L=VLV^\dagger$, and $\widehat L=\sum_{y:p_y>0}(a_y/p_y)P_y$. Then
\begin{equation}J_Q-J_C=\Tr[\widetilde\sigma(\widetilde L-\widehat L)^2]\ge0.\tag{189}\end{equation}
\begin{proof}
On the real space of Hermitian operators, use the symmetric semidefinite form $\langle X,Y\rangle_{\widetilde\sigma}=\tfrac12\Tr\widetilde\sigma(XY+YX)$. The projectors $P_y$ are orthogonal in this form with squared norms $p_y$, and (188) gives $\langle\widetilde L,P_y\rangle_{\widetilde\sigma}=a_y$. Thus $\widehat L$ is the orthogonal projection onto their span, modulo null directions. Its squared norm is $\sum a_y^2/p_y$, while that of $\widetilde L$ is $\Tr\sigma L^2$. Pythagoras proves the identity.
\end{proof}
\statement{Remark}{XLIV.7}{Exact preservation criterion, distinct from Helstrom's}
Zero loss means $\widetilde L=\widehat L$ in the seminorm induced by $\widetilde\sigma$. This is optimality for quadratic contrast, analogous to SLD/Braunstein--Caves, not a general Helstrom criterion for maximizing total variation. Both have classical literatures and are not claimed as a new universal quantum-information discovery.
\section{The Fisher tensor identity and a false dimensional inference}
\statement{Theorem}{XLIV.8}{Fisher-information dissipation under a parameter-independent channel}
Let $p_\theta(x)$ be a regular dominated parametric model with finite Fisher information and score $S_X=\nabla_\theta\log p_\theta(X)$. Let $Z$ result from a channel independent of $\theta$. When differentiation under the integral is valid,
\begin{equation}I_X(\theta)-I_Z(\theta)=\mathbb E_\theta[\operatorname{Cov}_\theta(S_X\mid Z)]\succeq0.\tag{190}\end{equation}
Equality holds exactly when $S_X$ is $\sigma(Z)$-measurable almost surely.
\begin{proof}
The output score is $S_Z=\mathbb E[S_X\mid Z]$. Apply total covariance and $\mathbb ES_X=0$. Equality means zero conditional covariance in every direction, equivalently reconstructibility of $S_X$ from $Z$.
\end{proof}
\statement{Proposition}{XLIV.9}{A scalar output can preserve two parameters}
Let $X_1\sim N(\mu,e^{2\eta})$ and $X_2\sim N(0,1)$ independently, with $X=(X_1,X_2)$ and scalar compression $Z=X_1$. For $\theta=(\mu,\eta)$,
\[I_Z(\theta)=I_X(\theta)=\begin{pmatrix}e^{-2\eta}&0\\0&2\end{pmatrix},\qquad\rank I_Z=2>\dim Z=1.\]
Thus $\rank I_Z\le\dim Z$ is false: one real random variable can carry Fisher information about several parameters.
\section{Counterexample to the proposed finite-horizon CTMC barrier}
\statement{Theorem}{XLIV.10}{A four-state chain refutes the spectral bound}
For $f(t)=\sin^2t$ on $[0,\pi]$, use a homogeneous continuous-time chain $1\to2\to3\to4$, all rates $\lambda=3/\pi$, with state 4 absorbing, start at 1, and effect $(0,1,1,0)$. Its response is
\begin{equation}g(t)=e^{-\lambda t}\left(\lambda t+\frac{(\lambda t)^2}{2}\right),\tag{191}\end{equation}
and rigorously
\begin{equation}\sup_{0\le t\le\pi}|f(t)-g(t)|<\frac9{20}=0.45.\tag{192}\end{equation}
The incoming proposal's false universal bound at $m=4$, $\omega=1$, $T=\pi$ would require error at least $(1-e^{-2\pi})/2>9/20$.
\begin{proof}
The number of jumps before absorption is a truncated Poisson process, giving (191) from probabilities of states 2 and 3. With $x=3t/\pi$, $\frac d{dx}[e^{-x}(x+x^2/2)]=e^{-x}(1-x^2/2)$ implies $|g'(t)|\le3/\pi<1$; also $|f'(t)|\le1$. The absolute error is 2-Lipschitz. On $t_k=k\pi/100$, $0\le k\le100$, an exact rational certificate using Machin's formula for $\pi$, alternating arctangent series, and Taylor expansions for cosine and exponential proves maximum grid error below 0.414284. Every point is within $\pi/200$ of a node, so the supremum is below $0.414284+\pi/100<0.45$. The code \path{checks/verify_v40.py} compares rational bounds, not rounded doubles. Finally, $e^{-2\pi}<e^{-6}<1/10$ proves the contradiction.
\end{proof}
\begin{figure}[ht]\centering\includegraphics[width=.9\textwidth]{assets/figure_133.pdf}
\caption{An explicit four-state CTMC refutes the proposed universal finite-horizon lower bound. A separate rigorous certificate uses rational arithmetic.}\end{figure}
\statement{Remark}{XLIV.11}{What failed in the spectral proof}
A bound on the ratio of imaginary and real parts of generator eigenvalues does not by itself imply a uniform error lower bound for all outputs $e^\top e^{tQ}\mu$. A finite-time response can be produced by transient real modes and need not contain a complex mode at the target frequency. This argument does not prove $m=\Omega(\omega T)$ for all the stated chains.
\section{Unspecified physical models and future research}
\begin{enumerate}
\item A branching quantum network needs a joint CPTP map from one input to $b$ outputs, a global initial state, measurements, and inter-branch correlations. The single-edge map $w\mapsto\rho e^{i\theta}w$ does not specify it. Arbitrary unknown states cannot be cloned. Without further assumptions, $b\rho^2\cos^2\theta>1$ cannot be declared necessary and sufficient for preserving all quantum information.
\item Cycle holonomies must be products of transport operators in a fixed gauge. Neither the sign of their TV effect nor a universal additive correction $\sum(1-\cos\Phi)$ follows from local edge phases alone; such formulas remain hypotheses.
\item Random AI-representation channels can depend on the learned parameter $\theta$, making parameter independence essential in Theorem XLIV.8. The proof applies to a fixed encoder on an input-distribution family, not automatically to a network whose weights are being changed.
\end{enumerate}
\section{Verification, methodological origins, and the error register}
\begin{itemize}
\item \path{checks/RESULTS_v40.json} records exact enumeration of 6, 30, and 140 section vertices; the explicit $3\ne0$ contradiction; 120 random full-rank quantum experiments with Naimark dilation; rational CTMC grid and Lipschitz certificates; and the Gaussian example with rank-two information matrix.
\item Predecessors are Kronecker--Weyl, stochastic asymptotic periodicity, Vandermonde identities, conditional score variance, the symmetric logarithmic derivative, Braunstein--Caves, and Naimark dilation. The formulas synthesize these mechanisms; independent worldwide priority is not established.
\item OpenAI Math connections: \#119 (information contraction), \#140 (memory and samples), \#144 (Koopman spectra), \#229 (tree reconstruction), \#276 (amplitude damping). Their theorems have much narrower domains than arbitrary quantum gravity or Transformers.
\item The original infinite normalization-premium hypothesis, the proposed POVM loss formula, quantum branching without a joint channel, the general finite-time CTMC inequality, and $\rank I_Z\le\dim Z$ are not accepted as proved results of this monograph.
\end{itemize}

\chapter{A Qubit Topological Specification, Finite-Sample Invariants, and an Audit of 42 Claims}
\note{Status.} This chapter separates (a) the normalized spinor space $S^3\subset\C^2$, (b) the parameter and pure-state sphere $S^2\simeq\mathbb{CP}^1$, and (c) hypothetical physical cosmological space $S_R^3$. The first two are linked by the Hopf fibration; the third does not follow. The results concern a specified two-level Hamiltonian family, not gravitational dynamics. Discrete Chern-number methods were studied by Fukui, Hatsugai, and Suzuki (2005), and topological numbers have been measured in superconducting qubits (Schroer et al., 2014). This chapter develops an experimental specification, an aliasing qualification, and controlled-compression applications; neither the Chern class nor its identities are declared new physics.
\section{A topological model with exactly computable charge}
Let $q\in\Z$, $\theta\in[0,\pi]$, $\phi\in\R/(2\pi\Z)$. For $q\ne0$, put
\begin{equation}r_q(\theta)=\tan^{|q|}(\theta/2),\qquad n_q(\theta,\phi)=\frac1{1+r_q^2}\begin{pmatrix}2r_q\cos(q\phi)\\2r_q\sin(q\phi)\\1-r_q^2\end{pmatrix}.\tag{193}\end{equation}
At the poles use continuous extension. This is a smooth sphere map of degree $q$, with conjugate orientation for negative $q$. For $q=0$, set $n_0=(1,0,0)$ constant. The Hamiltonian
\begin{equation}H_q(\theta,\phi)=-n_q(\theta,\phi)\cdot\sigma\tag{194}\end{equation}
has exact gap 2 everywhere.
\statement{Theorem}{XLV.1}{Chern number of this two-level family}
The energy-$-1$ eigenstates define a complex line bundle whose first Chern number, with the specified orientations, is
\begin{equation}c_1(H_q)=\frac1{2\pi}\int_{S^2}F_q=q.\tag{195}\end{equation}
For $q\ne0$, regular local representatives satisfy
\[\psi_N=\frac{(1,r_qe^{iq\phi})^T}{\sqrt{1+r_q^2}},\qquad\psi_S=e^{-iq\phi}\psi_N,\qquad\psi_N=e^{iq\phi}\psi_S.\]
\begin{proof}
With $A_N=-i\langle\psi_N,d\psi_N\rangle$ and similarly for $S$, direct differentiation gives $A_N=qr_q^2(1+r_q^2)^{-1}d\phi$, $A_S=-q(1+r_q^2)^{-1}d\phi$, and $A_N-A_S=q\,d\phi$. Hence $F_q=q\,d[r_q^2/(1+r_q^2)]\wedge d\phi$. As $\theta$ runs from 0 to $\pi$, the bracket runs from 0 to 1, giving integral $2\pi q$. The constant map at $q=0$ has zero curvature.
\end{proof}
\statement{Remark}{XLV.2}{Do not confuse the two different three-spheres}
$\psi$ and $e^{i\chi}\psi$ represent the same physical pure state. The unit $S^3\subset\C^2$ parameterizes amplitudes; the charge belongs to a bundle over the control-parameter sphere. Neither the radius or metric of cosmological $S_R^3$ nor observed spacetime curvature follows from $c_1(H_q)$.
\section{Finite sampling and guaranteed reconstruction}
At $\theta=\pi/2$, $r_q=1$ and
\begin{equation}z_q(\phi):=\langle\sigma_x\rangle+i\langle\sigma_y\rangle=e^{iq\phi}.\tag{196}\end{equation}
This is observable: the spinor's global $U(1)$ phase cancels. For $\phi_j=2\pi j/M$, $0\le j<M$, $M\ge3$, and nonzero measured $\widehat z_j$, define
\begin{equation}\widehat q=\frac1{2\pi}\sum_{j=0}^{M-1}\operatorname{Arg}\left(\frac{\widehat z_{j+1}}{|\widehat z_{j+1}|}\frac{\overline{\widehat z_j}}{|\widehat z_j|}\right),\qquad\widehat z_M=\widehat z_0.\tag{197}\end{equation}
Here $\operatorname{Arg}\in(-\pi,\pi]$. In exact arithmetic the argument sum is an integer multiple of $2\pi$.
\statement{Theorem}{XLV.3}{A robust finite-grid certificate}
Assume the Hamiltonian belongs to the prespecified family (194), with unknown integer $|q|\le Q$, and every measurement obeys $|\widehat z_j-z_q(\phi_j)|\le\delta<1$. If
\begin{equation}\frac{2\pi Q}{M}+2\arcsin\delta<\pi,\tag{198}\end{equation}
then $\widehat q=q$, without probabilistic assumptions.
\begin{proof}
Tangent geometry of the unit circle gives $|\operatorname{Arg}(\widehat z_j/z_q(\phi_j))|\le\arcsin\delta$. Denote these errors by $\epsilon_j$, cyclically $\epsilon_M=\epsilon_0$. The true increment is $2\pi q/M$; (198) keeps $2\pi q/M+\epsilon_{j+1}-\epsilon_j$ inside $(-\pi,\pi)$, preventing spurious $2\pi$ jumps. Summing telescopes to $2\pi\widehat q=2\pi q$.
\end{proof}
\statement{Corollary}{XLV.4}{Finite statistical cost of the certificate}
At each $\phi_j$, independently measure each of $\sigma_x,\sigma_y$ $N$ times with outcomes $\pm1$. Set $\widehat z_j=\overline X_j+i\overline Y_j$. For $\alpha\in(0,1)$,
\begin{equation}N\ge\left\lceil\frac4{\delta^2}\log\frac{4M}{\alpha}\right\rceil\tag{199}\end{equation}
ensures $\max_j|\widehat z_j-z_q(\phi_j)|\le\delta$ with probability at least $1-\alpha$. Under (198), the algorithm therefore returns the correct integer Chern number with that probability.
\begin{proof}
Hoeffding gives $\Pr(|\overline X-x|>t)\le2e^{-Nt^2/2}$. Coordinate errors at most $\delta/\sqrt2$ imply complex error at most $\delta$. Union over two coordinates and $M$ nodes gives $4Me^{-N\delta^2/4}$.
\end{proof}
\note{Numerical specification.} For $Q=3$, $M=16$, $\delta=0.2$, $\alpha=0.05$, (198) has margin above 1.56 radians and $N=716$ per axis per node suffices. This is $2MN=22\,912$ binary outcomes for one full check of unknown $q$. It is a conservative budget for ideal independent statistics; SPAM errors, drift, and uncalibrated noise need separate estimates.
\begin{figure}[ht]\centering\includegraphics[width=.9\textwidth]{assets/figure_135.pdf}
\caption{Equatorial phases for $q=0,1,2,3$ and finite measurement nodes for $\langle\sigma_x\rangle,\langle\sigma_y\rangle$. In family (194), these determine the first Chern number (source label: Figure XLI.1).}\end{figure}
\section{Two mandatory negative controls}
\statement{Proposition}{XLV.5}{Exact aliasing with too few nodes}
On the $M=4$ grid, $q=3$ and $q=-1$ have identical equatorial data. Thus reconstruction without (198) cannot be guaranteed even with infinitely many measurements per node.
\begin{proof}
At $\phi_j=\pi j/2$, $e^{i3\phi_j}=e^{-i\phi_j}e^{i4\phi_j}=e^{-i\phi_j}$.
\end{proof}
\statement{Proposition}{XLV.6}{Global $c_1$ cannot be recovered from one local region}
Let $U\subset S^2$ leave outside its closure a closed disk with nonempty interior. Smooth two-level Hamiltonian families can agree throughout $U$ while having different Chern numbers. A local spectrum or single point determines neither global Chern class nor physical-space topology.
\begin{proof}
Place inside the unobserved disk a smooth map to $S^2$, constant on the boundary, of any integer degree after collapsing the boundary to a point, using $\pi_2(S^2)=\Z$. Keep the map constant outside. Different degrees give the same data on $U$ but different curvature integrals. $H=-n\cdot\sigma$ retains gap 2.
\end{proof}
\statement{Remark}{XLV.7}{Why one equatorial loop is insufficient in general}
Theorem XLV.3 obtains a global charge only because family (194) is assumed. An arbitrary Hamiltonian can acquire an additional wrapping inside one hemisphere while leaving the equator unchanged. Then $\widehat q$ is unchanged but global $c_1$ differs. Independent verification of an arbitrary model needs surface tomography or gauge-invariant discrete Berry loops, such as Fukui--Hatsugai--Suzuki (2005).
\section{A testable consequence for AI compression}
Encode $Z_{\rm full}(\phi)=(\cos q\phi,\sin q\phi)$. An interface preserving both features with complex error below $\delta$ preserves the charge within this family when (198) holds. If it transmits only $Z_{\rm real}(\phi)=\cos q\phi$, the signals $q$ and $-q$ are indistinguishable at every observation depth. Preserving oriented complex phase is thus a task-dependent interface requirement, not a universal dimensional lower bound for Transformers or evidence of their cosmological nature.
\section{Critical audit of the submitted ``42 laws'' program}
The 42 items are matched to exact specifications in \path{results/V41_42_CLAIMS_AUDIT.csv}. Four corrections are especially significant.
\begin{enumerate}
\item Probability normalization: for $d\ge3$, $\sum_j\sin^2(\pi k\alpha_j)$ may exceed one and cannot be a fixed binary-POVM probability. The physical version is $d^{-1}\sum_j\sin^2(\pi k\alpha_j)$, realized by an averaged local effect on $d$ qubits. For rationally independent frequencies, its exact infinite-horizon minimax error against every finite homogeneous classical chain is $1/2$, by Chapter XL after normalization.
\item False Fisher-rank inference: scalar $Z\sim N(\mu,\sigma^2)$ with unknown $\theta=(\mu,\sigma)$ has $I_Z=\diag(1/\sigma^2,2/\sigma^2)$ and rank two, exceeding $\dim Z=1$. The correct identity $I_X-I_Z=\mathbb E[\operatorname{Cov}(S_X\mid Z)]$ remains valid under differentiation conditions.
\item False universal graph spectral bound: with degree $b+1=5$, no phases, and edge amplitude $\rho=3/10$, the nonbacktracking operator of a finite regular graph has $\rho_{\rm spec}(B)=b\rho=6/5>1/\sqrt b=1/2$, while $b\rho^2=9/25<1$. The proposed equivalence is false even as a translation of the standard tree criterion. On a fixed finite graph, a depth limit $R\to\infty$ also requires a new definition of a graph sequence.
\item Neither Chern class nor Hankel rank determines a cosmic metric. Pure-state spaces, cyclic Markov processes, and $S_R^3$ geometry require independent physical transport operators and data. A finite autonomous classical CTMC has no persistent nonzero frequencies, whereas a quantum GKLS generator may have purely imaginary eigenvalues from its Hamiltonian term.
\end{enumerate}
\section{Proof status and prospects for a physical experiment}
The chapter gives four analytically provable properties of a specified two-level family: exact gap, integer charge, robust finite-grid reconstruction, and a nonasymptotic sample bound. \path{checks/verify_v41_hopf_passport.py} reproduces seven charges $-3\le q\le3$, checks spinor global-phase invariance and exact aliasing, and simulates independent binary measurements. This is simulation, not hardware experimentation.

A real instrument needs preparation of the lower eigenstate of (194), independent Pauli series at every equatorial node, gap and model checks along meridians, preparation/readout calibration, certificate (198), and independent full-surface Chern reconstruction. Only then can physical realization of the invariant be discussed; spatial cosmology $S_R^3$ requires separate data.
\note{Bibliographical predecessors.} T. Fukui, Y. Hatsugai, H. Suzuki, \emph{Chern Numbers in Discretized Brillouin Zone}, J. Phys. Soc. Jpn. 74, 1674--1677 (2005), \url{https://doi.org/10.1143/JPSJ.74.1674}; M. D. Schroer et al., \emph{Measuring a Topological Transition in an Artificial Spin-1/2 System}, Phys. Rev. Lett. 113, 050402 (2014), \url{https://doi.org/10.1103/PhysRevLett.113.050402}. We claim neither discovery of the Chern class nor the first qubit experiment.
\setcounter{equation}{199}


% 11_open_systems.tex
\chapter{Independent Measurement of Global Topology under Unknown Deformation and Hardware Errors}
\note{Result and scope.} Unlike Chapter XLI, this chapter assumes neither the analytical family $H_q$ nor a prescribed bound $Q$ on the unknown charge. Measurements cover the entire closed parameter surface $S^2$. Instead of a known-form family, it requires (i) an independent uniform lower gap bound, (ii) a measured or hardware-bounded rate of eigenprojector variation, and (iii) known preparation, positioning, and readout error bounds. These global conditions cannot be obtained from a finite sample list alone; this is an exact identifiability limit. Nothing identifies $S^3\subset\C^2$ with cosmological physical space.
\section{An unknown smooth field and an unambiguous charge orientation}
Let $\mathcal P\simeq S^2$ be a controlled compact oriented surface with geodesic distance $d$. Consider an unknown Hermitian family
\begin{equation}H(x)=E(x)I_2-\mathbf d(x)\cdot\sigma,\qquad\|\mathbf d(x)\|\ge\gamma/2>0,\qquad n(x)=\frac{\mathbf d(x)}{\|\mathbf d(x)\|}.\tag{200}\end{equation}
The ground-state projector is $P(x)=(I+n(x)\cdot\sigma)/2$. Fix $A=-i\langle\psi,d\psi\rangle$, so
\begin{equation}C(H)=\frac1{2\pi}\int_{\mathcal P}dA=\frac1{4\pi}\int_{\mathcal P}n\cdot(\partial_1n\times\partial_2n)\,dx_1dx_2=\deg n\in\Z.\tag{201}\end{equation}
The opposite convention $A=+i\langle\psi,d\psi\rangle$ reverses the sign; below we use (201).
\statement{Lemma}{XLVI.1}{A hardware regularity certificate from the gap}
If $\|D\mathbf d\|_{\rm op}\le D_*$ along unit tangent vectors, then
\begin{equation}\|n(x)-n(y)\|_{\R^3}\le Ld(x,y),\qquad L=2D_*/\gamma.\tag{202}\end{equation}
\begin{proof}
Differentiating normalization gives $D(\mathbf d/\|\mathbf d\|)=(I-nn^\top)D\mathbf d/\|\mathbf d\|$, of norm at most $D_*/(\gamma/2)$. Integrate along a shortest geodesic.
\end{proof}
\section{A global lattice-degree theorem}
Let $\mathcal T$ triangulate all of $\mathcal P$ with its orientation, with $V$ vertices and $F$ triangles. Let $h$ bound the distance from any point of a triangle to each of its vertices. For geodesic triangles with sides below $\pi/2$, the maximum edge length is admissible: the radius-$h$ ball is geodesically convex and contains all three vertices. At vertex $v_i$, an unnormalized measured Bloch vector $\widehat n_i\in\R^3$ need not lie on the unit sphere. Suppose
\begin{equation}\max_{1\le i\le V}\|\widehat n_i-n(v_i)\|_2\le\eta.\tag{203}\end{equation}
Within each topological triangle use standard barycentric coordinates $\lambda_j$, and set $p(x)=\sum_{j=1}^3\lambda_j(x)\widehat n_{i_j}$, $\widetilde n(x)=p(x)/\|p(x)\|$ when nonzero.
\statement{Theorem}{XLVI.2}{Full-surface topological certification}
If
\begin{equation}Lh+\eta<1,\tag{204}\end{equation}
then $p$ never vanishes, $n$ and $\widetilde n$ are homotopic, and the finite-data integer is
\begin{equation}C(H)=\deg n=\frac1{4\pi}\sum_{(i,j,k)\in\mathcal T}\Omega(\widehat n_i,\widehat n_j,\widehat n_k)\in\Z.\tag{205}\end{equation}
Here $\Omega$ is the branch of the solid angle continuous on a nondegenerate oriented spherical triangle, computed without spinor gauge choices:
\begin{equation}\Omega(a,b,c)=2\operatorname{atan2}\bigl(\det(\widehat a,\widehat b,\widehat c),1+\widehat a\cdot\widehat b+\widehat b\cdot\widehat c+\widehat c\cdot\widehat a\bigr),\quad\widehat a=a/\|a\|.\tag{206}\end{equation}
\begin{proof}
Convexity and the bounds give
\[\|n(x)-p(x)\|\le\sum_j\lambda_j[\|n(x)-n(v_{i_j})\|+\|n(v_{i_j})-\widehat n_{i_j}\|]\le Lh+\eta<1.\]
Hence $\|p(x)\|\ge1-Lh-\eta>0$. The straight interpolation $(1-t)n+tp$ remains within distance less than one of $n$ and cannot vanish; normalization gives a homotopy. Each triangle maps by radial projection of a nonzero affine triangle and has oriented spherical area (206), with a consistent branch. The sum on the closed triangulation is $4\pi\deg\widetilde n=4\pi\deg n$. No one-parameter image family is assumed; arbitrary unknown deformations are allowed within the independent bound (204).
\end{proof}
\statement{Remark}{XLVI.3}{What is actually integer valued}
Expression (205) is mathematically integer valued in exact arithmetic, not through statistical rounding. Machine rounding is justified only after separately bounding accumulated numerical error and excluding nearly antipodal triangles. Mesh size alone is insufficient without $L$. Fukui--Hatsugai--Suzuki (2005) uses related but distinct gauge-invariant link variables rather than this Bloch solid-angle formula.
\section{Measurement, calibration, and preparation errors}
At every vertex, record $X_{i\alpha k}\in\{-1,+1\}$ in three independent batches of $N$ measurements for $\alpha=x,y,z$. The hardware specification is:
\begin{enumerate}
\item SPAM/readout: systematic mean-vector error has norm at most $\varepsilon_{\rm sys}$, obtained from external calibration, not the same data.
\item Positioning: the actual control setting is within geodesic distance $\xi$ of the nominal setting.
\item Preparation: an independent energy test $\Tr(\rho_xH(x))-E_0(x)\le\zeta$ and gap at least $\gamma$ imply Bloch-vector preparation error at most $\varepsilon_{\rm prep}=2\sqrt{\zeta/\gamma}$.
\end{enumerate}
\statement{Theorem}{XLVI.4}{A nonasymptotic experimental specification}
For independent, conditionally identically distributed binary measurements at confidence $1-\alpha$, take
\begin{equation}\eta=\underbrace{\sqrt{\frac6N\log\frac{6V}{\alpha}}}_{\varepsilon_{\rm stat}}+\varepsilon_{\rm sys}+\varepsilon_{\rm prep}+L\xi.\tag{207}\end{equation}
If $L,\gamma,\varepsilon_{\rm sys},\zeta,\xi$ are valid independently justified bounds and $Lh+\eta<1$, (205) returns the correct global charge with probability at least $1-\alpha$. For a chosen statistical allowance $\delta>0$, it suffices that
\begin{equation}N\ge\left\lceil\frac6{\delta^2}\log\frac{6V}{\alpha}\right\rceil,\qquad Lh+\varepsilon_{\rm sys}+\varepsilon_{\rm prep}+L\xi+\delta<1.\tag{208}\end{equation}
The total is $3VN$ independent binary outcomes.
\begin{proof}
Hoeffding at coordinate threshold $\delta/\sqrt3$, followed by a union bound over $3V$ means, gives failure probability at most $6Ve^{-N\delta^2/6}\le\alpha$. Systematic, preparation, and positioning errors add by the triangle inequality; apply Theorem XLVI.2. For preparation, ground-state fidelity $F_x=\Tr(\rho_xP_x)\ge1-\zeta/\gamma$. Fuchs--van de Graaf gives $\frac12\|\rho_x-P_x\|_1\le\sqrt{1-F_x}$, and for a qubit $\|r_x-n_x\|_2=\|\rho_x-P_x\|_1$.
\end{proof}
\statement{Corollary}{XLVI.5}{Homotopy stability under Hamiltonian deformation}
If $H_s(x)=E_s(x)I-\mathbf d_s(x)\cdot\sigma$ is continuous in $s\in[0,1]$ and $\min_{s,x}\|\mathbf d_s(x)\|>0$, all charges agree. For $\mathbf d_s=\mathbf d_0+s\delta\mathbf d$, it suffices that $\sup_x\|\delta\mathbf d(x)\|<\inf_x\|\mathbf d_0(x)\|$.
\begin{proof}
The normalized nonvanishing family is an explicit homotopy, and degree is homotopy invariant.
\end{proof}
\section{Rigorous negative controls}
\statement{Proposition}{XLVI.6}{Finite tomography without a continuity modulus is nonidentifiable}
For any finite nodes $x_1,\ldots,x_V\in S^2$ and integers $q_1,q_2$, smooth fields $n_1,n_2:S^2\to S^2$ can agree at all nodes and have degrees $q_1,q_2$. Both Hamiltonians $H_a=-n_a\cdot\sigma$ have gap 2 everywhere. Infinite statistics at fixed finite nodes therefore cannot determine degree without independent inter-node regularity.
\begin{proof}
Choose a disk without nodes. Keep both fields constant outside, and insert inside a smooth map constant near its boundary of degree $q_1$ or $q_2$ after boundary collapse. Existence follows from $\pi_2(S^2)=\Z$.
\end{proof}
\statement{Proposition}{XLVI.7}{Uncalibrated spatial orientation}
If the measurement basis is reflected by $R\in O(3)$ with $\det R=-1$, then $\deg(R\circ n)=-\deg n$, while local vector lengths and probability errors are unchanged. Charge sign requires independently established orientation of the Pauli-axis triple.
\begin{proof}
Degrees multiply under composition; a sphere reflection has degree $-1$.
\end{proof}
\statement{Remark}{XLVI.8}{Hardware domain of validity}
Confidence $1-\alpha$ does not cover uncontrolled inter-run drift, calibration with unknown orientation sign, an unknown upper bound $L$, or correlated readout errors. Independent-measurement statistics do not replace separate physical systematic-error and gap estimates.
\section{Reproducible numerical illustration and sources}
The independent code \path{checks/verify_v42_unknown_topology.py} uses an oriented geodesic sphere mesh with $V=1026$, $F=2048$, and maximum edge $h\simeq0.153$. Fourteen simulations cover $q=-3,\ldots,3$ for ideal and anisotropically deformed families, without passing deformation parameters to the degree algorithm. Each uses $N=1800$ measurements per component per node at $\alpha=0.05$, totaling $5\,540\,400$ outcomes per model. All certified margins $1-Lh-\eta$ are positive, with minimum above 0.124. This is idealized tomography simulation with bounded systematic bias, not an independent-instrument physical experiment.
\begin{figure}[ht]\centering\includegraphics[width=\textwidth]{assets/figure_140.pdf}
\caption{Full-sphere degree tomography instead of a known equatorial family: 14 simulations and certified threshold margins.}\end{figure}
\note{Independent analysis of the measured table.} The archive also includes \path{checks/certify_v42_from_csv.py}. It receives only Bloch-component estimates at all nodes and external bounds $L,N,\alpha,\varepsilon_{\rm sys},\varepsilon_{\rm prep},\xi$; it is not given $q$. It works for arbitrary unknown fields satisfying the theorem. On a control CSV with $q=2$ known only to the synthetic-data generator, independently calculated $L=2.60$, $N=1800$, and $\varepsilon_{\rm sys}=0.018$, it returns $C=2$ with margin above 0.387. This is reproducible simulation, not real hardware calibration certification. Solid angles use floating-point arithmetic; a rigorous machine proof additionally requires interval checking of rounding and $\operatorname{atan2}$ branches.

\note{Origins of the methods.} Mapping degree, nonvanishing-field homotopies, oriented spherical areas, Hoeffding bounds, and the connection between Fisher information and measurements are established mathematics. Discrete Chern methods include Fukui--Hatsugai--Suzuki, J. Phys. Soc. Jpn. 74 (2005), 1674--1677, \url{https://doi.org/10.1143/JPSJ.74.1674}. This revision contributes a unified testable specification with separate gap, regularity, preparation, systematic-error, and sampling bounds, plus negative controls. Historical originality of this combination is unestablished. No result proves nonlocal quantum gravity, cosmological $S_R^3$ topology, spacetime curvature, or an advantage for a real LLM.

\chapter{Two Spectral Gaps: Stationary-State Topology and the Impossibility of Stationary-Mode Charge}
\note{Scientific setting.} New user proposals on non-Hermitian stationary GKLS connections, nonlinear tree reconstruction, and a streaming AR(1) limit are checked directly. Three distinct objects must be separated: the Liouvillian zero eigenspace in operator space, a density-matrix eigenspace in qubit Hilbert space, and the measurable Bloch-map degree. Their topological classes generally differ. We provide a concrete GKLS family, a finite-sample certificate, and exact negative controls.
\section{Zero topology of the unique stationary Liouvillian mode}
Let $\mathcal L_x$, $x\in S^2$, be a smooth finite-dimensional trace-preserving GKLS family with a unique smoothly varying normalized stationary state $\rho_*(x)$ and a simple zero eigenvalue. The left zero eigenoperator is always $I$.
\statement{Theorem}{XLVII.1}{Triviality of the unique stationary eigenline}
The line $\ker\mathcal L_x=\C\rho_*(x)$ is a topologically trivial complex line bundle and has Chern number zero. With $\Tr\rho_*=1$, the proposed biorthogonal connection vanishes:
\[A_{\rm steady}=\Tr(I\,d\rho_*)=d\Tr\rho_*=0,\qquad c_1(\ker\mathcal L)=0.\]
\begin{proof}
$\rho_*(x)$ is a global smooth nowhere-zero section because its trace is one. Such a section trivializes a complex line bundle. $I$ is a global left biorthogonal section, $\Tr(I\rho_*)=1$; the normalized connection and its curvature are zero.
\end{proof}
\note{Boundary.} This concerns precisely the simple zero mode of a unique normalized stationary state. Nonzero modes, whether simple or degenerate, higher-dimensional stationary subspaces, and eigenlines of $\rho_*$ as a matrix are different objects; their topology is not determined by this theorem.
\section{A global GKLS generator with arbitrary eigenspace charge}
Let $n_q:S^2\to S^2$ be smooth of degree $q\in\Z$. In place of the potentially pole-nonsmooth expression $\sin\theta\cos(q\phi)$, use
\[d_q(x,y,z)=\bigl(\operatorname{Re}(x+i\operatorname{sgn}(q)y)^{|q|},\operatorname{Im}(x+i\operatorname{sgn}(q)y)^{|q|},z\bigr),\quad n_q=d_q/\|d_q\|\quad(q\ne0),\]
and the constant $(0,0,1)$ for $q=0$. The polynomial never vanishes on the sphere and has degree $q$, by equatorial winding and continuous pole interpolation. Set
\[\rho_{r,q}(x)=\tfrac12(I+rn_q(x)\cdot\sigma),\quad |r|\le1,\qquad\mathcal L_{x,r,q}(X)=\gamma[\Tr(X)\rho_{r,q}(x)-X],\quad\gamma>0.\]
\statement{Theorem}{XLVII.2}{Separating dissipative and density gaps}
This is a valid globally smooth GKLS generator with unique stationary state, and
\[\spec\mathcal L_{x,r,q}=\{0,-\gamma,-\gamma,-\gamma\},\qquad\Delta_{\rm diss}=\gamma,\qquad\Delta_{\rm purity}=|r|.\]
For $r>0$, the leading eigenline of $\rho_{r,q}$ has Chern number $q$; for $r<0$, it has $-q$. At $r=0$, $I/2$ has degenerate spectrum and no defined leading eigenline, although the dissipative gap remains $\gamma$. Throughout, $c_1(\ker\mathcal L)=0$.
\begin{proof}
In a fixed computational basis, let $R(x)=\sqrt{\rho_{r,q}(x)}$ and define four jumps $J_{ab}=\sqrt\gamma R|a\rangle\langle b|$. Then $\sum_{ab}J_{ab}XJ_{ab}^\dagger=\gamma\Tr(X)\rho_{r,q}$ and $\sum_{ab}J_{ab}^\dagger J_{ab}=\gamma I$, giving Lindblad form without local eigenspinors. On $\C\rho_{r,q}\oplus\{X:\Tr X=0\}$, the generator is $0\oplus(-\gamma I_3)$. For $r>0$, density projectors are $(I\pm n_q\cdot\sigma)/2$, with
\[F_+=-i\Tr(P_+dP_+\wedge dP_+)=\tfrac12n_q\cdot(\partial_\theta n_q\times\partial_\phi n_q)\,d\theta\wedge d\phi.\]
Its integral divided by $2\pi$ is $q$. For negative $r$, the leading projector is $P_-$ and the number reverses. At zero, density eigenvalues coincide. The zero Liouvillian line is trivial by Theorem XLVII.1.
\end{proof}
\statement{Corollary}{XLVII.3}{A topological transition without dissipative-gap closure}
For $q\ne0$, crossing $r=0$ changes the leading density-bundle charge from $-q$ to $q$, with the entire Liouvillian spectrum and exponential relaxation rate unchanged. Thus a positive dissipative gap does not certify $c_1$; the purity gap must be controlled independently.
\begin{figure}[ht]\centering\includegraphics[width=\textwidth]{assets/page_143_1.png}
\caption{Two independent gaps: relaxation toward a unique stationary state has constant rate $\gamma$, but the density matrix's eigenline loses topological definition at $r=0$. The diagram uses exact analytical formulas, not laboratory data.}\end{figure}
\section{Exact dynamics and an independent finite-point specification}
For arbitrary $X$,
\[e^{t\mathcal L}(X)=e^{-\gamma t}X+(1-e^{-\gamma t})\Tr(X)\rho_{r,q}.\]
For an initial physical state with Bloch vector $v_0$, $v_t=(1-e^{-\gamma t})rn_q+e^{-\gamma t}v_0$, so $\|v_t-rn_q\|\le(1+|r|)e^{-\gamma t}\le2e^{-\gamma t}$.
\statement{Theorem}{XLVII.4}{A full-sphere experimental certificate for two gaps}
Suppose the unknown stationary Bloch field $v:S^2\to\R^3$ has independently confirmed bounds $\|v(x)\|\ge r_{\min}>0$ and Lipschitz constant $L_v$. Use a full oriented sphere triangulation with cell diameter at most $h$ and $V$ vertices. Estimate each component at each vertex from $N$ independent $\pm1$ measurements, with failure budget $\alpha$ and external systematic, preparation, and positioning bounds $\varepsilon_{\rm sys},\varepsilon_{\rm prep},\xi$. For preparation by reset relaxation for time $t$, with $\gamma\ge\gamma_{\min}>0$, set
\[\eta_N=\sqrt{\frac{6\log(6V/\alpha)}N}+\varepsilon_{\rm sys}+\varepsilon_{\rm prep}+L_v\xi+2e^{-\gamma_{\min}t}.\]
If $L_vh+\eta_N<r_{\min}$, then with probability at least $1-\alpha$, the leading stationary-density projector's Chern number equals the integer lattice degree computed from normalized $\widehat v_i$. The procedure uses neither the analytical form $n_q$ nor true $q$.
\begin{proof}
Hoeffding and a union bound over $3V$ means give simultaneous Euclidean error $\sqrt{6\log(6V/\alpha)/N}$. Add other errors by the triangle inequality. A nonnegative barycentric interpolant $\widetilde v(x)=\sum_iw_i(x)\widehat v_i$ obeys $\|\widetilde v-v\|\le\eta_N+L_vh<r_{\min}\le\|v\|$. The straight homotopy avoids zero, so normalizations have equal degree. Nonvanishing on each triangle permits degree computation by oriented spherical solid angles with the correct $2\operatorname{atan2}$ branch. For a qubit, this is the leading density-eigenbundle Chern number.
\end{proof}
\statement{Remark}{XLVII.5}{}
External bounds $L_v,r_{\min},\gamma_{\min},\varepsilon_{\rm sys}$ require confirmation outside charge reconstruction. Hoeffding does not establish them. Even a large dissipative gap does not exclude $r_{\min}=0$. This is a mathematical specification of a simulated experiment, not a laboratory report.
\section{Three mandatory negative controls from the new proposals}
\statement{Proposition}{XLVII.6}{Symmetric pairwise Potts statistics are uninformative}
For symmetric $q$-state broadcasting on a $b$-ary tree with $M=\lambda I+(1-\lambda)\one\one^\top/q$, define
\[Q_t=\sum_{u\ne v\in\partial T_t}(q\one_{X_u=X_v}-1).\]
Its conditional law is independent of the root label. More generally, every statistic $F(X_{\partial T_t})$ invariant under simultaneous permutation of all $q$ labels satisfies $I(X_\varnothing;F)=0$ for a uniform root. Moreover,
\[\mathbb E[Q_t\mid X_\varnothing=k]=(q-1)(b-1)\sum_{h=0}^{t-1}b^{2t-h-1}\lambda^{2(t-h)}\]
for every $k$.
\begin{proof}
A simultaneous label permutation maps the process conditioned on root $k$ to that conditioned on any other label while preserving pair equalities. Thus conditional laws coincide. A leaf pair with common ancestor at depth $h$ has expected contribution $(q-1)\lambda^{2(t-h)}$, and there are $(b-1)b^{2t-h-1}$ ordered pairs of this type. Summing gives the formula.
\end{proof}
\note{Arithmetic control.} $b^2\lambda^4=(b\lambda^2)^2$, so $b^2\lambda^4>1$ is equivalent to $b\lambda^2>1$ and cannot establish purported subcritical reconstruction. This error does not replace the serious literature on below-KS reconstruction in other models and using other statistics.
\statement{Remark}{XLVII.7}{An incorrect AR(1) information-budget estimate}
If the unknown parameter is exactly one of two hypotheses $B\in\{0,1\}$, then $I(B;M_T)\le H(B)\le1$ bit for every $\varepsilon$. A requirement of $\Omega(\log(1/\varepsilon))$ information bits is therefore impossible as $\varepsilon\downarrow0$. The submitted proof of $T=\Omega(\log(1/\varepsilon)/(\varepsilon S))$ is invalid. The standalone nonasymptotic inequality is neither proved here nor declared false in every formulation.
\section{Exact PSD factorization as an independent physical-realizability check}
The six-label table $W^{(6)}_{ij}=\sin^2[\pi(i-j)/6]$ has a real PSD factorization of size two:
\[A_j=(I+\cos\theta_j\sigma_z+\sin\theta_j\sigma_x)/2,\quad B_i=I-A_i,\quad\theta_j=2\pi j/6,\quad\Tr(A_jB_i)=W^{(6)}_{ij}.\]
Since $\rank W^{(6)}=3>1$, its PSD rank is exactly two. This is an independent instance of classical PSD-factorization theory, not a proof of new complexity classes, and does not change the preceding geometric theorem about $\mu_{\rm norm}$.
\section{Reproducibility, priority limits, and catalog connections}
\path{checks/verify_v43_steady_topology.py} checks the global GKLS specification, operator spectra, real PSD factorization, tree symmetry, and full-sphere mixed-state charge simulations, including mandatory refusal to certify at $r=0$. True charge is used only in an external accuracy test, not by the degree estimator. No independent hardware data are available.

Mixed-state topology and dissipative preparation of topological phases predate this revision: A. Rivas, O. Viyuela, M. A. Martin-Delgado, Phys. Rev. B 88, 155141 (2013), DOI 10.1103/PhysRevB.88.155141; J. C. Budich, P. Zoller, S. Diehl, Phys. Rev. A 91, 042117 (2015), DOI 10.1103/PhysRevA.91.042117. These finite-dimensional reset constructions and certification conditions claim neither worldwide priority nor Millennium-Problem solutions. OpenAI Math 119, 229, 276, and 365 are catalog reference points, not direct proofs of our theorems.

\chapter{Topological Distinguishability and GKLS Stability under Hardware Perturbations}
\note{Setting and origin.} This chapter examines the user's 42 hypotheses concerning non-Hermitian topology, Markov memory, and spectral geometry. It isolates one mathematically closed package: topological separation of density families in trace norm, norm-controlled stability of stationary GKLS states, an exact phase boundary for additional reset, orientation calibration, and finite-sample certification. Degree, trace norm, and contractive semigroups have established predecessors; global priority of their particular combination is unestablished.
\section{Topological separation permits perfect qubit discrimination}
Let $a,b:S^2\to\R^3\setminus\{0\}$ be continuous on an oriented sphere, with normalized directions $\widehat a,\widehat b$ and positive minima $r_a,r_b$. Set $\rho_a=(I+a\cdot\sigma)/2$, $\rho_b=(I+b\cdot\sigma)/2$; physical states additionally require $\|a\|,\|b\|\le1$.
\statement{Theorem}{XLVIII.1}{Antipodal necessity and topological distance}
If $\deg\widehat a\ne\deg\widehat b$, some $x_*$ satisfies $\widehat a(x_*)=-\widehat b(x_*)$. Consequently,
\begin{equation}\sup_{x\in S^2}\|\rho_a(x)-\rho_b(x)\|_1\ge r_a+r_b.\tag{209}\end{equation}
For constant equal polarizations $r$, the supremum is exactly $2r$.
\begin{proof}
Without antipodal directions, normalized straight interpolation between the two unit fields is a global homotopy: a convex combination can vanish only at $t=1/2$ with opposite vectors. This contradicts unequal degrees. At the antipodal point, $a=r'_au$, $b=-r'_bu$, so the difference has trace norm $r'_a+r'_b\ge r_a+r_b$. For constant radii $r$, trace norm is everywhere at most $2r$ and attains it there.
\end{proof}
\statement{Corollary}{XLVIII.2}{Helstrom bound and an orthogonal quantum section}
For equally likely states at $x_*$,
\begin{equation}P_{\rm succ}(x_*)=\frac12+\frac14\|\rho_a(x_*)-\rho_b(x_*)\|_1\ge\frac12+\frac14(r_a+r_b).\tag{210}\end{equation}
If both states are pure, they are orthogonal at $x_*$ and perfectly distinguishable.
\note{Operational meaning and restriction.} The families must share the same physical parameter base and calibrated correspondence of settings. Chern numbers from unrelated experiments cannot be compared without it. The result asserts existence of a setting, not the ability to locate unknown $x_*$ without regularity and sufficient measurements. Its classical predecessor is homotopy invariance of sphere-map degree.
\statement{Corollary}{XLVIII.3}{Exact distance between reset-generator topological sectors}
For reset generators $\mathcal L_{a,x}(X)=\gamma[\Tr(X)\rho_a(x)-X]$ and $\mathcal L_{b,x}$ at common rate $\gamma$, common constant polarization $r>0$, and different leading-eigenprojector Chern numbers,
\begin{equation}\sup_{x\in S^2}\|\mathcal L_{a,x}-\mathcal L_{b,x}\|_{1\to1}=2\gamma r.\tag{211}\end{equation}
\begin{proof}
Their difference is $X\mapsto\gamma\Tr(X)(\rho_a-\rho_b)$. Since $|\Tr X|\le\|X\|_1$, with equality on density matrices, its induced norm is exactly $\gamma\|\rho_a-\rho_b\|_1$. Apply Theorem XLVIII.1.
\end{proof}
\section{Stationary GKLS states and an explicitly bounded perturbation}
For smooth $\rho_x=(I+a(x)\cdot\sigma)/2$, consider
\begin{equation}\mathcal L_{0,x}(X)=\gamma[\Tr(X)\rho_x-X],\qquad\gamma>0.\tag{212}\end{equation}
Let $\mathcal E_x$ be a smooth Hermiticity-preserving, trace-annihilating superoperator family such that $\mathcal L_x=\mathcal L_{0,x}+\mathcal E_x$ remains a CPTP-semigroup generator. Write $\kappa=\sup_x\|\mathcal E_x\|_{1\to1}$.
\statement{Theorem}{XLVIII.4}{Stationary-topology robustness against arbitrary GKLS perturbations}
If $\kappa<\gamma$, each generator has a unique stationary state $\widetilde\rho_x$. For normalized states $X,Y$,
\begin{align}\|e^{t\mathcal L_x}(X-Y)\|_1&\le e^{-(\gamma-\kappa)t}\|X-Y\|_1,\tag{213}\\
\sup_x\|\widetilde\rho_x-\rho_x\|_1&\le\kappa/\gamma.\tag{214}\end{align}
If $r_0=\min_x\|a(x)\|_2>\kappa/\gamma$, the leading eigenbundles of $\widetilde\rho_x$ and $\rho_x$ have the same degree, and the density gap is at least $r_0-\kappa/\gamma$.
\begin{proof}
On traceless $Z$, $\mathcal L_{0,x}Z=-\gamma Z$, and $\mathcal E_x$ preserves tracelessness. Thus $e^{t\mathcal L_x}Z=e^{-\gamma t}e^{t\mathcal E_x}Z$, giving (213). A finite-dimensional CPTP semigroup has a stationary state, for example via a subsequential Ces\`aro-average limit; contraction makes it unique. Stationarity gives $\gamma(\widetilde\rho_x-\rho_x)=\mathcal E_x(\widetilde\rho_x)$, and $\|\widetilde\rho_x\|_1=1$ proves (214). For qubits, trace-norm state difference equals Euclidean Bloch-vector difference. The perturbed field and its straight homotopy from $a$ stay at distance at least $r_0-\kappa/\gamma$ from zero; degree is preserved.
\end{proof}
\statement{Proposition}{XLVIII.5}{Lipschitz stability under control-parameter drift}
If $\|\rho_x-\rho_y\|_1\le L_0d(x,y)$ and $\|\mathcal E_x-\mathcal E_y\|_{1\to1}\le Kd(x,y)$, then for $\kappa<\gamma$,
\begin{equation}\|\widetilde\rho_x-\widetilde\rho_y\|_1\le L_*d(x,y),\qquad L_*=\frac{\gamma L_0+K}{\gamma-\kappa}.\tag{215}\end{equation}
\begin{proof}
Subtract stationary equations and move $\mathcal E_x(\widetilde\rho_x-\widetilde\rho_y)$ to the left. The right side has norm at most $(\gamma L_0+K)d(x,y)$, while the left is bounded below by $(\gamma-\kappa)\|\widetilde\rho_x-\widetilde\rho_y\|_1$.
\end{proof}
\section{A phase transition under additional physical reset}
For $q>0$, $0<r\le1$, use
\begin{equation}n_q(x)=\frac{(\operatorname{Re}(x_1+ix_2)^q,\operatorname{Im}(x_1+ix_2)^q,x_3)}{\sqrt{(x_1^2+x_2^2)^q+x_3^2}}.\tag{216}\end{equation}
For negative $q$, use exponent $|q|$ and reverse the imaginary component. The map is smooth of degree $q$: poles are fixed and the equator winds $q$ times. With $\tau_+=(I+\sigma_z)/2$, set
\begin{equation}\mathcal L_{x,\delta}(X)=\gamma[\Tr(X)\rho_x-X]+\delta[\Tr(X)\tau_+-X],\quad\rho_x=\tfrac12(I+rn_q(x)\cdot\sigma),\quad\delta\ge0.\tag{217}\end{equation}
\statement{Theorem}{XLVIII.6}{Exact density-topology transition boundary}
For all $\delta\ge0$, this GKLS generator has unique stationary state and spectrum
\begin{equation}\rho_{*,x,\delta}=\frac{\gamma\rho_x+\delta\tau_+}{\gamma+\delta},\qquad\spec\mathcal L_{x,\delta}=\{0,-(\gamma+\delta),-(\gamma+\delta),-(\gamma+\delta)\}.\tag{218}\end{equation}
Its leading density projector has Chern number
\begin{equation}C(\delta)=\begin{cases}q,&0\le\delta<\gamma r,\\\text{undefined},&\delta=\gamma r,\\0,&\delta>\gamma r.\end{cases}\tag{219}\end{equation}
The minimum density gap is exactly $g_\rho(\delta)=|\gamma r-\delta|/(\gamma+\delta)$, while the dissipative gap increases as $g_{\mathcal L}=\gamma+\delta$.
\begin{proof}
The sum is again a reset generator, at rate $\gamma+\delta$ and target (218). Its Bloch vector is $(\gamma rn_q+\delta e_z)/(\gamma+\delta)$. Below threshold, adding the offset through $s\delta e_z$ never crosses zero and preserves degree $q$. Above threshold, the third component is positive everywhere, so the map lies in a hemisphere and has degree zero. At equality it vanishes where $n_q=-e_z$, so the leading projector is undefined. The minimum numerator length is $|\gamma r-\delta|$, attained at that same point; division gives the density gap.
\end{proof}
\begin{figure}[ht]\centering\includegraphics[width=\textwidth]{assets/page_148_1.png}
\caption{Exact phase threshold for two physical reset channels: the density gap closes at $\delta=\gamma r$ while the relaxation rate $\gamma+\delta$ continues growing. Analytical model, without experimental data.}\end{figure}
\section{Calibration orientation and an experimental certificate}
\statement{Proposition}{XLVIII.7}{Affine calibration and charge sign}
For continuous nonzero $v:S^2\to\R^3$, $A\in GL(3,\R)$, and $c\in\R^3$, if $\|c\|_2<\sigma_{\min}(A)\min\|v\|_2$, then
\begin{equation}\deg\frac{Av+c}{\|Av+c\|}=\operatorname{sgn}(\det A)\deg\frac v{\|v\|}.\tag{220}\end{equation}
\begin{proof}
Straight removal of $c$ avoids zero. In the polar decomposition $A=OP$, deform positive $P$ to $I$. The orthogonal factor changes degree by exactly $\operatorname{sgn}\det O$.
\end{proof}
\statement{Corollary}{XLVIII.8}{Finite samples under unknown GKLS drift}
Use an oriented $S^2$ triangulation whose vertices are within geodesic distance $h$ of every point of their face. At each node take $N$ independent measurements of each Pauli component. For total systematic error $\eta_{\rm sys}$ and preparation error $\eta_{\rm prep}$, put
\begin{equation}\eta_N=\sqrt{\frac6N\log\frac{6V}\alpha}+\eta_{\rm sys}+\eta_{\rm prep}.\tag{221}\end{equation}
Under Theorem XLVIII.4 and Proposition XLVIII.5, if
\begin{equation}L_*h+\eta_N<r_0-\kappa/\gamma,\tag{222}\end{equation}
then with probability at least $1-\alpha$, the leading-eigenspace Chern number of the actual unknown stationary state equals the degree of the radially normalized piecewise-affine tomography interpolant. Its oriented face solid-angle sum divided by $4\pi$ gives that integer. After finite relaxation time $t$, the conservative bound $\eta_{\rm prep}\le2e^{-(\gamma-\kappa)t}$ is admissible.
\begin{proof}
Hoeffding and a union bound over $3V$ components give simultaneous three-dimensional statistical error at most the square-root term. Within every face, affine interpolation differs from the true field by at most $L_*h+\eta_N$. Condition (222) makes the straight homotopy nonvanishing. Radial projection computes degree and excludes irreducible solid-angle branch ambiguity. The relaxation estimate follows from (213) and initial trace distance at most 2.
\end{proof}
\note{Critical safeguards against incorrect generalizations in the proposed register.}
\begin{enumerate}
\item For the unique normalized stationary mode, $\Tr(I\,d\rho_*)=0$; isolating exceptional points cannot give it a nonzero Liouvillian Chern number.
\item A spectral-radius condition on a weighted Hashimoto operator cannot imply a universal TV-reconstruction threshold without a quantum broadcasting channel and joint output correlations. A necessarily negative ``loop correction'' is also unproved.
\item Pairwise leaf coincidences in symmetric Potts broadcasting are invariant under global label permutation and cannot recover the absolute root symbol. The proposed expectation formula also omitted lowest-common-ancestor depth.
\item A dissipative gap alone does not guarantee positive topological margin. Density gap, a global continuity modulus, and axis-orientation calibration are also required.
\item The Fibonacci identity uses $\beta=1-\varphi^{-4}=(3\sqrt5-5)/2\approx0.854101966$. The register's expression $(7\sqrt5-15)/2$ is incorrect.
\item For $d>1$, $\sum_j\sin^2(\pi k\alpha_j)$ is not a single binary probability; divide by $d$ or specify a real multichannel signal.
\end{enumerate}
\note{Reproducibility and physical interpretation.} Mathematical certificates and independent exact/numerical tests are in \path{checks/verify_v44_topological_robustness.py}; the graph is \path{FIGURE_v44_PHASE_TRANSITION.png}; the 42 incoming proposals are preserved unedited in \path{sources/}. The model is finite-dimensional quantum theory, not laboratory measurement, spacetime-$S^3$ topology, or a Millennium-Problem solution. Theorems use classical degree, operator-norm bounds, and quantum distinguishability; independent priority requires separate assessment.
\setcounter{equation}{222}


% 12_horizon_gravity.tex
\chapter{Source Audit, Horizon Modular Energy, and a Verifiable Information Certificate}
\section{A unified source: correcting omitted chapters}
This revision corrects a release error: in the separate sources for v38--v44, new Chapters XXXVIII--XLIV were connected through \verb|\input| rather than included in the body of the file delivered to the user. The original full v42 archive also contained an outdated main source that omitted Chapter XLII, although the separately published source included it. In this release, \path{META_CAUSAL_MASTER.tex} is an expanded main text containing all these chapters directly, without \verb|\input|, and compiles with the images supplied in the archive. The PDF page count and source metadata were checked after recompiling that exact file. Correcting file composition is not a new mathematical theorem.
\section{The Dorau--Much result and its actual scope}
This section uses P. Dorau and A. Much, \emph{From Quantum Relative Entropy to the Semiclassical Einstein Equations}, Physical Review Letters 136, 091602 (2026), \url{https://doi.org/10.1103/lmq8-nsty}, published 2 March 2026. The published six-page version was examined, especially equations (3)--(13), (18)--(20), and (23)--(28). It considers a real Klein--Gordon field on a local Rindler horizon, a quasifree Hadamard/KMS vacuum, and coherent Weyl-algebra excitations. In this setting, modular theory yields
\begin{equation}S_{\rm rel}(\omega_0\|\omega_\phi)=-2\pi\int_{H_B^R}U[\partial_U\phi(U,s)]^2\,dU\,d\operatorname{vol}_S,\qquad U<0.\tag{223}\end{equation}
Nonnegativity relies on the sign of $-U$ and admissible profiles. This is neither a formula for arbitrary noncoherent states nor finite-dimensional von Neumann entropy of a local type-III algebra. The article explicitly restricts itself to coherent states, with QFT used in a local Minkowski approximation.

The geometric step additionally identifies area change with entropy. In the article's notation,
\begin{equation}\delta A=\frac{\alpha}{2\pi}S_{\rm rel},\tag{224}\end{equation}
where $\alpha=8\pi$ in units $G=\hbar=c=k_B=1$, adopting Bekenstein--Hawking normalization $\delta A=4S_{\rm rel}$. In particular, its equation (19) posits an induced-metric variation ansatz yielding (20); (223) alone does not prove that ansatz. Under linearized Raychaudhuri assumptions and area identification for all local null directions, the authors obtain $R_{ab}-\tfrac12Rg_{ab}+\Lambda g_{ab}=\alpha\langle:T_{ab}:\rangle$. An approximation-free proof would require control of curvature corrections, vanishing boundary terms, horizon-generator normalization, and extension beyond coherent excitations. None follows from a single META--CAUSAL observable information invariant.
\section{An exact finite-dimensional specification of quadratic modular energy}
Let $U\in[-1,0]$, $s\in S$, and let real $h(s)$ satisfy $\int_Sh^2\,d\operatorname{vol}_S=1$. Choose real smooth longitudinal profiles $f_1,\ldots,f_d$ of compact support and define
\[\phi_c(U,s)=h(s)\sum_{j=1}^dc_jf_j(U),\qquad G_{ij}=\int_{-1}^0(-U)f'_i(U)f'_j(U)\,dU.\]
$G$ is positive semidefinite, and Dorau--Much equation (12) directly gives the exact family identity
\begin{equation}S_{\rm rel}(c)=2\pi c^TGc.\tag{225}\end{equation}
For a nondegenerate mode family, $G\succ0$ defines $\|c\|_G^2=c^TGc$. An invertible real basis change $\widetilde f=Af$ gives $\widetilde c=A^{-T}c$, $\widetilde G=AGA^T$, leaving (225) unchanged. This is a testable basis-independent quadratic invariant, not a consequence of global spatial topology $S^3$.
\statement{Theorem}{XLIX.1}{A certificate from incomplete modular observations}
Let $G\in\R^{d\times d}$ be symmetric positive definite, $C\in\R^{r\times d}$ have full row rank, and the linear measurement data be $y=Cc$. Then
\begin{equation}S_{\rm rel}(c)\ge2\pi y^T(CG^{-1}C^T)^{-1}y.\tag{226}\end{equation}
Equality occurs at $c_*=G^{-1}C^T(CG^{-1}C^T)^{-1}y$, and
\begin{equation}c^TGc=y^T(CG^{-1}C^T)^{-1}y+(c-c_*)^TG(c-c_*).\tag{227}\end{equation}
\begin{proof}
Lagrange multipliers for minimizing $c^TGc/2$ subject to $Cc=y$ give $Gc=C^T\lambda$, yielding $c_*$ and $Cc_*=y$. For $v=c-c_*$, $Cv=0$, hence $v^TGc_*=v^TC^T\lambda=0$. Expand $\|c_*+v\|_G^2$.
\end{proof}
For noisy $\widetilde y=Cc+e$ with independently certified $\|e\|_{H^{-1}}\le\eta$, $H=CG^{-1}C^T\succ0$, the reverse triangle inequality gives
\begin{equation}S_{\rm rel}(c)\ge2\pi[\max\{0,\|\widetilde y\|_{H^{-1}}-\eta\}]^2.\tag{228}\end{equation}
Here $\|v\|_{H^{-1}}^2=v^TH^{-1}v$. If only $\|e\|_2\le\eta_2$ is known, take $\eta=\eta_2/\sqrt{\lambda_{\min}(H)}$. Independent error control is essential; without it, the lower bound can be false.
\section{A reproducible rational-check protocol}
For $x=-U\in[0,1]$, choose $f_j(-x)=x^{j+2}(1-x)^2$, $j=0,1,2,3$. Then $G_{ij}=\int_0^1x\,\partial_xf_i\partial_xf_j\,dx$ is rational. The program \path{checks/verify_v45_horizon.py} computes it symbolically and checks Sylvester's criterion, minimum (226), residual orthogonality (227), rational basis-change stability, and nonidentifiable directions $\ker C\ne\{0\}$.

For $C=\left(\begin{smallmatrix}1&0&1&0\\0&1&1&0\end{smallmatrix}\right)$ and $c=(1,-1,2,1/2)^T$, exact arithmetic gives
\[c^TGc=\frac{142}{9009},\qquad y^T(CG^{-1}C^T)^{-1}y=\frac{10789}{836755425}.\]
The strict lower bound is much smaller than the actual energy. Its deficit reflects unreconstructed modes and cannot be declared gravitational energy obtained from observations alone.
\section{A valid bridge to geometry and information compression}
Positivity and monotonicity of relative entropy under restricting the observable algebra or applying a physical data channel imply
\[0\le S_{\rm rel}^{\rm measured}\le S_{\rm rel}^{\rm horizon}.\]
Thus (226) certifies information reconstructibility of part of the horizon excitation. Only with the additional assumption (224) does it imply the conditional area bound $\delta A\ge\alpha y^T(CG^{-1}C^T)^{-1}y$. It does not derive $\alpha$, prove semiclassical Einstein equations, recover global topology from a local horizon, or identify the Hopf bundle's Chern class with physical-space geometry.

\note{Negative controls.} (1) If $\ker C$ is nontrivial, $c$ and $c+v$, $v\in\ker C$, give identical data but may have different horizon energies. (2) Globally different spacetimes can have locally isometric Rindler neighborhoods and equal (223); a local signal does not distinguish $S^3$ from other topologies. (3) Changing the sign or normalization of the null coordinate generator without consistently transforming modular flow destroys the interpretation. The article separately emphasizes nonunique Rindler boost-field normalization. (4) Without (224), the energy identity gives no quantitative physical area change. These are parts of the theorem's specification, not minor experimental details.

\note{Result status.} The modular relative-entropy and energy-flux formulas belong to Dorau--Much and their predecessors, including Araki, Uhlmann, Casini, Kay--Wald, and Longo. Theorem XLIX.1 is an elementary but useful independently formulated operational specialization of a known quadratic minimization principle. No priority for a universal theorem or solution of quantum gravity is claimed.

\chapter{A Constructive Bridge between Quantum Distinguishability, the Hopf Triad, and Gravitational Dynamics}
\section{Exact scope and limitations}
Earlier chapters established operational-memory quantum--classical separations, topological certificates for $S^2\to S^2$ maps, and a local-horizon flux information certificate. None alone turns statistical distinguishability into a stress--energy tensor. Here we solve a narrower, fully specified problem: construct an explicit local action connecting quantum Fisher information to classical matter, and find an exact self-consistent solution of the Einstein and field equations on compact $S^3$ slices. The action, metric coupling, and matter fields are additional physical model hypotheses, neither derived from one probability table nor experimentally established.
\statement{Definition}{L.1}{The unit three-sphere and Hopf triad}
Let $\mathbb H$ be the quaternions, $S^3=\{q\in\mathbb H:|q|=1\}$ have unit round metric $\gamma$, and $e_1,e_2,e_3$ be an oriented orthonormal imaginary-quaternion basis with $e_ie_j=-\delta_{ij}+\epsilon_{ijk}e_k$. Define
\begin{equation}n_a(q)=qe_aq^{-1}\in S^2,\qquad a=1,2,3.\tag{229}\end{equation}
The left orthonormal frame $E_i(q)=qe_i$ satisfies $[E_i,E_j]=2\epsilon_{ijk}E_k$.
\statement{Theorem}{L.2}{The exact tensor identity of the Hopf triad}
Each $n_a$ is a standard Hopf fibration with $|\operatorname{Hopf}(n_a)|=1$, its sign depending on orientation conventions. Moreover,
\begin{equation}\sum_{a=1}^3dn_a\otimes dn_a=8\gamma,\qquad\sum_{a=1}^3|dn_a|_\gamma^2=24,\tag{230}\end{equation}
and each map is harmonic:
\begin{equation}\Delta_\gamma n_a=-8n_a,\qquad|dn_a|_\gamma^2=8.\tag{231}\end{equation}
Here $\Delta_\gamma$ acts componentwise in $\R^3$.
\begin{proof}
$E_i(n_a)=q[e_i,e_a]q^{-1}=2q(e_i\times e_a)q^{-1}$, so
\[\sum_aE_i(n_a)\cdot E_j(n_a)=4\sum_a(\delta_{ij}-\delta_{ia}\delta_{ja})=8\delta_{ij}.\]
The total trace is 24 and each fixed-$a$ trace is 8. Bi-invariance gives $\nabla_{E_i}E_i=0$; differentiating again gives $\Delta_\gamma n_a=4q\sum_ie_i\times(e_i\times e_a)q^{-1}=-8n_a$. The map $q\mapsto qe_3q^{-1}$ is the standard Hopf projection. For dual left forms $\theta_i$, $d\theta_3=-2\theta_1\wedge\theta_2$ and $n_3^*\omega_{S^2}=4\theta_1\wedge\theta_2=d(-2\theta_3)$, with $\int_{S^2}\omega_{S^2}=4\pi$. Since $\operatorname{Vol}(S^3)=2\pi^2$,
\[\left|(4\pi)^{-2}\int_{S^3}(-2\theta_3)\wedge n_3^*\omega_{S^2}\right|=1.\]
Target rotation replaces $e_3$ by $e_a$ without changing the absolute invariant.
\end{proof}
\statement{Theorem}{L.3}{The single-fibration obstruction}
If smooth $n:(S^3,\gamma)\to S^2$ has $dn\otimes dn$ pointwise proportional to $\gamma$, then $n$ is constant. Thus one nontrivial Hopf map in an ordinary positive $S^2$ sigma model cannot produce nonzero strictly isotropic spatial stress on the round three-sphere.
\begin{proof}
$dn:T_qS^3\to T_{n(q)}S^2$ has rank at most two, while nonzero $c\gamma$ has rank three. Hence $c=0$, $dn=0$, and connectedness makes $n$ constant. In the sigma model the traceless part of $T_{ij}$ equals that of $\partial_in\cdot\partial_jn$, so isotropy requires precisely this condition.
\end{proof}
\section{Operational six-label memory on every spatial slice}
Let $|\psi_j\rangle=(|0\rangle+e^{2\pi ij/6}|1\rangle)/\sqrt2$, $j=0,\ldots,5$, and let $U(q)$ be the fundamental qubit action of a unit quaternion, consistent with $\operatorname{Ad}_q$. Set
\begin{equation}R_j(q)=U(q)|\psi_j\rangle\langle\psi_j|U(q)^\dagger,\qquad M_i(q)=I-U(q)|\psi_i\rangle\langle\psi_i|U(q)^\dagger.\tag{232}\end{equation}
Trace cyclicity gives, independently of $q$,
\begin{equation}\Tr[R_j(q)M_i(q)]=\sin^2\frac{\pi(j-i)}6=W^{(6)}_{ji}.\tag{233}\end{equation}
By the central theorem of Chapter XXXVIII, the ranks are 3 and 5 and normalized memory is 6, while one qubit realizes the full table. These numbers persist in the spatial family but do not determine $U(q)$'s spatial derivative.
\statement{Proposition}{L.4}{A single table does not identify the gravitational source}
The families $U(q)\equiv I$ and $U(q)=q$ have the same table (233) everywhere but respectively zero and nonzero gradient tensors. Thus $\rank W$, $\rank_+W$, and $\mu_{\rm normalized}(W)$ cannot recover a local gravitational source without a spatial state specification in a fixed physical frame.
\begin{proof}
Simultaneous unitary conjugation of preparation and effect preserves their trace pairing, but generally not their spatial derivatives. Those vanish for constant $U$ and are nonzero for $U=q$ by (230).
\end{proof}
\section{A quantum-information metric as matter: an explicitly chosen action}
For constant $0<r\le1$, define three qubit states
\begin{equation}\rho_a(x)=\tfrac12(I+rn_a(x)\cdot\sigma).\tag{234}\end{equation}
Here $r$ is constant in space and time, not the cosmological sphere radius. For a smooth family, define the symmetric SLD quantum Fisher metric $I^{(a)}_{\mu\nu}=\tfrac12\Tr\rho_a\{L_{a,\mu},L_{a,\nu}\}$, with $\partial_\mu\rho_a=\tfrac12\{\rho_a,L_{a,\mu}\}$.
\statement{Lemma}{L.5}{An exact qubit SLD--Fisher identity}
For constant $r$,
\begin{equation}I^{(a)}_{\mu\nu}=r^2\partial_\mu n_a\cdot\partial_\nu n_a.\tag{235}\end{equation}
It also holds at the pure boundary $r=1$, with the natural support-based SLD definition.
\begin{proof}
Since $n_a\cdot\partial_\mu n_a=0$, choose $L_{a,\mu}=r\partial_\mu n_a\cdot\sigma$. The Pauli relation $\{u\cdot\sigma,v\cdot\sigma\}=2(u\cdot v)I$ gives both the SLD equation and (235).
\end{proof}
\statement{Definition}{L.6}{A minimally coupled Einstein--QFI--$S^2$ model}
On a four-dimensional Lorentzian manifold of signature $(-,+,+,+)$, with fixed $r$ and independent maps $n_a:M\to S^2$, postulate
\begin{equation}S[g,n_a]=\int_Md^4x\sqrt{-g}\left[\frac{R-2\Lambda}{16\pi G}-\frac\beta2g^{\mu\nu}\sum_{a=1}^3I^{(a)}_{\mu\nu}\right],\qquad\beta>0.\tag{236}\end{equation}
Write $\kappa=\beta r^2>0$. Memory theorems do not determine $\beta$. The three maps are three matter sigma fields/sectors, not three simultaneous states of one physical qubit.
\statement{Theorem}{L.7}{Covariant conservation and a local null-direction bridge}
Varying (236) gives
\begin{align}G_{\mu\nu}+\Lambda g_{\mu\nu}&=8\pi GT_{\mu\nu},\tag{237}\\
T_{\mu\nu}&=\beta(I_{\mu\nu}-\tfrac12g_{\mu\nu}I^\alpha{}_{\alpha}),\qquad I_{\mu\nu}=\sum_aI^{(a)}_{\mu\nu},\tag{238}\\
\Box_gn_a+(\nabla_\mu n_a\cdot\nabla^\mu n_a)n_a&=0.\tag{239}\end{align}
On shell, $\nabla^\mu T_{\mu\nu}=0$. For null $k$,
\begin{equation}R_{\mu\nu}k^\mu k^\nu=8\pi G\beta\sum_{a=1}^3I^{(a)}_{\mu\nu}k^\mu k^\nu\ge0.\tag{240}\end{equation}
For a hypersurface-orthogonal null-geodesic congruence with affine parameter $\lambda$,
\begin{equation}\frac{d\vartheta}{d\lambda}=-\tfrac12\vartheta^2-\sigma_{\mu\nu}\sigma^{\mu\nu}-8\pi G\beta I_{\mu\nu}k^\mu k^\nu.\tag{241}\end{equation}
\begin{proof}
By (235), the matter action is three ordinary $S^2$ sigma models with coupling $\kappa$. Metric variation gives (238); constrained variation $n_a\cdot n_a=1$ gives (239). Stress divergence vanishes on these equations. Null contraction of (237) removes the cosmological and trace terms, and $I(k,k)=r^2\sum_a|k^\mu\partial_\mu n_a|^2\ge0$. The last identity is standard zero-twist Raychaudhuri. These are model identities, not a derivation of gravity from an information axiom.
\end{proof}
\section{Exact spatially closed solutions of the Einstein system}
\statement{Theorem}{L.8}{A self-consistent $S^3$-FLRW cosmology}
For
\begin{equation}ds^2=-dt^2+a(t)^2\gamma_{ij}\,dq^idq^j,\qquad a(t)>0,\tag{242}\end{equation}
and the time-independent fields (229), the field equations hold exactly. Stress is homogeneous and isotropic, with
\begin{equation}\rho=\frac{12\kappa}{a^2},\quad p=-\frac{4\kappa}{a^2},\quad w=-\frac13,\quad\dot\rho+3H(\rho+p)=0.\tag{243}\end{equation}
The full Einstein system is equivalent to
\begin{equation}\dot a^2=32\pi G\kappa-1+\frac{\Lambda a^2}3,\qquad\ddot a=\frac\Lambda3a.\tag{244}\end{equation}
Every positive solution of these equations therefore solves both dynamical systems exactly.
\begin{proof}
For static fields, $\Box_gn_a=a^{-2}\Delta_\gamma n_a$ and $\nabla n_a\cdot\nabla n_a=8/a^2$, so (231) gives (239). The information tensor has $I_{00}=I_{0i}=0$, $I_{ij}=8r^2\gamma_{ij}$, and trace $24r^2/a^2$. Thus $T_{00}=12\kappa/a^2$ and $T_{ij}=-4\kappa\gamma_{ij}$. Substitute into $3(H^2+a^{-2})=8\pi G\rho+\Lambda$ and $\ddot a/a=-4\pi G(\rho+3p)/3+\Lambda/3$. Continuity follows by differentiating $\rho\propto a^{-2}$.
\end{proof}
\statement{Corollary}{L.9}{Complete analytical regime table}
Write $C=32\pi G\kappa-1$.
\begin{enumerate}[label=(\alph*)]
\item If $\Lambda>0$, let $h=\sqrt{\Lambda/3}$. For $C<0$, there is a nonsingular scale-factor bounce
\begin{equation}a(t)=\frac{\sqrt{-C}}h\cosh[h(t-t_0)],\qquad a_{\min}=\frac{\sqrt{1-32\pi G\kappa}}h.\tag{245}\end{equation}
For $C=0$, $a=a_0e^{\pm h(t-t_0)}$. For $C>0$, $a=(\sqrt C/h)\sinh[h(t-t_0)]$ on its positive interval.
\item If $\Lambda=0$ and $C>0$, $a=\sqrt C(t-t_0)$ on $a>0$. For $C=0$, every positive constant is a static solution; for $C<0$, none exists.
\item If $\Lambda<0$, let $\omega=\sqrt{-\Lambda/3}$. Positive solutions exist only for $C>0$: $a=(\sqrt C/\omega)\sin[\omega(t-t_0)]$ on a positive open interval.
\end{enumerate}
Every solution satisfies
\begin{equation}R=4\Lambda+\frac{192\pi G\kappa}{a^2},\qquad E_{S^3}=\int_{S^3}\rho\,dV=24\pi^2\kappa a,\tag{246}\end{equation}
and $dE_{S^3}=-p\,d(2\pi^2a^3)$.
\begin{proof}
Integrate (244) for each sign combination. The curvature follows from the Einstein trace and $T^\mu{}_{\mu}=-\rho+3p=-24\kappa/a^2$. Energy uses $\operatorname{Vol}(S_a^3)=2\pi^2a^3$.
\end{proof}
\statement{Remark}{L.10}{A bounce without NEC violation}
For $0<32\pi G\kappa<1$, $\Lambda>0$, the bounce has $\rho+p=8\kappa/a^2>0$. Acceleration is due to positive $\Lambda$, not negative energy. This is a model solution, not an observational conclusion about our Universe.
\statement{Proposition}{L.11}{Local null focusing in the FLRW solution}
In an orthonormal FLRW frame, a null vector with time component $k^{\widehat0}$ satisfies
\begin{equation}I_{\mu\nu}k^\mu k^\nu=\frac{8r^2}{a^2}(k^{\widehat0})^2,\qquad R_{\mu\nu}k^\mu k^\nu=\frac{64\pi G\kappa}{a^2}(k^{\widehat0})^2.\tag{247}\end{equation}
\begin{proof}
$I_{\widehat i\widehat j}=8r^2\delta_{ij}/a^2$, and nullness gives $\sum_i(k^{\widehat i})^2=(k^{\widehat0})^2$. Apply (240).
\end{proof}
\section{Triad rigidity: isotropic sources do not imply arbitrary topological charge}
\statement{Theorem}{L.12}{Constant-density rigidity of compatible triads}
For smooth $F:S^3\to S^3$, let $n_a^F(q)=F(q)e_aF(q)^{-1}$. Then
\begin{equation}\sum_a dn_a^F\otimes dn_a^F=8F^*\gamma.\tag{248}\end{equation}
If this equals constant $c\gamma$, $c\ge0$, either $c=0$ and $F$ is constant, or $c=8$ and $F$ is a round-sphere isometry. Thus in this compatible homogeneous class, component Hopf-charge moduli are one; arbitrary large charges cannot be obtained by substituting $F$ while preserving homogeneity and isotropy.
\begin{proof}
Pull back (230). At $c=0$, $dF=0$. For $c>0$, $F^*\gamma=(c/8)\gamma$ makes $F$ a local isometry from the rescaled sphere to the unit sphere. Comparing sectional curvatures $8/c$ and 1 gives $c=8$. Compactness and simple connectedness make the local isometry global. Its degree is $\pm1$, so component Hopf invariants have modulus one.
\end{proof}
\section{A three-dimensional finite-sample topological laboratory}
Chapters XLI--XLIV mainly concern the parameter sphere $S^2$ and Chern class. Here the physical spatial slice is $S^3$. A Chern class on $S^3$ must not be identified with a Hopf invariant: $H^2(S^3;\Z)=0$, whereas $\pi_3(S^2)\cong\Z$.
\statement{Theorem}{L.13}{A finite-grid Hopf-invariant certificate}
Let smooth $v:S^3\to\R^3$ satisfy $\|v(q)\|\ge r_{\min}>0$, with independently known Lipschitz constant $L_v$. Let a tetrahedral triangulation have geodesic diameter at most $h$, and suppose all node estimates obey $\|\widehat v_i-v(q_i)\|\le\eta$. Let $u$ be componentwise barycentric interpolation. If
\begin{equation}L_vh+\eta<r_{\min},\tag{249}\end{equation}
then $u$ never vanishes and $\widehat n=u/\|u\|$ is homotopic to $n=v/\|v\|$. In particular,
\begin{equation}\operatorname{Hopf}(\widehat n)=\operatorname{Hopf}(n).\tag{250}\end{equation}
\begin{proof}
Within a simplex,
\[\|u(q)-v(q)\|\le\sum_i\lambda_i[\|\widehat v_i-v(q_i)\|+\|v(q_i)-v(q)\|]\le\eta+L_vh<\|v(q)\|.\]
The straight interpolation never vanishes; normalization gives a homotopy $S^3\to S^2$ preserving Hopf invariant. The integer itself should be computed by a calibrated discrete preimage-linking or certified cochain algorithm, not by rounding an uncontrolled numerical integral.
\end{proof}
\statement{Corollary}{L.14}{A sufficient sampling cost for three Pauli axes}
With $N$ independent measurements per axis at each of $V$ nodes, additional calibration error $\eta_{\rm sys}$, positioning error $\xi$, and failure budget $\alpha$, a sufficient condition at confidence $1-\alpha$ is
\begin{equation}\sqrt{\frac6N\log\frac{6V}\alpha}+\eta_{\rm sys}+L_v(h+\xi)<r_{\min}.\tag{251}\end{equation}
For the exact triad, $L_v\le2r$ and $r_{\min}=r$. With no additional errors, $2h+\sqrt{6\log(6V/\alpha)/N}/r<1$ suffices.
\begin{proof}
Hoeffding bounds each mean error by failure probability $2e^{-Nt^2/2}$. Union over $3V$ events and conversion to Euclidean error gives the square-root term. Other errors add by the triangle inequality; $\|dn_a\|_{\rm op}\le2$ follows from Theorem L.2.
\end{proof}
\section{Local basis versus physical field: a mandatory negative control}
\statement{Remark}{L.15}{A spurious gravitational source from a frame change}
Equation (235) describes statistically distinguishable states relative to a fixed physical identification of qubit frames across points. If $q(x)$ is only a local $SU(2)$ fiber-coordinate change, ordinary derivatives are not gauge invariant. With $D_\mu n=\partial_\mu n+[A_\mu,n]$ and the pure-gauge connection $A_\mu=-(\partial_\mu q)q^{-1}$, one has $D_\mu(qe_aq^{-1})=0$. Stress built from $\|Dn_a\|^2$ then vanishes. The nonzero model therefore assumes physical matter textures in an operationally fixed frame, not merely an arbitrary gauge. This condition cannot be removed when asserting universality.
\section{Compatibility with the horizon certificate of Chapter XLV}
Equation (240) directly links null matter flux and SLD distinguishability within the stated model. Chapter XLV treats a different algebraic-QFT object: relative entropy of coherent free-field excitations on a local horizon. They cannot be automatically identified:
\[\text{QFI sigma-model flux }\beta I_{kk}\ \not\equiv\ \text{horizon relative entropy }S_{\rm rel}.\]
Extending the horizon formula to these interacting sigma fields would require new modular theory, a horizon-state algebra, and additional entropy--area conditions. This agrees with the limits of Dorau--Much, PRL 136, 091602 (2026), \url{https://doi.org/10.1103/lmq8-nsty}. Gravitating sigma models have also been studied, for example by F. Canfora, N. Dimakis, A. Paliathanasis, \url{https://arxiv.org/abs/1707.02270}, and European Physical Journal C 77, 909 (2017), \url{https://doi.org/10.1140/epjc/s10052-017-5477-2}. This construction claims no priority for the basic sigma-model equations or solution type.
\section{Limitations, controls, and scientific status}
\begin{enumerate}
\item Proved within the action: (230)--(246), the null energy condition, the isotropy obstruction, compatible-triad rigidity, and a conditional Hopf certificate.
\item Assumed physically: the Einstein action term with specified $G$, coupling $\beta$, three independent matter fields, and a fixed operational frame. No cosmological-data test was performed.
\item Not proved: universal Einstein dynamics from memory alone, determination of $G$ or $\Lambda$ by topology, stability against all anisotropic perturbations, gravity quantization, a physical neutrino-mass connection, the primordial origin of a cosmological three-sphere, or worldwide priority.
\item Essential negative controls: one Hopf map is anisotropic; $W$ cannot distinguish zero from nonzero spatial source; pure-gauge texture has no covariant source; a Chern class on $S^2$ is not a Hopf invariant on $S^3$.
\end{enumerate}
\note{A reproducible example without parameter fitting.} In units $G=1$, take $\kappa=1/(64\pi)$ and $\Lambda=3$. Then $C=-1/2$, $h=1$, and
\begin{equation}a(t)=\frac1{\sqrt2}\cosh t,\quad\rho(t)=\frac3{8\pi\cosh^2t},\quad p(t)=-\frac1{8\pi\cosh^2t},\quad R(t)=12+\frac6{\cosh^2t}.\tag{252}\end{equation}
Direct substitution gives $3(H^2+a^{-2})=8\pi\rho+\Lambda$, $\ddot a/a=\Lambda/3$, and continuity exactly. Since the existing article does not determine $G,\kappa,\Lambda$ observationally, this illustrates a solution, not a numerical cosmological prediction.
\setcounter{equation}{252}


% 13_effective_gravity.tex
\chapter{The Origin of an Effective Gravitational Coupling, a Negative Hopf Mode, and Testable Cosmography}
\section{A rigorous obstruction: statistical geometry does not fix Newton's constant}
\statement{Theorem}{LI.1}{The preceding axioms cannot uniquely determine the coupling}
The previous chapter's assumptions---a smooth metric, physically distinguishable qubit families, a positive SLD metric $I_{\mu\nu}$, locality, general covariance, and on-shell stress conservation---together with the triad's topological invariants determine neither $G$, the stiffness $\beta$, nor their product.
\begin{proof}
For every $M^2>0$, $\beta>0$,
\begin{equation}S_{M,\beta}=\int d^4x\sqrt{-g}\left[\frac{M^2}2R-\Lambda M^2-\frac\beta2g^{\mu\nu}I_{\mu\nu}\right]\tag{253}\end{equation}
satisfies those kinematic assumptions. Yet null focusing is $R_{\mu\nu}k^\mu k^\nu=(\beta/M^2)I_{\mu\nu}k^\mu k^\nu$; varying either parameter changes the observed amplitude while QFI and topology remain fixed. Unique determination therefore does not follow. $M^2=(8\pi G)^{-1}$ is an independent action parameter.
\end{proof}
\statement{Remark}{LI.2}{What a microscopic theory actually adds}
Computing the coefficient requires additional physical carriers, their curvature couplings, a regularization scale, and a rule for the admissible bare $R$ term. This is dynamical closure, not a new algebraic operation on the six-label table.
\section{An independent microscopic branch: proper-time induction}
\statement{Definition}{LI.3}{A conditionally specified UV model}
Let $(M_E,g_E)$ be a smooth compact boundaryless Riemannian four-manifold, with independent real scalars $\phi_A$, $A=1,\ldots,N$, and
\begin{equation}P_A=-\Delta_{g_E}+m_A^2+\xi_AR_E,\qquad W_{\rm one\ loop}=\frac12\sum_A\Tr\log P_A.\tag{254}\end{equation}
Introduce a proper-time cutoff $s\ge\ell_*^2$, assume $\ell_*^2|R_E|\ll1$, and define the relative effective functional by subtracting a fixed reference metric. Neither $\ell_*$ nor $(m_A,\xi_A)$ is determined by Hopf topology.
\statement{Theorem}{LI.4}{Explicit Einstein--Hilbert term in the chosen regularization}
To first local curvature order,
\begin{align}\Tr e^{-sP_A}&=\frac1{(4\pi s)^2}\int d^4x\sqrt{g_E}\,e^{-m_A^2s}\left[1+s(\tfrac16-\xi_A)R_E+O(s^2R^2)\right],\tag{255}\\
W_R&=-\frac1{32\pi^2}\sum_A(\tfrac16-\xi_A)\int_{\ell_*^2}^\infty\frac{e^{-m_A^2s}}{s^2}\,ds\int d^4x\sqrt{g_E}R_E+O(R^2).\tag{256}\end{align}
Comparison with $-\int\sqrt{g_E}R_E/(16\pi G)$ gives
\begin{equation}\frac1{G_{\rm eff}}=\frac1{G_{\rm bare}}+\frac1{2\pi}\sum_A(\tfrac16-\xi_A)\int_{\ell_*^2}^\infty\frac{e^{-m_A^2s}}{s^2}\,ds.\tag{257}\end{equation}
For $m_A=0$, $\xi_A=\xi$,
\begin{equation}\frac1{G_{\rm eff}}=\frac1{G_{\rm bare}}+\frac{N(1-6\xi)}{12\pi\ell_*^2}.\tag{258}\end{equation}
\begin{proof}
Use $\frac12\Tr\log P=-\frac12\int_{\ell_*^2}^\infty(ds/s)\Tr e^{-sP}$, with a metric-independent subtraction. The standard first local heat coefficient is $a_1=(1/6-\xi)R_E$, producing $s^{-2}$ in the integral and coefficient $-\frac12(4\pi)^{-2}\sum_A(1/6-\xi_A)\int s^{-2}e^{-m_A^2s}ds$. For zero mass the integral is $\ell_*^{-2}$.
\end{proof}
\statement{Corollary}{LI.5}{A predictive but conditional branch}
If a separate UV hypothesis imposes $G_{\rm bare}^{-1}=0$, $m_A=0$, and $\xi_A=0$, then
\begin{equation}G_{\rm ind}=\frac{12\pi\ell_*^2}{N},\qquad8\pi G_{\rm ind}\beta=\frac{96\pi^2\beta\ell_*^2}{N}.\tag{259}\end{equation}
This predicts the coefficient of $I(k,k)$ given $\ell_*,N,\beta$. QFI alone still does not determine $\beta$.
\statement{Remark}{LI.6}{Mandatory limits of the calculation}
\begin{enumerate}[label=(\roman*)]
\item The quadratically divergent $R$ contribution is scheme dependent: changing the UV cutoff or an allowed finite counterterm changes $G_{\rm eff}$.
\item At conformal coupling $\xi=1/6$, these massless scalars give no quadratic $R$ contribution in this scheme; for $\xi>1/6$ its sign reverses.
\item A volume term of order $\ell_*^{-4}$ and local $R^2$, $R_{\mu\nu}^2$, and other corrections arise simultaneously. Without solving the cosmological-constant problem, closure is incomplete.
\item This is a one-loop low-energy expansion of a specified microscopic sector, not universal emergence of GR or a derivation of sigma-sector masses and couplings.
\end{enumerate}
Predecessors include Sakharov's program as reviewed by M. Visser, \url{https://arxiv.org/abs/gr-qc/0204062}, and T. Jacobson's thermodynamic equation of state, \url{https://doi.org/10.1103/PhysRevLett.75.1260}. In the latter, $G$ is determined given a horizon entropy density, which is not derived from our six probabilities.
\section{A separate UV definition of matter stiffness: a regulated state lattice}
\statement{Proposition}{LI.7}{The QFI coefficient from specified microscopic exchange}
On a four-dimensional Euclidean hypercubic lattice of spacing $\ell_m$, take three physically observable unit-vector fields $n_a$ and positive dimensionless exchange $J_{\rm lat}$. Postulate
\begin{equation}S_{\rm lat}=\frac{J_{\rm lat}}2\sum_{x,\mu,a}\|n_a(x+\ell_m\widehat\mu)-n_a(x)\|^2.\tag{260}\end{equation}
If fields are smooth on that scale and a rotationally invariant low-energy limit exists, its leading term is
\begin{equation}S_{m,E}=\frac\kappa2\int d^4x\sqrt{g_E}\sum_{a,\mu}|\partial_\mu n_a|^2+O(\ell_m^2\partial^4),\qquad\kappa=\frac{J_{\rm lat}}{\ell_m^2}.\tag{261}\end{equation}
In a spatially fixed physical frame and at constant mixedness $r$,
\begin{equation}\beta=\frac{J_{\rm lat}}{r^2\ell_m^2}.\tag{262}\end{equation}
With the additional UV assumptions $G_{\rm bare}^{-1}=0$ and $m_A=\xi_A=0$, the null-focusing coefficient becomes
\begin{equation}8\pi G_{\rm ind}\beta=\frac{96\pi^2J_{\rm lat}}{Nr^2}\left(\frac{\ell_*}{\ell_m}\right)^2.\tag{263}\end{equation}
\begin{proof}
Each squared finite difference is $\ell_m^2|\partial_\mu n|^2+O(\ell_m^3)$, and the cell sum becomes $\ell_m^{-4}\int d^4x$. Lattice symmetry removes odd integrated terms, leaving the displayed higher-derivative correction. Equation (235) gives $I=r^2\sum_a dn_a\otimes dn_a$, hence $\kappa=\beta r^2$. Combine with (259).
\end{proof}
\statement{Remark}{LI.8}{A microscopic parameter is not automatically a constant of nature}
$J_{\rm lat},\ell_m,\ell_*,N,r$ are physical inputs to a specific UV scenario, not topological deductions. Extending scalar proper time to a dynamical lattice and obtaining reflection positivity, general covariance, and Lorentzian causality remain open. Equation (263) conditionally matches two models; it is not a universal derivation of Newton's constant.
\section{An unexpected negative result: an explicit energy-lowering conformal mode}
\statement{Definition}{LI.9}{Nontrivial deformations at fixed topology}
Embed unit $S^3$ in $\R^4$ and set $f(q)=q_4$. Let $\phi_s$ be the global flow of the conformal vector field $V=\nabla_\gamma f$. It is a smooth diffeomorphism isotopic to the identity, with
\begin{equation}\phi_s^*\gamma=\omega_s^2\gamma,\qquad\omega_s(q)=\frac1{\cosh s+f(q)\sinh s}.\tag{264}\end{equation}
Consider $n_a^{(s)}=n_a\circ\phi_s$.
\statement{Theorem}{LI.10}{An exact negative energy mode of the Hopf triad}
On fixed round $S^3$, the positive sigma-model energy $E[n]=(\kappa/2)\sum_a\int|dn_a|_\gamma^2dV_\gamma$ obeys
\begin{equation}E[n^{(s)}]=12\kappa\operatorname{Vol}(S^3)\operatorname{sech}^2(s/2).\tag{265}\end{equation}
Consequently,
\begin{equation}E'(0)=0,\qquad E''(0)=-6\kappa\operatorname{Vol}(S^3)<0.\tag{266}\end{equation}
Every finite-$s$ component retains Hopf-invariant modulus one. The original triad is therefore an energy saddle under admissible smooth joint deformations at fixed geometry; topological charge does not ensure energetic stability.
\begin{proof}
Pulling back (230) gives $\sum_a d(n_a\circ\phi_s)\otimes d(n_a\circ\phi_s)=8\omega_s^2\gamma$, hence $E(s)=12\kappa\int\omega_s^2dV_\gamma$. The coordinate $f$ has normalized density $(2/\pi)\sqrt{1-f^2}$ on $[-1,1]$, and direct integration yields
\[\frac1{\operatorname{Vol}(S^3)}\int\omega_s^2dV_\gamma=\frac2\pi\int_{-1}^1\frac{\sqrt{1-f^2}}{(\cosh s+f\sinh s)^2}\,df=\operatorname{sech}^2(s/2).\]
Differentiate at zero. Isotopy of $\phi_s$ preserves the maps' homotopy classes.
\end{proof}
\statement{Corollary}{LI.11}{An explicit dangerous-mode equation on a prescribed FLRW background}
For fixed $g=-dt^2+a(t)^2\gamma$ and ansatz $n_a(t,q)=n_a(\phi_{s(t)}(q))$, the quadratic matter Lagrangian is
\begin{equation}L_m^{(2)}=12\kappa\mathcal I[a^3\dot s^2+as^2],\qquad\mathcal I=\int_{S^3}f^2dV_\gamma=\tfrac14\operatorname{Vol}(S^3).\tag{267}\end{equation}
Its equation is
\begin{equation}\ddot s+3H\dot s-\frac{s}{a^2}=0.\tag{268}\end{equation}
This is an actual negative matter second-variation mode on a fixed background and excludes an unsupported assertion of automatic Hopf--FLRW stability. It is not the final spectrum of the coupled Einstein--matter system; linear gravitational constraints and mixed perturbations must also be solved.
\begin{proof}
At $s=0$, $\partial_tn_a=dn_a(V)\dot s$ and $\sum_a|dn_a(V)|^2=8|V|^2$. Since $\int|\nabla f|^2=3\int f^2=3\mathcal I$, the kinetic term is $12\kappa a^3\mathcal I\dot s^2$. Expansion of (265) gives $E(s)=12\kappa\operatorname{Vol}(S^3)-3\kappa\operatorname{Vol}(S^3)s^2+O(s^4)$, yielding $+12\kappa a\mathcal I s^2$ in the Lagrangian. Euler--Lagrange gives (268).
\end{proof}
\statement{Remark}{LI.12}{The mode on a bounce}
For $a=a_{\min}\cosh(ht)$, the equation is $\ddot s+3h\tanh(ht)\dot s-a_{\min}^{-2}\operatorname{sech}^2(ht)s=0$. At the bounce, $s(0)>0$, $\dot s(0)=0$ imply $\ddot s(0)=a_{\min}^{-2}s(0)>0$. On the late de Sitter branch, the negative coefficient tends to zero, so growth need not remain exponential. Full nonlinear or gravitationally coupled stability is not established.
\section{Observable identities and falsification without fitting}
\statement{Theorem}{LI.13}{Background curvature degeneracy and an independent geometric test}
For the model of Chapter XLVI with conservatively added noninteracting dust, radiation, and $\Lambda$, let $a_0$ be today's curvature radius and define
\begin{align}\Omega_K&=-\frac1{a_0^2H_0^2}<0,\tag{269}\\\Omega_H&=\frac{32\pi G\kappa}{a_0^2H_0^2}>0,\tag{270}\\\Omega_2&=\Omega_K+\Omega_H.\tag{271}\end{align}
Then
\begin{equation}\frac{H^2(z)}{H_0^2}=\Omega_r(1+z)^4+\Omega_m(1+z)^3+\Omega_2(1+z)^2+\Omega_\Lambda.\tag{272}\end{equation}
Because both terms scale as $a^{-2}$, $H(z)$ alone cannot separate $\Omega_H$ from $\Omega_K$. In closed geometry, however,
\begin{equation}D_M(z)=a_0\sin\left[\frac1{a_0}\int_0^z\frac{dz'}{H(z')}\right],\qquad D_L(z)=(1+z)D_M(z),\tag{273}\end{equation}
before the antipodal caustic, and for $z>0$, $D_M\ne0$,
\begin{equation}\frac{[H(z)D'_M(z)]^2-1}{H_0^2D_M(z)^2}=\Omega_K<0.\tag{274}\end{equation}
Consequently,
\begin{equation}\Omega_H=\Omega_2-\Omega_K,\qquad\frac{\Omega_H}{|\Omega_K|}=32\pi G\kappa.\tag{275}\end{equation}
These are conditional background consistency tests, not a new universal curvature law.
\begin{proof}
Insert $\rho_H=12\kappa/a^2$ into the first Friedmann equation with $k=+1$. The transverse length in a closed slice is the sine of the radial angle, giving (273). Differentiation gives $HD'_M=\cos(\chi/a_0)$, hence $(HD'_M)^2-1=-D_M^2/a_0^2$. Divide to obtain (274); definitions give (275).
\end{proof}
\statement{Proposition}{LI.14}{Bounce restrictions with ordinary matter}
With dust and radiation,
\begin{equation}\frac{\ddot a}a=\frac\Lambda3-\frac{4\pi G}3(\rho_m+2\rho_r).\tag{276}\end{equation}
At $H=0$, positive acceleration requires
\begin{equation}\Lambda>4\pi G(\rho_m+2\rho_r).\tag{277}\end{equation}
The $w=-1/3$ Hopf sector does not contribute to acceleration. Adding ordinary early matter does not automatically preserve the vacuum bounce of Chapter XLVI.
\begin{proof}
Use the acceleration equation with $w_m=0$, $w_r=1/3$, $w_H=-1/3$, and require $\ddot a>0$ at the scale minimum.
\end{proof}
\statement{Remark}{LI.15}{An observational protocol cannot be replaced by synthetic data}
Testing requires joint treatment of independent radial $H(z)$ or $D_H/r_d$, transverse $D_M/r_d$, supernova distances, and CMB acoustic scale, lensing, and perturbation spectra. Geometry and small $|\Omega_K|$ alone do not establish compatibility: the Hopf sector itself has the preceding perturbations. DESI DR2 BAO chains are public at \url{https://data.desi.lbl.gov/doc/papers/dr2/}. Planck 2018 plus BAO gave $\Omega_K=0.0007\pm0.0019$ in its adopted parameter model, not a universal limit on all geometrically possible theories: \url{https://doi.org/10.1051/0004-6361/201833910}. Parameter selection and model admissibility need the original likelihood. No such fit was performed in this revision.
\section{Proved results, open tasks, and explicit completion conditions}
\begin{enumerate}
\item Proved within the model: inability to fix constants from prior axioms; the $R$ coefficient for a specified scalar UV sector; exact conformal deformation energy and its negative second variation; a negative linear matter mode on fixed FLRW; background kinematic identities and a necessary bounce condition.
\item Unproved: UV justification of $G_{\rm bare}=0$ and $\ell_*$; quantum origin of $\beta$; cancellation of large induced $\Lambda$; the complete coupled Einstein--Hopf linear spectrum; DESI/Planck/SN compatibility; and a physically observable triad.
\item Next rigorous targets: full Bianchi IX Hamiltonian reduction and scalar/tensor spectra with gravitational constraints; microscopic matrix-model calculation of $\beta$ with an operationally fixed frame; likelihood analysis with explicit priors and controlled $R^2$ corrections; and an independent gauge invariant distinguishing matter texture from a local $SU(2)$ basis choice.
\end{enumerate}
Negative harmonic-map modes, induced gravity, and cosmographic curvature tests have substantial prior literature. Novelty of their META--CAUSAL combination needs independent systematic review and is not declared a fundamental discovery.

\chapter{A Renormalization-Consistent Spectral Bridge and the Barrier to Predicting $G$}
\section{Physical assumptions and exact normalization}
All dimensional expressions use $\hbar=c=1$. Real scalars live on a smooth compact boundaryless Riemannian manifold. The operator
\begin{equation}P_A=-\nabla^2+m_A^2+\xi_AR,\qquad\Gamma_{\rm matter}=\tfrac12\sum_A\Tr'\log P_A\tag{278}\end{equation}
need not be strictly positive: at $m_A=\xi_A=0$, the constant function is an exact zero mode on a connected compact manifold. The prime excludes the kernel; alternatively use a positive IR regulator and separately define its removal. A zero mode cannot remain in an unregulated $\int^\infty ds/s$.

Local UV coefficients use small proper time. One cannot integrate small-time asymptotics over the entire positive half-line without qualification. Let a dimensionless cutoff profile have finite positive moments
\begin{equation}C_1[F]=\int_0^\infty\frac{F(u)}{u^2}\,du,\qquad C_0[F]=\int_0^\infty\frac{F(u)}{u^3}\,du.\tag{279}\end{equation}
Take $F(0)=0$, $F(\infty)=1$, and use $F(s/\ell_*^2)$ with a proper IR procedure. Both moments depend on the profile.
\statement{Theorem}{LII.1}{Scheme dependence of induced local coefficients}
At leading local UV order, $N$ massless real scalars with common $\xi$ give
\begin{equation}\Gamma_{\rm UV,loc}\supset-\frac{NC_0[F]}{32\pi^2\ell_*^4}\int\sqrt g\,d^4x-\frac{N(1-6\xi)C_1[F]}{192\pi^2\ell_*^2}\int\sqrt gR\,d^4x,\tag{280}\end{equation}
and under $S^E_{\rm EH}=-M^2\int\sqrt gR/2$,
\begin{equation}M_{\rm eff}^2=M_{\rm bare}^2+\frac{N(1-6\xi)C_1[F]}{96\pi^2\ell_*^2}.\tag{281}\end{equation}
If the bare term is zero and $N(1-6\xi)>0$, the coefficient is positive but cutoff-shape dependent.
\begin{proof}
The coefficients $a_0=1$, $a_1=(1/6-\xi)R$ produce $s^{-3}$ and $s^{-2}$. Substituting $s=\ell_*^2u$ gives $\ell_*^{-4}C_0$ and $\ell_*^{-2}C_1$. Multiply by $-N/[2(4\pi)^2]$ and compare the $R$ coefficient with $-M^2/2$.
\end{proof}
\statement{Proposition}{LII.2}{Three admissible profiles with different normalization}
For $F_{\rm sharp}(u)=\one_{u\ge1}$, $(C_1,C_0)=(1,1/2)$; for $F_{\rm exp}(u)=e^{-1/u}$, they are $(1,1)$. More generally, for $F_a(u)=e^{-a/u}$, $a>0$,
\begin{equation}C_1[F_a]=a^{-1},\qquad C_0[F_a]=a^{-2}.\tag{282}\end{equation}
Thus equal $C_1$ need not mean equal induced vacuum density.
\begin{proof}
The sharp profile gives elementary power integrals. Substitution $t=a/u$ gives $C_1=\Gamma(1)/a$ and $C_0=\Gamma(2)/a^2$.
\end{proof}
\section{A theorem of renormalization underdetermination}
\statement{Theorem}{LII.3}{Physically equivalent UV-action parameterizations}
For fixed observable low-energy $M_{\rm ren}^2$, changing $F\to\widetilde F$ at fixed $\ell_*$ requires and suffices to change the bare coefficient by
\begin{equation}M_{\rm bare}^2(\widetilde F)=M_{\rm bare}^2(F)+\frac{N(1-6\xi)}{96\pi^2\ell_*^2}[C_1[F]-C_1[\widetilde F]].\tag{283}\end{equation}
Neither geometric invariants of fixed $S^3\to S^2$ maps nor the free-operator spectrum on a fixed background distinguish these parameterizations by themselves. Thus $M_{\rm bare}^2=0$ is not a scheme-invariant physical law.
\begin{proof}
In each scheme, $M_{\rm ren}^2$ is the sum of the bare and local one-loop coefficients. Subtract (281). Shifting a purely metric counterterm leaves matter fields and their fixed-background spectrum unchanged, as well as independently defined map topology. This freedom prevents unique determination without a physical matching condition.
\end{proof}
\statement{Remark}{LII.4}{Scope of the assertion}
This does not say Newton's constant is fundamentally unpredictable; it excludes derivation from the limited inputs listed. A UV-complete theory with an independent physical condition on the contact term could give a unique answer.
\statement{Corollary}{LII.5}{Quantitative spectral control of the vacuum term}
With $M_{\rm bare}=0$ and no additional volume-density contribution,
\begin{equation}\frac{|\rho_{\rm vac,ind}|}{M_{\rm ind}^4}=\frac{288\pi^2}{N(1-6\xi)^2}\frac{C_0[F]}{C_1[F]^2}.\tag{284}\end{equation}
The cutoff scale cancels. The profile ratio is $1/2$ for the sharp cutoff and 1 for the exponential. This is a regulator-dependent consistency test, not the observed cosmological constant; an independent cosmological counterterm is allowed.
\begin{proof}
Divide $NC_0/(32\pi^2\ell_*^4)$ by the square of $N(1-6\xi)C_1/(96\pi^2\ell_*^2)$. The scale cancels and $96^2/32=288$.
\end{proof}
\section{Mass thresholds and the correct infrared procedure}
\statement{Proposition}{LII.6}{The exact proper-time mass function}
For sharp cutoff and $m>0$, the coefficient multiplying $1/6-\xi$ is
\begin{equation}I(m,\ell_*)=\int_{\ell_*^2}^\infty\frac{e^{-m^2s}}{s^2}\,ds=\frac{e^{-m^2\ell_*^2}}{\ell_*^2}-m^2E_1(m^2\ell_*^2),\tag{285}\end{equation}
where $E_1(z)=\int_z^\infty e^{-t}dt/t$. As $m\ell_*\to0$, $I\to\ell_*^{-2}$; for $m\ell_*\gg1$, it is exponentially suppressed. The massive formula does not remove the need to subtract zero modes in the strictly massless compact case.
\begin{proof}
Integration by parts gives $\int_a^\infty e^{-m^2s}s^{-2}ds=e^{-m^2a}/a-m^2\int_a^\infty e^{-m^2s}ds/s$. Substitute $t=m^2s$.
\end{proof}
On the compact product $S_T^1\times S_a^3$, where $R=6/a^2$, the scalar eigenvalues and multiplicities are
\begin{equation}\lambda_{n\ell}=(2\pi n/T)^2+\ell(\ell+2)/a^2+m^2+6\xi/a^2,\quad d_{n\ell}=(\ell+1)^2,\quad n\in\Z,\ \ell\ge0.\tag{286}\end{equation}
For $m=\xi=0$, $(n,\ell)=(0,0)$ must be excluded from $\Tr'\log P$. Global geometry determines the IR spectrum and finite geometric corrections, but local leading UV coefficients do not uniquely fix $\ell_*$.
\section{Measuring the coupling: the stress-correlator slope}
\statement{Proposition}{LII.7}{A two-point coefficient as a matching condition}
Suppose the full Euclidean metric effective action near flat space has a local analytical expansion through two derivatives,
\begin{equation}\Gamma[g]=-\frac{M_{\rm ren}^2}2\int\sqrt gR+\rho_{\rm ren}\int\sqrt g+O(\partial^4),\tag{287}\end{equation}
with separately normalized volume density. Let $g_{\mu\nu}=e^{2\sigma}\delta_{\mu\nu}$ for compactly supported $\sigma$, and let $\mathcal K(p^2)$ be the Fourier-space second-variation kernel. Then
\begin{equation}\mathcal K(p^2)=\mathcal K(0)-6M_{\rm ren}^2p^2+O(p^4),\qquad M_{\rm ren}^2=-\frac16\left.\frac{d\mathcal K}{d(p^2)}\right|_{0}.\tag{288}\end{equation}
For microscopic quantum matter, the kernel contains both the connected two-point function of $T^\mu{}_{\mu}$, with sign fixed by $-\log Z$, and mandatory contact terms from the action's second variation. This full coefficient, not a formal integral over scattered spectral levels, must be matched to induced $G$.
\begin{proof}
In four dimensions, $\sqrt gR=e^{2\sigma}[-6\Box\sigma-6(\partial\sigma)^2]$. Integration by parts gives $\int\sqrt gR=6\int e^{2\sigma}(\partial\sigma)^2$, hence $\Gamma_{\rm EH}^{(2)}=-3M_{\rm ren}^2\int(\partial\sigma)^2$. With quadratic normalization $\frac12\int\sigma\mathcal K\sigma$, the kernel is $-6M_{\rm ren}^2p^2+\cdots$. The volume term has no derivatives. Differentiating $\Gamma=-\log Z$ gives a contact term minus the connected correlator of action variations.
\end{proof}
\statement{Remark}{LII.8}{Microscopic closure}
Equation (288) is in the spirit of the Adler--Zee program, expressing induced gravitational coupling through stress--energy correlations. Measuring the slope alone does not forbid a finite local $R$ counterterm. Determining $G$ requires a complete UV theory, its metric-matching rule, one contact-term scheme, and independently measured matter constants. See S. Adler, Phys. Lett. B 95 (1980), 241; J. Donoghue and G. Menezes, Phys. Rev. D 97 (2018), 056022, \url{https://doi.org/10.1103/PhysRevD.97.056022}.
\section{An independent dynamical scale instead of an arbitrary UV cutoff}
\statement{Proposition}{LII.9}{A conditional route through dimensional transmutation and Adler--Zee}
Assume a complete asymptotically free four-dimensional quantum theory with one dynamical mass scale, positive mass gap, no other explicit dimensional parameters, and $\mu\,dg/d\mu=-b_0g^3/(16\pi^2)+O(g^5)$, $b_0>0$. At one loop,
\begin{equation}\Lambda_{\rm dyn}=\mu\exp\left[-\frac{8\pi^2}{b_0g^2(\mu)}\right],\qquad\mu\frac{d\Lambda_{\rm dyn}}{d\mu}=0.\tag{289}\end{equation}
If the full conformal response, including contact matching, is finite, has only its intrinsic scale, and permits no independent local gravitational term, dimensional analysis gives
\begin{equation}M_{\rm ind}^2=C_{\rm AZ}\Lambda_{\rm dyn}^2,\quad G_{\rm ind}=\frac1{8\pi C_{\rm AZ}\Lambda_{\rm dyn}^2},\quad C_{\rm AZ}=-\frac{\mathcal K'(0)}{6\Lambda_{\rm dyn}^2}.\tag{290}\end{equation}
$C_{\rm AZ}$ is an as-yet uncomputed dimensionless coefficient of that full microscopic theory. If the first physical excitation has independently determined mass $m_{\rm gap}=c_{\rm gap}\Lambda_{\rm dyn}$, then $M_{\rm ind}^2/m_{\rm gap}^2=C_{\rm AZ}/c_{\rm gap}^2$ needs no separately postulated cutoff. Computing these two dimensionless numbers is an explicit research problem, not a claimed solution for our world.
\begin{proof}
The logarithmic derivative of $-8\pi^2/(b_0g^2)$ is $-1$ at one loop, proving (289). With no other scale and finite response, the slope has mass dimension two and is proportional to $\Lambda_{\rm dyn}^2$, normalized by (288). Eliminate the scale using the gap.
\end{proof}
\statement{Remark}{LII.10}{Assumptions cannot substitute for proof}
Positivity of $C_{\rm AZ}$, absence of an allowed local $R$ term, finiteness, and a contact-term definition are independent conditions; general covariance does not establish them. The one-scale assumption is not valid for the real Standard Model's massive sectors without further justification.
\section{Physical domain of a cosmological bounce}
\statement{Corollary}{LII.11}{A stronger UV-parameter and curvature condition}
For the Chapter XLVI model with $\kappa,\Lambda>0$, zero inverse bare $G$, massless scalars, and $N(1-6\xi)C_1>0$, define
\begin{equation}\eta=32\pi G_{\rm ind}\kappa=\frac{4\kappa}{M_{\rm ind}^2}=\frac{384\pi^2\kappa\ell_*^2}{N(1-6\xi)C_1[F]}.\tag{291}\end{equation}
The exact vacuum bounce has $a_{\min}=\sqrt{3(1-\eta)/\Lambda}$ for $0<\eta<1$. Validity of the local UV expansion near it additionally requires
\begin{equation}\frac{\ell_*^2}{a_{\min}^2}=\frac{\Lambda\ell_*^2}{3(1-\eta)}\ll1,\tag{292}\end{equation}
and smallness of all other dimensionless curvature and derivative invariants. The limit $\eta\to1^-$ is therefore not automatically a reliable effective-model regime.
\begin{proof}
From $\dot a^2=\eta-1+\Lambda a^2/3$, obtain $a_{\min}^2=3(1-\eta)/\Lambda$ and substitute in the spatial-curvature smallness condition. It is necessary for higher-derivative control, not sufficient for full quantum stability.
\end{proof}
\section{A falsifiable completion plan}
\begin{enumerate}
\item Specify a real one-scale microscopic theory with unitarity, symmetries, an exact stress--energy operator, and zero-mode removal procedures.
\item Compute the renormalized slope $\mathcal K'(0)$ from the full microscopic correlator, including contact terms, independently of the heat kernel; match both methods in one scheme.
\item Test regulator changes: $M_{\rm ren}$ and physical amplitudes must be invariant after (283). Requiring zero bare mass in every scheme immediately produces a contradiction.
\item Simultaneously check the vacuum term, induced $R^2$, $R_{\mu\nu}R^{\mu\nu}$, and nonanalytic form factors against (292).
\item Study coupled stability---Einstein constraints, matter modes, and anisotropy---before DESI/CMB likelihood comparison. Fixing $G$ does not remove the negative matter direction of Chapter XLVII.
\end{enumerate}
The results establish structure of the chosen effective scenario and its numerical underdetermination without extra physical inputs. Universal $G$, physical scale origins, and full stability remain unproved. Renaming an unknown numerical constant does not compute it.
\setcounter{equation}{292}


% 14_stability_barriers.tex
\chapter{Hopf Energy Concentration, Faddeev--Skyrme Stabilization, and Rigorous Cosmological Bounce Bounds}
\section{Novelty status and testable assumptions}
This chapter retains the locally covariant model of XLVI--XLVIII but examines its incomplete physical-stability argument. We prove theorems within an explicit effective model, not universal emergence of gravity. Variational properties of a single Hopf fibration in the Faddeev--Skyrme model were studied by R. S. Ward, Nonlinearity 12 (1999), 241, DOI 10.1088/0951-7715/12/2/005; J. M. Speight and M. Svensson, math/0605516; and A. Guerra, X. Lamy, K. Zemas, Communications in Mathematical Physics 407 (2026), article 171, DOI 10.1007/s00220-026-05702-5. We do not present those results as first discoveries. Calculations new to this construction concern three maps, their FLRW source, compatibility of bounce and matter second-variation sign, and tomography resources. Worldwide priority of this combination is unestablished.

Let $\gamma$ be the unit round $S^3$ metric, $n_a(q)=qe_aq^{-1}$, $\kappa=\beta r^2>0$, and $\omega_{S^2}$ the unit-sphere area form with integral $4\pi$. Put $F_a=n_a^*\omega_{S^2}$. The chosen extension is
\begin{equation}S=\int d^4x\sqrt{-g}\left[\frac{R-2\Lambda}{16\pi G}-\frac\kappa2\sum_{a=1}^3g^{\mu\nu}\partial_\mu n_a\cdot\partial_\nu n_a-\frac\lambda4\sum_{a=1}^3F_{a,\mu\nu}F_a^{\mu\nu}\right],\quad\lambda>0.\tag{293}\end{equation}
The quartic coupling is a new free action input, not automatically derived from $\kappa$ or Sakharov's formula. Signature is $(-,+,+,+)$, and $F_a$ denotes the pulled-back area two-form.
\section{First obstruction: topological charge without an energy gap}
\statement{Theorem}{LIII.1}{Collapse of the pure quadratic-energy infimum in every Hopf class}
For every integer $Q$, there are smooth $n_\varepsilon:S^3\to S^2$ with Hopf invariant $Q$ such that
\begin{equation}\inf_{\operatorname{Hopf}(n)=Q}\int_{S^3}|dn|_\gamma^2dV_\gamma=0,\qquad\int|dn_\varepsilon|_\gamma^2dV_\gamma=O(\varepsilon)\quad(\varepsilon\downarrow0).\tag{294}\end{equation}
For three maps with any prescribed charges $(Q_1,Q_2,Q_3)$, their total quadratic energy can simultaneously be made arbitrarily small while retaining each charge.
\begin{proof}
Represent each class by a smooth based map $h_Q:S^3\to S^2$. A smooth degree-one collapse $c_\varepsilon:S^3\to S^3$, constant outside a geodesic ball $B_\varepsilon$, has derivative $O(\varepsilon^{-1})$ inside; it smooths $B_\varepsilon/\partial B_\varepsilon\simeq S^3$. Then $n_\varepsilon=h_Q\circ c_\varepsilon$ retains Hopf invariant $Q$, has zero gradient outside the ball and gradient $O(\varepsilon^{-1})$ inside. Volume $O(\varepsilon^3)$ gives energy $O(\varepsilon)$. Use the same collapse for all three maps; nonnegativity completes the proof.
\end{proof}
\statement{Remark}{LIII.2}{What the theorem does not assert}
It excludes a universal positive bound $E_2\ge C(|Q|)>0$ in the pure three-dimensional sigma model. It does not keep arbitrary localized fields inside a strictly homogeneous FLRW ansatz: localization generally makes them inhomogeneous. Nor does it prove instability of the full Einstein--matter system after gravitational constraints.
\section{An exact tensor identity for the quartic stabilizer}
\statement{Lemma}{LIII.3}{Hopf two-forms and isotropy of their quadratic tensor}
In the unit sphere's orthonormal left frame,
\begin{equation}F_a=4\theta_b\wedge\theta_c\quad\text{for cyclic }(a,b,c),\qquad\sum_a|F_a|_\gamma^2=48,\qquad\sum_aF_{a,ik}F_{a,jk}=32\gamma_{ij}.\tag{295}\end{equation}
The triad is stationary for both quadratic and quartic matter functionals on fixed round $S^3$.
\begin{proof}
From $E_i(n_a)=2q(e_i\times e_a)q^{-1}$, each area pullback has horizontal coefficient 4 and squared norm 16. Its contracted tensor is $16(\gamma_{ij}-\theta_{ai}\theta_{aj})$; summing yields $32\gamma_{ij}$. Harmonicity was proved in XLVI. For the quartic variation, writing $\alpha_a=n_a^*\omega_{S^2}$ and integrating by parts gives a pairing of the codifferential $\delta\alpha_a$ with $n_a^*(\iota_{\delta n_a}\omega_{S^2})$. The former is vertical, along $\theta_a$, by $d\theta_a=-2\theta_b\wedge\theta_c$; the latter is horizontal and annihilates $\ker dn_a=\operatorname{span}(E_a)$. The first variation vanishes.
\end{proof}
\statement{Theorem}{LIII.4}{The exact FLRW system with a topological stabilizing term}
For $g=-dt^2+a(t)^2\gamma$, $\Lambda>0$, and spatial-frame-static fields, (293) satisfies the matter equations exactly and has perfect-fluid stress
\begin{equation}\rho=\frac{12\kappa}{a^2}+\frac{24\lambda}{a^4},\qquad p=-\frac{4\kappa}{a^2}+\frac{8\lambda}{a^4}.\tag{296}\end{equation}
With $h^2=\Lambda/3$, $\eta=32\pi G\kappa$, $D=64\pi G\lambda$, $C=1-\eta$,
\begin{equation}\dot a^2=h^2a^2-C+\frac D{a^2},\qquad\ddot a=h^2a-\frac D{a^3}.\tag{297}\end{equation}
A regular expanding branch with a nondegenerate bounce exists exactly when
\begin{equation}0<\eta<1,\qquad0<4h^2D<C^2.\tag{298}\end{equation}
Its exact trajectory is
\begin{equation}a(t)^2=\frac{C+\sqrt{C^2-4h^2D}\cosh[2h(t-t_0)]}{2h^2}.\tag{299}\end{equation}
At equality, the root is degenerate: a stationary branch exists, but there is no nonzero bounce acceleration.
\begin{proof}
For spatial two-forms, $T^{(4)}_{00}=(\lambda/4)(96/a^4)=24\lambda/a^4$, while the spatial tensor from (295) is $8\lambda a^{-2}\gamma_{ij}$, so $p_4=8\lambda/a^4$. Add the quadratic contribution from XLVI. Both separately obey continuity. The Friedmann equation gives (297). For $x=a^2$, $\dot x^2=4(h^2x^2-Cx+D)$; two positive nondegenerate roots exist exactly under (298). The upper root $x_+=[C+\sqrt{C^2-4h^2D}]/(2h^2)$ has positive acceleration. Substitution verifies (299) and both equations.
\end{proof}
\section{The exact conformal mode and impossibility of permanent protection}
\statement{Theorem}{LIII.5}{Triad energy along a conformal deformation}
For the conformal flow of XLVII, $\phi_s^*\gamma=\omega_s^2\gamma$, $\omega_s=(\cosh s+f\sinh s)^{-1}$, and $n_a^{(s)}=n_a\circ\phi_s$, at fixed $a>0$,
\begin{equation}E(s,a)=\operatorname{Vol}(S^3)\left[12\kappa a\operatorname{sech}^2(s/2)+\frac{24\lambda}a\cosh s\right],\tag{300}\end{equation}
with exact second variation
\begin{equation}\partial_s^2E(0,a)=\operatorname{Vol}(S^3)\left(-6\kappa a+\frac{24\lambda}a\right).\tag{301}\end{equation}
The conformal matter direction is nonnegative exactly when
\begin{equation}\frac{4\lambda}{\kappa a^2}\ge1.\tag{302}\end{equation}
If $a(t)\to\infty$ with fixed finite positive couplings, this eventually fails. A single standard Faddeev--Skyrme term cannot keep this second variation nonnegative throughout an unbounded expanding branch.
\begin{proof}
The quadratic gradient sum is $24\omega_s^2$ and the quartic two-form norm sum is $48\omega_s^4$. Exact averages on $S^3$ are $\langle\omega_s^2\rangle=\operatorname{sech}^2(s/2)$ and $\langle\omega_s^4\rangle=\cosh s$. For the latter, use $\phi_s^*dV=\omega_s^3dV$ to obtain
\[\int\omega_s^4(q)dV(q)=\int\omega_s(\phi_s^{-1}(y))dV(y)=\int[\cosh s-f(y)\sinh s]dV(y),\]
whose odd part vanishes. Substitute and differentiate. The last claim uses unbounded $a$ and constant couplings. It concerns one spatial matter-energy direction at fixed metric, not full dynamical GR instability.
\end{proof}
\statement{Theorem}{LIII.6}{Exact compatibility window at the bounce minimum}
Define
\begin{equation}\eta=32\pi G\kappa,\qquad\zeta=64\pi G\lambda h^2,\qquad h^2=\Lambda/3.\tag{303}\end{equation}
A nondegenerate bounce whose conformal matter second variation is nonnegative at the minimum exists exactly when
\begin{equation}0<\eta<\tfrac12,\qquad\frac\eta2-\frac{3\eta^2}4\le\zeta<\frac{(1-\eta)^2}4.\tag{304}\end{equation}
This is a condition for one specified fixed-metric variation, not stability of the full Einstein--matter system.
\begin{proof}
Let $b=h^2a_{\min}^2$, $C=1-\eta$. The upper root is $b=(C+\sqrt{C^2-4\zeta})/2$, with $\zeta=b(C-b)$ and $C/2<b<C$. Condition (302) is $\zeta\ge(\eta/2)b$, equivalently $b\le C-\eta/2$. The allowed interval is nonempty exactly when $\eta<C$, or $\eta<1/2$. On $b>C/2$, $b(C-b)$ decreases. Its allowed minimum is $(\eta/2)(C-\eta/2)=\eta/2-3\eta^2/4$, and its supremum is $C^2/4$, unattained for a nondegenerate bounce.
\end{proof}
\statement{Corollary}{LIII.7}{No weak quartic stabilization: exact equation-of-state cost}
Every triad configuration with nonnegative conformal matter Hessian requires
\begin{equation}\frac{\rho_4}{\rho_2}=\frac{2\lambda}{\kappa a^2}\ge\frac12,\qquad w_m=\frac{-1+2\lambda/(\kappa a^2)}{3[1+2\lambda/(\kappa a^2)]}\ge-\frac19.\tag{305}\end{equation}
Thus correcting the negative conformal mode cannot occur with quartic density negligible relative to quadratic density. This concerns the ratio $E_4/E_2$, not UV completeness of other microscopic models.
\begin{proof}
Equation (296) gives the density ratio, and (301) gives $\lambda/(\kappa a^2)\ge1/4$. The displayed equation-of-state function increases for nonnegative argument and equals $-1/9$ at $1/4$.
\end{proof}
\statement{Corollary}{LIII.8}{A finite time for loss of the conformal sign}
For every solution (299) in window (304), the Hessian is negative whenever $|t-t_0|>t_{\rm crit}$, where
\begin{equation}t_{\rm crit}=\frac1{2h}\operatorname{arcosh}\left[\frac{8h^2\lambda/\kappa-(1-\eta)}{\sqrt{(1-\eta)^2-4\zeta}}\right].\tag{306}\end{equation}
At the lower window boundary $t_{\rm crit}=0$; in its interior it is positive.
\begin{proof}
Set $a^2=4\lambda/\kappa$ in (299) and solve for the hyperbolic cosine. Monotonicity after the bounce gives uniqueness.
\end{proof}
\statement{Corollary}{LIII.9}{Double degeneracy of background cosmological observations}
With separately conserved dust and standard radiation,
\begin{equation}\frac{H(z)^2}{H_0^2}=\Omega_m(1+z)^3+(\Omega_{r,\rm std}+\Omega_{H4})(1+z)^4+(\Omega_K+\Omega_{H2})(1+z)^2+\Omega_\Lambda.\tag{307}\end{equation}
Here $\Omega_{H2}=32\pi G\kappa/(H_0^2a_0^2)$, $\Omega_{H4}=64\pi G\lambda/(H_0^2a_0^4)$, and $\Omega_K=-1/(H_0^2a_0^2)$. $H(z)$ alone cannot separately recover either pair $(\Omega_{r,\rm std},\Omega_{H4})$ or $(\Omega_K,\Omega_{H2})$. Angular distances, linear perturbations, anisotropic stress, and other observables can break this exact background degeneracy.
\begin{proof}
The quadratic and quartic sources scale as $a^{-2}$ and $a^{-4}$, dust as $a^{-3}$, and radiation as $a^{-4}$. Substitute in Friedmann with $1+z=a_0/a$. Only the parameter sums appear, leaving their separate inversion underdetermined.
\end{proof}
\section{Resource consequence: finite tomography versus a weak gravitational source}
\statement{Theorem}{LIII.10}{Three Hopf-charge certificates and a measurement budget}
Let $\rho_a(q)=(I+rn_a(q)\cdot\sigma)/2$, $0<r\le1$, for the exact Hopf triad. Suppose $M$ nodes have a continuous partition of unity locally supported within geodesic radius $h_{\rm mesh}<1/2$. Measure all three Pauli components for all three sectors at each node, $N_{\rm shots}$ independent times each, totaling $9MN_{\rm shots}$ outcomes. If
\begin{equation}N_{\rm shots}>\frac{6\log(18M/\alpha)}{r^2(1-2h_{\rm mesh})^2},\qquad0<\alpha<1,\tag{308}\end{equation}
then with probability at least $1-\alpha$, all three reconstructed continuous normalized fields retain the original Hopf classes, assuming the known bound $\|dn_a\|_{\rm op}\le2$. Equivalently, for $\rho_2=12\beta r^2/a^2$,
\begin{equation}N_{\rm shots}\rho_2>\frac{72\beta\log(18M/\alpha)}{a^2(1-2h_{\rm mesh})^2}\tag{309}\end{equation}
is sufficient when comparing fixed $\beta,a,M,h_{\rm mesh},\alpha$.
\begin{proof}
$v_a=rn_a$ is $2r$-Lipschitz. Hoeffding and union over $9M$ coordinates give simultaneous coordinate error $t=\sqrt{2\log(18M/\alpha)/N_{\rm shots}}$, hence vector error $\eta_N=\sqrt{6\log(18M/\alpha)/N_{\rm shots}}$. A continuous convex local interpolant obeys $|\widehat v_a(q)-rn_a(q)|\le2rh_{\rm mesh}+\eta_N<r$. The straight homotopy avoids zero, preserving normalized Hopf class. Substitute $\rho_2$ to get (309).
\end{proof}
\statement{Remark}{LIII.11}{A statistical threshold is not Newton's law}
Equation (309) concerns a particular measurement protocol for an already specified field with known regularity and action. It neither makes gravity depend on measurement count nor lower-bounds $G_N$. It is sufficient, not necessary, and excludes correlated noise, preparation errors, and unknown calibration. A GKLS dissipative gap may stay constant as $r\downarrow0$ while $N_{\rm shots}\sim r^{-2}$ diverges.
\section{The program's fundamental boundary}
Four precise restrictions have been checked: conserved Hopf charge alone fixes neither positive $E_2$ nor gravitational coupling; positive quartic interaction gives exact cosmological trajectories and a testable nondegenerate-bounce criterion; compatibility with a nonnegative conformal direction requires window (304), inevitably lost on an infinite expanding branch at constant couplings; and finite tomography assigns an operational cost to distinguishing topology and the information-geometric source.

A fundamental physical law still requires the fully coupled scalar/vector/tensor perturbations with Einstein constraints, microscopic derivation of $\beta,\lambda$, independent computation of $G$ with all counterterms, realistic origins of cosmic field content, and likelihood fits to real BAO/CMB/lensing data. Symbolic tests replace none of these.

\chapter{A Robust Dimensional Barrier, Polynomial Witnesses, and a Spectral Audit}
\section{Result status and inclusion rules}
This section continues the unified manuscript; it does not replace Chapters I--XLIX. It distinguishes rigorously proved conditional statements, results explicitly correcting the submitted revision L, and open questions. Global physical novelty and independently verified publication priority are not claimed.
\section{Theorem L.2+: a stochastic coarse-graining barrier with imperfect calibration}
Use column-stochastic matrices, $\|A\|_{1\to1}=\max_j\sum_i|A_{ij}|$, and $\operatorname{TV}(p,q)=\|p-q\|_1/2$. On microstates $e_1,\ldots,e_n$, $n>d\ge2$, define
\begin{equation}Pe_1=e_2,\quad Pe_j=e_{j+1}\ (2\le j<n),\quad Pe_n=e_2,\quad Qe_1=Qe_2=e_1,\quad Qe_j=e_j\ (j\ge3).\tag{310}\end{equation}
Let $L\in\R_{\ge0}^{d\times n}$ and $\bar P,\bar Q\in\R_{\ge0}^{d\times d}$ be column stochastic, with macro-effect $\bar\alpha\in[0,1]^d$. Define
\begin{equation}\varepsilon=\tfrac12\max\{\|LP-\bar PL\|_{1\to1},\|LQ-\bar QL\|_{1\to1}\},\qquad\delta=\|\bar\alpha^TL-e_1^T\|_{\max}<1.\tag{311}\end{equation}
For rows, the maximum norm means maximum absolute coordinate.
\statement{Theorem}{LIV.1}{The robust dimensional barrier L.2+}
Put $\gamma_\delta=\delta+\delta/(1-\delta)$. Then for every $n>d\ge2$,
\begin{equation}\varepsilon\ge\max\left\{0,\frac{1/d-\gamma_\delta}{d},\frac{1-d\gamma_\delta}{d\sqrt{(d+1)(2d+1)/6}}\right\}.\tag{312}\end{equation}
At exact calibration,
\begin{equation}\varepsilon\ge\max\left\{\frac1{d^2},\frac{\sqrt6}{d\sqrt{(d+1)(2d+1)}}\right\}>0,\tag{313}\end{equation}
independently of $n$.
\begin{proof}
Let $v_0=e_2$ and $v_s=e_{n-s+1}$ for $1\le s\le d-1$. These indices are distinct. The cycle and $Q$ give
\begin{equation}e_1^TQP^rv_s=\one_{r=s},\qquad0\le r,s<d.\tag{314}\end{equation}
Set $u=Le_1$, $\beta=\bar\alpha^Tu\ge1-\delta>0$, and $t_s=\bar\alpha^TLv_s/\beta\in[0,\delta/(1-\delta)]$. The vectors $w_s=Lv_s-t_su$ belong to $\ker\bar\alpha^T$, of dimension at most $d-1$, so
\begin{equation}R_{rs}=\bar\alpha^T\bar Q\bar P^rw_s,\qquad0\le r,s<d,\tag{315}\end{equation}
has rank at most $d-1$. For $G_{rs}=\bar\alpha^T\bar Q\bar P^rLv_s$, telescoping one-step defects gives $|G_{rs}-\bar\alpha^TLQP^rv_s|\le(r+1)\varepsilon$. Calibration and (314) imply $|G_{rs}-\delta_{rs}|\le(r+1)\varepsilon+\delta$. Since $G_{rs}-R_{rs}=t_s\bar\alpha^T\bar Q\bar P^ru$ and the last factor lies in $[0,1]$,
\begin{equation}|R_{rs}-\delta_{rs}|\le(r+1)\varepsilon+\gamma_\delta.\tag{316}\end{equation}
A singular $R$ must satisfy $\|I_d-R\|_{\max}\ge1/d$; otherwise its induced infinity norm is below one and the Neumann series makes $R$ invertible. Thus $1/d\le d\varepsilon+\gamma_\delta$. Eckart--Young also gives $\|I_d-R\|_F\ge1$, while (316) yields
\[1\le\|I_d-R\|_F\le\varepsilon\sqrt{d\sum_{j=1}^dj^2}+d\gamma_\delta=\varepsilon d\sqrt{(d+1)(2d+1)/6}+d\gamma_\delta.\]
These are the two nontrivial bounds.
\end{proof}
\statement{Remark}{LIV.2}{The theorem's boundary}
For fixed $d$ and this transition pair with the specified effect restriction, error remains positive as $n\to\infty$. The theorem neither gives the optimal constant, excludes effective models with growing $d$, nor automatically extends to arbitrary microdynamics.
\section{Corrected theorem L.3: a hierarchy of finite-time witnesses}
Let $x_k=e^TM^k\mu$ be an $m$-state stochastic response, with admissible effect and preparation. Put $p=m-1\ge1$, $\Delta_k=x_{k+1}-x_k$. Exactly $2m-1$ samples $x_0,\ldots,x_{2m-2}$ define $H=[\Delta_{i+j}]_{0\le i,j<p}$ and $b=[\Delta_{i+p}]_{0\le i<p}$. Let $D=\det H$, $A=\det H_A$ where the last column is replaced by $b$, and, for $p\ge2$, $B=\det H_B$ where the penultimate column is replaced by $-b$. For $p=1$, set $B=0$.
\statement{Theorem}{LIV.3}{A polynomial trace witness}
Every such model satisfies
\begin{equation}T_m=(m-1)D^2-2DA+(m-1)A^2-2mDB\ge0.\tag{317}\end{equation}
The separate condition $\det[\Delta_{i+j}]_{0\le i,j<m}=0$ needs at least $2m$ samples $x_0,\ldots,x_{2m-1}$, not $2m-1$.
\begin{proof}
The characteristic polynomial contains $(z-1)$, so differences satisfy a recurrence of order at most $p$, determined by the remaining eigenvalues $\lambda_j$. Write $e_1=\sum_j\lambda_j$, $e_2=\sum_{i<j}\lambda_i\lambda_j$. If $D\ne0$, Cramer's rule gives $A=De_1$, $B=De_2$. If $D=0$, the minimal Hankel rank is below $p$, so all $p\times p$ minors of $[H\mid b]$, including $A,B$, vanish. Nonnegativity gives $\Tr(M^2)\ge\sum_iM_{ii}^2\ge(\Tr M)^2/m$. Since $\Tr M=1+e_1$ and $\Tr(M^2)=1+e_1^2-2e_2$, one gets $(m-1)-2e_1+(m-1)e_1^2-2me_2\ge0$. Multiply by $D^2$.
\end{proof}
\statement{Corollary}{LIV.4}{The corrected formula for coherent poles}
For $m=2K+1$, suppose the $2K$ visible nonzero difference poles are $\lambda_{j,\pm}=\rho_je^{\pm i\theta_j}$, with $D\ne0$. If $E_1=2\sum_{j=1}^K\rho_j\cos\theta_j$, then
\begin{equation}T_{2K+1}=D^2\left[2K-2E_1-E_1^2+2(2K+1)\sum_{j=1}^K\rho_j^2\cos(2\theta_j)\right].\tag{318}\end{equation}
For every $K\ge1$, distinct phases sufficiently close to $\pi/2$ make this negative. The real bounded signal $x_k=1/2+(2K)^{-1}\sum_j\cos(k\theta_j)$ then has a finite-dimensional quantum realization but no $(2K+1)$-state classical realization of its first $4K+1$ samples.
\begin{proof}
Use $e_1=E_1$ and $e_1^2-2e_2=\sum_j\lambda_j^2=2\sum_j\rho_j^2\cos2\theta_j$. At unit radii and phases tending to $\pi/2$, the bracket tends to $2K-2(2K+1)K=-4K^2<0$. Distinct nearby phases give nonzero $D$ by Vandermonde, since poles are distinct and amplitudes nonzero. A probabilistic mixture of $K$ equatorially initialized phase-rotating qubits with a common binary effect realizes the signal.
\end{proof}
\statement{Remark}{LIV.5}{Two different sampling requirements}
Witness (317) uses $2m-1$ values, while the zero $m\times m$ difference-Hankel determinant uses $2m$. For $m=2$, the polynomial is identically nonnegative for every sequence and distinguishes nothing.
\section{Corrected spectral energy on a finite abelian group}
Let $G$ have order $N$, $\varnothing\ne A\subset G$, $a=|A|$, and $r_k=\one_A^{*k}$ under unnormalized convolution. Use $\widehat f(\chi)=\sum_{x\in G}f(x)\overline{\chi(x)}$, and define $E_{2k}(A)=\sum_xr_k(x)^2$, $p_\chi=|\widehat{\one_A}(\chi)|^2/(Na)$.
\statement{Proposition}{LIV.6}{An exact entropy--energy identity}
For $k\ge2$ and $H_k=(1-k)^{-1}\log\sum_\chi p_\chi^k$,
\begin{equation}E_{2k}(A)=\frac1N\sum_{\chi\in\widehat G}|\widehat{\one_A}(\chi)|^{2k}=N^{k-1}a^ke^{-(k-1)H_k},\tag{319}\end{equation}
and
\begin{equation}\frac{|kA|}{|A|}\ge\left(\frac aN\right)^{k-1}e^{(k-1)H_k}.\tag{320}\end{equation}
\begin{proof}
Plancherel and $\widehat r_k=\widehat{\one_A}^{\,k}$ give the first equality; substitute $p_\chi$. Cauchy--Schwarz gives $a^{2k}=(\sum_xr_k(x))^2\le|kA|E_{2k}$. The factor $N^{-(2k-1)}$ in one submitted formula is wrong for this unnormalized convention.
\end{proof}
\section{Checking Berger geometry and spectral asymmetry}
For $g_\tau=\sigma_1^2+\sigma_2^2+\tau^2\sigma_3^2$, under the submitted normalization, $a_0=2\pi^2\tau$ and $a_1/a_0=(4-\tau^2)/3$. For $0<\tau<2$, Hitchin's known formula in this convention is $\eta_D(0)=-(\tau^2-1)^2/6$. Algebraically,
\begin{equation}\frac{a_1}{a_0}+\frac{\operatorname{sgn}(\tau^2-1)}3\sqrt{-6\eta_D(0)}=1,\qquad\tau\ne1.\tag{321}\end{equation}
Both sides extend continuously at one. Indeed, the square root is $|\tau^2-1|$, making the second term $(\tau^2-1)/3$. This is not an independent new physical identity. The earlier proposed expression involving $3(8-3\tau^2+|\tau^2-1|)$ is false: at $\tau=1$, it gives 15 instead of 5.
\section{Corrected laboratory version L.4: calibrated binary observation}
Let $v:X\to\R^3$, $X=S^2$ or $S^3$, satisfy $|v|\ge r_{\min}>0$ and Lipschitz constant $L_v$. At each of $M$ nodes, component $Y_\alpha\in\{-1,+1\}$ has mean $v_\alpha(x)$. An independent calibrated binary symmetric channel flips the sign with known probability $q_\alpha<1/2$. Write $c_\alpha=1-2q_\alpha$, $c_{\min}=\min c_\alpha>0$, and observed output $Z_\alpha$.
\statement{Proposition}{LIV.7}{A topological certificate after calibrated information loss}
From $N_s$ independent measurements per component per node, use $\widehat v_\alpha=(N_sc_\alpha)^{-1}\sum_{j=1}^{N_s}Z_{\alpha,j}$. With probability at least $1-\alpha_0$, simultaneously at all nodes,
\begin{equation}\|\widehat v(x_i)-v(x_i)\|_2\le\frac{\sqrt{6\log(6M/\alpha_0)}}{c_{\min}\sqrt{N_s}}.\tag{322}\end{equation}
If this plus $L_vh$ is strictly below $r_{\min}$, continuous interpolation with local convex weights on a radius-$h$ covering yields homotopic normalized fields: equal degree/Chern number on $S^2$, equal Hopf invariant on $S^3$. For a balanced-label binary experiment with the same symmetric channel, $\mathcal J(Z)=c_\alpha^2\mathcal J(Y)$; the budget grows by $c_{\min}^{-2}$.
\begin{proof}
$\mathbb EZ_\alpha=c_\alpha v_\alpha$, so the estimator is unbiased. Hoeffding gives $\Pr(|\widehat v_\alpha-v_\alpha|>t/c_\alpha)\le2e^{-N_st^2/2}$. Union over $3M$ coordinates with $t=\sqrt{2\log(6M/\alpha_0)/N_s}$ proves (322). Local convex interpolation adds at most $L_vh$; the straight homotopy avoids zero. In the balanced experiment, $\mathbb E[B\mid Y]=v_\alpha Y$ and $\mathbb E[B\mid Z]=c_\alpha v_\alpha Z$, giving the squared contrast ratio.
\end{proof}
\statement{Remark}{LIV.8}{An invalid step in the original L.4}
$\mathbb E[B\mid Z]$ is random, not the deterministic component $v(x)$. The proposed $\mathbb E[Y\mathbb E(B\mid Z)]$ requires paired $(Y,Z)$ data unavailable after deleting $Y$. Original L.4 is therefore unproved without calibration or extra data; the present version explicitly states accessible observations and measurement procedure.
\section{Open questions not promoted to theorems}
\begin{enumerate}
\item The full staircase $E_m(f_{n,a,b})$ from L.1 remains a conditional candidate: reduction to independent residue classes for arbitrary periods $d_j$ is unproved. Attaining constructions provide upper bounds for particular partitions, not a universal equality.
\item Minimizing $\sum_{m\le M}(m\|m\omega\|_\Z)^{-2}/\log M$ at $\omega=\varphi$ does not follow merely from Hurwitz's theorem and the Lagrange constant. The claimed unique minimizer is unproved.
\item The exact optimum $\varepsilon(d)$, normalization-premium scaling on hypercubes, and a physically independent derivation of $G_N$ remain open.
\end{enumerate}
\setcounter{equation}{322}


% 15_ricci_clocks.tex
\chapter{An Exact Spectral Chronometer for Ricci Flow on Berger $S^3$}
\note{Status and priority.} This section gives a self-contained curvature calculation on $SU(2)$ and exactly integrates the ratio of specified scalar eigenvalues under homogeneous Ricci flow. Berger geometry, Hitchin's formula, Ricci-flow analysis, and spectral evolution have prior literatures. These combinations may be new forms of known equations within this program; worldwide priority is not claimed. Auxiliary Ricci-flow time must not be confused with cosmological time in Einstein--matter dynamics.
\section{Normalization, commutators, and an independent Ricci-tensor derivation}
On $SU(2)\cong S^3$, use left-invariant fields normalized by
\begin{equation}[X_1,X_2]=2X_3,\quad[X_2,X_3]=2X_1,\quad[X_3,X_1]=2X_2.\tag{323}\end{equation}
The dual forms obey $d\sigma_1=-2\sigma_2\wedge\sigma_3$ cyclically, with sign reversed by orientation convention. Fix
\begin{equation}g_{a,\tau}=a^2(\sigma_1^2+\sigma_2^2+\tau^2\sigma_3^2),\qquad a,\tau>0.\tag{324}\end{equation}
In the orthonormal frame $e_1=X_1/a$, $e_2=X_2/a$, $e_3=X_3/(a\tau)$, at $a=1$ the structure constants are $c^3_{12}=2\tau$, $c^1_{23}=c^2_{31}=2/\tau$. Koszul's formula is
\begin{equation}2\langle\nabla_{e_i}e_j,e_k\rangle=c^k_{ij}-c^i_{jk}+c^j_{ki}.\tag{325}\end{equation}
Together with $R(e_i,e_j)e_k=\nabla_i\nabla_je_k-\nabla_j\nabla_ie_k-\nabla_{[e_i,e_j]}e_k$, it yields $\nabla_1e_2=\tau e_3$, $\nabla_2e_1=-\tau e_3$, $\nabla_1e_3=-\tau e_2$, $\nabla_2e_3=\tau e_1$, $\nabla_3e_1=(2/\tau-\tau)e_2$, $\nabla_3e_2=(\tau-2/\tau)e_1$. At general $a$, connection coefficients gain $a^{-1}$.
\statement{Theorem}{LV.1}{Berger Ricci tensor and curvature anisotropy}
In the orthonormal frame,
\begin{align}\operatorname{Ric}_{g_{a,\tau}}&=a^{-2}\diag(4-2\tau^2,4-2\tau^2,2\tau^2),\tag{326}\\
R_g&=(8-2\tau^2)/a^2,\qquad\operatorname{Vol}(g)=2\pi^2a^3\tau,\tag{327}\\
|\operatorname{Ric}^{\circ}|_g^2&=\frac{32}{3a^4}(\tau^2-1)^2.\tag{328}\end{align}
\begin{proof}
Substitute the structure constants into (325); the six coefficients above are the only nonzero ones. Compute curvature and contract $\operatorname{Ric}(e_j,e_k)=\sum_i\langle R(e_i,e_j)e_k,e_i\rangle$. Metric scaling gives $a^{-2}$ for orthonormal Ricci components. Volume uses the frame determinant and round volume $2\pi^2$. Finally, $|\operatorname{Ric}^\circ|^2=\Tr(\operatorname{Ric}^2)-R^2/3$. An independent symbolic Koszul/curvature implementation is in \path{checks/verify_v51_su2_spectral_flow.py}.
\end{proof}
\section{The scalar spectrum and a scale-invariant anisotropy sensor}
Unimodularity makes the $X_i$ divergence free, so the nonnegative Laplace--Beltrami operator is
\begin{equation}-\Delta_g=-a^{-2}(X_1^2+X_2^2+\tau^{-2}X_3^2).\tag{329}\end{equation}
Peter--Weyl gives Casimir eigenvalues $n(n+2)$ and $-X_3^2$ weights $m^2$, where $m=-n,-n+2,\ldots,n$, each with left multiplicity $n+1$. Thus
\begin{equation}\lambda_{n,m}=a^{-2}[n(n+2)+(\tau^{-2}-1)m^2].\tag{330}\end{equation}
Choose actual weights $m=0,2$ in shell $n=2$, avoiding the previously noted use of nonexistent $m=0$ in odd shells. Define
\begin{equation}\mathcal A(g)=\frac{\lambda_{2,2}-\lambda_{2,0}}{\lambda_{2,0}}=\frac{\tau^{-2}-1}2,\qquad\tau^2=\frac1{1+2\mathcal A}.\tag{331}\end{equation}
This uses symmetry-labeled levels, not necessarily the two lowest unlabeled eigenvalues.
\section{Exactly integrable spectral evolution under Ricci flow}
For ordinary $\partial_tg=-2\operatorname{Ric}$ or volume-normalized $\partial_tg=-2\operatorname{Ric}+\tfrac23Rg$, homogeneity and $U(2)$ symmetry retain the Berger family. In the ordinary flow, let $A=a^2$, $B=a^2\tau^2$, and $y=B/A$:
\begin{equation}\dot A=-8+4y,\qquad\dot B=-4y^2,\qquad\dot y=\frac8A y(1-y).\tag{332}\end{equation}
The normalized term adds the same relative scaling to $\dot A$ and $\dot B$, canceling in $\dot y$.
\statement{Theorem}{LV.2}{An exact spectral chronometer for homogeneous Ricci flow}
On any smooth existence interval, set $s(t)=\int_0^t8/a(u)^2\,du$. For both flows,
\begin{equation}\frac{d\mathcal A}{ds}=-\mathcal A,\qquad\mathcal A(s)=\mathcal A(0)e^{-s},\qquad\tau^2(s)=\frac1{1+2\mathcal A(0)e^{-s}}.\tag{333}\end{equation}
If $\mathcal A(0)\ne0$, two symmetry-identified levels alone recover elapsed geometric time:
\begin{equation}s=\log\frac{|\mathcal A(0)|}{|\mathcal A(s)|}.\tag{334}\end{equation}
The sign is preserved and $\tau\to1$ as $s\to\infty$.
\begin{proof}
Equation (332) gives $dy/ds=y(1-y)$. Since $2\mathcal A=y^{-1}-1$, $2\mathcal A'=-y^{-2}y(1-y)=-2\mathcal A$. Integrate. The denominator is positive because $1+2\mathcal A(0)=\tau(0)^{-2}>0$ and $e^{-s}\in(0,1]$ for forward time.
\end{proof}
\section{Curvature rigidity and spectral asymmetry}
Restrict to $0<\tau<4$ and the standard spin structure. In this Dirac-zero-mode-free region, Hitchin's formula is $\eta(D_\tau)=-(\tau^2-1)^2/6$ with the stated operator orientation. At $\tau\ge4$, spectral-flow corrections occur and the following formulas do not apply as written.
\statement{Corollary}{LV.3}{An exact local--global relation within the Berger family}
For $0<\tau<4$,
\begin{equation}a^4|\operatorname{Ric}^{\circ}|^2=-64\eta(D_\tau)=\frac{32}3\left(\frac{2\mathcal A}{1+2\mathcal A}\right)^2.\tag{335}\end{equation}
Along the flow starting in this interval,
\begin{equation}\frac{d\eta(D_{\tau(s)})}{ds}=\tfrac13\tau(s)^2[\tau(s)^2-1]^2\ge0.\tag{336}\end{equation}
Equality at finite time occurs only for the round metric.
\begin{proof}
Substitute the Ricci and Hitchin formulas. Differentiate $\eta=-(y-1)^2/6$ using $y'=y(1-y)$ to obtain $y(y-1)^2/3$.
\end{proof}
\note{Physical interpretation.} This is a Ricci-flow relation, not Einstein dynamics or a universal gravitational field. Cosmology needs an independent evolution system with constraints, a source tensor, and observational tests. These identities do not determine $G_N$. The spectral indicator is a function of homogeneous anisotropy, not a new independent topological charge; $\eta$ is also a function of $\mathcal A$, not an additional first integral.
\section{Checks and scientific reproducibility}
\begin{enumerate}
\item The orthonormal symbolic algorithm starts only from (323), constructs connection, curvature, and Ricci, and does not substitute the final formula.
\item Both flow normalizations are derived directly, checking cancellation of the common $\tfrac23Rg$ contribution from $B/A$.
\item Actual weights $n=2$, $m=0,2$ and scale invariance of their ratio are checked.
\item Equation (335) is tested only for $0<\tau<4$; $\tau=4$ is excluded from the formula without spectral flow.
\end{enumerate}
Primary-source context: N. Hitchin, \emph{Harmonic Spinors}, Adv. Math. 14 (1974); S. Goette, \emph{Computations and Applications of $\eta$ Invariants} (2011); M. Capoferri et al., \emph{Beyond the Hodge Theorem: Curl and Asymmetric Pseudodifferential Projections}, J. Lond. Math. Soc. (2026), Appendix D, \url{https://doi.org/10.1112/jlms.70431}; \emph{Harmonic Spinors in the Ricci Flow}, J. Geom. Anal. (2024), \url{https://doi.org/10.1007/s12220-024-01665-y}. Familiar starting identities do not justify worldwide-discovery claims without independent priority research.

\chapter{Overlapping Cyclic Clocks, Exact Memory, and a Gauge-Invariant Hankel Pencil}
\note{Scientific correction.} The incoming v52 text declared the universal full periodic-memory staircase L.1 proved. An exact counterexample below invalidates that claim. In its place we establish an exact finite linear-programming reduction, an optimal tail step for an infinite period class, and a strictly conditional theorem on uncalibrated spectral reconstruction. Valid identities do not establish worldwide priority.
\section{Counterexample to L.1: period 42 with 41 states}
For $n\ge2$, let $f_n(k)=\one_{n\mid k}$ and define
\begin{equation}E_m(f_n)=\inf_{\substack{M\ge0,\ \one^TM=\one^T\\\mu\in\Delta_{m-1},\ e\in[0,1]^m}}\sup_{k\ge0}|f_n(k)-e^TM^k\mu|.\tag{337}\end{equation}
Chains are homogeneous, transitions column stochastic, and measurements physically normalized.
\statement{Theorem}{LVI.1}{An explicit counterexample to the universal full staircase L.1}
For $n=42$, $m=41$,
\begin{equation}E_{41}(f_{42})=\frac14,\tag{338}\end{equation}
whereas submitted v52 formula L.1 requires $1/3$ and is false in general.
\begin{proof}
Use disjoint deterministic cycles of lengths 6, 14, and 21, totaling 41 states. Give each initial phase mass $1/3$, effect $3/4$ at phase zero and zero elsewhere. Then
\begin{equation}g(k)=\tfrac14(\one_{6\mid k}+\one_{14\mid k}+\one_{21\mid k}).\tag{339}\end{equation}
At multiples of 42 it is $3/4$. Elsewhere at most one indicator is one, because any pair has lcm 42. The uniform error is $1/4$. Theorem LVI.3 below supplies the matching lower bound for $\omega(42)=3$ and every $m<42$. In the incorrect L.1, prime-power factors are $2,3,7$, $K_2(42)=\min\{6+7,14+3,21+2\}=13$, and $K_1(42)=42$, predicting $\max\{1/4,1/3\}=1/3$ at 41 states.
\end{proof}
\note{Locating the proof defect.} Periods 6, 14, and 21 overlap in prime factors, each containing two of $2,3,7$. Their additive clock sum cannot be represented without loss as a sum on disjoint prime-power blocks. The Chinese remainder theorem does not justify that step. L.1 loses its universal-theorem status; its block-partition constructions remain admissible upper bounds.
\section{A rigorous finite-dimensional reduction to cyclic models}
\statement{Theorem}{LVI.2}{Exact LP reduction of periodic Markov minimax error}
Let $f:\Z_{\ge0}\to[0,1]$ have period $n$. For each nonempty set of distinct positive divisors $D\subseteq\{d:d\mid n\}$ with $\sum_{d\in D}d\le m$, define $\operatorname{LP}(D;f)$ in variables $w_d\ge0$, $h_{d,s}\ge0$ for $0\le s<d$, and $\varepsilon\ge0$:
\begin{equation}\begin{aligned}
&\min\varepsilon,\qquad\sum_{d\in D}w_d=1,\qquad0\le h_{d,s}\le w_d,\\
&-\varepsilon\le f(k)-\sum_{d\in D}h_{d,k\bmod d}\le\varepsilon\qquad(0\le k<n).
\end{aligned}\tag{340}\end{equation}
Then
\begin{equation}E_m(f)=\min_{\substack{\varnothing\ne D\subseteq\{d:d\mid n\}\\\sum_{d\in D}d\le m}}\operatorname{LP}(D;f).\tag{341}\end{equation}
For rational $f(k)$, the minimax value is rational and admits a finite combinatorial-linear exact algorithm.
\begin{proof}
An LP solution gives cycles of lengths $d$, with initial phase-zero mass $w_d$ and effect $h_{d,s}/w_d$ (zero if $w_d=0$). Effects lie in $[0,1]$ and masses sum to one, proving one inequality.

Conversely, any $m$-state model with error $\varepsilon_0$ has an asymptotically periodic response $g_\infty=\sum_jg_j$, with periods $d_j$, $0\le g_j\le w_j$, $\sum_jw_j=1$, and $\sum_jd_j\le m$. This follows by separating transient states and closed irreducible classes; $w_j$ are eventual class-entry masses. Since the target is periodic and $g-g_\infty\to0$, the asymptotic response retains error at most $\varepsilon_0$ on every common-period residue.

Let $L=\lcm(n,d_1,\ldots,d_t)$ and $T=L/n$. Average target-preserving shifts:
\[\bar g(k)=T^{-1}\sum_{h=0}^{T-1}g_\infty(k+hn),\qquad\bar g_j(k)=T^{-1}\sum_{h=0}^{T-1}g_j(k+hn).\]
Error does not increase, and each component now has period $\delta_j=\gcd(d_j,n)$, since shifts by $n$ cover exactly the corresponding residue cosets. Each $\delta_j\mid n$ and their sum is at most $m$. Merge equal periods while retaining $0\le h_{d,s}\le w_d$. This produces a feasible LP with no larger error. Taking the infimum and the minimum over finitely many compact LPs proves equality.
\end{proof}
\section{An exact tail-step law when $\sum_i1/q_i<1$}
Write $n=\prod_{i=1}^rp_i^{\alpha_i}=\prod_iq_i$, $r=\omega(n)$, and
\begin{equation}M(n)=\sum_{i=1}^r\frac n{q_i}.\tag{342}\end{equation}
\statement{Theorem}{LVI.3}{A universal lower bound when a full period is unavailable}
For $f_{n,b}(k)=b\one_{n\mid k}$, $b\in[0,1]$, every homogeneous model with $m<n$ states satisfies
\begin{equation}E_m(f_{n,b})\ge\frac b{\omega(n)+1}.\tag{343}\end{equation}
\begin{proof}
By LVI.2, use nonnegative periodic components $h_d$ with proper divisor periods $d\mid n$. For each $d$, some $i$ has $v_{p_i}(d)<\alpha_i$, hence $d\mid n/p_i$ and $h_d(n/p_i)=h_d(0)$. All other values are nonnegative, so
\begin{equation}g(0)\le\sum_{i=1}^rg(n/p_i).\tag{344}\end{equation}
None of $n/p_i$ is divisible by $n$. Accuracy gives $g(0)\ge b-\varepsilon$ and each term at most $\varepsilon$, hence $b-\varepsilon\le r\varepsilon$.
\end{proof}
\statement{Theorem}{LVI.4}{The exact tail step from overlapping clocks}
If $M(n)\le m<n$, then
\begin{equation}E_m(f_{n,b})=\frac b{r+1},\qquad r=\omega(n).\tag{345}\end{equation}
This interval is nonempty exactly when $\sum_i1/q_i<1$.
\begin{proof}
Use $r$ disjoint cycles of lengths $n/q_i$, each with mass $1/r$, phase-zero effect $rb/(r+1)\le1$, and zero elsewhere:
\begin{equation}g(k)=\frac b{r+1}\sum_{i=1}^r\one_{n/q_i\mid k}.\tag{346}\end{equation}
At multiples of $n$, all indicators are one. Otherwise choose $j$ with $q_j\nmid k$; all indicators except possibly $i=j$ vanish. Error is exactly $b/(r+1)$. Pad unused states with zero mass and apply LVI.3. Finally, $M(n)<n$ is equivalent to the reciprocal-sum condition.
\end{proof}
\statement{Corollary}{LVI.5}{An infinite counterexample family and the actual scope of L.1}
For prime $p\ge7$ and $0<b\le1$,
\begin{equation}E_m(f_{6p,b})=b/4\qquad(6+5p\le m<6p).\tag{347}\end{equation}
At $m=6+5p$, the false L.1 staircase predicts $b/3$, since its block number is two.
\section{Numerical checking and the actual staircase for $n=42$}
The full LP implementation gives the following rational-value table, consistent with the analytical thresholds:
\begin{center}\begin{tabular}{lll}\toprule Budget $m$&$E_m(f_{42})$&Attaining periods\\\midrule
$1\le m\le11$&$1/2$&1\\12&$2/5$&2, 3, 7\\$13\le m\le40$&$1/3$&6, 7\\41&$1/4$&6, 14, 21\\$m\ge42$&0&42\\\bottomrule\end{tabular}\end{center}
This finite numerical enumeration checks computation, but the preceding proofs do not rely on numerical optimization. Fully automated rational dual certificates for every intermediate LP are not yet attached; treat the table as numerically established except for analytically proved rows.
\section{Stable uncalibrated spectral reconstruction of a Hankel pencil}
\note{Notation distinction.} In v50, L.4 denoted an unproved uncalibrated topology--information certificate. The new submission reuses L.4 for the unrelated problem of gauge-invariant shift-pencil reconstruction. We give it a new label to avoid confusion.
\statement{Theorem}{LVI.6}{Hankel-pencil spectra under constant diagonal calibration}
Suppose the series has exact Hankel rank $r$, with invertible $r\times r$ base panel $H_0=OR$, where observability and reachability matrices $O,R$ are square and invertible. For all words $w$, let $H_w=OT_wR$. If the same invertible diagonal $D_L,D_R$ apply to every panel,
\[\widetilde H_0=D_LH_0D_R,\qquad\widetilde H_w=D_LH_wD_R,\]
then $P_w=\widetilde H_0^{-1}\widetilde H_w$ satisfies
\begin{equation}P_w=(RD_R)^{-1}T_w(RD_R).\tag{348}\end{equation}
Characteristic polynomials, eigenvalues, traces of powers, and traces of all products $P_{w_1}\cdots P_{w_t}$ equal the corresponding minimal-linear-realization invariants. Eigenvectors and entries of $T_w$ in the original hidden basis are not thereby identified.
\begin{proof}
Direct cancellation gives $(D_LORD_R)^{-1}D_LOT_wRD_R=(RD_R)^{-1}T_w(RD_R)$. The same similarity applies to every word, hence to products and their traces.
\end{proof}
\statement{Proposition}{LVI.7}{Quantitative finite-measurement stability}
Let $A=\widetilde H_0$, $B=\widetilde H_w$, $P=A^{-1}B$, and $\sigma=\sigma_{\min}(A)>0$. If $\|E\|_2,\|F\|_2\le\eta<\sigma$, then
\begin{equation}\|(A+E)^{-1}(B+F)-P\|_2\le\frac{\eta(1+\|P\|_2)}{\sigma-\eta}.\tag{349}\end{equation}
For diagonalizable $P=V\Lambda V^{-1}$, Bauer--Fike puts each perturbed eigenvalue within $\kappa_2(V)$ times this bound of some eigenvalue of $P$. Without diagonalizability, use a pseudospectral inclusion instead.
\begin{proof}
$\widehat P-P=(A+E)^{-1}(F-EP)$, and $\|(A+E)^{-1}\|_2\le(\sigma-\eta)^{-1}$ by the smallest-singular-value bound.
\end{proof}
\statement{Remark}{LVI.8}{Common calibration is necessary}
If either calibration changes between base and shifted panels, failure occurs even at rank one: true $H_0=H_a=[1]$ but measured $\widetilde H_0=[1]$, $\widetilde H_a=[2]$ give false $P_a=[2]$ instead of $T_a=[1]$. ``Uncalibrated'' means unknown but stable calibration, not arbitrary drift. The result is related to classical Kalman--Ho spectral realization and should not be declared wholly new physics without a literature review.

\note{Computational commitments.} \path{verify_v52_exact.py} uses exact rational cycle checks, symbolic SymPy pencil calculations, and finite LP enumeration for $n=42$. It replaces neither the proofs nor formal verification in Lean. Independent worldwide priority and a continuous gravitational interpretation of these operator invariants remain open.
\setcounter{equation}{349}


% 16_tail_berger_symplectic.tex
\chapter{The Necessary Cost of the Tail Optimum, a Quantitative Gap, and Exact Dual Certificates}
\note{Continuity and status.} This chapter continues the correction of L.1 in LVI; it does not restore the refuted disjoint-block-clock formula. Proofs concern autonomous finite stochastic models with normalized preparations and effects in $[0,1]$. Worldwide priority is not established; these are exact results within this research program.
\section{Optimal state count for the lower bound $b/(r+1)$}
Let $n=\prod_{i=1}^rq_i$, $q_i=p_i^{\alpha_i}$ with distinct primes, $r=\omega(n)$, and $f_{n,b}(k)=b\one_{n\mid k}$, $0<b\le1$. Recall the homogeneous stochastic minimax definition of $E_m$. By LVI.2, infinite-horizon optima are attained by finite LP models with nonnegative periodic contributions $h_d$, $d\mid n$, and $\sum_{d\in D}d\le m$. For $m<n$, all periods are proper divisors. Set
\[k_i=n/p_i,\quad I(d)=\{i:d\mid k_i\},\quad c(d)=|I(d)|,\quad M_*(n)=\sum_{i=1}^rn/q_i.\]
Every proper divisor has $c(d)\ge1$.
\statement{Lemma}{LVII.1}{The witness's exact nonnegative remainder}
For $g(k)=\sum_{d\in D}h_d(k)$,
\begin{equation}\sum_{i=1}^rg(k_i)-g(0)=\sum_{d\in D}\left[(c(d)-1)h_d(0)+\sum_{i\notin I(d)}h_d(k_i)\right]\ge0.\tag{350}\end{equation}
Equality requires every component with $h_d(0)>0$ to have $c(d)=1$, and $h_d(k_i)=0$ for $i\notin I(d)$.
\begin{proof}
For $i\in I(d)$, periodicity gives $h_d(k_i)=h_d(0)$. Sum over $i$ and subtract $h_d(0)$. Every remaining term is nonnegative.
\end{proof}
\statement{Theorem}{LVII.2}{The exact threshold for attaining the tail level}
For $m<n$,
\begin{equation}E_m(f_{n,b})=\frac b{r+1}\quad\Longleftrightarrow\quad m\ge M_*(n)=\sum_{i=1}^r\frac n{q_i}.\tag{351}\end{equation}
If $\sum_iq_i^{-1}\ge1$, no model with $m<n$ attains that level. If the sum is below one, $M_*(n)$ is necessary as well as sufficient.
\begin{proof}
Sufficiency uses $r$ cycles of lengths $n/q_i$, masses $1/r$, and phase-zero effects $rb/(r+1)$. Their output is $b/(r+1)\sum_i\one_{n/q_i\mid k}$ and has exact error $b/(r+1)$, because distinct lengths have lcm $n$. The matching lower bound is LVI.3.

For necessity, the finite LP reduction attains its minimum. At $\varepsilon=b/(r+1)$, accuracy and Lemma LVII.1 force $g(0)=r\varepsilon$, $g(k_i)=\varepsilon>0$ for all $i$, and zero remainder. Thus every positive-origin component has $c(d)=1$, and all off-$I(d)$ contributions to $g(k_i)$ vanish. For each $i$, a distinct period $d_i$ must have $I(d_i)=\{i\}$ and positive origin value. For $j\ne i$, $d_i\nmid n/p_j$ forces full valuation $v_{p_j}(d_i)=\alpha_j$. Hence $d_i$ is a multiple of $n/q_i$ and at least that large. Summing the distinct cycle costs gives $m\ge M_*(n)$.
\end{proof}
\section{An explicit gap below the required memory budget}
\statement{Theorem}{LVII.3}{The quantitative cost of a missing exclusive clock}
If $m<\min\{n,M_*(n)\}$, then
\begin{equation}E_m(f_{n,b})\ge\frac{2b}{2r+1}>\frac b{r+1}.\tag{352}\end{equation}
\begin{proof}
For a feasible response with error $\varepsilon$, put $S=\sum_ig(k_i)-g(0)$. Then $0\le S\le(r+1)\varepsilon-b$. If every $i$ had a used period with $I(d)=\{i\}$, the budget would be at least $M_*$. Thus some $i_0$ lacks one. Each contribution to $g(k_{i_0})$ is either from a component with $c(d)\ge2$ and $i_0\in I(d)$, bounded by its $(c(d)-1)h_d(0)$ contribution to $S$, or is explicitly an off-$I(d)$ term of $S$. Hence $g(k_{i_0})\le S$. Also
\[g(k_{i_0})=\sum_ig(k_i)-\sum_{j\ne i_0}g(k_j)\ge(b-\varepsilon)-(r-1)\varepsilon=b-r\varepsilon.\]
Combine to get $\varepsilon\ge2b/(2r+1)$; strict comparison follows from $2(r+1)>2r+1$.
\end{proof}
\statement{Corollary}{LVII.4}{The complete two-prime-divisor problem}
For $n=p^\alpha q^\beta$, distinct primes $p,q$, write $Q=p^\alpha$, $R=q^\beta$. Then $M_*=Q+R<n$, and
\[E_m(f_{n,b})=\begin{cases}b/2,&1\le m<Q+R,\\b/3,&Q+R\le m<n,\\0,&m\ge n.\end{cases}\]
\begin{proof}
The distinct prime powers satisfy $QR>Q+R$, including $(Q,R)=(2,3)$ where $6>5$. The middle region follows from LVII.2. The first-improvement threshold $K(n)=Q+R$ and error $b/2$ below it use LVI.3 together with the earlier threshold theorem (LVI and historical revision v37). One $n$-cycle gives exactness at $m\ge n$.
\end{proof}
\section{Rational dual certificates: the complete period-42 staircase}
The proper divisors are $\{1,2,3,6,7,14,21\}$. Of $2^7-1=127$ nonempty subsets, exactly 104 have total cost below 42. For each, the program constructs feasible rational primal and dual vectors with equal objective. The verifier rechecks
\[A_Dx\le b_D,\quad C_Dx=1,\quad x\ge0,\quad A_D^Ty+C_D^Tz\le c,\quad y\le0,\]
and $c^Tx=b_D^Ty+z$. The previously approximate LP table therefore has an exact reproducible computer certificate:
\[E_m(f_{42,1})=\begin{cases}1/2,&1\le m\le11,\\2/5,&m=12,\\1/3,&13\le m\le40,\\1/4,&m=41,\\0,&m\ge42.\end{cases}\]
The files \path{checks/LP42_EXACT_RATIONAL_CERTIFICATES.json} and \path{checks/verify_v53_rational.py} form one certificate. Rational arithmetic removes rounding from certificate verification, but the program and finite-LP reduction remain outside Lean, and independent priority is unconfirmed.
\section{A conditional interface to information geometry and the status of 42 claims}
The pencil reconstruction of LVI.6 remains valid with the same unknown diagonal calibration on both panels. General drift destroys that invariance. Periodic approximation errors and spectral numerical invariants cannot automatically become a measurable gravitational source. The v52 register marks the full disjoint-block staircase as refuted; subsequent catalogs of 42 proposed laws must retain that status. Testable identities and proved barriers are neither worldwide discoveries nor solutions of Clay open problems.

\chapter{Berger Spectral Reconstruction, Exact Ricci Clocks, and the Distinction between Helicities}
\note{Status and priority.} This chapter derives exact consequences of the $SU(2)$ left-invariant frame and a special Ricci flow. The complete Berger curl spectrum and its eta invariant are already known (Gibbons; Capoferri--Vassiliev and collaborators, 2026), and the round-sphere zeta determinant was previously published. Consequences new to this program are separated from those known results and unresolved Navier--Stokes or gravitational questions.
\section{Direct derivation of three curl eigenmodes on anisotropic $S^3$}
Normalize left-invariant forms by $d\sigma_1=2\sigma_2\wedge\sigma_3$ cyclically, without doubled index summation. Let
\begin{equation}g_{a,\tau}=a^2(\sigma_1^2+\sigma_2^2)+a^2\tau^2\sigma_3^2,\qquad a,\tau>0,\tag{353}\end{equation}
and use the positively oriented orthonormal coframe $\theta_1=a\sigma_1$, $\theta_2=a\sigma_2$, $\theta_3=a\tau\sigma_3$. Then
\begin{equation}d\theta_1=\frac2{a\tau}\theta_2\wedge\theta_3,\quad d\theta_2=\frac2{a\tau}\theta_3\wedge\theta_1,\quad d\theta_3=\frac{2\tau}a\theta_1\wedge\theta_2.\tag{354}\end{equation}
\statement{Theorem}{LVIII.1}{Two-frequency inversion of Berger geometry}
The operator $\operatorname{curl}_{a,\tau}=*_{g_{a,\tau}}d$ preserves the left-invariant one-form space $\mathcal L$, with
\begin{equation}\operatorname{curl}_{a,\tau}|_{\mathcal L}=\diag(h,h,v),\qquad h=\frac2{a\tau},\quad v=\frac{2\tau}a.\tag{355}\end{equation}
From labeled horizontal double mode $h$ and vertical mode $v$,
\begin{equation}\tau^2=v/h,\qquad a^2=4/(hv),\qquad\operatorname{Vol}(g_{a,\tau})=\frac{16\pi^2}{h^2v}.\tag{356}\end{equation}
Equality $h=v$ identifies roundness within this family. For unequal frequencies, recognition of these modes from the full unlabeled spectrum is not asserted.
\begin{proof}
Apply the Hodge star to (354). Volume is $2\pi^2a^3\tau$, since $\int\sigma_1\wedge\sigma_2\wedge\sigma_3=2\pi^2$. Substitute $a=2/\sqrt{hv}$ and $\tau=\sqrt{v/h}$.
\end{proof}
\statement{Theorem}{LVIII.2}{Curvature and anisotropy directly from frequencies}
In this frame,
\begin{align}\operatorname{Ric}_{11}=\operatorname{Ric}_{22}&=hv-\tfrac12v^2,\tag{357}\\
\operatorname{Ric}_{33}&=\tfrac12v^2,\tag{358}\\
R&=2hv-\tfrac12v^2,\tag{359}\\
K_{12}&=hv-\tfrac34v^2,\qquad K_{13}=K_{23}=\tfrac14v^2.\tag{360}\end{align}
Thus the three left-invariant modes determine all sectional curvatures in this family, not arbitrary three-manifolds.
\begin{proof}
Koszul gives $K_{12}=(4-3\tau^2)/a^2$, $K_{13}=K_{23}=\tau^2/a^2$. Substitute $hv=4/a^2$, $v^2=4\tau^2/a^2$. The attached script also computes Ricci directly from structure constants without these final formulas.
\end{proof}
\section{An exact spectral-clock law along Ricci flow}
The symmetric-family flow satisfies
\begin{equation}\frac d{dt}a^2=-8+4\tau^2,\qquad\frac d{dt}(a^2\tau^2)=-4\tau^4,\qquad\frac{d\tau^2}{dt}=\frac{8\tau^2(1-\tau^2)}{a^2}.\tag{361}\end{equation}
Set $s(t)=\int_0^t8/a(u)^2\,du$ and $q=h/v-1=\tau^{-2}-1$.
\statement{Theorem}{LVIII.3}{Exact exponential decay of the frequency ratio}
Before metric degeneration,
\begin{equation}\frac{dq}{ds}=-q,\qquad q(s)=q(0)e^{-s},\qquad\frac d{dt}\left(\frac hv\right)=-2hv\left(\frac hv-1\right).\tag{362}\end{equation}
One measured pair gives anisotropy; ratios on two slices give accumulated internal time. At $q(0)=0$, the round branch does not define a clock this way.
\begin{proof}
Insert $q=\tau^{-2}-1$ into (361), obtaining $\dot q=-(8/a^2)q$. Use $hv=4/a^2$. Nothing automatically applies to Lorentzian Einstein dynamics or cosmological time.
\end{proof}
\statement{Corollary}{LVIII.4}{Closed Ricci-flow solution in spectral time}
For $q_0=\tau_0^{-2}-1>-1$ and $a_0=a(0)$,
\begin{equation}\tau^2(s)=\frac{e^s}{e^s+q_0},\qquad a^2(s)=a_0^2e^{-s}\sqrt{\frac{e^s+q_0}{1+q_0}},\qquad s\ge0.\tag{363}\end{equation}
Substitution in (355) gives all three frequencies in closed form. The flow parameter $t$ is recovered by the strictly positive quadrature $dt/ds=a^2(s)/8$.
\begin{proof}
Use $q=q_0e^{-s}$ and integrate $d\log a^2/ds=-1+\tau^2(s)/2$.
\end{proof}
\section{Known spectral asymmetry and a zeta-determinant check}
The full Berger curl spectrum, including negative levels, and the following eta invariant were published by Capoferri et al., J. London Math. Soc. (2026), Appendix D:
\begin{equation}\eta_{\rm curl}(0)=\tfrac23(\tau^2-1)^2.\tag{364}\end{equation}
For comparison with that literature, the complete published spectrum is reproduced without a novelty claim. For $k=2,3,\ldots$, the first two positive series are
\begin{equation}\lambda_k^{\rm(I)}=\frac{k}{a\tau}\quad[\text{multiplicity }2k-2],\qquad
\lambda_k^{\rm(II)}=\frac{k+2(\tau^2-1)}{a\tau}\quad\begin{cases}1,&k=2,\\2k-2,&k\ge3.\end{cases}\tag{365}\end{equation}
For $1\le j\le\lfloor k/2\rfloor$, the remaining series are
\begin{equation}\lambda_{k,j}^{(\pm)}=\frac{\tau\pm\sqrt{\tau^2+k(k+2)+(\tau^{-2}-1)(k-2j)^2}}a.\tag{366}\end{equation}
Multiplicity is $2k+2$ for odd $k$, or even $k$ with $j<k/2$; for even $k$ and $j=k/2$, it is $k+1$. These restate the cited Theorem D.1. At $k=2$, the first two series give $h$ with multiplicity two and $v$ with multiplicity one. They may coincide with other full-spectrum levels; reconstruction requires identifying this geometric sector.

Hitchin's formula on the branch without the relevant crossing is $\eta_D(0)=-(\tau^2-1)^2/6$, hence $\eta_{\rm curl}=-4\eta_D$. This compares known formulas; it is neither a new index nor a new proof of Hitchin's result.
\statement{Corollary}{LVIII.5}{Curl spectral asymmetry, Ricci anisotropy, and exact decay}
With the specified orientation,
\begin{equation}a^4\|\operatorname{Ric}^\circ\|^2=16\eta_{\rm curl}(0)=\frac{32}3(\tau^2-1)^2,\qquad\frac{d\eta_{\rm curl}}{ds}=-2\tau^2\eta_{\rm curl}\le0.\tag{367}\end{equation}
In particular, $\eta_{\rm curl}(s)=\frac23q_0^2e^{-2s}/(1+q_0e^{-s})^2$. This follows exactly from the published eta formula and Ricci flow, not an independent asymmetry theorem for arbitrary manifolds.
\begin{proof}
The Ricci eigenvalues give squared traceless norm $32(\tau^2-1)^2/(3a^4)$. Differentiate (364) using $d\tau^2/ds=\tau^2(1-\tau^2)$.
\end{proof}
\statement{Proposition}{LVIII.6}{The corrected zeta determinant of $S^3$}
For the unit round scalar Laplacian with constant mode removed,
\begin{equation}\zeta'_\Delta(0)=-\log\pi-\frac{\zeta_R(3)}{2\pi^2},\qquad\det'\Delta=\pi\exp\left[\frac{\zeta_R(3)}{2\pi^2}\right].\tag{368}\end{equation}
This agrees with the published sphere-determinant result in Advances in Difference Equations (2013), formula (4.18). The submitted exponent had the opposite sign on $\zeta(3)$.
\begin{proof}
Start from $\zeta_\Delta(s)=\sum_{k=2}^\infty k^2(k^2-1)^{-s}$, continue via expansion of $(1-k^{-2})^{-s}$, and use $\zeta'_R(0)=-\tfrac12\log(2\pi)$ and $\zeta'_R(-2)=-\zeta_R(3)/(4\pi^2)$. The known value is independently sourced; the computational module checks 30 digits.
\end{proof}
\section{Critical corrections to hydrodynamic laws 63--64}
On a fixed closed Riemannian three-sphere, let $\omega=\operatorname{curl}u$, and implement viscosity with the positive Hodge Laplacian, equal to $\operatorname{curl}^2$ on divergence-free one-forms. Then
\begin{equation}\partial_tu+\nu\operatorname{curl}^2u=P(u\times\omega),\qquad P\nabla p=0.\tag{369}\end{equation}
Distinguish
\begin{equation}H_{\rm kin}=\int u\cdot\operatorname{curl}u\,dV,\qquad H_{\rm pot}=\int(\operatorname{curl}^{-1}u)\cdot u\,dV.\tag{370}\end{equation}
\statement{Proposition}{LVIII.7}{The actually conserved Euler invariants}
For smooth inviscid solutions, $E=\tfrac12\|u\|_2^2$ and $H_{\rm kin}$ are constant because $\langle u,u\times\omega\rangle=0$ and $\langle\omega,u\times\omega\rangle=0$. No such cancellation holds for $H_{\rm pot}$:
\[\frac d{dt}H_{\rm pot}=2\langle\operatorname{curl}^{-1}u,P(u\times\omega)\rangle.\]
The geometrically valid inequality $E\ge|H_{\rm pot}|/R$ on round $S_R^3$ therefore does not automatically give a time-invariant Euler-energy lower bound: potential helicity is not the usual conserved velocity invariant.
\statement{Proposition}{LVIII.8}{Exact single-chirality eigenvalue shells}
If $\operatorname{curl}u_0=\mu u_0$, $\mu\ne0$, then $u\times\omega=0$ and $u(t)=e^{-\nu\mu^2t}u_0$ solves (369) without nonlinear mode transfer. This holds for one eigenvalue, not an arbitrary mixture of opposite chiralities in one absolute shell. For $u=u_++u_-$ with $\operatorname{curl}u_\pm=\pm\lambda u_\pm$, $u\times\omega=-2\lambda u_+\times u_-$ is generally nonzero.
\statement{Proposition}{LVIII.9}{A Beltrami triad: the correct quadratic invariants}
Let $\Gamma=\int(\Phi_1\times\Phi_2)\cdot\Phi_3\,dV$ and $\operatorname{curl}\Phi_i=\mu_i\Phi_i$. A closed isolated real-amplitude triad with this orientation obeys
\begin{equation}\dot c_1=\Gamma(\mu_3-\mu_2)c_2c_3,\quad\dot c_2=\Gamma(\mu_1-\mu_3)c_3c_1,\quad\dot c_3=\Gamma(\mu_2-\mu_1)c_1c_2.\tag{371}\end{equation}
It preserves $\sum_ic_i^2$ and $\sum_i\mu_ic_i^2$, but not generally $\sum_ic_i^2/\mu_i$. Physical kinetic helicity is thus distinct from vector-potential helicity.
\begin{proof}
Real symmetrization of $u\times\omega$ gives coefficient $(\mu_2-\mu_1)\Gamma$ for $c_1c_2$ in the third equation, and cyclic counterparts. Multiplication by $c_i$ or $\mu_ic_i$ makes the weight sums vanish; multiplication by $c_i/\mu_i$ generally does not.
\end{proof}
\section{What remains unproved, and why this is not a Clay solution}
The complete curl spectrum for arbitrary anisotropic Berger spheres is not derived here; it is known and much richer than three left-invariant modes. No general three-dimensional hydrodynamic a priori $H^1$ estimate, absence of mixed-chirality nonlinear cascades, global regularity of arbitrary Navier--Stokes solutions, or spectral derivation of gravitational constants has been proved. These are precisely restricted mathematical developments, not a Millennium-Problem solution.

\chapter{Berger--Hopf Symplectic Geometry and a Spectral Ricci-Flow Integral}
\note{Subject and status.} We construct a specified phase space $T^*SU(2)$ rather than postulating one smooth manifold of all operational specifications. Symplectic reduction and Ricci flow are classical; we express their joint description through two labeled curl frequencies in forms useful to this program. Independent priority of that combination is unestablished. Normalizations match LV and LVIII.
\section{The cotangent bundle and an explicit Poisson structure}
On $G=SU(2)$, take $[X_i,X_j]=2\epsilon_{ijk}X_k$ and $d\sigma_1=-2\sigma_2\wedge\sigma_3$ cyclically. In left trivialization $T^*G\simeq G\times\mathfrak{su}(2)^*$, put
\begin{equation}\Theta=\sum_{i=1}^3p_i\sigma_i,\qquad\Omega_{\rm can}=-d\Theta=\sum_i\sigma_i\wedge dp_i+2(p_1\sigma_2\wedge\sigma_3+p_2\sigma_3\wedge\sigma_1+p_3\sigma_1\wedge\sigma_2).\tag{372}\end{equation}
The sign gives $\{q,p\}=1$ in canonical local coordinates.
\statement{Proposition}{LIX.1}{The exact momentum-bracket structure}
For momentum-only functions,
\begin{equation}\{p_i,p_j\}=-2\epsilon_{ijk}p_k,\qquad\{F,G\}(p)=-2p\cdot(\nabla_pF\times\nabla_pG).\tag{373}\end{equation}
The opposite symplectic or group convention reverses the sign. This is ordinary Lie--Poisson reduction on $\mathfrak{su}(2)^*$, not a universal bracket for arbitrary memory models.
\begin{proof}
For $J_X(q,p)=p(X(q))$, canonical coordinates give $\{J_X,J_Y\}=-J_{[X,Y]}$. Substitute the frame fields and extend by Leibniz and the chain rule.
\end{proof}
\section{Geodesic precession and Hopf magnetic reduction}
For $g=A(\sigma_1^2+\sigma_2^2)+B\sigma_3^2$, $A=a^2>0$, $B=a^2\tau^2>0$,
\begin{equation}H_{\rm geo}(p)=\frac{p_1^2+p_2^2}{2A}+\frac{p_3^2}{2B}.\tag{374}\end{equation}
\statement{Theorem}{LIX.2}{Exact momentum precession and spectral clocks}
For fixed or smoothly time-dependent $A,B$, both $p_3$ and $p_1^2+p_2^2$ are conserved, with
\begin{align}\dot p_1&=-\Omega_{\rm prec}p_2,\qquad\dot p_2=\Omega_{\rm prec}p_1,\qquad\dot p_3=0,\tag{375}\\
\Omega_{\rm prec}(t)&=2p_3[A(t)^{-1}-B(t)^{-1}]=\frac{p_3}2\omega_\perp(t)[\omega_\parallel(t)-\omega_\perp(t)].\tag{376}\end{align}
Here $\omega_\perp=2/\sqrt B$, $\omega_\parallel=2\sqrt B/A$ are the absolute left-invariant curl frequencies. Thus $P=p_1+ip_2$ obeys $P(t)=P(0)\exp[i\int_0^t\Omega_{\rm prec}(u)du]$.
\begin{proof}
Insert (374) into (373), giving $\dot p_1=2p_2p_3(B^{-1}-A^{-1})$, $\dot p_2=2p_1p_3(A^{-1}-B^{-1})$, and $\dot p_3=0$. The spectral form uses $A^{-1}=\omega_\perp\omega_\parallel/4$ and $B^{-1}=\omega_\perp^2/4$.
\end{proof}
\statement{Proposition}{LIX.3}{Hopf magnetic reduction and Weyl volume}
Let the right $U(1)$ action be generated by $X_3$, with Hopf connection $\alpha=\sigma_3$. At momentum $p_3=\mu$, the reduced space is diffeomorphic to $T^*S^2$, with
\begin{equation}\Omega_\mu=\Omega_{T^*S^2}-\mu\pi^*F,\qquad F=d\alpha=-2\sigma_1\wedge\sigma_2.\tag{377}\end{equation}
$F$ is basic; its sign depends on fiber orientation. The six-dimensional symplectic volume of $\{H_{\rm geo}\le E\}$ is
\begin{equation}V_E=\frac{4\pi}3(2E)^{3/2}\operatorname{Vol}(S^3,g)=\frac{8\pi^3}3(2E)^{3/2}A\sqrt B.\tag{378}\end{equation}
It is constant at fixed $E$ under volume-normalized homogeneous Ricci flow, but not under ordinary flow.
\begin{proof}
Decompose the covector as $\eta+\mu\alpha$, so $\Theta=\Theta_{T^*S^2}+\mu\alpha$. Restrict $-d\Theta$ to the momentum level and quotient by $U(1)$. The momentum ellipsoid has volume $(4\pi/3)(2E)^{3/2}A\sqrt B$ at each base point; integrate against invariant volume $2\pi^2$. Normalized flow preserves Riemannian volume by definition.
\end{proof}
\section{A corrected Ricci-flow first integral in frequency space}
\statement{Theorem}{LIX.4}{The Berger--Ricci spectral integral}
For the ordinary homogeneous flow
\begin{equation}\dot A=-8+4B/A,\qquad\dot B=-4B^2/A^2,\tag{379}\end{equation}
the signed quantity
\begin{equation}I_{\rm BR}=\omega_\perp^3(\omega_\perp-\omega_\parallel)=\frac{16(A-B)}{AB^2},\qquad\frac{dI_{\rm BR}}{dt}=0\tag{380}\end{equation}
is conserved on every smooth existence interval. It vanishes on $A=B$; off that branch an equivalent integral is $B/\sqrt{|1-B/A|}=4/\sqrt{|I_{\rm BR}|}$. For precession with $p_3\ne0$,
\begin{equation}\frac{\Omega_{\rm prec}(t)\omega_\perp(t)^2}{p_3}=-\tfrac12I_{\rm BR}=\text{constant}.\tag{381}\end{equation}
The last relation assumes Hamiltonian motion with the same time-dependent $A,B$; it does not derive physical Einstein--matter dynamics from Ricci flow.
\begin{proof}
Substitute the frequency definitions to get the rational expression. Direct differentiation gives
\[\frac d{dt}\frac{A-B}{AB^2}=\frac{(\dot A-\dot B)AB^2-(A-B)(\dot AB^2+2AB\dot B)}{A^2B^4}=0\]
using (379). Rearrangement gives the equivalent integral, and (376) gives (381). Volume normalization adds a conformal term, so (380) is generally not conserved there.
\end{proof}
\statement{Remark}{LIX.5}{Required correction to preliminary meta-law 11}
The proposed $(\omega_\perp-\omega_\parallel)/(\omega_\perp^2\omega_\parallel)=(A-B)/4$ is not a first integral: its derivative is $-2+B/A+(B/A)^2$, generally nonzero. Subsequent models should use (380), not that simpler combination.
\section{Spectral dispersion of two helicities: exact applicability}
\statement{Proposition}{LIX.6}{Nonnegative spectral dispersion}
Let a smooth divergence-free field on fixed round $S_R^3$ contain only positive chiral modes, $\operatorname{curl}\Phi_j=\lambda_j\Phi_j$, $\lambda_j>0$, and let $a_j=|c_j|^2$. Then
\begin{equation}H_{\rm pot}H_{\rm kin}-(2E)^2=\sum_{i<j}a_ia_j\frac{(\lambda_i-\lambda_j)^2}{\lambda_i\lambda_j}\ge0.\tag{382}\end{equation}
For a nonzero field, equality holds exactly when all active modes share one curl eigenvalue. This is algebraic; the difference is not declared conserved under full Euler dynamics.
\begin{proof}
Use $2E=\sum_ja_j$, $H_{\rm pot}=\sum_ja_j/\lambda_j$, and $H_{\rm kin}=\sum_j\lambda_ja_j$. Pairing cross terms gives $\sum_{i<j}a_ia_j(\lambda_i/\lambda_j+\lambda_j/\lambda_i-2)$.
\end{proof}
\note{Limits of meta-generalization.} Classical reduction of $T^*SU(2)$ does not, without additional structure, make GKLS-generator spaces, nonnegative-realization polytopes, and AI Grassmannians one canonical symplectic manifold. Such a claim needs an object category, smooth structure, morphisms, compatible forms, and a proof of compatibility. New notation alone is not a theorem.

\note{Published foundations.} Cotangent reduction and magnetic terms are discussed by J. E. Marsden and T. S. Ratiu, \emph{Mechanical Systems: Symmetries and Reduction}, Springer (2012), \url{https://doi.org/10.1007/978-1-4614-1806-1_59}. Berger metrics and homogeneous Ricci flow belong to classical Hamilton--Milnor theory. Formula (380) is checked as a consequence of a known homogeneous ODE, not claimed as a solution of an open geometric-analysis problem.
\setcounter{equation}{382}


% 17_phase_three_primes.tex
\chapter{Exact Topological Phase Boundaries of Competing GKLS Resets}
\note{Status.} This is a rigorous theorem for an explicitly specified finite-dimensional GKLS family, not a new universal law of nature. Mixed stationary-state topology and the distinction between purity and damping gaps were discussed previously (Budich--Diehl 2015; Budich--Zoller--Diehl 2015). This audit combines a multiple-reset parameterization, an exact superoperator perturbation boundary, and a finite-sample certificate in one model. Independent priority for the combined formulas is unestablished.
\section{An explicit physical model and complete spectral decomposition}
Let $X=S^2$ have fixed orientation, and let smooth $n:X\to S^2$ have nonzero degree $q$. For $r\in(0,1]$, $\gamma>0$, $\kappa_j\ge0$, and constant $b_j\in\R^3$ with $\|b_j\|\le1$, set
\[\rho_n(x)=\tfrac12(I+rn(x)\cdot\sigma),\qquad\sigma_j=\tfrac12(I+b_j\cdot\sigma),\]
\[\mathcal L_x(Z)=\gamma[\Tr(Z)\rho_n(x)-Z]+\sum_{j=1}^m\kappa_j[\Tr(Z)\sigma_j-Z].\]
Write $A=\gamma r>0$, $w=\sum_j\kappa_jb_j$, $B=\|w\|$, and $\Gamma=\gamma+\sum_j\kappa_j>0$. Then $\mathcal L_x(Z)=\Gamma[\Tr(Z)\bar\rho_x-Z]$, with $\bar\rho_x=(I+v_x\cdot\sigma)/2$ and $v_x=[An(x)+w]/\Gamma$.
\statement{Lemma}{LX.1}{Complete positivity and exact relaxation}
The generator has standard Gorini--Kossakowski--Sudarshan--Lindblad form, and
\[e^{t\mathcal L_x}(Z)=e^{-\Gamma t}Z+(1-e^{-\Gamma t})\Tr(Z)\bar\rho_x.\]
Its unique stationary state is $\bar\rho_x$, its operator-space spectrum is $\{0,-\Gamma,-\Gamma,-\Gamma\}$, and $\Delta_{\rm Liouv}=\Gamma$.
\begin{proof}
If $\bar\rho_x=\sum_{a=1}^2p_a|a\rangle\langle a|$, the jumps $L_{ab}=\sqrt{\Gamma p_a}|a\rangle\langle b|$ give $\sum L_{ab}ZL_{ab}^\dagger=\Gamma\Tr(Z)\bar\rho_x$ and $\sum L_{ab}^\dagger L_{ab}=\Gamma I$. The generator acts as $-\Gamma$ on traceless operators and as zero on $\bar\rho_x$. Decompose $Z=\Tr(Z)\bar\rho_x+[Z-\Tr(Z)\bar\rho_x]$ to exponentiate.
\end{proof}
\statement{Theorem}{LX.2}{A multichannel topological boundary with a nonclosing dissipative gap}
The minimum stationary-density gap is exactly
\[g_{\rm pur}=\inf_x[\lambda_+(\bar\rho_x)-\lambda_-(\bar\rho_x)]=\frac{|A-B|}{\Gamma}.\]
For $A\ne B$, the globally defined normalized field satisfies
\[\deg\widehat v=\begin{cases}q,&A>B,\\0,&0<A<B.\end{cases}\]
At $A=B>0$, some $x_*$ has $\bar\rho_{x_*}=I/2$, so its leading projector is undefined. The relaxation gap remains positive for all parameters. The leading eigenbundle's first Chern number is $\pm\deg\widehat v$, according to Berry-curvature sign convention.
\begin{proof}
A qubit density gap is its Bloch-vector length. Nonzero degree makes $n$ surjective: omission of a point would factor it through a contractible punctured sphere. Therefore $\min_x\|An(x)+w\|=|A-\|w\||$. For $A>B$, normalized $An+tw$ gives a homotopy to $n$. For $A<B$, normalized $(1-t)An+w$ gives a homotopy to the constant $w/B$. At equality, choose $n(x_*)=-w/B$. The Liouvillian spectrum follows from the lemma.
\end{proof}
\statement{Corollary}{LX.3}{Signed counting of local critical points}
If $A=B>0$ and $-w/B$ is a regular value, the critical points are finite and $\sum_{x_*}\operatorname{ind}_{x_*}(n)=q$. There are at least $|q|$ of them.
\begin{proof}
This is the oriented regular-preimage degree formula; compactness and regularity give finiteness.
\end{proof}
\statement{Corollary}{LX.4}{Topological control without changing relaxation rate}
For two resets with fixed positive rates $\kappa_1,\kappa_2$ and Bloch moduli $r_1,r_2$, let $\vartheta$ be the angle between their vectors and put $u=\kappa_1r_1$, $v=\kappa_2r_2$. Then $B(\vartheta)^2=u^2+v^2+2uv\cos\vartheta$, while $\Gamma=\gamma+\kappa_1+\kappa_2$ is constant. Rotation reaches the transition exactly when $|u-v|\le A\le u+v$. In the strict interior, the unique critical angle is
\[\vartheta_* =\arccos\frac{A^2-u^2-v^2}{2uv}.\]
Thus relative reset direction alone causes $q\leftrightarrow0$ at unchanged Liouvillian gap. At the endpoints the critical angle is 0 or $\pi$.
\begin{proof}
Use the cosine rule and its monotonicity on $[0,\pi]$. The range of $B$ is $[|u-v|,u+v]$, and rotating Bloch vectors does not change rates.
\end{proof}
\section{Robustness under arbitrary smooth GKLS perturbations}
\statement{Theorem}{LX.5}{A precise superoperator certificate}
For $A\ne B$, let $\widetilde{\mathcal L}_x=\mathcal L_x+\Delta_x$ be smooth trace-preserving GKLS generators with $\varepsilon=\sup_x\|\Delta_x\|_{1\to1}<|A-B|$. They have unique stationary states $\widetilde\rho_x$, and
\[\sup_x\|\widetilde\rho_x-\bar\rho_x\|_1\le\varepsilon/\Gamma,\qquad\inf_x[\lambda_+(\widetilde\rho_x)-\lambda_-(\widetilde\rho_x)]\ge\frac{|A-B|-\varepsilon}{\Gamma}>0.\]
Hence normalized perturbed and original fields have equal degree. The norm is induced by trace norm.
\begin{proof}
On traceless operators, $\mathcal L_x=-\Gamma I$. Since $\varepsilon<|A-B|\le\Gamma$, Neumann inversion makes $-\Gamma I+\Delta_x$ invertible. CPTP dynamics ensures existence of a stationary state; invertibility ensures uniqueness. Stationarity gives $\Gamma(\widetilde\rho_x-\bar\rho_x)=\Delta_x(\widetilde\rho_x)$, yielding the norm bound. Qubit trace distance equals Euclidean Bloch difference, so the straight vector homotopy avoids zero.
\end{proof}
\section{Statistically reproducible recovery of integer charge}
\statement{Theorem}{LX.6}{A sufficient finite-sample certificate}
Let $L_v$ be a known Lipschitz constant for the stationary Bloch field. Triangulate $X$ so every point of a simplex is within $h$ of each vertex. At $M$ vertices, take $N$ independent measurements of each Pauli operator. Define
\[g=|A-B|/\Gamma,\qquad d_N=\sqrt{6\log(6M/\alpha)/N}.\]
If $L_vh+d_N<g$, then with probability at least $1-\alpha$, the piecewise-linear Bloch interpolant is everywhere nonzero and homotopic to the true field, recovering degree and the corresponding Chern number exactly. A sufficient per-component, per-vertex budget is
\[N>\frac{6\log(6M/\alpha)}{(g-L_vh)^2},\qquad g>L_vh.\]
\begin{proof}
Hoeffding gives $\Pr(|\widehat v_j-v_j|>t)\le2e^{-Nt^2/2}$. Union over $3M$ components at $t=\sqrt{2\log(6M/\alpha)/N}$ bounds simultaneous Euclidean node error by $d_N$. Convex barycentric interpolation adds at most $L_vh$. Error smaller than $\inf\|v\|=g$ prevents zero along the straight homotopy, preserving degree.
\end{proof}
\statement{Remark}{LX.7}{Physical scope}
A positive Liouvillian gap gives exponential forgetting of the initial state, but alone does not protect leading-projector topology. As $A\to B$, the stated sufficient tomography budget diverges as $|A-B|^{-2}$ at fixed mesh. This is not a global lower bound on optimal complexity of every quantum measurement protocol.
\section{Counterchecks of proposed laws 65--114}
\begin{enumerate}
\item Law 74 is confirmed and generalized above: the multichannel threshold depends on the norm of the vector sum of resets, not the sum of norms.
\item Law 85: a bipartite $k$-regular graph has normalized transition eigenvalue $-1$, so its second-largest eigenvalue modulus is 1, even if nontrivial modes satisfy a Ramanujan condition. Universal TV contraction does not follow without excluding bipartiteness or using a lazy chain.
\item Law 86: ``if and only if $M\ge\exp(O(n^{1/3}))$'' is not an exact necessary-and-sufficient threshold. Big-O does not specify one, and an upper bound is not a lower bound without proof.
\item Law 94: preservation of Vlasov phase volume and Casimir norms is not equivalent to global continuation of smooth solutions for arbitrary data. OpenAI Math 362 requires separate checking of geometric, functional, and algorithmic hypotheses.
\item Law 100: exact nodal area $2\pi R^2\ell$ for an arbitrary ``sectoral'' mode does not follow from Donnelly--Fefferman without defining the mode and nodal-set regularity.
\item Law 113: the mean-Casimir identity $[\sum_{k=1}^Kk^2(k^2-1)]/[\sum_{k=1}^Kk^2]=(3K^2+3K-6)/5$ is correct, but is a known application of Faulhaber formulas, not a new physical law.
\item Methodology: five tests in the submitted script cannot check all 50 laws. Its content does not support the claim of ``all 50 rigorous audits.''
\end{enumerate}
\section{Verifiable primary sources and novelty limits}
Budich and Diehl, \emph{Topology of density matrices}, Phys. Rev. B 91, 165140 (2015), DOI 10.1103/PhysRevB.91.165140; Budich, Zoller and Diehl, \emph{Dissipative preparation of Chern insulators}, Phys. Rev. A 91, 042117 (2015), DOI 10.1103/PhysRevA.91.042117. Nothing here implies protection solely by a Liouvillian gap or computation of a fundamental gravitational coefficient.

\chapter{Complete Analytical Classification of Three-Prime Periodic Memory and a Standalone Manuscript}
\note{Scientific specification.} This section establishes the exact minimax staircase for a periodic spike with three distinct prime divisors. It completes this special case of L.1 without restoring the refuted universal block formula. Unlike a preliminary proof, the $2b/5$ lower bound follows directly from divisor incidences without implicitly assuming pairwise coprime periods. All results use the physically normalized homogeneous Markov class of (337).
\section{Canonical thresholds and the complete theorem}
Let
\begin{equation}n=q_1q_2q_3,\qquad2\le q_1<q_2<q_3,\qquad q_i=p_i^{\alpha_i},\quad\alpha_i\ge1,\tag{383}\end{equation}
with distinct primes. For $f_{n,b}(k)=b\one_{n\mid k}$, $0<b\le1$, define
\begin{equation}K_\Sigma=q_1+q_2+q_3,\qquad K_{\rm pair}=q_1q_2+q_3,\qquad M_*=q_1q_2+q_1q_3+q_2q_3.\tag{384}\end{equation}
Always $K_\Sigma<K_{\rm pair}<n$ and $K_{\rm pair}<M_*$. Depending on arithmetic, $M_*$ may lie below or at/above $n$.
\statement{Theorem}{LXI.1}{The complete three-prime-power staircase}
For every integer $m\ge1$,
\begin{equation}E_m(f_{n,b})=\begin{cases}
b/2,&1\le m<K_\Sigma,\\
2b/5,&K_\Sigma\le m<K_{\rm pair},\\
b/3,&K_{\rm pair}\le m<\min\{M_*,n\},\\
b/4,&M_*\le m<n\quad\text{if }M_*<n,\\
0,&m\ge n.
\end{cases}\tag{385}\end{equation}
No levels between $b/3$ and $b/4$ occur when $\omega(n)=3$; the last nonzero level is absent if $q_1^{-1}+q_2^{-1}+q_3^{-1}\ge1$.
\section{Incidence witnesses for lower bounds}
Use LVI.2 to reduce, for $m<n$, to
\begin{equation}g(k)=\sum_{d\in D}h_d(k),\qquad h_d(k+d)=h_d(k),\quad h_d(k)\ge0,\quad\sum_{d\in D}d\le m,\tag{386}\end{equation}
where $D$ consists of distinct proper divisors. Put $t_i=n/p_i$ and
\begin{equation}I(d)=\{i\in\{1,2,3\}:d\mid t_i\},\qquad h_d(t_i)=h_d(0)\quad(i\in I(d)).\tag{387}\end{equation}
Every $I(d)$ is nonempty because a proper divisor lacks at least one full prime power.
\statement{Lemma}{LXI.2}{A double incidence cover}
If every period contains at most one full prime-power factor $q_i$, then $|I(d)|\ge2$ and
\begin{equation}2g(0)\le g(t_1)+g(t_2)+g(t_3).\tag{388}\end{equation}
Its spike approximation error is therefore at least $2b/5$.
\begin{proof}
$d\mid n/p_i$ is equivalent to $v_{p_i}(d)<\alpha_i$, that is, absence of the full factor. At least two factors are absent. Thus $\sum_ih_d(t_i)\ge|I(d)|h_d(0)\ge2h_d(0)$. Sum over periods, then use $g(t_i)\le\varepsilon$ and $g(0)\ge b-\varepsilon$, giving $2(b-\varepsilon)\le3\varepsilon$.
\end{proof}
\statement{Lemma}{LXI.3}{A two-point cover when an exclusive cycle is absent}
If $m<M_*$, some index $i_0$ has no used period with $I(d)=\{i_0\}$. For the other two indices $j,k$,
\begin{equation}g(0)\le g(t_j)+g(t_k),\qquad\varepsilon\ge b/3.\tag{389}\end{equation}
\begin{proof}
A singleton incidence $\{i\}$ requires $d\ge q_jq_k=n/q_i$. All three singleton types would therefore cost at least $M_*$. For a missing type $i_0$, every nonempty $I(d)$ intersects $\{j,k\}$, so $h_d(0)\le h_d(t_j)+h_d(t_k)$. Sum and use $b-\varepsilon\le2\varepsilon$.
\end{proof}
\statement{Lemma}{LXI.4}{The budget for the first full arithmetic cover}
If $n\mid\lcm(D)$, then $\sum_{d\in D}d\ge K_\Sigma$. If this sum is below $K_{\rm pair}$, no period contains two distinct full prime powers.
\begin{proof}
Assign each covered $q_i$ to one period that contains it. This partitions indices into blocks and lower-bounds the total cost by the sum of their products. For distinct prime powers $a,b\ge2$, $ab\ge a+b$, so the cost is at least $K_\Sigma$. If a period contains $q_iq_j$, it costs at least that product. Since $m<n$, another period must cover $q_k$, costing at least $q_k$. The sum is at least $q_iq_j+q_k\ge q_1q_2+q_3=K_{\rm pair}$, by the ordering of $(q_i-1)(q_j-1)$.
\end{proof}
\section{Proof of optimality of all five steps}
\begin{proof}[Proof of Theorem LXI.1]
For $m<n$, if $n\nmid L=\lcm(D)$, then $g(jL)=g(0)$ while the target takes both 0 and $b$ on that subsequence, so error is at least $b/2$. A constant response $b/2$ always attains that bound.

If the periods cover $n$, Lemma LXI.4 makes $m<K_\Sigma$ impossible. For $K_\Sigma\le m<K_{\rm pair}$, all periods contain at most one full factor, and LXI.2 gives $2b/5$. Cycles $q_1,q_2,q_3$, each with initial mass $1/3$ and phase-zero effect $3b/5$, attain it through $g=(b/5)\sum_i\one_{q_i\mid k}$. At multiples of $n$, error is $2b/5$; elsewhere at most two indicators are active.

For $K_{\rm pair}\le m<\min(M_*,n)$, LXI.3 gives $b/3$. Cycles $q_1q_2,q_3$, each of mass $1/2$ with phase-zero effect $2b/3$, give $g=(b/3)(\one_{q_1q_2\mid k}+\one_{q_3\mid k})$. Their lcm is $n$, giving exact error $b/3$.

For $M_*\le m<n$, LVI.3 gives $b/4$. Cycles $q_1q_2,q_1q_3,q_2q_3$, each of mass $1/3$ and phase-zero effect $3b/4$, give $g=(b/4)(\one_{q_1q_2\mid k}+\one_{q_1q_3\mid k}+\one_{q_2q_3\mid k})$. Any pair has lcm $n$, so at most one is active away from multiples and the error is exactly $b/4$.

For $m\ge n$, one $n$-cycle with phase-zero effect $b$ is exact. Unused states can be unreachable in every construction.
\end{proof}
\statement{Corollary}{LXI.5}{Arithmetic criterion for the fourth step}
A $b/4$ layer below $n$ exists exactly when $q_1^{-1}+q_2^{-1}+q_3^{-1}<1$. It is absent for $n=30$, begins at $m=41$ for 42, at 47 for 60, and at 71 for 105.
\section{Reproducibility and scope}
Independent scalar LP reduction enumerates all admissible divisor sets, including overlapping periods. \path{checks/verify_v56_three_primes_staircase.py} checks all budgets $1\le m\le n$ for $n=30,42,60,66,70,78,84,90,105$, using SciPy's independent HiGHS solver. \path{checks/VERIFICATION_V56.json} stores problem counts and steps. Numerical optima do not replace the analytical proof. International priority has not been checked; the full staircase for more distinct prime divisors remains separate.
\section{One file: removing external directives}
In this revision, historical Chapters LIV and LV are included directly in the main \LaTeX{} source instead of two \verb|\input| commands. Their original texts are also retained separately for reproducibility, but compiling \path{META_CAUSAL_MASTER.tex} no longer depends on them. The proposed alternative LIV swapping kinetic and potential helicity was rejected: inviscid Euler on a fixed closed sphere conserves $\int u\cdot\operatorname{curl}u\,dV$, not generally $\int(\operatorname{curl}^{-1}u)\cdot u\,dV$.
\setcounter{equation}{389}


% 18_hypergraphs.tex
\chapter{The Divisor Hypergraph, Exact Memory for Period 210, and the Corner--Entropy--Flux P0 Protocol}
\note{Scientific status.} The following reduction is an exact theorem within the class of finite homogeneous stochastic realizations with an arbitrary normalized measured effect $0\le e_j\le1$. The result for $n=210$ is accompanied by an exhaustive, rational, computer-checkable certificate; it is neither a Lean proof nor an establishment of scientific priority. The physical Corner--Entropy--Flux P0 program is studied separately as an open test, not as a consequence of Markov combinatorics.
\section{The normalized operational setting}
Let $n\ge2$, $0<b\le1$, and $f_{n,b}(k)=b\one_{\{n\mid k\}}$ for $k\in\N_0$, and define
\begin{equation}
E_m(f_{n,b})=\inf_{\substack{P\in\R_{\ge0}^{m\times m},\ \one^\top P=\one^\top\\\pi\ge0,\ \one^\top\pi=1,\ 0\le e\le1}}\ \sup_{k\ge0}|e^\top P^k\pi-f_{n,b}(k)|.\tag{390}
\end{equation}
All limits and averages below refer to the tail of the homogeneous chain, so that transients cannot produce spurious periodic witnesses.

Let $p_1,\ldots,p_r$ be the distinct prime divisors of $n$, and let $q_i=p_i^{\alpha_i}$ be their maximal powers in $n$. Write
\begin{equation}
I(d):=\{i\in\{1,\ldots,r\}:d\mid n/p_i\},\qquad d\mid n,\ d<n.\tag{391}
\end{equation}
This set is nonempty: every proper divisor of $n$ lacks at least one maximal prime power. The canonical divisor with the same signature is $\bar d=\prod_{i\notin I(d)}q_i\mid d$. Thus replacing $d$ by $\bar d$ never increases the cost of a period and leaves its incidence unchanged.

\statement{Theorem}{LXII.1}{Exact reduction of spike Markov memory to a fractional hypergraph game}
For any $1\le m<n$, consider all nonempty sets of distinct proper divisors $D\subset\{d:d\mid n,\ d<n\}$ with cost $\sum_{d\in D}d\le m$. Define
\begin{equation}
c(D):=\min_{\substack{x_d\ge0\\\sum_{d\in D}x_d=1}}\ \max_{1\le i\le r}\sum_{\substack{d\in D\\i\in I(d)}}x_d.\tag{392}
\end{equation}
Then the following exact identity holds:
\begin{equation}
E_m(f_{n,b})=b\min_{\substack{\varnothing\ne D\subset\{d\mid n,\ d<n\}\\\sum_{d\in D}d\le m}}\frac{c(D)}{1+c(D)}.\tag{393}
\end{equation}
The dual form of the fractional game is
\begin{equation}
c(D)=\max_{\substack{y_i\ge0\\\sum_i y_i=1}}\ \min_{d\in D}\sum_{i\in I(d)}y_i.\tag{394}
\end{equation}
It suffices to minimize over canonical divisors $\bar d$ and supports $D$ of cardinality at most $r$. For $m\ge n$, the error is exactly zero.

\begin{proof}
For a finite chain, the periodic part of its asymptotics is formed by closed recurrent classes of periods $d_j$ and normalized nonnegative profiles $h_j(k\bmod d_j)$ bounded by initial masses $w_j$, with $\sum_jw_j=1$. Passing to this asymptotic part does not increase the upper bound on the uniform error. Averaging it over shifts $k\mapsto k+tn$ preserves the $n$-periodic target and the error bound; each period $d_j$ is thereby replaced by $\gcd(d_j,n)\le d_j$. Equal resulting periods are combined, without increasing the total cost.

Because the target vanishes outside $k\equiv0\pmod n$, every profile $h_d$ can be set to zero at all residues other than zero, leaving $g(0)$ unchanged and not increasing deviations at $k\not\equiv0\pmod n$. This gives $g(k)=\sum_{d\in D:d\mid k}t_d$, where $t_d\ge0$ and $T:=\sum_dt_d\le1$. For any $k\not\equiv0\pmod n$, the divisor $a=\gcd(k,n)$ is proper. Choose a prime $p_i$ such that $a\mid n/p_i$; the set of periods active at $k$ is then contained in the set active at $k_i=n/p_i$. Hence
\[
\max_{k\not\equiv0\ (\mathrm{mod}\ n)}g(k)=\max_i\sum_{d:i\in I(d)}t_d.
\]
For $T>0$, write $t_d=Tx_d$. Minimization over $x$ gives $c(D)$, and the amplitude problem reduces to $\min_{0\le T\le1}\max\{|b-T|,Tc(D)\}$. Since $0<b\le1$ and $c(D)>0$, its optimum occurs at $T=b/(1+c(D))$ and equals $bc(D)/(1+c(D))$. This optimum is physically realizable: prepare a direct sum of deterministic cycles of periods $d\in D$ with initial probabilities $x_d$, and set the observed effect to $b/(1+c(D))$ only at the zero phase of each cycle. The normalized measured response is $g(k)=T\sum_{d\mid k}x_d$. For $m\ge n$, use one exact cycle of length $n$.

Duality (394) is the finite-dimensional minimax theorem for a bilinear game on two simplices. Replacing $d$ by its canonical divisor preserves $I(d)$ and decreases the cost. Finally, the linear program (392) has $r$ load inequalities and one simplex equality; an optimal basic feasible solution has at most $r$ positive $x_d$, since $c(D)>0$. Larger supports are therefore unnecessary.
\end{proof}

\section{Compression of nested cycles and strict divisor antichains}
\statement{Proposition}{LXII.2}{Exact reduction to antichains of periods}
Consider a periodic realization $g(k)=\sum_{d\in D}h_d(k\bmod d)$ with $0\le h_d(s)\le w_d$ and $\sum_{d\in D}w_d=1$. If $d_1,d_2\in D$ and $d_1\mid d_2$, the contribution of the $d_1$-cycle can be absorbed into the $d_2$-cycle without changing $g(k)$ for any $k\ge0$, without violating probability normalization, and with a budget reduction of $d_1$. Consequently, the optimum in (390) admits a realization whose periods form an antichain under divisibility. For canonical periods $d_S=\prod_{i\notin S}q_i$,
\[
d_{S_b}\mid d_{S_a}\quad\Longleftrightarrow\quad S_a\subseteq S_b.
\]
Thus, if $S_a\subsetneq S_b$, it is the shorter cycle $d_{S_b}$ that is removed, not $d_{S_a}$.
\begin{proof}
Replace the two profiles $h_{d_1}$ and $h_{d_2}$ by the single $d_2$-periodic profile $h(s)=h_{d_1}(s\bmod d_1)+h_{d_2}(s)$. At every phase, $0\le h(s)\le w_{d_1}+w_{d_2}$. It is therefore realized by a deterministic $d_2$-cycle with initial mass $w_{d_1}+w_{d_2}$ and effect $h(s)/(w_{d_1}+w_{d_2})\in[0,1]$. The sum of the responses is preserved identically, while the number of required states decreases by $d_1$. Repeating this operation for all comparable periods yields an antichain. The divisibility equivalence follows from pairwise coprimality of the maximal prime powers $q_i$.
\end{proof}

\statement{Proposition}{LXII.3}{Corrected number of simultaneous false activations}
Let $\mathcal F$ be a nonempty family of distinct nonempty missing sets $S\subseteq[r]$. Each cycle has period $d_S=n/\prod_{i\in S}q_i$, and the output spikes at zero phases have the same amplitude $u\ge0$. Set $M=|\mathcal F|$ and
\[
\nu(\mathcal F)=\max_{1\le i\le r}|\{S\in\mathcal F:i\in S\}|.
\]
Then $\max_{k\not\equiv0\ (\mathrm{mod}\ n)}g(k)=\nu(\mathcal F)u$, $g(0)=Mu$, and the optimal amplitude $u=b/(M+\nu)$ gives uniform error
\[
\varepsilon_{\mathcal F}=\frac{b\nu(\mathcal F)}{M+\nu(\mathcal F)}.
\]
This is the exact error of the stated equal-amplitude construction, not the universal optimum over unequal weights.
\begin{proof}
At $k_i=n/p_i$, we have $d_S\mid k_i\Longleftrightarrow i\in S$. Every false time $k\not\equiv0\pmod n$ admits an index $i$ such that $\gcd(k,n)\mid n/p_i$; all cycles active at $k$ are therefore also active at $k_i$. The maximum false response is $\nu u$, and the spike at $k=0$ is $Mu$. Minimizing $\max\{|b-Mu|,\nu u\}$ gives $u=b/(M+\nu)$. Physical normalization is achieved with initial mass $1/M$ on each of the $M$ cycles and effect $Mu=Mb/(M+\nu)\le1$ at its distinguished phase.
\end{proof}

\section{The complete optimal staircase for $n=210$}
\statement{Theorem}{LXII.4}{A rationally certified seven-step staircase for period 210}
For $n=210=2\cdot3\cdot5\cdot7$ and $0<b\le1$, the minimax error in (390) is
\begin{equation}
\frac{E_m(f_{210,b})}{b}=\begin{cases}
1/2,&1\le m\le16,\\
3/7,&m=17,\\
2/5,&18\le m\le28,\\
1/3,&29\le m\le99,\\
2/7,&100\le m\le106,\\
1/4,&107\le m\le209,\\
0,&m\ge210.
\end{cases}\tag{395}
\end{equation}
In particular, the optimal four-cycle and three-cycle constructions in the final intermediate layers are
\begin{equation}
g_{100}(k)=\frac b7\bigl(\one_{14\mid k}+\one_{21\mid k}+\one_{35\mid k}\bigr)+\frac{2b}{7}\one_{30\mid k},\tag{396}
\end{equation}
\begin{equation}
g_{107}(k)=\frac b4\bigl(\one_{30\mid k}+\one_{35\mid k}+\one_{42\mid k}\bigr).\tag{397}
\end{equation}
\begin{proof}[Finite, rationally checkable proof]
Theorem LXII.1 reduces verification to all canonical supports of at most four periods. For $n=210$, there are $2^4-1=15$ types of nonempty missing sets $I(d)$; exactly 1905 supports of cardinality $1,2,3,4$ have total cost below 210. For each support $D$, the file \path{checks/CERTIFICATES_V57_210.json} contains rational $x_d\ge0$, $y_i\ge0$, and $c\in\Q_{>0}$ satisfying exactly
\[
\sum_dx_d=1,\qquad \sum_iy_i=1,\qquad
\sum_{d:i\in I(d)}x_d\le c\ (\forall i),\qquad
\sum_{i\in I(d)}y_i\ge c\ (\forall d).
\]
By weak duality, these two feasible solutions establish $c=c(D)$ without rounding. The self-contained program \path{checks/verify_v57_rational.py}, using only \texttt{fractions.Fraction}, re-enumerates all 1905 supports, checks every certificate, and takes the exact lower cost envelope. This yields the only undominated thresholds $m=17,18,29,100,107$ and the exact values (395). Before $m=17$, the error is the constant $b/2$; at $m=210$, an exact cycle exists. For computational independence, certificate generation and the main verification script are provided separately.
\end{proof}

\statement{Remark}{LXII.5}{Checking a certificate versus solving an optimization problem}
The generation script may use SciPy/HiGHS to find candidates, but the verifier checks the primal and dual rational witnesses, exhaustive support enumeration, and final staircase using exact integer operations alone. The proof does not establish worldwide priority, the physical origin of a particular Markov process, or a global classification for all $r$ by a single closed threshold formula.

\section{Exact expansion of hypergraph certificates into the full temporal LP}
\statement{Proposition}{LXII.6}{A compressed primal--dual certificate over the entire period}
Let $n=\prod_{i=1}^rp_i^{\alpha_i}$, let $D$ be a nonempty support of proper divisors, and let $x_d,y_i,c\in\Q$ satisfy the rational conditions in the proof of Theorem LXII.4. For $b=1$, set
\begin{equation}
\varepsilon=\frac c{1+c},\qquad w_d=x_d,\qquad h_{d,0}=\frac{x_d}{1+c},\qquad h_{d,s}=0\quad(s\ne0).\tag{398}
\end{equation}
For the dual of the original phase LP, written as $A_DX\le b_D$, $C_DX=1$, $X\ge0$, with dual multipliers $Y\le0$ and free $z$, only the following nonzero coordinates are needed:
\begin{equation}
Y_{\mathrm{peak},0}=-\frac c{1+c},\qquad
Y_{\mathrm{false},n/p_i}=-\frac{y_i}{1+c},\qquad z=0,\qquad Y_{\mathrm{cap}}=0.\tag{399}
\end{equation}
All primal constraints throughout $k=0,\ldots,n-1$ and all dual constraints then hold exactly. Both objective values equal $\varepsilon$, so the duality gap is zero.
\begin{proof}
For $k\not\equiv0\pmod n$, choose $i$ such that $\gcd(k,n)\mid n/p_i$. Every cycle active at $k$ is also active at $n/p_i$. Thus $g(k)\le c/(1+c)=\varepsilon$, while $g(0)=1/(1+c)=1-\varepsilon$. The bound $h_{d,0}\le w_d$ and normalization $\sum_dw_d=1$ are immediate. In the dual problem, the $\varepsilon$ column satisfies $-Y_{\mathrm{peak},0}-\sum_iY_{\mathrm{false},n/p_i}=1$. For the $h_{d,0}$ column, use $\sum_{i:d\mid n/p_i}y_i\ge c$ to obtain $-Y_{\mathrm{peak},0}+\sum_{i:d\mid n/p_i}Y_{\mathrm{false},n/p_i}\le0$. For $h_{d,s}$ with $s\ne0$, the sum of dual contributions is nonpositive; for $w_d$, it is zero. The dual objective is $-Y_{\mathrm{peak},0}=\varepsilon$, equal to the primal objective.
\end{proof}

\statement{Remark}{LXII.7}{Compact exhaustive certification for period 210}
The 15 proper divisors of period 210 have exactly 166 nonempty divisibility antichains, of which 164 have cost below 210. Among all $2^{15}-1=32767$ nonempty subsets, exactly 20684 satisfy a budget $m<210$. Absorbing comparable cycles reduces each of them to one of the 164 antichains of no greater cost, preserving the response. The file \path{checks/LP210_164_ANTICHAIN_EXACT_CERTIFICATES.json} contains 164 compact rational triples $(x,y,c)$; the standalone \path{checks/verify_v62_full_lp_210.py} checks the entire period and exact agreement of the dual objectives. Nine additional explicit control supports, including the critical cost 100 and a separate relaxation of cost 261, which is not a model with budget 209, are stored in \path{checks/LP210_9_EXPLICIT_FULL_LP_WITNESSES.json}. The catalogue of 1905 historically independent supports of cardinality at most four is retained for reproducible cross-checking; it is not required to use the shorter antichain certificate.

\section{The independent P0 direction: Corner--Entropy--Flux and its scope}
As of 14 October 2026, three recent preprints have been compared:
\begin{enumerate}[label=(\alph*)]
\item J. L. V. Cerdeira and T. Ort\'in, \emph{On the quantum statistical relation and the covariant phase space Hamiltonian}, arXiv:2610.13377 (first submitted 9 October 2026): the relation of the full boundary--gauge variational principle to the covariant Hamiltonian and Euclidean action. \url{https://arxiv.org/abs/2610.13377}
\item N. Bobev, J. Boruch, L. V. Iliesiu, and D. Maskalaniec, \emph{Proving the OSV conjecture with the gravitational path integral}, arXiv:2610.12550 (8 October 2026): localization of gravitational saddle contributions to a special black-hole index in four-dimensional $\mathcal N=2$ supersymmetric theory. \url{https://arxiv.org/abs/2610.12550}
\item W. Arata et al., \emph{Asymptotic AEther Misalignment and Universal-Horizon Mechanics}, arXiv:2610.12559 (8 October 2026): nonintegrability of the energy under independent variation of two boundary data and an additional aether-flux contribution to universal-horizon thermal balance. \url{https://arxiv.org/abs/2610.12559}
\end{enumerate}
These works do not prove the identity $Z^{\mathrm{ren}}_{\mathrm{corner/Type\text{-}II}}=Z^{\mathrm{ren}}_{\mathrm{Euclidean}}$. Here it is formulated solely as a test with identical physical boundary data, ensemble, normalization, gauge-fixing/ghost/edge measures, counterterms, and quantum corrections. The first calculation supplies a testable classical boundary bridge; the second concerns a protected index rather than general Type-II entropy; the third rules out an uncritical universal identification of temperature with the area coefficient in all modified theories. P0 status: open.

\chapter{Universal Determinant Signatures of Fractional Hypergraph Covers}
\note{Scientific status and scope.} An exact, finite classification, unconditional within definition (390), is obtained for all optimal stochastic approximations of a pure periodic spike at arbitrary $r=\omega(n)$ through the rational vertices of one universal polytope. This is a classification by a finite determinant formula, not a proof that a short threshold list of polynomial size in $r$ exists. Fractional matching/cover duality and the standard vertex theory of linear polyhedra are known; the new element within this manuscript is their joint application to the budget of physically normalized Markov cycles. Independent worldwide priority has not been established.

\section{Exact identification of the memory game with fractional matching}
Let $n=\prod_{i=1}^rq_i$, where $q_i=p_i^{\alpha_i}$, the $p_i$ are pairwise distinct, and $0<b\le1$. For each nonempty $S\subseteq[r]$, define its canonical period and cost by
\begin{equation}
d_S(q):=\prod_{i\notin S}q_i,\qquad C_q(\mathcal F):=\sum_{S\in\mathcal F}d_S(q).\tag{400}
\end{equation}
The set $S$ equals $I(d_S)$ in (391). Let $A_r=(\one_{i\in S})_{i\in[r],\ \varnothing\ne S\subseteq[r]}$ be the full incidence matrix of all nonempty subsets. Consider the universal polytope
\begin{equation}
\mathcal P_r:=\{z\in\R_{\ge0}^{2^r-1}:A_rz\le\one_r\},\qquad
\nu(z)=\sum_Sz_S,\qquad C_q(z)=\sum_{S:z_S>0}d_S(q).\tag{401}
\end{equation}
Unlike the original Markov matrices, $\mathcal P_r$ does not depend on the numerical values of $q_i$. It is bounded: every hyperedge is nonempty, so $0\le z_S\le1$.

\statement{Theorem}{LXIII.1}{Exact correspondence between the hypergraph game, fractional matching, and memory}
For a nonempty set of canonical cycles $\mathcal F$, let $\nu^*(\mathcal F)$ be the maximum total mass of a fractional matching in the hypergraph $([r],\mathcal F)$:
\[
\nu^*(\mathcal F)=\max\left\{\sum_{S\in\mathcal F}z_S:z_S\ge0,\ \sum_{S\ni i}z_S\le1\ (i\in[r])\right\}.
\]
Then the minimax load $c(\mathcal F)$ in (392) and the optimal error are exactly
\begin{equation}
c(\mathcal F)=\frac1{\nu^*(\mathcal F)},\qquad
E_m(f_{n,b})=\frac b{1+\max_{\mathcal F:C_q(\mathcal F)\le m}\nu^*(\mathcal F)}\quad(1\le m<n).\tag{402}
\end{equation}
The second formula admits an explicit optimal stochastic realization from any maximum fractional matching $z$: with $\nu=\sum_Sz_S$,
\begin{equation}
g_z(k)=\frac b{1+\nu}\sum_{S:d_S\mid k}z_S,\qquad
\sup_{k\ge0}|g_z(k)-f_{n,b}(k)|=\frac b{1+\nu}.\tag{403}
\end{equation}
\begin{proof}
Set $x_S=z_S/\nu$ and $c=1/\nu$. Then $\sum_Sx_S=1$ and $\sum_{S\ni i}x_S\le c$. Conversely, from any feasible $x$ in game (392) with $c>0$, set $z_S=x_S/c$ to obtain a feasible fractional matching of mass $1/c$. The optimizations are reciprocal; combine this with the rigorously proved reduction LXII.1. For an explicit construction, take independent deterministic cycles $d_S$, initial masses $z_S/\nu$, and effect $b\nu/(1+\nu)$ at the zero phase of each. Normalization holds because $b\nu/(1+\nu)\le1$. At multiples of $n$, the value is $b\nu/(1+\nu)$; at other times it is at most $b/(1+\nu)$ by $A_rz\le\one$, since the active periods are contained in one set $\{d_S:S\ni i\}$ for a suitable $i$ (Theorem LXII.1).
\end{proof}

\section{Universal classification by integer determinants}
\statement{Theorem}{LXIII.2}{A complete determinant catalogue of optima for arbitrary $r$}
Let $\mathcal V_r$ denote the nonzero vertices of (401). Each vertex is specified by a finite combinatorial certificate:
\begin{enumerate}[label=(\roman*)]
\item Choose $1\le s\le r$, pairwise distinct nonempty hyperedges $\mathcal F=\{S_1,\ldots,S_s\}$, and $s$ distinct row indices $I\subseteq[r]$.
\item The square zero--one matrix $B=A_r[I,\mathcal F]$ is invertible.
\item $z_{\mathcal F}=B^{-1}\one_s$ is strictly positive and $A_r[:,\mathcal F]z_{\mathcal F}\le\one_r$; the other coordinates of $z$ vanish.
\end{enumerate}
Conversely, every such construction gives a vertex, with repeated certificates identified. Cramer's rule expresses the coordinates through integer determinants:
\begin{equation}
z_{S_j}=\frac{\det B_j}{\det B},\qquad
\nu(z)=\frac{\sum_{j=1}^s\det B_j}{\det B}.\tag{404}
\end{equation}
For any $1\le m<n$, the minimax error has the following closed finite formula, containing no optimization variables, stochastic matrices, or linear programs:
\begin{equation}
E_m(f_{n,b})=\frac b{1+\max\{\nu(v):v\in\mathcal V_r,\ C_q(v)\le m\}}.\tag{405}
\end{equation}
For $m\ge n$, the error is zero. The set of possible values $\nu(v)$ depends only on $r$, all values are rational, and at most $r$ active optimal cycles are needed.
\begin{proof}
On each finite feasible support $\mathcal F$, the linear function $\sum_Sz_S$ attains its maximum at a vertex of the corresponding face $\{z\in\mathcal P_r:z_S=0\ (S\notin\mathcal F)\}$. Since this is a face of $\mathcal P_r$, its vertex is also a vertex of $\mathcal P_r$. Passing to a vertex does not enlarge the support or increase $C_q$. A set of at most $r$ positive coordinates specifies a vertex if and only if the active constraints $A_rz\le\one_r$ contain $s$ equations linearly independent in those coordinates. These are precisely conditions (i)--(iii); (404) follows from Cramer's rule. Substitution in (402) proves (405). Boundedness and rationality of $\mathcal P_r$ guarantee a finite catalogue. For $s\le r$, Hadamard's inequality gives $|\det B|\le s^{s/2}$, so the denominators of the masses have universal bounds for fixed $r$.
\end{proof}

\statement{Corollary}{LXIII.3}{Universal thresholds and arithmetic chambers}
Let $\Lambda_r=\{\nu(v):v\in\mathcal V_r\}$ and $\lambda\in\Lambda_r$. The exact minimum budget for attaining a load of at least $\lambda$ is
\begin{equation}
K_r(\lambda;q)=\min_{\substack{v\in\mathcal V_r\\\nu(v)\ge\lambda}}C_q(v).\tag{406}
\end{equation}
Every $C_q(v)$ is a polynomial in the $q_i$ with coefficients zero or one, composed of square-free monomials. The ordering of thresholds and the set of realized steps are therefore constant in each chamber of a finite semialgebraic partition of $q_i>1$, defined by the signs of $C_q(v)-C_q(w)$. On the walls, optimal models may coincide.
\begin{proof}
Interchange the finite maximum and minimum operations in (405); polynomiality follows directly from (400).
\end{proof}

\section{Complete exact classification for small $r$: 15 and 94 vertices}
An independent rational algorithm enumerates all admissible $\mathcal F,I$, inverts the corresponding matrices by Gaussian elimination over $\Q$, and identifies duplicates. The exact numbers of vertices and permutation orbits under $S_r$ are:
\begin{center}
\begin{tabular}{cccc}\toprule
$r$&$|\mathcal V_r|$&$|\mathcal V_r/S_r|$&Values of $\nu(v)$\\\midrule
2&4&3&$1,2$\\
3&15&7&$1,\frac32,2,3$\\
4&94&18&$1,\frac43,\frac32,\frac53,2,\frac52,3,4$\\\bottomrule
\end{tabular}
\end{center}
For $r=4$, the frequencies of the listed levels among the 94 vertices are, respectively, $15,1,22,4,37,4,10,1$. The original machine-checkable catalogue contains all 18 isomorphism types with rational weights, not merely their orbit counts. For example, the type $\nu=5/3$ has four permutation representatives and consists of four edges:
\begin{equation}
\begin{gathered}
(S_1,S_2,S_3,S_4)=(\{1,2,3\},\{1,4\},\{2,4\},\{3,4\}),\\
z=(2/3,1/3,1/3,1/3).
\end{gathered}\tag{407}
\end{equation}
Its cost is
\begin{equation}
C_{5/3}(q;4)=q_4+q_2q_3+q_1q_3+q_1q_2,\tag{408}
\end{equation}
with analogous variants for distinguished index 1, 2, or 3. By (405), this configuration gives feasible error $3b/8$ at the stated budget, but an actual plateau occurs only if no type with a higher load is cheaper.

\statement{Proposition}{LXIII.4}{An explicit optimal $3b/8$ plateau absent for $n=210$}
For $n=1050=2\cdot3\cdot7\cdot25$ and the pure spike $f_{1050,b}$, exactly
\begin{equation}
\frac{E_m(f_{1050,b})}{b}=\begin{cases}
1/2,&1\le m\le36,\\
3/7,&m=37,\\
2/5,&38\le m\le65,\\
3/8,&m=66,\\
1/3,&67\le m\le341,\\
2/7,&342\le m\le366,\\
1/4,&367\le m\le1049,\\
0,&m\ge1050.
\end{cases}\tag{409}
\end{equation}
At $m=66$, the optimal normalized response is given by the exact formula
\begin{equation}
\begin{gathered}
g_{66}(k)=\frac b4\one_{25\mid k}+\frac b8\bigl(\one_{21\mid k}+\one_{14\mid k}+\one_{6\mid k}\bigr),\\
\sup_{k\ge0}|f_{1050,b}(k)-g_{66}(k)|=\frac{3b}{8}.
\end{gathered}\tag{410}
\end{equation}
\begin{proof}
For (407) and $q=(2,3,7,25)$, the cost is $25+21+14+6=66$, and $\nu=5/3$. Formula (403) gives coefficients $b/(1+5/3)\times(2/3,1/3,1/3,1/3)=(b/4,b/8,b/8,b/8)$. Full minimality of each of the seven thresholds follows from Theorem LXIII.2 and exhaustive rational enumeration of all 94 vertices with their costs $C_q$; the exact data and verifier are included in the archive. No floating-point approximations are required for this statement.
\end{proof}

\section{Two infinite families of rational certificates}
\statement{Proposition}{LXIII.5}{Cyclic uniform $k$-hyperedges}
For $1\le k<r$, let $S_j=\{j,j+1,\ldots,j+k-1\}\pmod r$ for $j=1,\ldots,r$, and assign $z_{S_j}=1/k$ to each of these $r$ hyperedges. Then $A_rz=\one_r$, $\nu(z)=r/k$, and an explicit feasible construction is
\begin{equation}
C_{r,k}(q)=\sum_{j=1}^r\prod_{i\notin S_j}q_i,\qquad
E_{C_{r,k}(q)}(f_{n,b})\le\frac{bk}{r+k}.\tag{411}
\end{equation}
If $\gcd(r,k)=1$, the corresponding circulant incidence matrix is invertible and $z$ is a vertex of $\mathcal P_r$. This does not assert optimality at this budget compared with other types in (405).
\begin{proof}
Every hypergraph vertex belongs to exactly $k$ of the stated hyperedges, so its total weight is one. The circulant spectrum consists of the sums $1+\omega^h+\cdots+\omega^{(k-1)h}$, where $\omega=e^{2\pi\ii/r}$; a nontrivial zero eigenvalue exists if and only if $\gcd(r,k)>1$. In the coprime case, this gives a nonsingular matrix $B$ in Theorem LXIII.2.
\end{proof}

\statement{Proposition}{LXIII.6}{Fractional projective-plane signatures}
Suppose a finite projective plane of order $t$ exists, for example when $t$ is a prime power. On its $r=t^2+t+1$ points, take the $r$ lines, each containing $t+1$ points. Assigning weight $1/(t+1)$ to each edge yields a vertex of the universal $\mathcal P_r$ with
\begin{equation}
\nu=\frac{t^2+t+1}{t+1},\qquad
E_{C_q}(f_{n,b})\le\frac{b(t+1)}{t^2+2t+2},\qquad
C_q=\sum_{S\text{ a line}}\prod_{i\notin S}q_i.\tag{412}
\end{equation}
For the Fano plane ($t=2$, $r=7$), this gives the rational certificate $\nu=7/3$ and error at most $3b/10$, provided $C_q<n$.
\begin{proof}
Each point lies on exactly $t+1$ lines, so all $r$ constraints are saturated. The incidence matrix $B$ is square and $BB^\top=tI_r+J_r$: every two distinct points lie on exactly one common line, and each point lies on $t+1$ lines. For $t\ge1$, $BB^\top$ is positive definite, so $B$ is invertible and the certificate is a vertex. Substitute $\nu$ into (403).
\end{proof}

\section{The limits of closed-form classification and independent checkability}
All formulas in this chapter concern the exact minimax problem for a pure spike and a stationary classical Markov realization. For a spike with baseline $a>0$, a time-varying calibration, nonstationary chains, or continuous time, this catalogue alone is not a complete classification. The determinant formula is closed for each $r$, but the number of hyperedges $2^r-1$ and the number of vertices may grow superexponentially: it does not imply a computationally efficient algorithm, a constant-size list, or a solution to any Millennium Prize Problem. The literature already contains classical LP duality for fractional matchings and covers; independent scientific priority of this specialization has not been checked. The verifier \path{checks/verify_v58_universal_vertices.py} uses \texttt{Fraction}, independently generates the vertices for $r\le4$, and reproduces the staircases for $n=210$ and $n=1050$.
\setcounter{equation}{412}


% 19_structural_thresholds.tex
\chapter{The Structural Skeleton of Optimal Hypergraph Memory without Vertex Enumeration}
\note{Scope.} This chapter considers only the pure $n$-periodic spike $f(k)=b\one_{\{n\mid k\}}$, $0<b\le1$, finite homogeneous classical Markov realizations, and the hypergraph reduction of Chapter LXIII. It proves a short structural classification of the upper boundary of the spectrum of optimal loads for arbitrary $r=\omega(n)$, not a polynomial classification of all lower levels. The vertex characterization for fractional matchings used here is well known; the new element is the explicit arithmetic thresholds for the specific cost $d_S=\prod_{i\notin S}q_i$.

Write $n=\prod_{i=1}^rq_i$, where $2\le q_1<q_2<\cdots<q_r$ are pairwise coprime maximal prime powers. For $\varnothing\ne S\subseteq[r]$, set $d_S=n/\prod_{i\in S}q_i$. Consider the fractional polytope
\[
\mathcal P_r=\left\{z_S\ge0:\sum_{S\ni i}z_S\le1\ (1\le i\le r)\right\},\qquad \nu(z)=\sum_Sz_S.
\]
The support cost $C_q(z)=\sum_{S:z_S>0}d_S$ is the discrete cost of including a cycle, not a linear function of $z$. For budget $1\le m<n$, Theorem LXIII.1 gives minimax error $E_m=b/(1+\max_{C_q(z)\le m}\nu(z))$.

\statement{Lemma}{LXIV.1}{Exact load deficit}
For every $z\in\mathcal P_r$,
\begin{equation}
r-\nu(z)=\sum_{i=1}^r\left(1-\sum_{S\ni i}z_S\right)+\sum_{S\ne\varnothing}(|S|-1)z_S.\tag{413}
\end{equation}
In particular, $\nu(z)\le r$, with equality only when $z_{\{i\}}=1$ for all $i$ and all other coefficients vanish.
\begin{proof}
Expand the brackets. Each $z_S$ is subtracted $|S|$ times in the first sum and added $|S|-1$ times in the second. Every term on the right is nonnegative. Their simultaneous vanishing requires the absence of hyperedges of size at least two and saturation of all rows.
\end{proof}

\statement{Theorem}{LXIV.2}{Complete classification of extremal structures with forbidden singleton hyperedges}
Let $J\subseteq[r]$ be the indices for which singleton hyperedges $\{i\}$ are forbidden, and let $k=|J|$. On the corresponding face $\mathcal P_{r,J}=\{z\in\mathcal P_r:z_{\{i\}}=0\ (i\in J)\}$, the maximum load is
\begin{equation}
\max_{z\in\mathcal P_{r,J}}\nu(z)=\begin{cases}
r,&k=0,\\r-1,&k=1,\\r-k/2,&k\ge2.
\end{cases}\tag{414}
\end{equation}
For $k\ge2$, all optimal points have $z_{\{i\}}=1$ for $i\notin J$ and form fractional perfect matchings on the complete graph with vertex set $J$: positive hyperedges within $J$ have size exactly two, and $\sum_{j\in J\setminus\{i\}}z_{\{i,j\}}=1$ for every $i\in J$. Extremal optima are disjoint unions of single edges of weight one and odd cycles with weight $1/2$ on every edge.
\begin{proof}
This follows directly from LP duality. For $k=0$, the dual price vector $y_i=1$ covers all hyperedges and gives $r$. For $k=1$, put $y_i=0$ on the unique index in $J$ and $y_i=1$ outside $J$; each permitted hyperedge satisfies $\sum_{i\in S}y_i\ge1$, yielding $\nu\le r-1$. For $k\ge2$, put $y_i=1/2$ on $J$ and $y_i=1$ outside $J$. Every permitted hyperedge satisfies $\sum_{i\in S}y_i\ge1$, giving $\nu\le r-k/2$.

To attain the bound, assign weight one to the singletons outside $J$. Inside $J$, for even $k$ use a matching with weights one; for odd $k\ge3$, use one triangle with weights $1/2$ and a matching of the remaining vertices with weights one. Equality for $k\ge2$, with a strictly positive dual vector, requires all rows to be saturated. Complementary slackness permits only hyperedges whose dual prices sum to exactly one: singletons outside $J$ and pairs inside $J$. Finally, vertices of a graph's fractional perfect matching polytope consist of single edges with weight one and odd cycles with weight $1/2$: leaves force weight one and are removed, while a remaining component with independent columns cannot contain an even cycle or two independent cycles.
\end{proof}

\statement{Theorem}{LXIV.3}{The final three structural thresholds for an arbitrary number of prime powers}
Let $r\ge3$. Define three explicit integer-valued functions:
\begin{equation}
U(q)=\sum_{i=4}^r\frac n{q_i}+\frac n{q_1q_2}+\frac n{q_1q_3}+\frac n{q_2q_3},\tag{415}
\end{equation}
\begin{equation}
V(q)=\sum_{i=3}^r\frac n{q_i}+\frac n{q_1q_2},\tag{416}
\end{equation}
\begin{equation}
M(q)=\sum_{i=1}^r\frac n{q_i}.\tag{417}
\end{equation}
Then $U(q)<V(q)<M(q)$. For all budgets $m<n$ lying in the following nonempty intervals, the exact optimal errors are
\begin{equation}
E_m(f_{n,b})=\begin{cases}
\displaystyle\frac{2b}{2r-1},&U(q)\le m<V(q),\\[5pt]
\displaystyle\frac b r,&V(q)\le m<M(q),\\[5pt]
\displaystyle\frac b{r+1},&M(q)\le m<n.
\end{cases}\tag{418}
\end{equation}
Moreover, for $m<U(q)$ the optimal load is strictly below $r-3/2$, and for $m<M(q)$ it is at most $r-1$. Thus the last three thresholds are found by sorting the $q_i$ and performing $O(r)$ arithmetic operations, without enumerating the $2^r-1$ hyperedges or the vertices of the polytope.
\begin{proof}
For the level $\nu=r$, Lemma LXIV.1 requires all $r$ singleton hyperedges, of total cost exactly $M$. If any singleton is absent, Theorem LXIV.2 gives $\nu\le r-1$.

To attain $\nu>r-3/2$, at most two singleton edges can be forbidden: for $k\ge3$, the upper bound is $r-k/2\le r-3/2$. When $k\le1$, the singleton cost is at least $M-n/q_i$, which exceeds the minimum cost $V$. For $k=2$, with missing singletons $\{i\},\{j\}$, absence of the pair $\{i,j\}$ gives the dual certificate $y_i=y_j=0$, $y_h=1$ for the other indices. All permitted hyperedges are covered, so $\nu\le r-2$. Therefore $\nu>r-3/2$ requires the pair $\{i,j\}$, and the cost is at least
\[
V_{ij}=\sum_{h\ne i,j}\frac n{q_h}+\frac n{q_iq_j}
=M-n\left(\frac1{q_i}+\frac1{q_j}-\frac1{q_iq_j}\right).
\]
Since $x+y-xy$ is strictly increasing for $0<x,y<1$, $V_{ij}$ is minimized at the two smallest prime powers $q_1,q_2$ and equals $V$. This threshold is attained by the $r-2$ singleton edges $\{3\},\ldots,\{r\}$, each of weight one, and the edge $\{1,2\}$ of weight one; their load is $r-1$.

If $\nu\ge r-3/2$ and $m<V$, the preceding paragraph excludes $k\le2$. For $k\ge4$, Theorem LXIV.2 gives $\nu\le r-2$. Hence $k=3$ and the maximum load $r-3/2$ is attained. By the structural part of Theorem LXIV.2, all singletons outside the triple $J=\{i,j,k\}$ and all three pairs within $J$ are present, the latter each of weight $1/2$. Their total cost is
\[
U_{ijk}=M-n(x_i+x_j+x_k)+n(x_ix_j+x_ix_k+x_jx_k),\qquad x_i=q_i^{-1}.
\]
Replacing one selected index $i$ by an unselected $h$ changes the cost by $n(x_h-x_i)(x_j+x_k-1)$. For pairwise coprime maximal prime powers, any two reciprocals satisfy $x_j+x_k\le1/2+1/3<1$. Thus the triple with the largest reciprocals, $J=\{1,2,3\}$, has the least cost, namely $U$. A triangle of weights $1/2$, together with the singletons outside $J$, attains load $r-3/2$.

Finally, $V-U=(n/q_3)(1-1/q_1-1/q_2)>0$, and $M-V=n/q_1+n/q_2-n/(q_1q_2)>0$. All interval formulas now follow from the exact identity $E_m=b/(1+\nu_m)$ proved in Chapter LXIII.
\end{proof}

\statement{Corollary}{LXIV.4}{Control examples and the limits of a full structural description}
For $n=210=2\cdot3\cdot5\cdot7$, we obtain $(U,V,M)=(100,107,247)$ and the exact steps $E_{100\ldots106}=2b/7$, $E_{107\ldots209}=b/4$. For $n=1050=2\cdot3\cdot7\cdot25$, we obtain $(U,V,M)=(342,367,1067)$, reproducing exactly the two rightmost steps of Chapter LXIII. For $r=3$, the functions $U$ and $V$ agree with the thresholds $K_{\rm sum}$ and $K_{\rm pair}$ of Chapter LVI.

A complete classification of all budget-constrained optima for arbitrary $r$ without exponential search has not been proved here. It would require solving the fixed support-cost problem, which permits complex irregular hypergraphs. Even fractional graph matchings contain infinite families of independent odd-cycle structures; arbitrary hyperedges introduce projective-plane types as well. The three thresholds displayed here are a genuinely closed structural classification of the upper edge for all $r$, not a disguised enumeration.
\setcounter{equation}{418}


% 20_references.tex
\chapter*{References and Origins of the Methods}
\addcontentsline{toc}{chapter}{References and Origins of the Methods}
\note{Comparison with existing research.} L. B. Vieira, S. Milz, G. Vitagliano, and C. Budroni, \emph{Witnessing environment dimension through temporal correlations}, Quantum \textbf{8}, 1224 (2024), DOI:10.22331/q-2024-01-10-1224. This work belongs to the related program of dimension witnesses but describes a different multitime measurement scheme.

\begin{enumerate}[label={[\arabic*]},leftmargin=*,itemsep=7pt]
\item S. Kiefer, \emph{Notes on Equivalence and Minimization of Weighted Automata} (2020). \url{https://arxiv.org/abs/2009.01217}.
\item R. Miclo, \emph{On contraction properties of Markov kernels}; information on the Dobrushin coefficient. \url{https://perso.math.univ-toulouse.fr/miclo/files/2012/04/contraction.pdf}.
\item S. M. Rump, \emph{Bounds for the Componentwise Distance to the Nearest Singular Matrix}, SIAM J. Matrix Anal. Appl. \url{https://doi.org/10.1137/S0895479895289170}.
\item OpenAI, \emph{Mathematical manuscript catalogue}; sources and Lean checks differ in scope. \url{https://github.com/openai/math}.
\item J. A. Tropp, \emph{User-Friendly Tail Bounds for Sums of Random Matrices}, Found. Comput. Math. \textbf{12} (2012), 389--434. \url{https://doi.org/10.1007/s10208-011-9099-z}.
\item OpenAI Math: scope documents \url{https://github.com/openai/math/tree/main/lean/docs}; families 003, 095, 144, 229, 269, 287, 350, 362, 365, 374, 376. Statements in the catalogue are not all formalized to the same extent.
\item M. C. Gonzalez-Garcia and R. Wendell, \emph{Neutrino Masses, Mixing, and Oscillations}, Review of Particle Physics (PDG, 2025), PDG review (PDF).
\item T. Ohlsson and H. Snellman, \emph{Neutrino oscillations and mixings with three flavors}, Phys. Rev. D \textbf{60}, 093007 (1999). \url{https://doi.org/10.1103/PhysRevD.60.093007}.
\item \emph{A Tutorial on the Spectral Theory of Markov Chains}, Neural Computation \textbf{35} (2023), 1713--1751, spectral review. In particular, peripheral roots of stochastic matrices are classical results.
\item A. Naumovs, M. Dimitrijevs, and A. Yakary\i lmaz, \emph{The minimal probabilistic and quantum finite automata recognizing uncountably many languages with fixed cutpoints}, DMTCS \textbf{22}:1 (2020), doi:10.23638/DMTCS-22-1-13. Quantum--classical separations are known; priority of the proposed three-flavor signature synthesis has not been checked.
\item A. Bendory, N. Boumal, C. Ma, Z. Zhao, and A. Singer, \emph{Bispectrum Inversion with Application to Multireference Alignment}, IEEE Trans. Signal Process. \textbf{66} (2018), 1037--1050. \url{https://doi.org/10.1109/TSP.2017.2775591}. Status of the preceding theory.
\item \emph{Difference set}, Encyclopedia of Mathematics. \url{https://encyclopediaofmath.org/index.php?title=Difference_set}. Parameters of cyclic $(v,k,\lambda)$-sets.
\item J. Singer, \emph{A Theorem in Finite Projective Geometry and Some Applications to Number Theory}, Trans. Amer. Math. Soc. \textbf{43} (1938), 377--385. \url{https://doi.org/10.2307/1990067}. The classical construction of cyclic difference sets.
\item L. Benvenuti, \emph{Minimal positive realizations: A survey}, Automatica \textbf{143} (2022), 110422. \url{https://doi.org/10.1016/j.automatica.2022.110422}.
\item M. Zonnios, A. Boyd, and F. C. Binder, \emph{Memory-minimal quantum generation of stochastic processes: spectral invariants of quantum hidden Markov models}, arXiv:2412.12812 (2024). \url{https://arxiv.org/abs/2412.12812}. Earlier quadratic memory in a different operational setting.
\item T. Heinosaari, O. Kerppo, L. Lepp\"ajarvi, and M. Pl\'avala, \emph{Simple information-processing tasks with unbounded quantum advantage}, Phys. Rev. A \textbf{109}, 032627 (2024). \url{https://doi.org/10.1103/PhysRevA.109.032627}. Direct prior work on single-qubit antidistinguishing matrices with sine-squared entries.
\item A. P. Goucha, J. Gouveia, and P. M. Silva, \emph{On Ranks of Regular Polygons}, SIAM J. Discrete Math. \textbf{31} (2017). \url{https://doi.org/10.1137/16M1105608}. Previous literature on nonnegative ranks and polygonal factorizations.
\item L. B. Vieira, S. Milz, G. Vitagliano, and C. Budroni, \emph{Witnessing environment dimension through temporal correlations}, Quantum \textbf{8}, 1224 (2024). \url{https://doi.org/10.22331/q-2024-01-10-1224}.
\item R. Kakarala, \emph{The Bispectrum as a Source of Phase-Sensitive Invariants for Fourier Descriptors: A Group-Theoretic Approach}, J. Math. Imaging Vis. (2012). \url{https://doi.org/10.1007/s10851-012-0330-6}. Earlier general bispectral completeness.
\item OpenAI Math, \emph{Three-State Reconstruction} (family 229), selected Lean scope: \url{https://github.com/openai/math/blob/main/lean/docs/229.md}; \emph{Forced Navier--Stokes Computation} (family 376), \url{https://github.com/openai/math/blob/main/lean/docs/376.md}. For both families, the scope of each formalization is strictly narrower than the surrounding physical and mathematical questions.
\item META--CAUSAL VIII, IX, X, XI (11 October 2026): previously prepared sections are now integrated into one continuous manuscript; an independent priority audit remains open.
\item D. Aldous and J. A. Fill, \emph{Reversible Markov Chains and Random Walks on Graphs}, Chapter 3, spectral representation of reversible chains. \url{https://www.stat.berkeley.edu/~aldous/RWG/Book_Ralph/Ch3.S4.html}.
\item J. Berstel and C. Reutenauer, \emph{Rational Series and Their Languages}, EATCS Monographs \textbf{12} (1988), Hankel minimal linear realization. \url{https://www-igm.univ-mlv.fr/~perrin/Livres/RationalSeries9April2010.pdf}.
\end{enumerate}

\end{document}
